\documentclass[10pt,a4paper]{article}
\usepackage{fontspec}
\usepackage{polyglossia}
\setmainlanguage{finnish}
\setmainfont{TeX Gyre Pagella}[Ligatures=TeX]
\newfontfamily\sansfont{TeX Gyre Heros}[Ligatures=TeX]
\newfontfamily\monofont{TeX Gyre Cursor}
\usepackage[a4paper,top=20mm,bottom=20mm,left=22mm,right=22mm]{geometry}
\usepackage{microtype}\usepackage{booktabs}\usepackage{tabularx}\usepackage{longtable}
\usepackage{enumitem}\usepackage{titlesec}\usepackage{fancyhdr}\usepackage{xcolor}
\usepackage[hidelinks]{hyperref}
\definecolor{meri}{HTML}{2C4051}\definecolor{leima}{HTML}{B8511F}\definecolor{harmaa}{HTML}{5C7385}
\setlength{\parindent}{0pt}\setlength{\parskip}{0.5em}\linespread{1.05}
\emergencystretch=3em \tolerance=2000 \hbadness=10000
\titleformat{\section}{\sansfont\bfseries\large\color{meri}}{\thesection.}{0.5em}{}
\titleformat{\subsection}{\sansfont\bfseries\normalsize\color{meri}}{\thesubsection}{0.5em}{}
\pagestyle{fancy}\fancyhf{}
\renewcommand{\headrulewidth}{0.4pt}
\renewcommand{\headrule}{\color{harmaa}\hrule height 0.4pt}
\fancyhead[L]{\sansfont\footnotesize\color{harmaa} Heimolehdistön teemat 1925--1939 · Juri Jukola}
\fancyhead[R]{\sansfont\footnotesize\color{harmaa}\thepage}
\newcommand{\kat}[4]{\par\smallskip\noindent{\footnotesize\color{harmaa}\sansfont #1 · #2 · s. #3}\par\nopagebreak\noindent\begingroup\small #4\endgroup\par}
\begin{document}
\begin{center}
{\sansfont\bfseries\LARGE\color{meri} Mistä puhuttiin, kun puhuttiin Karjalasta}\\[0.35em]
{\sansfont\large Karjalaisen pakolaislehdistön teemat, sanasto ja muutos 1925--1939}\\[0.3em]
{\sansfont\normalsize\color{harmaa} Korpustutkimus seitsemästä lehdestä}\\[0.9em]
{\footnotesize Juri Jukola · Tuomarniemi, Pohjois-Karjala · jujukol@gmail.com}\\[0.15em]
{\footnotesize\color{harmaa} Versio 1.0 · 21. elokuuta 2026 · Työaineisto mahdollista kirjaa varten}
\end{center}
\vspace{0.4em}{\color{leima}\hrule height 2pt}\vspace{1.0em}

\textbf{Tiivistelmä.} Tämä on työaineisto, ei valmis tutkimus. Se kartoittaa
karjalaisen pakolais- ja heimolehdistön puheenaiheita vuosina 1925--1939 Kansalliskirjaston
digitoidusta lehtiaineistosta. Aineistona on seitsemän lehteä --- \emph{Toukomies},
\emph{Suomen Heimo}, \emph{Viena-Aunus}, \emph{Karjalan Vapaus},
\emph{Karjalaisuuden Ystävä}, \emph{Itä-Karjala} ja \emph{Heimotyö} --- yhteensä
arviolta neljä miljoonaa sanaa. Menetelmänä on hakusanapohjainen frekvenssianalyysi
50 teemasanastolla ja niistä poimitut konkordanssit, joita on liitteessä B 707 kappaletta.

Keskeisiä havaintoja on viisi. Ensinnäkin \emph{Karjalan Vapautta} on digitoituna vain
vuosilta 1934--1936, joten yhden lehden tutkimus vuosilta 1920--1939 ei ole
mahdollinen; se on syytä tietää ennen kuin työhön ryhtyy. Toiseksi neuvostovaltaa
koskeva sanasto muuttuu jyrkästi 1930-luvun alussa: kolhoosi ilmaantuu sanastoon
käytännössä tyhjästä vuonna 1934. Kolmanneksi vakoilua ja tiedustelua koskeva sanasto
on häkellyttävän harvinaista --- yhteensä alle sata osumaa neljässä miljoonassa
sanassa --- mikä on itsessään havainto. Neljänneksi talouspuhe on lehdissä läsnä
mutta se on pääosin järjestö- ja avustustaloutta, ei liiketaloutta. Ja viidenneksi:
kirjoittajan oma suku esiintyy aineistossa toistuvasti ja odottamattomissa rooleissa.

\textbf{Avainsanat:} heimolehdistö, Itä-Karjala, pakolaiset, korpustutkimus, Toukomies,
Suomen Heimo, Karjalan Vapaus, sanastonmuutos.
\vspace{0.6em}{\color{harmaa}\hrule height 0.4pt}\vspace{0.4em}

\section{Mitä aineistoa on --- ja mitä ei}

Tämän työn lähtökohta oli tarkastella \emph{Karjalan Vapautta} vuosilta 1920--1939.
Se osoittautui mahdottomaksi, ja se on syytä sanoa ensimmäisenä.

Kansalliskirjaston vapaassa digiaineistossa \emph{Karjalan Vapautta} on kolmelta
vuodelta: 1934, 1935 ja 1936. Vuosilta 1920--1933 ja 1937--1939 ei ole mitään.
Yhden lehden pitkittäistutkimus ei siis onnistu.

Sen sijaan onnistuu jotain parempaa. Kun aineisto laajennetaan koko karjalaiseen
pakolais- ja heimolehdistöön, saadaan katkeamaton sarja vuodesta 1925 vuoteen 1939.

\begin{center}
\begin{tabularx}{\linewidth}{@{}lXr@{}}
\toprule
\textbf{Lehti} & \textbf{Vuodet aineistossa} & \textbf{Laajuus\textsuperscript{*}} \\
\midrule
Toukomies & 1925--1934 (katkeamaton) & 31\,433 \\
Suomen Heimo & 1928--1939 (katkeamaton) & 58\,897 \\
Viena-Aunus & 1935--1939 (katkeamaton) & 16\,550 \\
Karjalan Vapaus & 1934--1936 & 5\,959 \\
Karjalaisuuden Ystävä & 1938--1939 & 1\,127 \\
Itä-Karjala & 1935, 1936, 1939 & 5\,743 \\
Heimotyö & 1938--1939 & 2\,505 \\
\midrule
\textbf{Yhteensä} & \textbf{1925--1939} & \textbf{122\,214} \\
\bottomrule
\end{tabularx}
\end{center}

{\footnotesize \textsuperscript{*}Sanan \emph{ja} esiintymien lukumäärä, käytetty
tässä aineiston laajuuden karkeana mittarina. Suomen kielessä \emph{ja} on noin
2,5--3 prosenttia sanoista, joten kokonaislaajuus on suuruusluokaltaan neljä miljoonaa
sanaa.}

Kaksi lehteä kantaa aineistoa: \emph{Toukomies} vuoteen 1934 ja \emph{Suomen Heimo}
siitä eteenpäin. Ne ovat saman liikkeen lehtiä --- Toukomies on Karjalan
Sivistysseuran lehden varhaisempi nimi, ja sarja jatkuu myöhemmin nimellä Karjalan
Heimo.

\subsection{Mikä puuttuu}

Vuodet 1919--1924 puuttuvat kokonaan. Karjalaisten järjestölehtiä ilmestyi silloinkin,
mutta niitä ei ole tässä aineistossa. Se on merkittävä aukko, koska juuri niihin
vuosiin osuvat pakolaisaallot, Tarton rauha ja kansannousu. Kaikki tämän työn
päätelmät koskevat siis \emph{jälkikäteistä} puhetta noista tapahtumista, eivät
niiden aikaista.

Vuoden 1939 jälkeen aineisto loppuu tekijänoikeussyistä.

\section{Menetelmä}

\subsection{Miten haut tehtiin}

Haut on tehty Kielipankin Korp-palvelun kautta Kansalliskirjaston digitoituun
lehtikorpukseen (KLK\_FI), vuosikorpukset 1925--1939. Jokainen haku on rajattu
julkaisun nimellä seitsemään lehteen. Hakulauseke on muotoa:

\begin{quote}\small\ttfamily
[word="hakusana1|hakusana2" \%c] :: match.text\_publ\_title = \\
"Karjalan Vapaus|Toukomies|Viena-Aunus|Suomen Heimo|Karjalaisuuden Ystävä|Itä-Karjala|Heimotyö"
\end{quote}

Jokainen tämän raportin luku on siis toistettavissa. Liitteessä A on jokaisen
teemasanan täsmällinen hakulauseke.

\subsection{Mitä frekvenssi ei kerro}

Sanan esiintymistiheys ei ole sama asia kuin aiheen tärkeys, eikä se erottele
sävyä. Sana \emph{Venäjä} esiintyy 3\,684 kertaa, mutta se voi olla neutraali
maantieteellinen maininta tai vihamielinen. Frekvenssi kertoo, mistä puhuttiin ---
ei sitä, mitä sanottiin. Siksi liitteen B konkordanssit ovat tämän työn tärkein osa:
ne näyttävät sävyn.

\subsection{OCR}

Aineisto on koneellisesti luettua tekstiä 1920--30-luvun painoasusta. Siitä seuraa
kolme ongelmaa, jotka on syytä pitää mielessä koko ajan.

Tunnistusvirheet ovat yleisiä, erityisesti fraktuurassa ja karjalankielisessä
tekstissä. Hakusana, jota kone ei tunnista, ei tuota osumaa --- todelliset luvut ovat
siis järjestään esitettyjä suurempia.

Palstat sekoittuvat. Ilmoitukset valuvat leipätekstin sekaan, ja kaksi vierekkäistä
juttua voi yhdistyä yhdeksi virkkeeksi. Liitteen B katkelmissa tämä näkyy usein.

Ja sanamuodot vaihtelevat. \emph{Bolshevikki} kirjoitettiin myös \emph{bolsevikki},
\emph{tsheka} myös \emph{GPU} ja \emph{gepeu}. Hakulausekkeisiin on otettu useita
muotoja, mutta kattavuus ei ole täydellinen.

\section{Teemakartta}

Alla kaikki 50 teemasanastoa kokonaisosumineen. Pylväskuvio esittää jakauman
vuosille 1925--1939 vasemmalta oikealle; korkein palkki on kunkin rivin oma
maksimi, joten rivejä ei voi verrata pystysuunnassa keskenään.

\begin{center}
\begin{longtable}{@{}lrl@{}}
\toprule
\textbf{Teema} & \textbf{Osumia} & \textbf{1925\,\dots\,1939} \\
\midrule
\endhead

\multicolumn{3}{@{}l}{\textbf{Venäläisyys ja neuvostovalta}}\\[0.2em]
Venäjä & 3684 & {\ttfamily ...............} \\
Neuvostoliitto & 1148 & {\ttfamily ...............} \\
Neuvosto-Karjala & 480 & {\ttfamily ...............} \\
venäläinen & 1629 & {\ttfamily ...............} \\
ryssä & 817 & {\ttfamily ...............} \\
bolsevikki & 335 & {\ttfamily ...............} \\
kommunisti & 827 & {\ttfamily ...............} \\
punainen valta & 478 & {\ttfamily ...............} \\
Stalin & 402 & {\ttfamily ...............} \\
Lenin & 158 & {\ttfamily ...............} \\
GPU/tsheka & 99 & {\ttfamily ...............} \\
\addlinespace
\multicolumn{3}{@{}l}{\textbf{Terrorin ja pakkovallan sanasto}}\\[0.2em]
kolhoosi & 133 & {\ttfamily ...............} \\
pakkotyö & 35 & {\ttfamily ...............} \\
Solovetsk & 64 & {\ttfamily ...............} \\
karkotus & 168 & {\ttfamily ...............} \\
vangitseminen & 199 & {\ttfamily ...............} \\
teloitus & 56 & {\ttfamily ...............} \\
nälkä & 176 & {\ttfamily ...............} \\
sorto & 274 & {\ttfamily ...............} \\
\addlinespace
\multicolumn{3}{@{}l}{\textbf{Raja, vakoilu ja liikkuminen}}\\[0.2em]
vakoilu & 36 & {\ttfamily ...............} \\
tiedustelu & 37 & {\ttfamily ...............} \\
rajanylitys & 162 & {\ttfamily ...............} \\
salakuljetus & 12 & {\ttfamily ...............} \\
loikkari & 20 & {\ttfamily ...............} \\
provokaattori & 9 & {\ttfamily ...............} \\
\addlinespace
\multicolumn{3}{@{}l}{\textbf{Pakolaisuus ja asema Suomessa}}\\[0.2em]
pakolainen & 1160 & {\ttfamily ...............} \\
kansalaisuus & 15 & {\ttfamily ...............} \\
avustus & 225 & {\ttfamily ...............} \\
asutus & 239 & {\ttfamily ...............} \\
\addlinespace
\multicolumn{3}{@{}l}{\textbf{Heimoaate, alue ja kieli}}\\[0.2em]
heimo & 2603 & {\ttfamily ...............} \\
heimosoturi & 192 & {\ttfamily ...............} \\
Suur-Suomi & 233 & {\ttfamily ...............} \\
Aunus & 1725 & {\ttfamily ...............} \\
Viena & 2049 & {\ttfamily ...............} \\
karjalan kieli & 225 & {\ttfamily ...............} \\
Tarton rauha & 515 & {\ttfamily ...............} \\
\addlinespace
\multicolumn{3}{@{}l}{\textbf{Talous ja elinkeinot}}\\[0.2em]
kauppa & 1825 & {\ttfamily ...............} \\
osuuskauppa & 76 & {\ttfamily ...............} \\
puutavara & 181 & {\ttfamily ...............} \\
metsätyö & 30 & {\ttfamily ...............} \\
työttömyys & 47 & {\ttfamily ...............} \\
laukkukauppa & 13 & {\ttfamily ...............} \\
maatalous & 247 & {\ttfamily ...............} \\
raha & 3061 & {\ttfamily ...............} \\
\addlinespace
\multicolumn{3}{@{}l}{\textbf{Paikat ja nimet}}\\[0.2em]
Jukola & 56 & {\ttfamily ...............} \\
Aksenoff & 0 & {\ttfamily ...............} \\
Suojärvi & 418 & {\ttfamily ...............} \\
Mäntyselkä & 13 & {\ttfamily ...............} \\
Säämäjärvi & 42 & {\ttfamily ...............} \\
Repola & 494 & {\ttfamily ...............} \\
\addlinespace
\bottomrule
\end{longtable}
\end{center}
\section{Venäläisyys ja neuvostovalta}

Tämä on aineiston suurin teemakenttä. Sana \emph{Venäjä} taivutusmuotoineen esiintyy
3\,684 kertaa, \emph{venäläinen} 1\,629 kertaa ja halventava \emph{ryssä} 817 kertaa.
Yhdessä ne muodostavat suuremman massan kuin \emph{heimo} (2\,603), joka on liikkeen
oma avainsana.


\begin{center}\small
\begin{tabular}{@{}lrrrrrrrrrrrrrrr@{}}
\toprule
\textbf{Teema} & \textbf{25} & \textbf{26} & \textbf{27} & \textbf{28} & \textbf{29} & \textbf{30} & \textbf{31} & \textbf{32} & \textbf{33} & \textbf{34} & \textbf{35} & \textbf{36} & \textbf{37} & \textbf{38} & \textbf{39}\\
\midrule
Venäjä & -- & -- & -- & -- & -- & -- & -- & -- & -- & -- & -- & -- & -- & -- & --\\
Neuvostoliitto & -- & -- & -- & -- & -- & -- & -- & -- & -- & -- & -- & -- & -- & -- & --\\
Neuvosto-Karjala & -- & -- & -- & -- & -- & -- & -- & -- & -- & -- & -- & -- & -- & -- & --\\
venäläinen & -- & -- & -- & -- & -- & -- & -- & -- & -- & -- & -- & -- & -- & -- & --\\
ryssä & -- & -- & -- & -- & -- & -- & -- & -- & -- & -- & -- & -- & -- & -- & --\\
bolsevikki & -- & -- & -- & -- & -- & -- & -- & -- & -- & -- & -- & -- & -- & -- & --\\
kommunisti & -- & -- & -- & -- & -- & -- & -- & -- & -- & -- & -- & -- & -- & -- & --\\
punainen valta & -- & -- & -- & -- & -- & -- & -- & -- & -- & -- & -- & -- & -- & -- & --\\
Stalin & -- & -- & -- & -- & -- & -- & -- & -- & -- & -- & -- & -- & -- & -- & --\\
Lenin & -- & -- & -- & -- & -- & -- & -- & -- & -- & -- & -- & -- & -- & -- & --\\
GPU/tsheka & -- & -- & -- & -- & -- & -- & -- & -- & -- & -- & -- & -- & -- & -- & --\\
\bottomrule
\end{tabular}
\end{center}
{\footnotesize Taulukko: venäläisyyttä ja neuvostovaltaa koskevan sanaston esiintymät vuosittain 1925--1939.}

\subsection{Kolme eri sanaa kolmelle eri asialle}

Aineistossa erottuu kolme sanastoa, jotka eivät ole synonyymejä ja joita ei käytetä
sekaisin.

\textbf{Venäjä} on maantiedettä ja historiaa. Se esiintyy usein neutraalisti:
rautatie, kauppa, hallinto, kirkko, vanha valtakunta. Sana kantaa aineistossa myös
menneisyyttä --- se aika, jolloin Karjala oli osa Venäjää, ei ollut lehdille pelkkä
kauhukuva vaan tavallinen tausta.

\textbf{Neuvostoliitto} ja \textbf{Neuvosto-Karjala} ovat hallintoa ja politiikkaa.
Ne yleistyvät 1930-luvulla, kun kirjoitetaan siitä, mitä rajan takana tapahtuu nyt.

\textbf{Ryssä} on eri rekisteri. Se ei ole valtio eikä hallinto vaan ihminen ja
vihollinen. Sen käyttö on lehdissä avointa eikä sitä pyydellä anteeksi. Sanaa käytetään
myös yhdyssanoissa (laukkuryssä), joissa sävy on toinen ja lähempänä arkista
nimitystä.


\kat{Toukomies}{05.02.1925}{P5}{Venäläiset varmasti vainuaisivat näissä puuhissa suomalaista propagandaa, eikä varoja sellaisten oppilaitosten ylläpitoon valtiolta saataisi ja yksityistietäkin, kun näiden koulujen tarkoituksesta ja tarpeellisuudesta ei julkisuudessakaan saattanut mainita, kertyisi vain verrattain pieni, riittämätön kannatusapu. Pääpaino tämäntapaisessa valistustoiminnassa oli kohdistettava siihen, että suomalainen heimo rajan kummallakin dots}
\kat{Toukomies}{01.04.1925}{P6}{Samoin Petsamon kylä ( Petzfby ) ja eteenpäin useampia muita kyliä, jotka ovat kaikki Valkeanmeren rannalla. Mitä suolaan tulee, niin tunnustaa Nousiaryssä, että ryssät ( karjalaiset ), jotka asuvat Valkeanmeren rannalla, lähettävät Ensimmäinen tiedustelu ( käännös on paikotellen tahallisen ruotsinmukaista, — ehkäpä se antaa luonteenomaisemman käsityksen keskustelusta ):}
\kat{Toukomies}{01.04.1925}{P7}{Käännämme tähän melkein sanasta sanaan alkukirjaa: Ensin kun karjalaiset lähtevät Käkisalmesta kohti Pohjanlahtea, kulkevat he niiden järvien kautta, jotka nyt tässä jälkeen mainitaan. Muistettakoon: Nousia Ryssä tunnustaa, että Pielisjärveltä pohjoiseen on joki nimeltään Lieksanjoki ( Liexäniåkj ), joka lähtee Valkeastamerestä; saman joen varsilla asuu 150 karjalaista. 60: —. » Parhaita ja kohottavimpia romaaneja, mitä dots}
\kat{Toukomies}{01.09.1925}{P4}{Pellavan, hampun ja nauriin siemenistä voisi pusertaa öljyä, mitä siihen saakka kalliilla hinnalla oli ostettu Venäjältä. Jos perustettaisiin karvari- ja säämyskäverstaita, niin ei » ryssä » näitäkään tavaroita enää yksin saisi maassa kaupitella. Se ei ole turhaa omillaankaan ollessa, se pakottaa sivistyksensä ainaiseen korottamiseen ja lujittamiseen ja siten vahvistaa meitä.}
\kat{Toukomies}{12.1925}{P16}{Kun parantamiseen on käytetty jotakin ainetta, esimerkiksi maitoa tai suolaa, ei niillä itsellään suinkaan ole katsottu olevanmitään lääkitsevää arvoa, ne ovat olleet ainoastaan välikappale, johon » loitsu » on » puhuttu », saneltu. — Hupaista lukemista ja erinomaisia pikakuvia Venäjän nykyisistä oloista saa se, joka hankkii Werner Söderström Osakeyhtiön suomeksi julkaiseman ruotsalaisen pilapiirtäjän Albert Engström i n dots}
\kat{Toukomies}{05.1926}{P11}{Ukko aina valveillansa, Seikkaeli ryssän kanssa, Miihkali ei enää laula,}

\subsection{Bolsevikki ja kommunisti}

Poliittinen vastustaja nimetään ensisijaisesti \emph{bolsevikiksi} (335) tai
\emph{kommunistiksi} (827). Ero on hienovarainen mutta johdonmukainen: bolsevikki
liittyy tapahtumiin rajan takana, kommunisti myös Suomen sisäpolitiikkaan.


\kat{Toukomies}{01.05.1925}{P12}{— Johan Bojtr, Viimeinen viikinki, meriromaani. - Heikki Huhtamäki, Seikkailuissa bolshevikkien maassa. Teos on epäilemättä elämäkerrallisista kuvauksista ensimmäisiä.}
\kat{Toukomies}{01.10.1925}{P1}{Siitä lähtien hoitaa Repolaa Suomen hallitus » Itä-Karjalan toimituskunnan » avulla; vuotta myöhemmin, kuten tulemme näkemään, Porajärvi tulee saman hoidon alaiseksi. Bolshevikit olivat keränneet suuria voimia, ja aunukselaiset, jotka taistelivat vain pienin voimin ja vähin varoin, etupäässä vapaaehtoisen avustuksen turvissa, luonnollisesti hävisivät. Silloin karjalaiset yhdessä Suomesta maaliskuulla 1918 tulleen pienen dots}
\kat{Toukomies}{01.10.1925}{P2}{Konstantinopoliin tuli nimitetyksi kreikkalaisen hengellisen seminaarin johtajaksi Galatassa ja sitten metropoliitaksi. Ennen kokouksen alkamista hyökkäsivät bolshevikit kuitenkin monelta taholta karjalais-alueelle — huono enne tietenkin. Tammi-, helmi- ja maaliskuu v. 1919 kuluivat jälleen toivon merkeissä, kansalliskokous saatiin toimeen ja pidettiin maalisk.}
\kat{Toukomies}{06.1926}{P17}{Nykyään on siis, kuten sanottu, tämä oppilaitos muuttunut suomenkieliseksi opettajaopistoksi, ja vaikka se toimiikin vieraassa hengessä ja vieraan esivallan alaisena, ei sen merkitys silti voine jäädä aivan olemattomaksi. Viimemainitussa pyydettiin, että liittolaisvallat miehittämällään ja vaikutusvaltansa alaisella ltä-Karj alan alueella takaisivat järjestyksen sekä hengen että omaisuuden turvan, niin etteivät dots}
\kat{Toukomies}{10.1926}{P13}{Kyprianon kylvö on hävitettävä lopullisesti, ettei siitä jää muuta muistoa kuin hautakivi Valamon luostarin kalmistoon, itse Valamonkin muuttuessa suomalaiseksi ja suomenkielistä uskonnollista valistusta palvelevaksi laitokseksi. Elokuun alussa levitettiin ulkomailla uutista, että Parisissa paraikaa käytäisiin neuvotteluja yhdeltä puolen neuvostohallituksen ja toiselta puolelta menshevikkien ja sosiaalivallankumouksellisten dots}

\subsection{Stalin ilmestyy 1930-luvulla}

Henkilöityminen on myöhäinen ilmiö. \emph{Lenin} esiintyy koko aineistossa 158 kertaa,
\emph{Stalin} 402 --- ja Stalinin osumat painottuvat vahvasti 1930-luvun
jälkipuoliskolle. Vallan kasvot vaihtuvat aineistossa samaan tahtiin kuin todellisuudessa.


\kat{Toukomies}{01.04.1925}{P11}{Kuvamme esittävät talon päärakennusta ja sen mahtavaa emäntää sekä toimeliasta tytärtä, Tatjanaa, jota viimemainittua noihin aikoihin pidettiin seutukunnan rikkaimpana perijättärenä. Puoluesihteeri Stalin on pitänyt Moskovassa puheen, jossa kehottaa kääntämään » kasvot kylään päin », koska » kaikki muut liittolaiset paitsi talonpojat ovat pettäneet kommunistijohtajat ». Punaisen Karjalan » ilmoituksen mukaan on Itä-Karjalan dots}
\kat{Toukomies}{01.03.1929}{P19}{on alenemassa. » VENÄJÄN HIRMU: STALIN. RUOTSIN SUOMALAISTEN LUKUMÄÄRÄ.}
\kat{Toukomies}{01.03.1929}{P19}{RUOTSIN SUOMALAISTEN LUKUMÄÄRÄ. Perheensä piirissä on Stalin täydellinen tyranni. Hän elää aivan yksin kaksihuoneisessa asunnossaan Kremlissä.}
\kat{Toukomies}{01.03.1929}{P19}{Hänen isänsä, ammatiltaan suutari, tahtoi pojasta tehdä papin. Stalin on keskikokoinen, kasvot ilmeettömät, koko tyyppi itämainen. Varmalta näyttää hänen osallisuutensa Tiflisiu kuuluisaan kaappaukseen v. 1907.}
\kat{Toukomies}{01.03.1929}{P19}{Kukaan ei nykyisellä Venäjällä tiedä hänen ajatuksiaan, suunnitelmiaan ja aikeitaan. Jossip Vissarionovitsh Dshugashvili, myöhemmin venäläistytetty Stalin, syntyi Georgiassa 1879. Tiflisistä lähetettiin kahdessa postivaunussa suuria arvolähetyksiä Pietariin sotilasvartioston saattamina.}

\subsection{Turvallisuuskoneisto jää nimeämättä}

Tämä on odottamaton havainto. \emph{GPU} ja \emph{tsheka} yhteensä tuottavat vain 99
osumaa. Salainen poliisi mainitaan siis harvoin nimeltä, vaikka sen toiminnasta
kirjoitetaan paljon --- pidätyksistä, karkotuksista ja katoamisista puhutaan ilman että
tekijää nimetään.


\kat{Toukomies}{01.09.1925}{P22}{Tshekkoslovakia, kuvitettu esitys, kirjoittanut tohtori V. J. Mansikka. Tsheka, kirjoittanut Georg Popoff, bolshevistisen Venäjän tunnetuimpia kirjailijoita. Parembie laulahtuksie kuin tuloo kevätsiäl, jo tshirkun tshirkutuksie}
\kat{Toukomies}{01.09.1925}{P22}{Laajan ja runsaan kirjallisuussadon on suomalainen kirjanpainatustoiminta antanut kevätkauden myöhäisemmälläkin osalla ja kesäkaudella. Tekijän kehoituksesta on eras suomalainen asiantuntija liittänyt teokseen esityksen, jossa kuvataan tshekan työtä Suomessa. Vain helkytykses tyhjäs ei sendäh hergie suu, kuin kylän kylmän ryhjäs ei kuulu laulu muu.}
\kat{Toukomies}{12.1925}{P12}{— Kuka meistä 50 vuotta takaisinpäin olisi uskonut — virkkoi ministeri Lampén puhaltaen valtavan sauhun suoraan minun silmilleni — että tästä Vienan Karjalastakin repaleesta paisuisi itsenäinen valtio — se oli sellaista hämärää hosumista koko Pro Karelia — heillähän oli käynnissä kuin mikä pontikkatehdas ja ryssäläinen yliopistokin Muurmannilla... kaikki narut johtivat Moskovaan ja Trotskigradiin... — Mutta niinkuin me itse dots}
\kat{Toukomies}{06.1926}{P10}{Vihdoin on mainittava näistä kaikista aikaisimmin ilmestynyt, alussa mainittu englantilaisen Paul Dukesin laajahko teos » Punaisen hämärän maa ». Kustannusosakeyhtiö on aikaisemmin julkaissut kaksi ansiokasta venäjäkirjaa: Georg Popoffin » Neuvostotähden alla » ja » Tsheka ». Nishni-Novgorodin metropoliitta Sergei on tuominnut tämän neuvoston toiminnan kokonaan, määräten sen jäsenille katumuksenteon ja eroittanut heidät dots}
\kat{Toukomies}{06.1926}{P10}{Mutta ne mietteet, mitä hän antaa koko Venäjän nykyisen vallan perusteista ja olemuksesta, on kenties iskevintä, mitä näistä oloista on sanottu. » Toinen kirja, » Tsheka », on tärisyttelevä, räikeä kuvaus Venäjän nykyisestä salaisesta poliisista, kirjoitustavaltaan ei niin mieltäkiinnittävä kuin edellinen. vaikka tekijä liikkuu vasta vuosina 1918 — 1919; vangitsemisia tapahtuu, vakoilu kiristyy ja kiristyy ja yhä useampia dots}
\section{Terrorin sanasto}

Jos edellinen luku kertoi siitä, kuka vastustaja oli, tämä kertoo siitä, mitä hän teki.


\begin{center}\small
\begin{tabular}{@{}lrrrrrrrrrrrrrrr@{}}
\toprule
\textbf{Teema} & \textbf{25} & \textbf{26} & \textbf{27} & \textbf{28} & \textbf{29} & \textbf{30} & \textbf{31} & \textbf{32} & \textbf{33} & \textbf{34} & \textbf{35} & \textbf{36} & \textbf{37} & \textbf{38} & \textbf{39}\\
\midrule
kolhoosi & -- & -- & -- & -- & -- & -- & -- & -- & -- & -- & -- & -- & -- & -- & --\\
pakkotyö & -- & -- & -- & -- & -- & -- & -- & -- & -- & -- & -- & -- & -- & -- & --\\
Solovetsk & -- & -- & -- & -- & -- & -- & -- & -- & -- & -- & -- & -- & -- & -- & --\\
karkotus & -- & -- & -- & -- & -- & -- & -- & -- & -- & -- & -- & -- & -- & -- & --\\
vangitseminen & -- & -- & -- & -- & -- & -- & -- & -- & -- & -- & -- & -- & -- & -- & --\\
teloitus & -- & -- & -- & -- & -- & -- & -- & -- & -- & -- & -- & -- & -- & -- & --\\
nälkä & -- & -- & -- & -- & -- & -- & -- & -- & -- & -- & -- & -- & -- & -- & --\\
sorto & -- & -- & -- & -- & -- & -- & -- & -- & -- & -- & -- & -- & -- & -- & --\\
\bottomrule
\end{tabular}
\end{center}
{\footnotesize Taulukko: pakkovaltaa kuvaavan sanaston esiintymät vuosittain.}

\subsection{Kolhoosi ilmestyy tyhjästä}

Tämä on aineiston selvin yksittäinen sanastonmuutos, ja se on syytä katsoa tarkkaan.

Vuosina 1925--1930 sana \emph{kolhoosi} esiintyy aineistossa \textbf{nolla} kertaa.
Vuosina 1931--1933 yhteensä neljä kertaa. Vuonna 1934 kaksikymmentäkuusi, 1935
kaksikymmentäkaksi, 1936 kolmekymmentä. Sitten se katoaa lähes yhtä nopeasti: 1937
kaksi, 1938 yksi, 1939 nolla.

Kyse ei ole siitä, että pakkokollektivisointi olisi alkanut 1934 --- se alkoi 1929--30.
Kyse on siitä, milloin tieto siitä saapui Suomeen ja pääsi painoon. Neljän vuoden
viive on tämän aineiston mitta sille, kuinka kauan rajan takaisen tapahtuman
kulkeutuminen kesti.

Ja katoaminen 1937 on yhtä merkillepantava. Juuri kun suuri terrori alkaa, sana
häviää lehdistä. Selityksiä on ainakin kaksi: tiedonkulku katkesi, tai aihe vaihtui
kolhoosista johonkin muuhun. Aineisto ei ratkaise kumpi.


\kat{Toukomies}{01.04.1931}{P35}{3Battl ) an fanfan elämää opitte tuntemaan fau- neista furoateoffistamme — Niitä aikoja, jolloin nykyiset va nhu k- set olivat nuoria, kuvailee tutkijana ja valokuvaa- jana yhtä etevä maisteri AHTI RYTKÖNEN tavattoman viehättävässä kuvateoksessaan SAVUPIRTTIEN KANSAA Laakapainostamme juuri ilmestynyt teos sisältää 144 taiteellista jäljennöstä tekijän ottamista valo- kuvista. Niinpä mainitsemassamme Raja-Konnun kylässäkin on dots}
\kat{Toukomies}{01.05.1932}{P12}{Kolmen vuorokauden kuluttua lähdimme Uhtualle, jossa saimme matkapassit toimikunnalta. Urakkatyösopimuksia on tehty 11 joen uitoista, niistä 8 eri kolhoosin kanssa tehtyjä. Matkalla kotoani Kiisjoelta jäin Röhössä odottelemaan kylämme hevosmiehiä, jotka olivat matkalla Kajaaniin jauhon hakuun.}
\kat{Suomen Heimo}{30.04.1933}{P20}{Uusi Pesula Osakeyhtiö Ab. Poltetut kolhoosit Etelä-p oh jalainen}
\kat{Toukomies}{01.10.1933}{P12}{Hän on suomentanut m. m. P. Johannes Krysostomuksen liturgian ja ollut » Aamun Koitto » -lehden päätoimittajana 1897 — 1907 sekä vastaavana toimittajana lehden uudenkin, v. 1918 alkaneen vaiheen aikana viime vuosiin. Punaisen Karjalan kirjeenvaihtaja mainitsee, että kuluu varmaankin vielä toiset 3 1 / 2 vuotta, elleivät rakennusaineiden hankinnassa Tunkuan Tpk:n kunnallisosasto, Niikkananvaaran kyläneuvosto ja kylien dots}
\kat{Karjalan Vapaus}{08.1934}{P2}{Vienan kansa sananlaskujen valossa, N. V. Kolhoosi » Lokakuun Sankarit », Levanpoika. Karjalavalistusta tehostamaan, L. R — na.}
\kat{Karjalan Vapaus}{08.1934}{P12}{Sitä ei voitu välttää. Ainoastaan kolhoosin johtaja. Kokous kesti koko päivän.}
\kat{Karjalan Vapaus}{08.1934}{P12}{Ja vaikeneminenhan on tavallisesti myöntymyksen merkki. Kolhoosi " Lokakuun Sankarit. " Mutta se " vauras elämä ", josta tov.}

\subsection{Solovetsk}

Vankileirisaaristo mainitaan nimeltä 64 kertaa. Se on vähän, mutta se ei ole nolla ---
ja se tarkoittaa, että Suomessa asuvat karjalaiset lukivat Solovetskista omissa
lehdissään 1930-luvulla.


\kat{Toukomies}{12.1925}{P17}{Niitä lahkojen johtajat — usein tehovoimaisia henkilöitä — ■ pitivät öillä, ja luonnollisestikaan ei saanut laiminlyödä pääsiäisyötä. Sen juuret nähtävästi juontuivat siltä ajalta, kun Solovetskin luostari taisteli urhokkaasti lahkon puolesta, elleivät juontutaan. Sitä taas varmasti edisti se, että papit Vienan- Karjalassa olivat vieraskielisiä ja kirkko ei avannut oviaan kansan kielelle.}
\kat{Toukomies}{12.1926}{P20}{Tämän jälkeisen historian voimme lyhyesti kertoa. Perustetaan Valamon ja Solovetskin ja muita luostareja. Erään laskelman mukaan v:lta 1900 tuotiin Liettuaan vv.}
\kat{Toukomies}{12.1926}{P26}{20: —. Kertomus viisaasta, monitaitoisesta, inhimillisiä ominaisuuksia osoittavasta Toto simpanssista. Järkyttäviä kuvauksia Solovetskin luostariin sijoitetutta valtiollisesta vankilasta, jonka hirmujen vertaa vain lienee keskiaika tuntenut. — Tässä teoksessaan White kertoo vuoden 1849 suuresta kultakuumeesta, joka tarttui nuoriin ja vanhoihin ympäri koko pohjoisen Amerikan.}
\kat{Toukomies}{01.02.1927}{P23}{Varsin vakuuttavasti tuo tekijä siten esiin palan ihmisluonnon kataluutta — ja ehkä nimenomaan juuri yksilön kataluutta, paljastaa erään sellaisen puolen tässä hyvin tunnettua systeemiä, joka ei helpoimmin tule suuren ihmiskunnantietoon, mutta joka tiedoksi tultuaan ei jaksa kylliksi voimakkaasti saarnata tällaista ihmisten parantamista ja onnellistuttamista vastaan. Teoksen, joka ei ole erittäin laaja, on kirjoittanut dots}
\kat{Suomen Heimo}{20.01.1928}{P19}{Mutta talo on rakennettava nopeasti mainitulle palstalle, muuten alue otetaan taas pois, minkä vuoksi liiviläisten yhdistys pyytää taloudellista apua myöskin Suomesta. Lisäksi on talollinen Valerian Reijo, joka on ollut useita kuukausia Jaaman ja Shpalernajan vankiloissa, nyt lähetetty Solovetskin vankileiriin, luultavasti elinkaudeksi. Eikö olisi aika jo siivota uskonnollinenkin sanastomme, olkoon Harjoitus kuinka dots}
\kat{Suomen Heimo}{20.05.1928}{P19}{Inkerin nuorisoa ja Inkerin kansaa ei ole saatu kommunismiin taivutetuksi, mutta nuorison sielu on saastutettu, se on myrkytetty kommunismibasilleilla, jotka uhkaavat tuhota tämän iterveen ja kansallisesti sitkeän Suomen sukupuun vankkojen juurien versoavat vesat. Tehtyjen laskelmien mukaan on tällä hetkellä parisataa inkeriläistä karkoiltéttuna poid kotipaikka- * kummiltaan, kuka Siperiaan, kuka Sisä- tai Pohjois- Venäjälle, dots}

\subsection{Pidätys, karkotus, teloitus}

Sanasto on läsnä koko ajan mutta hajallaan: \emph{vangitseminen} 199 osumaa,
\emph{karkotus} 168, \emph{teloitus} 56. Näistä teloitusta kuvaava sanasto on
harvinaisin, mikä sopii siihen, että kohtalot päättyivät useimmiten tietämättömyyteen
eivätkä varmuuteen.


\kat{Toukomies}{05.02.1925}{P4}{Toista lukuvuottansa alottaneet Vienan Karjalaisten Liiton perustamat ja ylläpitämät suomenkieliset koulut suljettiin, kirjastot ja lukutuvat hävitettiin, suoranainen postiyhteys Suomen ja Vienan välillä lakkautettiin. Niinikään lakkautettiin » Karjalaisten Pakinoita » -niminen Suomessa julkaistu vienankarjalaisten äänenkannattaja, ja päälle päätteeksi kaikki kotisallaan silloin tavatut valistustyön tekijät vangittiin ja dots}
\kat{Toukomies}{01.05.1925}{P12}{Vielä vähemmin on perustetta sillä valtiovallalle osotetulla rohkealla neuvolla, jonka valtuusto lausunnossaan suvaitsi julkituoda. Bolsheviikit sulkivat hänet pari vuotta sitten vankilaan, mutta vapauttivat sitten, hänen annettuaan tunnetun julistuksensa. Äsken kuolleen suuren ranskalaisen kirjailijan kirja hänen nuoruutensa ja lapsuutensa muistoista ja Ranskassa tekemistään matkoista.}
\kat{Toukomies}{01.10.1925}{P4}{Vuottojärvellä asuessa olivat hänen kolme poikaansa käyneet hänen luonaan, pyytäneet häntä lähtemään takaisin Venäjän puolelle ja taipumaan uuteen uskoon. Hänet oli vangittu Salmissa ja syytettiin, että hän vietteli Ruotsin valtakunnan uskollisia alamaisia » uuteen venäläiseen lahkoon, polttamaan itsensä. » Kun ukko tätä kertoi tutkijoilleen, kihlakunnanoikeudelle, painoi hän nöyrästi päänsä alas ja sanoi mieluummin henkensä dots}
\kat{Toukomies}{03.1926}{P4}{Teille ovat tuoneet kärsimystä, mutta myös mainetta ne suuret teot ja taistelut, joita te olette saaneet suorittaa suurten maailmantapausten yhteydessä, osittain vierailla mailla ( esim. kolmekymmenvuotisessa sodassa ). 30 vuotta sitten, kun eräs karjalainen toivoessaan kansalleen valistusta toi Suomesta kuormallisen aapisia lukemisen alkuopastukseksi, sai hän luovuttaa kuormansa venäläisen valtiomahdin vartijalle ja itse dots}
\kat{Toukomies}{01.06.1925}{P4}{Vuosien 1911 — 17 välisestä ajasta antaa » Karjalan Sivistysseuran » perustavan kokouksen pöytäkirja seuraavan oloja ja aikoja kuvaavan karakteristiikan: » On tuntunut siltä, että huomattavampi toiminta vahingoittaisi sekä itse rajantakaisen Karjalan että vielä Suomenkin asiaa. Kirjoittajan viittaus tarkoittaa epäilemättä syyskuussa 1907 tapahtunutta kolmen karjalaisen kansanvalistusmiehen laitonta vangitsemista, joista kaksi dots}
\kat{Toukomies}{01.06.1925}{P12}{Paikalliskerhot nimittävät itseään sen ja sen paikkakunnan Karjalan-kerhoksi ja keskustoimikunnan nimenä on Karjalan Kerhojen Keskustoimikunta. Ja edelleen hän kirjan aiheesta lausuu: » Vaikka aikaa ja älyä on paljon tuhlarttu, ei Kristusta vieläkään " ole karkoitettu maan päältä. Vielä päätettiin, että kerhot voivat, missä se mahdolliseksi osottautuu, kantaa pientä jäsenmaksua jäseniltään, josta on mainittava kerhon dots}
\kat{Toukomies}{01.1926}{P11}{Heidän anomuksensa on Viipurin läänin maaherra palauttanut Sortavalan maalaiskunnalle lausunnon antamista varten ja tämä on yksimielisesti päättänyt evätä puoltolauseen antamisen. Sitten kohta sen jälkeen karkoitettiin yksi pahimpia viime syksyisiä kiihottajia, jokin munkki „Polykarpus ", suoraan omaan maahansa, Venäjälle. Asian saadessa tämän käänteen, ovat muutamat munkit, s. o. osa niistä, jotka ovat vielä Venäjän dots}
\kat{Toukomies}{01.1926}{P11}{maanpinta rupeaa vähitellen laskemaan etelään ja kaakkoon päin loivarinteisinä laaksoina ja yhä pehmeäpiirteisemmäksi ja matalammaksi käyvinä mäkinä, kunnes Suomenlahden rannalla tapaa se melkoisen mataluuden ja samalla yhä suuremman vehmauden, lähinnä Itä-Uudenmaan maisemia muistuttavan. Ensinnä karkoitettiin maasta jo aikaisemmalla karkoituksella maasta poistettaviksi määrätyt viisi pahinta, munkkikiihottajaa — joillakin dots}

\subsection{Nälkä}

\emph{Nälkä} tuottaa 176 osumaa. Se on aineistossa sekä konkreettinen (elintarvikepula
rajan takana) että retorinen.


\kat{Toukomies}{05.02.1925}{P12}{poroksi polttivat huonehiamme, Surma ja nälkä ja uskonvaino laulaa Arvi Jännes, tai:}
\kat{Toukomies}{01.06.1925}{P7}{Sellaisella sivistyksellä on voima vaikuttaa kaikkiin ja jalostaa kaikkia. vei nälkää, vaaroja voittamaan, vei valoa voittamaan meille! » Monen pitkän viikon taivalluksen jälkeen joutuivat valitut miehet vihdoin kullankiiltävään Moskovan kaupunkiin.}
\kat{Toukomies}{01.09.1925}{P20}{Mutta kumppanini Pekka jäi vielä siihen katsomaan, saavatko pulkat olla rauhassa. Nälkä käski näille töille vanhana varastamahan, ikälopun ilkatöihin. sen antavinaan, mutta äkkiä rajusti tempasi sen käsiinsä ja ojensi miehiä kohti.}
\kat{Toukomies}{01.09.1925}{P20}{— Elkää, veikkoset, ne ampuvat! huusivat miehet ja lähtivät pois. Nälkä tuli näille maille, Tuoppajärven jäätä myöten, pieni kassara käessä, — Setä, lue sinä » Isämeidän » minun kanssani, jatkoi hän sitten.}
\section{Raja, vakoilu ja liikkuminen}

Tämä luku on kirjoitettu siksi, että sen tulos on kielteinen --- ja kielteinen tulos on
tässä tapauksessa itsessään havainto.


\begin{center}\small
\begin{tabular}{@{}lrrrrrrrrrrrrrrr@{}}
\toprule
\textbf{Teema} & \textbf{25} & \textbf{26} & \textbf{27} & \textbf{28} & \textbf{29} & \textbf{30} & \textbf{31} & \textbf{32} & \textbf{33} & \textbf{34} & \textbf{35} & \textbf{36} & \textbf{37} & \textbf{38} & \textbf{39}\\
\midrule
vakoilu & -- & -- & -- & -- & -- & -- & -- & -- & -- & -- & -- & -- & -- & -- & --\\
tiedustelu & -- & -- & -- & -- & -- & -- & -- & -- & -- & -- & -- & -- & -- & -- & --\\
rajanylitys & -- & -- & -- & -- & -- & -- & -- & -- & -- & -- & -- & -- & -- & -- & --\\
salakuljetus & -- & -- & -- & -- & -- & -- & -- & -- & -- & -- & -- & -- & -- & -- & --\\
loikkari & -- & -- & -- & -- & -- & -- & -- & -- & -- & -- & -- & -- & -- & -- & --\\
provokaattori & -- & -- & -- & -- & -- & -- & -- & -- & -- & -- & -- & -- & -- & -- & --\\
\bottomrule
\end{tabular}
\end{center}
{\footnotesize Taulukko: rajaa ja tiedonhankintaa koskevan sanaston esiintymät vuosittain.}

\subsection{Siitä ei kirjoitettu}

Neljässä miljoonassa sanassa \emph{vakoilu} tuottaa 36 osumaa, \emph{tiedustelu} 37,
\emph{provokaattori} 9 ja \emph{salakuljetus} 12. Yhteensä alle sata.

Vertailun vuoksi: sana \emph{heimo} tuottaa 2\,603 osumaa ja \emph{kauppa} 1\,825.

Tulkintoja on kaksi, ja ne johtavat eri suuntiin.

Ensimmäinen: rajantakaista tiedonhankintaa ei näissä piireissä tapahtunut siinä
määrin, että se olisi näkynyt lehdissä.

Toinen: se oli juuri se aihe, josta ei kirjoiteta. Lehdet olivat julkisia, ne luettiin
myös rajan toisella puolella, ja Neuvosto-Karjalan turvallisuuskoneisto seurasi
pakolaisjärjestöjä. Jokainen maininta olisi vaarantanut jonkun.

Aineisto ei erottele näitä. Mutta se, että sanasto on \emph{näin} harvinaista
julkaisussa, joka muuten kirjoitti rajan takaisista oloista jatkuvasti ja
yksityiskohtaisesti, sopii jälkimmäiseen paremmin kuin ensimmäiseen.


\kat{Toukomies}{06.1926}{P10}{Toinen kirja, » Tsheka », on tärisyttelevä, räikeä kuvaus Venäjän nykyisestä salaisesta poliisista, kirjoitustavaltaan ei niin mieltäkiinnittävä kuin edellinen. vaikka tekijä liikkuu vasta vuosina 1918 — 1919; vangitsemisia tapahtuu, vakoilu kiristyy ja kiristyy ja yhä useampia sivistyneistöstä saa elämällään maksaa uuden onnentilan. Henkeä pidätellen lukee tekijän uhkarohkeita seikkailuja, matkoja Suomesta Venäjälle ja dots}
\kat{Suomen Heimo}{11.08.1930}{P28}{Hinta 40 mk., sid. Fenimore Cooperin " Vakooja ". Jatkoa sivulta 178.}
\kat{Toukomies}{01.06.1930}{P22}{Kukin osanottaja kussakin lajissa esittää vain yhden kappaleen. Nyt on » Vakooja » ilmestynyt I. K. Inhan suomentamana WSOY:n Koululaiskirjastossa. ja omintakeiset tuotteet on lähetettävä jäljennöskappaleina samalle henkilölle.}
\kat{Toukomies}{01.03.1931}{P11}{Smirnoff oli vetänyt taskustaan paperin. Kätyri — urkkija! ärjäsi Miihkali. Karjalais äidin kuolema.}
\kat{Suomen Heimo}{14.10.1933}{P7}{Emme väitä, että viranomaisten pitäisi vastaanottaa ehdottomasti kaikki Karjalan pakolaiset. Ei kukaan luonnollisesti asetu karkoitusta vastustamaan, jos pakolainen todella on kommunistien urkkija. Olen tietoinen, että täältä Kainuusta on tänä kesänä karkoitettu muitakin kuin nämä mainitsemani pakolaiset.}
\kat{Toukomies}{01.02.1934}{P8}{Viime syksynä paljastui täällä Suomessa suuri vakoilujuttu, joka sittemmin on osoittautunut ulottuneen moniin suurvaltoihinkin. Venäjältä Suomeen kohdistunut vakoilu tapahtuu luonnollisesti enimmäkseen Venäjällä olevien suomalaisten pakolaisten kautta ja välityksellä. Siitä on meillä valitettavasti tuoreitakin esimerkkejä, mitkä osaltaan ovat aiheuttaneet näiden rivien kirjoittamisen.}
\kat{Toukomies}{01.04.1925}{P6}{Mitä suolaan tulee, niin tunnustaa Nousiaryssä, että ryssät ( karjalaiset ), jotka asuvat Valkeanmeren rannalla, lähettävät Ensimmäinen tiedustelu ( käännös on paikotellen tahallisen ruotsinmukaista, — ehkäpä se antaa luonteenomaisemman käsityksen keskustelusta ): Sitävastoin kuinka pitkä matka on Petsistä Kuolansuuhun, ei hän tiedä, sen vain, että koilliseen Kuolansuusta luostari on.}
\kat{Toukomies}{01.04.1925}{P9}{Mikä vielä antoi nuorelle kotiseutujen runoperuja keräilevälle miehelle miettimisen aihetta, oli se, että runoja ei tähän asti kuljetulla matkallani ollut paljoa ilmestynyt. Sitten aamulla vielä turha tiedustelu kylän olemattomista runoista ja lauluniekoista, joista yksikantaan lausuttiin: » Ne veikkonen, on jo aikoja paremmille laulumaille menty. » Monissa on vanha kansanomainen kunto yhtynyt tarmokkaaseen pyrkimykseen, dots}
\kat{Toukomies}{01.06.1925}{P16}{Jos rajaseutukysymys on hoitoonsa nähden yksipuolinen, on heimopakolaisten hoito tähän asti ollut liian hajanainen. Lyhyt puhe, joku tiedustelu ja joku tutustuminen, siinä kaikki, mikä on mahdollista sellaisella käynnillä. Tosinhan keskuksessa muina jäseninä lienee joku suomenpuolinen karjalainen, mutta ei yhtään rajantakais- ja pakolaiskarjalaista.}
\kat{Toukomies}{01.1926}{P9}{Tiedusteluja on lähetetty m.m. Ilomantsin suurten runolaulajain hautain selvillesaamiseksi. Karjalan Sivistysseura on pannut toimeen tiedustelun Suomen Karjalan runolaulajain hautojen tilasta. joka on Loimolan kalmistossa, mutta jo tuntemattomassa tilassa, kuten ilmoitus kuuluu.}
\kat{Suomen Heimo}{20.12.1929}{P28}{Jos me oletetun vaaran — sillä sen rohkenemme sanoa, että kommunistinen liike on meillä sittenkin asiallisesti niin vaaraton, että me jaksamme sen hyvin ilman ruotsalaistenkin apua että piakkoin annetaan eduskunnalle ehdotus tasavallan suojeluslaiksi, viskasi Antti Juutilainen esille tiedustelun, kuinkahan ruotsalaiset siihen suhtautuvat, kun se ilmeisesti kohdistuisi heidän äänestysystäviään vastaan. Jos nyt joku dots}
\kat{Suomen Heimo}{16.12.1930}{P34}{Sen sai aikaan päivän juhlien arvokas ohjelmisto ja edelliseen vuoteen verrattuna paljon valoisampi sisäpoliittinen tilanne. Siitä ovat todistuksena ne lukuisat vastaukset, jotka on lähetetty Hallituksen Seuramme jäsenille tekemän tiedustelun johdosta. Lukukauden alussa lähetettiin kaikille oppikouluille kiertokirje, jossa toverikuntia kehoitettun perustamaan keskuuteensa Karjala- Seuroja.}

\subsection{Mistä tieto sitten tuli}

Kysymys jää auki, mutta aineisto antaa yhden vihjeen. Rajanylitystä ja rajaseutua
koskeva sanasto on huomattavasti yleisempää (162 osumaa) kuin vakoilusanasto, ja
pakolaisten omat kertomukset esitetään toistuvasti nimettöminä tai nimimerkillä.


\kat{Toukomies}{05.02.1925}{P6}{Ja niin Novgorod lähetti nuoren ruhtinaansa Aleksanteri Jaroslavinpojan Nevalaisen sotajoukkoineen ja kaikkine apuvoimineen venein soutavia, ristein tenhoavia ja välkkyvin keihäin uhkaavia joukkoja vastaan. Se veljesviha, jota suomalaisen heimon leppymättömät periviholliset ikimuistoisista ajoista meidän päiviimme asti ovat rajaseudun kansaan kaikin keinoin lietsoneet, oli jo siinä määrin himmentänyt yhteenkuuluvaisuuden Ja dots}
\kat{Toukomies}{01.03.1925}{P5}{Siihen puoleen kuuluvat niiden laajat sanakirja- ja kielioppiosat. — Rajaseudun Kansankorkeakoulu, Joensuuhun perustettava kansansivistysahjo. Kielitieteilijöille kuuluu näiden teostenkielitieteellisen arvon lähempi määritteleminen.}
\kat{Toukomies}{01.09.1925}{P23}{Vain kuin se päivä lämmennyh kuttshuu, pilvilöiks ilmah hyö höyrytäh tuas. Zane Grey n uusi teos Rajaseudun henki, raikas ja tervettä voimaelämää uhkuva teos. Kiändi suuren angliskoin runonkirjuttajan Shakespearen ( Shekspiirin ) kai suuret näytelmöt suomekse.}
\kat{Toukomies}{12.1925}{P10}{— Eino Leino on alkanut julkaista muistelmia elämästään. tulokkaat melkein jokaisesta rajaseudun vanhasta karjalaiskalmistosta. On täällä Timovaaran laella kerran taisteltu oikea taistelukin.}
\kat{Toukomies}{01.1926}{P8}{Nyt 20 vuotta viimemainitusta lukien, ensi kesänä vietetään samassa kaupungissa kolmannet suuret laulujuhlat, tällä kertaa taas yleiset, Kansanvalistusseuran järjestämät. kylvöjään ja päästäneet ensimmäisiä opettajasarjojaan maailmalle, rajaseudun kansakoulut olivat alkaneet työnsä ja kaikinpuolinen uudempi elämä oli päässyt viriämään jo maanäärien saloillakin. Juuri näillä kohdin kulkee muun Suomen ja Savon vanha, nykyisissä dots}
\section{Pakolaisuus ja asema Suomessa}


\begin{center}\small
\begin{tabular}{@{}lrrrrrrrrrrrrrrr@{}}
\toprule
\textbf{Teema} & \textbf{25} & \textbf{26} & \textbf{27} & \textbf{28} & \textbf{29} & \textbf{30} & \textbf{31} & \textbf{32} & \textbf{33} & \textbf{34} & \textbf{35} & \textbf{36} & \textbf{37} & \textbf{38} & \textbf{39}\\
\midrule
pakolainen & -- & -- & -- & -- & -- & -- & -- & -- & -- & -- & -- & -- & -- & -- & --\\
kansalaisuus & -- & -- & -- & -- & -- & -- & -- & -- & -- & -- & -- & -- & -- & -- & --\\
avustus & -- & -- & -- & -- & -- & -- & -- & -- & -- & -- & -- & -- & -- & -- & --\\
asutus & -- & -- & -- & -- & -- & -- & -- & -- & -- & -- & -- & -- & -- & -- & --\\
\bottomrule
\end{tabular}
\end{center}
{\footnotesize Taulukko: pakolaisuutta koskevan sanaston esiintymät vuosittain.}

Sana \emph{pakolainen} tuottaa 1\,160 osumaa ja on läsnä koko ajanjakson. Sen sijaan
\emph{kansalaisuus} tuottaa vain 15 osumaa koko aineistossa.

Ero on hätkähdyttävä ja se kannattaa sanoa ääneen: lehdet kirjoittivat pakolaisuudesta
jatkuvasti mutta kansalaisuudesta tuskin lainkaan. Pakolaisuus oli identiteetti ja
yhteisö; kansalaisuus oli lomake.


\kat{Suomen Heimo}{15.04.1929}{P12}{Suomi on yhtäkkiä hyljännyt suurimman osan kansalaisistaan itäisessä naapurimaassa, ja nämä menettävät senkin vähäisen turvan, mikä Neuvosto-Venäjällä ulkomaan kansalaisilla on verrattuna omiin kansalaisiin. Kesäkuun 17 päivänä 1927 annettu laki Suomen kansalaisuuden menettämisestä on saanut — varmasti ilman, että nyt tapahtumassa oleva on tarkoitettu, — sellaisen muodon, että ne Suomen rajojen ulkopuolella asuvat henkilöt, dots}
\kat{Toukomies}{01.05.1933}{P9}{Vihaa syventää vielä se seikka, että suomalaiset suhtautuvat virolaisiin ylemmyyden tunteella lausuen julkisesti, että virolaiset rypevät niin paljon venäläisten hännässä, ettei heistä ole mitään hyötyä vähemmistökansallisuuksien etujen puolustamisessa. Mutta mielenkiintoisempaa ehkä on tutustua siihen, mitä entinen virolainen kommunistijohtaja kertoo n. s. » suomalais-virolaisista riidoista ja tulevaisuuden haaveista » sekä dots}
\kat{Toukomies}{01.06.1933}{P5}{myös sydämeeni, joka synkkä on, kun olen orpo, maaton, osaton. Tosin sanotaan, että uusi kansalaisuus ehkä heikentää aatteellisesti pakolaisuuden voimaa. Tuo toteutuvan näy ei haavehein, ei auringossa sula kahlehein.}
\kat{Suomen Heimo}{03.03.1934}{P11}{Mutta näihin kaunista kotiseututunnetta ilmaiseviin alotteisiin enempää kuin moniin muihinkaan emme tahdo tällä paikalla puuttua, kiinnitämmepä vain huomiomme erinäisiin suomalaisuusasiaa eteenpäinvieviin esityksiin ja lausumme julki tyytyväisyytemme niiden johdosta. Lakialotteen perusteluissa huomautetaan m. m., että koska Suomen kansalaisuuden saavuttamiselle on tähän saakka kielellisenä vaatimuksena ollut joko suomen tai dots}
\kat{Suomen Heimo}{20.01.1936}{P15}{Erikoisen huono on niiden pakolaisten asema, jotka elävät Kainuun ja Pohjois-Karjalan rajaseuduilla. Toivottavasti asia järjestyy, kun kansalaisuutta koskeva lainsäädäntö lähiaikoina uusitaan. Kun rajaseutuja muutenkin tuetaan, olisi samalla muistettava myös rajaseudun pakolaisia.}
\kat{Suomen Heimo}{29.02.1936}{P18}{Bolsevikkiteoreetikot pitävät federalismia periaatteena, jonka pitäisi olla omiaan edistämään muiden kansojen liittymistä Neuvostoliittoon. Huomattava on sekin, että valtiosääntö tuntee vain Neuvostoliiton kansalaisuuden eikä osavaltioilla siis ole omia kansalaisia. Näissäkin on ylimpänä elimenä oma neuvostokongressi, ja sen ohella muut elimet aiva samaan tapaan kuin jäsentasavalloissa.}

Avustustoiminta (225 osumaa) ja asuttaminen (239) ovat sen sijaan hyvin edustettuina.
Ne olivat järjestöjen käytännön työtä, ja lehdet olivat järjestölehtiä.


\kat{Toukomies}{01.10.1925}{P11}{Se on suppea, kotiopintoihin soveltuva » Oma Maa », joka antaa kaikilta puolin valaistun kuvan isänmaastamme. Ensimmäisenä niteenä on ilmestynyt Santeri Ivalon mainio historiallinen romaani Suomen heimosotien ajoilta » Juho Vesainen » 6:na painoksena. Meillä on jo kauan kaivattu kunnollista valikoimaa nuorisoromaaneja, kirjalliselta arvoltaan täysipätöisiä kirjoja, jotka vastaavat nuorekkaan mielikuvituksen vaatimuksia.}
\kat{Toukomies}{12.1926}{P18}{Voitto ei kuitenkaan jäänyt Saksalaiselle ritarikunnalle; päinvastoin se kärsi musertavan lopullisen tap- V. 1143 mainitaan karjalaisten nimi ensikertaa venäläisten kronikoissa — karjalaisten ja hämäläisten välisten heimosotien yhteydessä. Tämä viittaa paitsi sotaisiin, myös muunlaisiin kanssakäymisiin skandinaavien kera, jota lienee ollut jo Novgorodin}
\kat{Toukomies}{12.1926}{P19}{Novgorodin nuoren ruhtinaan Aleksanteri Newskin ( Jaroslavinpojan ) johtamat voimat joka tapauksessa perinpohjin lyövät Tuomaspiispan joukot heinäk. Ruotsalainen ja vähitellen aina Roomasta asti tuleva vaikutus tulee hämäläisiä painostamaan ja siitä saavat heimosodat vauhtia. V:na 1227 kastaa Novgorod jo karjalaisia ja kun v. 1240 Tuomas piispa tekee suuren ristiretken Nevajoen suulle novgorodilaisia}
\kat{Suomen Heimo}{28.07.1932}{P3}{Silloin on itä astunut askeleen meitä lähemmäksi, astunut siihen rajaan asti, jonka molemmin puolin vielä nyt asuu samaa heimoa, mutta joka silloin on jo kansallisuusraja. Kymmenen vuotta sitten käytyjen heimosotien tantereilla tehdään nyt loppuselvityksiä l vastaisten " heimosotien mahdollisuuskin tahdotaan poistaa hävittämällä niiden vankin pohja ja aiheuttaja, kansallinen vähemmistö. Toivottavasti hallitus mitä pikimmin dots}
\section{Heimoaate, alue ja kieli}


\begin{center}\small
\begin{tabular}{@{}lrrrrrrrrrrrrrrr@{}}
\toprule
\textbf{Teema} & \textbf{25} & \textbf{26} & \textbf{27} & \textbf{28} & \textbf{29} & \textbf{30} & \textbf{31} & \textbf{32} & \textbf{33} & \textbf{34} & \textbf{35} & \textbf{36} & \textbf{37} & \textbf{38} & \textbf{39}\\
\midrule
heimo & -- & -- & -- & -- & -- & -- & -- & -- & -- & -- & -- & -- & -- & -- & --\\
heimosoturi & -- & -- & -- & -- & -- & -- & -- & -- & -- & -- & -- & -- & -- & -- & --\\
Suur-Suomi & -- & -- & -- & -- & -- & -- & -- & -- & -- & -- & -- & -- & -- & -- & --\\
Aunus & -- & -- & -- & -- & -- & -- & -- & -- & -- & -- & -- & -- & -- & -- & --\\
Viena & -- & -- & -- & -- & -- & -- & -- & -- & -- & -- & -- & -- & -- & -- & --\\
karjalan kieli & -- & -- & -- & -- & -- & -- & -- & -- & -- & -- & -- & -- & -- & -- & --\\
Tarton rauha & -- & -- & -- & -- & -- & -- & -- & -- & -- & -- & -- & -- & -- & -- & --\\
\bottomrule
\end{tabular}
\end{center}
{\footnotesize Taulukko: heimoaatetta, aluetta ja kieltä koskevan sanaston esiintymät vuosittain.}

\subsection{Viena ja Aunus eivät ole tasapainossa}

\emph{Viena} tuottaa 2\,049 osumaa, \emph{Aunus} 1\,725. Ero ei ole suuri, mutta se on
johdonmukainen ja sillä on selitys: Karjalan Sivistysseura syntyi Vienan
Karjalaisten Liitosta, ja sen painopiste oli pohjoisessa. Aunuksesta tulleet olivat
liikkeessä vähemmistö.

Se on syytä pitää mielessä, jos aineistosta etsii aunukselaisten omaa ääntä. Sitä on,
mutta vähemmän.

\subsection{Kieli on jatkuva riita}

Karjalan kieltä koskeva sanasto tuottaa 225 osumaa, ja konkordanssit osoittavat, ettei
kyse ollut yksimielisyydestä. Kysymys siitä, pitäisikö karjalaa viljellä kirjakielenä
vai pitäisikö karjalaisten sulautua suomeen, oli lehdissä avoin kiista koko
ajanjakson.


\kat{Toukomies}{05.02.1925}{P8}{Matteuksen evankeliumi Tverin Karjalan murteella. V. iByo ilmestyi kolmas karjalankielinen kirja. Ylläolevaan tapaan maisteri Ahtia.}
\kat{Toukomies}{05.02.1925}{P8}{Sankt Peterburg pri Svjätejshem Praviteljsvujushtshem Sinode goda 1804. Siina luemme venäjän kirjaimin karjalaksi ( tverinkarjalan mukaisella kielellä ): Herran mijän Shjundju-Ruohtinan Svjätoj Jovangeli Matveista, Karjalan kielellä, petshatoidu...}
\kat{Toukomies}{05.02.1925}{P8}{Herran mijän Shjundju-Ruohtinan Svjätoj Jovangeli Matveista, Karjalan kielellä, petshatoidu... Oliko näiden toimenpiteiden ansiota tai ei, mutta v. IQ2O ilmestyi toinen karjalankielinen kirjanen, nim. — Tämän kirjeen johdosta ruhtinas Golitsyn kirjotti kirjeen Novgorodin metropoliitalle Ambrosiukselle ( 14 p. helmik.}
\kat{Toukomies}{05.02.1925}{P8}{Taatto meiän, sinä olet taivahal, i pyhitäh nimi sinun, i tulov tsarstva sinun, i lienöv valdu sinun, kui taivahal i maal, leiby meiän heitelematoi anna meile-nygöi, i jätä meile velgat meiän, kui i myö heitämmyö velguniekoin meiän, i elä vie meidä pahah, i päästä meidy ounahas. Tätä kirjaa tarkotetaan myös eräässä kirjeessä, jossa Aunuksen koulujen esimies v. 1868 käskee sitä lähettää 200 kpl. karjalaisille ja lausuu sen dots}
\kat{Toukomies}{01.04.1925}{P13}{Kibunoi tuas uusi tuhjo Otteita evankeliumeista karjalankielellä vain syväinpohjaspa vain hengävys tuas ylenöö.}
\kat{Toukomies}{01.04.1925}{P13}{2. Kuin on kirjutettu Jesai-prorokan ( profietan ) kirjas: » Katsho, minä työnnän miun anhelit siun silmies icl, ja hain valmistaa siun ties ( dorogas ) »; Koska karjalankielinen kirjallisuus pääasiassa sisältää evankeliumikirjain, rukousten ja muiden uskonnollisten lukemisten käännöksiä, julkaisemme silloin tällöin otteita ensinmainituista näytteiksi nykyisemmältä tavalla suoritetuista käännösyrityksistä. Kergeinä kiitelgäh dots}
\kat{Toukomies}{01.04.1925}{P16}{45. Vain mendyö pois mies rubei leviel kuulovoittamah ja azhies tieduo levittämäh, jotta lisus ei eniä suattannuh julgi mennä linnoih, vain oletteli niijen ulgopuolel tyhjis paikois; ja joga paikas tuldih hänen luo. 1. H.; Otteita evankeliuuneista karjalankielellä; Kibunoi, suomalaisia runoja aunukselaisille, sivut 9-16. — Kuvia: Kalevalan heimoa kansallisilla kesäjuhlillaan; Kaksi karjalaisia esittävää kuvaa Olaus Magnin dots}

Erimielisyys näkyy myös siinä, mitä lehdet julkaisivat. \emph{Karjalan Vapaus} ja
\emph{Karjalaisuuden Ystävä} painoivat karjalankielistä tekstiä; \emph{Suomen Heimo}
kirjoitti karjalasta enemmän kuin karjalaksi.

\subsection{Suur-Suomi}

Sana esiintyy 233 kertaa. Se on vähemmän kuin voisi olettaa, ja se kannattaa merkitä:
heimolehdistön keskeinen sanasto oli \emph{heimo} ja \emph{Karjala}, ei
\emph{Suur-Suomi}.


\section{Talous ja elinkeinot}


\begin{center}\small
\begin{tabular}{@{}lrrrrrrrrrrrrrrr@{}}
\toprule
\textbf{Teema} & \textbf{25} & \textbf{26} & \textbf{27} & \textbf{28} & \textbf{29} & \textbf{30} & \textbf{31} & \textbf{32} & \textbf{33} & \textbf{34} & \textbf{35} & \textbf{36} & \textbf{37} & \textbf{38} & \textbf{39}\\
\midrule
kauppa & -- & -- & -- & -- & -- & -- & -- & -- & -- & -- & -- & -- & -- & -- & --\\
osuuskauppa & -- & -- & -- & -- & -- & -- & -- & -- & -- & -- & -- & -- & -- & -- & --\\
puutavara & -- & -- & -- & -- & -- & -- & -- & -- & -- & -- & -- & -- & -- & -- & --\\
metsätyö & -- & -- & -- & -- & -- & -- & -- & -- & -- & -- & -- & -- & -- & -- & --\\
työttömyys & -- & -- & -- & -- & -- & -- & -- & -- & -- & -- & -- & -- & -- & -- & --\\
laukkukauppa & -- & -- & -- & -- & -- & -- & -- & -- & -- & -- & -- & -- & -- & -- & --\\
maatalous & -- & -- & -- & -- & -- & -- & -- & -- & -- & -- & -- & -- & -- & -- & --\\
raha & -- & -- & -- & -- & -- & -- & -- & -- & -- & -- & -- & -- & -- & -- & --\\
\bottomrule
\end{tabular}
\end{center}
{\footnotesize Taulukko: taloutta koskevan sanaston esiintymät vuosittain.}

\subsection{Kauppa on läsnä, liiketalous ei}

\emph{Kauppa} tuottaa 1\,825 osumaa ja rahaa koskeva sanasto 3\,061. Luvut ovat suuria,
mutta konkordanssit paljastavat, että suuri osa niistä on ilmoituksia, jäsenmaksuja,
tilausmaksuja ja avustussummia --- järjestötaloutta, ei liiketaloutta.

Se on itsessään havainto lehtien luonteesta. Ne olivat aatteellisia julkaisuja, joiden
talousaineisto koostui pääosin ilmoituksista ja rahankeruusta.


\kat{Suomen Heimo}{20.02.1928}{P29}{Sähkösanasto: erittäin edullisesti kaikkia rauta- Zebra Code 3:rd Edition, kauppa-alaan kuuluvia tavaroita. „....,.., Saha: Inkilän asema. Helsingin Osuuskauppa r. L 15 Siirtomaatavaramyymälää •a Hankinnan lisäys vuonna 1927 suu- T., °....... * Lmamyymälaa rempi kuin kolmella muulla maamme vanhimjmalla yhtiöllä yhteensä. 24,000,000 pilrustustakkeja y. m. Keskinäinen Vakuutusyhtiö Ét\textasciicircum{}\textasciicircum{}\textasciicircum{} KALEVA SUf\&ä ) Hyvä ja varma kotimainen dots}
\kat{Suomen Heimo}{20.12.1929}{P7}{243 OSUUSLIIKE ARINA rl. Tänään on itsenäisyysläivä.}
\kat{Suomen Heimo}{20.12.1929}{P7}{Jäsenmäärä yli 2,000. Vuosimyynti n. 15,000,000 mk. OULU Puolueeton, aina kuluttajien etuja valvova osuusliike. Siirtmnfi, a, -, ruoka- ja sekatavara-, liha- ja leipämyymälöitä.}
\kat{Toukomies}{01.02.1929}{P18}{Kalastus, joka on tullut käytäntöön verrattain myöhään, on nykyään melkein kokonaan lamassa eikä sen uudelleen käytäntöön tulosta ole suuria toiveita. ( Vehkaojan — Vehrutshei ) osuuskauppa on ryhtynyt myymään tahkoja ja kovasinkiviä niin lähiseuduille kuin kaukaisempiinkin seutuihin, jotka tarvitsevat tahkokiveä. Harvoissa kylissä enää harjoitetaan kalastusta ja se vähä, mitä saadaan, menee melkein kokonaan omaan talouteen; dots}
\kat{Toukomies}{01.12.1929}{P23}{Mutta kuitenkin se ansaitsee huomiotamme yhtenä ilmauksena siitä mielenkiinnosta, mitä venäläiset tutkijat viime aikoina ovat osoittaneet nimenomaan suomalaisia kansoja kohtaan, joiden pariin meikäläisillä tiedemiehillä pitkiin aikoihin ei ole ollut pääsyä. Selitetään petun valmistusta ja kerrotaan, että petun käyttö, joka vallankumouksen jälkeisinä nälkävuosina oli hyvin yleistä, nyttemmin on supistunut melkein olemattomiin, dots}

\subsection{Puutavara}

Puutavara-alan sanasto (181 osumaa) on merkityksellinen tämän työn kannalta erityisesti
siksi, että se oli se ala, jolle Suojärvelle asettuneet aunukselaiset pakolaiset
sijoittuivat.


\kat{Toukomies}{01.09.1925}{P17}{41. Ondrei, Simon Pedrin velli, oli töine niis kahtes, kudamat oldih kuuldu midä livan sanoi, da käveldy Isusan jälles. Tämän valtion sahan johtaja on monta kertaa myöntänyt karjalaisten ansion työssä, mutta näyttää hänellekin olevan vaikeata saada esim. kerhomme tarkoituksiin jotain pientä huonetta, jota olemme pyytäneet. 52. Da Hai sanoi hänele: » Tovesti, tovesti sanon teile: 'täs ielleh suatto nähtä taivahan olevan augi, dots}
\kat{Toukomies}{08.1926}{P18}{Huolimatta matkain pituudesta oli suurin osa jäseniä ansiopaikoiltaan tilaisuuteen saapunut. Kokouksessa päätettiin pitää samalla sahan vanhalla pirtillä illanvietto saman kuun 30 p:nä. Kerhomme taloudellinen asema ei anna myöten maksaa pienintäkään palkkaa yhdellekään toimihenkilölle.}
\kat{Toukomies}{01.02.1927}{P23}{Tässä onkin eräs sen ansioita. 9 p:ksi sahan vanhalle pirtille. N:ot 1 — 2 -}
\kat{Toukomies}{01.05.1927}{P13}{Runsaasta ohjelmasta, joka alkoi pianonsoitolla ja kerhon toimihenkilön Mikko Saariston tervehdyspuheella, on mainittava neiti Klaudia Karvosen karjalanmurteiset lausuntanumerot, tervehdysten lukeminen, nti Viitasen yksinlaulut, lehtori Tikanojan juhlapuhe ja kerholaisten 'esittämä nelikuvaelmainen esitys " Karjalan häät ". 20 p:nä Toppilan sahan vanhalla pirtillä ohjelmallisen illanvieton niinikään runsaalla ohjelmalla ja dots}
\kat{Toukomies}{01.09.1927}{P2}{1930 — 31: Povenetsiin rakennetaan puuseppätehdas, kust. 1929 — 31: Kemiin rakennetaan 6-raaminen saha, kust. 1931 — 33: Kemiin rakennetaan 4-raaminen saha, kust.}
\kat{Toukomies}{01.09.1927}{P2}{1929 — 31: Kemiin rakennetaan 6-raaminen saha, kust. 1931 — 33: Kemiin rakennetaan 4-raaminen saha, kust. Metsän kulutus vuodeksi 1927 — 28 lasketaan 650,000 kuutiosyleksi.}

\subsection{Laukkukauppa}

Vanha karjalainen laukkukauppaperinne mainitaan vain 13 kertaa. Se oli 1930-luvulla jo
menneisyyttä, mutta juuri se perinne selittää, miksi maansa menettäneet talolliset
siirtyivät nimenomaan kauppaan.


\kat{Toukomies}{08.1926}{P5}{Suomeenko aikojen kuluessa asettuneiden ja täällä oleskelevien karjalaisten keskuudessa ja toimesta, vai Vienassa kotosalla olevan kansan havahtumisesta? Toisaalta tuli sitä yhtä itsetiedottomasti karjalaismiesten omain matkojen kautta Suomeen — ikivanhoista ajoista tehtyjen laukkukauppa 3a muiden matkojen kautta. Sen, mikä syttyi täällä, voi sanoa syttyneen sielläkin, mitä siellä ajateltiin, se sai paremman toteutumisen ja dots}
\kat{Toukomies}{08.1926}{P6}{oli laukut täysinä nyt. Laukkuryssä ei ollut vielä häipynyt mielestä. Mut laukussa ei katinkulta,}
\kat{Toukomies}{01.03.1934}{P3}{Näiden kesken kaikilla aloilla ja lukuisissa piireissä on kehittynyt ja noussut myöskin yleisemmän edistys- ja sivistyselämän harrastajia, yhteiskunnallisia toimihenkilöitä, lahjoittajia, kunnallisiin ja yleisiin sivisty spyrkimyksiin innolla osallistuneita. Sitkeydessään melkein vertoja vetämätöntä on ollut Vienan miesten ikiajoista harjoittama laukkukauppa, joka toiselta puolen on tavaran hakuun ja ostoon houkutellut aina dots}
\kat{Toukomies}{01.06.1934}{P13}{Samana vuonna kuin Genetz- Jännes alkaa toinenkin Karjalan uuttera kuvaaja ja tutkija, valkjärveläinen /. M. Salenius laajan karjala- aiheisen tuotantonsa. * ) Lisäyksenä kirjoituksemme edell. osaan mainittakoon, että Aleksanteri Rahkonen v. 1869 julkaisi pienen näytelmämukailun » Laukkuryssä ». 1881 ) ja » Tutkimus Aunuksen kielestä » ( Suomi-kirja II 17, 1884, erikseen 1885 ), jotka aina näihin asti ovat olleet dots}
\kat{Karjalan Vapaus}{01.1935}{P8}{1 ) Kirjoitettava vain toiselle puolelle paperia; » Kuljeksin ta kuljeksin, a sumtsa selköä painau... » Idästä pilvi taas nousi, tuli raesade, oraan alut löi, ruhjoi kuin ruttotauti.}

\subsection{Pula-aika}

Työttömyyttä ja lamaa koskeva sanasto (47 osumaa) on niukkaa, vaikka 1930-luvun alun
lama osui suoraan siihen hetkeen, jolloin pakolaisten yritykset olivat hauraimmillaan.


\kat{Suomen Heimo}{18.11.1928}{P16}{Tällaisia talonpoikaismaita kutsuttiin nimellä obza, jopka suuruus arvioitiin siihen kylvetyn viljamäärän mukaan. Työttömyys ja köyhyys lienee tullut Nevan seutujen asukkaille, jotka alkoivat muuttaa työansioihin muualle! Mutta yleensä obzien suuruus vaihteli suuresti riippuen maanlaadusta, paikallisista olosuhteista y.m.}
\kat{Suomen Heimo}{15.02.1930}{P28}{Sokerista Sokeri ei ole vain mauste. Ennen oli sokeri harvinainen mauste, ja pula-aika opetti taas säästämään sitä. Sokeri on nimittäin parhaita ravintoainei- tamme, ja myöskin hintansa suhteen edullisimpia.}
\kat{Suomen Heimo}{11.08.1930}{P36}{42 haarakonttoria; \_ ",... Ennen oli sokeri harvinainen mauste, ja pula-aika opetti taas säästämään sitä. Sokeri on nimit- täin parhainta ravintoainettamme ja myöskin hintansa suhteen v edullisimpia.}
\kat{Suomen Heimo}{16.12.1930}{P38}{Sen I osa, " I s a ja t y t a r ", kuvastaa sisämaan pienen sahayhteiskunnan oloja venäläissorron aiheuttaman passiivisen vastarinnan vuosina ja kansanelämän sielunliikkeitä ennen sosialidemokratian tunkeutumista maahamme. Mielikuvituksellinen aihe: eräs huomattava suurliikemies vetäytyy syrjään tavallisesta liike-elämästä, rakennuttaa itselleen loistoasunnon pohjoiseen, erämaiden keskelle, ja luo siellä dots}
\section{Paikat ja nimet}


\begin{center}\small
\begin{tabular}{@{}lrrrrrrrrrrrrrrr@{}}
\toprule
\textbf{Teema} & \textbf{25} & \textbf{26} & \textbf{27} & \textbf{28} & \textbf{29} & \textbf{30} & \textbf{31} & \textbf{32} & \textbf{33} & \textbf{34} & \textbf{35} & \textbf{36} & \textbf{37} & \textbf{38} & \textbf{39}\\
\midrule
Suojärvi & -- & -- & -- & -- & -- & -- & -- & -- & -- & -- & -- & -- & -- & -- & --\\
Mäntyselkä & -- & -- & -- & -- & -- & -- & -- & -- & -- & -- & -- & -- & -- & -- & --\\
Säämäjärvi & -- & -- & -- & -- & -- & -- & -- & -- & -- & -- & -- & -- & -- & -- & --\\
Repola & -- & -- & -- & -- & -- & -- & -- & -- & -- & -- & -- & -- & -- & -- & --\\
Jukola & -- & -- & -- & -- & -- & -- & -- & -- & -- & -- & -- & -- & -- & -- & --\\
\bottomrule
\end{tabular}
\end{center}
{\footnotesize Taulukko: tässä työssä seurattujen paikkojen ja nimien esiintymät vuosittain.}

\subsection{Suojärvi}

Suojärvi tuottaa 418 osumaa. Se oli Raja-Karjalan pitäjä, jonne suuri osa Aunuksesta
tulleista asettui, ja se esiintyy aineistossa sekä kerhotoiminnan että liike-elämän
yhteydessä.


\kat{Toukomies}{05.02.1925}{P10}{Ken siltojen kera elon kauniimman tois. ruman häätäisi pois? ( Suojärven ensimmäisen nuorisoseuratalon vihkiäisiin Leppäniemellä 21 / g ig24. ) rajat maan, sulut heimon ja hengen yö, sun suurin ja voittoisin vartiotyö:}
\kat{Toukomies}{05.02.1925}{P10}{rajat maan, sulut heimon ja hengen yö, sun suurin ja voittoisin vartiotyö: Suojärvi, suurien salojen maa, Suojärvi, suurien vetten vyö, Suojärvi, min rantoja rajoittaa satu hohdollaan salot peittänyt maan, runon kulta ja kunnia välkkynyt, miss' oisi se hohto nyt hohtelemaan, joka uudestaan unisuudestaan}
\kat{Toukomies}{05.02.1925}{P10}{rajat maan, sulut heimon ja hengen yö, sun suurin ja voittoisin vartiotyö: Suojärvi, suurien salojen maa, Suojärvi, suurien vetten vyö, Suojärvi, min rantoja rajoittaa satu hohdollaan salot peittänyt maan, runon kulta ja kunnia välkkynyt, miss' oisi se hohto nyt hohtelemaan, joka uudestaan unisuudestaan}
\kat{Toukomies}{05.02.1925}{P10}{rajat maan, sulut heimon ja hengen yö, sun suurin ja voittoisin vartiotyö: Suojärvi, suurien salojen maa, Suojärvi, suurien vetten vyö, Suojärvi, min rantoja rajoittaa satu hohdollaan salot peittänyt maan, runon kulta ja kunnia välkkynyt, miss' oisi se hohto nyt hohtelemaan, joka uudestaan unisuudestaan}
\kat{Toukomies}{05.02.1925}{P10}{satu hohdollaan salot peittänyt maan, runon kulta ja kunnia välkkynyt, miss' oisi se hohto nyt hohtelemaan, joka uudestaan unisuudestaan Suojärvi, suurene kummuillas, sun palttees ja rantasi kultahan luo suur maailma olkohon suuntanas, mut omista juuristas voimasi juo, kuin malmit järvies pohjasta tuo, Maan äärille nouskaa te talot elämän, talot valkeuden nouskaa te paistelemaan, isot ikkunat illassa välkkymähän ja aamun dots}
\kat{Toukomies}{05.02.1925}{P17}{Karjalan Sivistysseura, Vienan Karjalaisten Liiton perijä, laskee täten syvimmällä kaipauksen, kiitollisuuden ja kunnioituksen tunteella tämän seppeleensä tohtori Benjamin Mitron haudalle. G. Gulin, os. Pitkäranta; Salmin jakopaikka, hoitajana vääpeli V. Jakobson, os. Salmi ja Suojärven jakopaikka, hoitajana vääpeli I. Halonen, os. Suojärvi, Suvilahti. Kajaanin jakopaikka, hoitajana V. Rasi, os. Kajaani; Kuhmoniemen dots}
\kat{Toukomies}{05.02.1925}{P17}{Karjalan Sivistysseura, Vienan Karjalaisten Liiton perijä, laskee täten syvimmällä kaipauksen, kiitollisuuden ja kunnioituksen tunteella tämän seppeleensä tohtori Benjamin Mitron haudalle. G. Gulin, os. Pitkäranta; Salmin jakopaikka, hoitajana vääpeli V. Jakobson, os. Salmi ja Suojärven jakopaikka, hoitajana vääpeli I. Halonen, os. Suojärvi, Suvilahti. Kajaanin jakopaikka, hoitajana V. Rasi, os. Kajaani; Kuhmoniemen dots}

\subsection{Mäntyselkä}

Kolmetoista osumaa koko aineistossa. Ne ovat kaikki liitteessä B. Pitäjä esiintyy
rautatieluettelossa, hallinnollisissa kuvauksissa ja muutamassa henkilöesittelyssä.


\kat{Toukomies}{06.1926}{P12}{Sinä et voinut palata rakastamillesi seuduille, vaikka niin hartaasti sitä toivoit, uusi järjestelmä ei herättänyt sinussa luottamusta. Suojärven radan päätekohtaa ja pohjoisin Kontiovaarasta, Äänis järven pohjoisimmasta perukasta, Mäntyselän ja Juustjärven kautta Poraj arvelle. Saamiemme varmojen tietojen mukaan on Venäjä päättänyt rakentaa kolme kapearaiteista rautatietä Muurmanin radalta länteen päin, Suomen rajan lähelle.}
\kat{Toukomies}{01.09.1927}{P13}{rautatie Mäntyselkä Njuhotskaja}
\kat{Suomen Heimo}{20.12.1931}{P26}{Sitäpaitsi uskonnollisestikin Karjala oli Venäjän vallankumoukseen saakka, ortodoksisen kirkon välityksellä, kiinteässä riippuvaisuussuhteessa mainittuun esivaltaan. Poventsan kihlakuntaan kuuluivat seuraavat' 10 kuntaa: Repola, Porajärvi, Rukajärvi, Lintajärvi, Mäntyselkä, Paatene, Aunuksen kihlakuntaan kuuluivat seuraavat 9 kuntaa: Tulemajärvi, Vieljärvi, Vitele, Kotkajärvi, Keskotjärvi, Riipuskala, Mätusova, Vaasheni ja dots}
\kat{Toukomies}{01.02.1931}{P18}{Koska Valtion Pakolaisavustuskeskus on kuluvaksi vuodeksi kieltäytynyt myöntämästä rahamäärää postimenojen peittämiseksi ja koska kansliaa hoitava papisto on lain mukaan oikeutettu perimään lunastukset todistuksista henkilöiltä, jotka postitse tilaavat papintodistukset tai muut asiakirjat, huomautetaan niille pakolaisille, jotka tulevat hakemaan postitse tarpeellisia todistuksia, että todistus ilman merkintää aviosuhteista dots}
\kat{Toukomies}{01.12.1934}{P18}{Kun alueen turvana oli pieni Pohjois-Aunuksen suojeluskunta, jonka muodostivat rohkeimmat miehet Repolan ja Porajärven pitäjistä päällikkönään Iyähes vuotta myöhemmin, kesäkuussa 1919 seurasivat esimerkkiä Porajärvi ja kohta sen jälkeen LÄndajärvi, s. o. Suomen rajalla oleva osa Mäntyselän pitäjää. Tämä miestemme sankariteko se oli, joka antoi voiman ja rohkeuden koko isänmaalliselle osalle kansastamme ryhtyä vapauttamaan dots}
\kat{Karjalan Vapaus}{05.1935}{P17}{Kauppias Ilja Zertijeff, Kauppias F. Ojamo on syntynyt Mäntyselän pit. 1914 — 1919 hän hoiti postinkuljettajan virkaa.}
\kat{Karjalan Vapaus}{05.1935}{P19}{Liikemies Pekka Jukola on synt. 29 / 1 1898 Mäntyselän pit. 1921 — 1922.}
\kat{Karjalan Vapaus}{05.1935}{P19}{Kun 1932 hänen toimesta perustettiin T:ni M. Pedas niminen liike, tuli hän sen johtajaksi. Toiminimi M. Pedaksen johtaja J. Dorojejeff on syntynyt Mäntyselän Karhumäellä 24 / 10 -94 mv. Joensuussa kangastavaroilla tukku- ja vähittäiskauppaa harjoittava kauppias S. Saharoff on syntynyt Rukajärvellä helmik.}

\subsection{Säämäjärvi}

Säämäjärvi (42 osumaa) oli Aunuksen pakolaisten suurin lähtökunta --- yli 400
pakolaista --- ja sieltä kotoisin olevien kertomuksia julkaistiin karjalaksi.


\kat{Toukomies}{05.1926}{P13}{Tarvitaan pitkä-aikaista, ennakkoluulotonta kansallista kasvatusta ja herätystyötä ennenkuin noiden seutujen asukkaat perinpohjin, juurta jaksain itsensä, » löytävät », tuntevat suomalaisiksi. Säämäjärven ja Pyhäniemen kylistä kotoisin olevat pakolaiset selittävät, että jo ennenkin he ovat käyneet Suoj arvella, mutta eivät koskaan Suomessa saakka — tarkoittaen Korpiselkää. Korpiselällä, Ilomantsissa, Pielis järvellä ei kansa dots}
\kat{Suomen Heimo}{25.05.1929}{P7}{Kevättalvi v. 1919 oli ollut erinomaisen leuto ja Ensimmäiset kelirikon oireet alkoivat ilmestyä jo neuvottelujen aikana maaliskuun puolivälissä. Orusjärven ryhmä hiihtäisi Tulema- ja Videljärven kautta Rihmakylään ( Prjäshaan ), yhtyen siellä Säämäjärven kautta hiihtävään Suojärven ryhmään. 5 miljoonaa kiväärin ja kk:n panosta, 600 shrapnellia, 200 granaattia, 11,000 käsigranaattia y.m. Yhteysvälineetkin olivat hyvin dots}
\kat{Suomen Heimo}{25.05.1929}{P9}{Pienempiä joukkoosastoja eteni vielä Aleksanteri Svirin luostarille ja Syvärin välittömään läheisyyteen. Yhtämittaa taistellen sitten edettiin Säämäjärven kautta Prjäshaan, jossa -jo II ryhmä sijaitsi. 2 " p! nä luovuttava Aunuksesta, mutta kaupunki vallattiin taas, ja. toisen kerran siitä luovuttiin toukok.}
\kat{Suomen Heimo}{20.06.1929}{P8}{Huhtikuun 6 päivänä saapuivat ensimmäiset 60 miestä, ja tämän jälkeen päivittäin 100 — 200 miestä, niin että 11. 4. Sortavalassa oli vasta n. 600 miestä. Etenemistä Suo järveltä Säämäjärven yli Rihmakylään oli pidettävä ainoastaan apuliikkeenä, jonka tarkoituksena oli suojata pohjoisen ryhmän vasenta sivustaa yllätyksiltä. Vihollinen, joka olisi uskaltanut jättäytyä molempien etenevien ryhmien väliin, olisi ehdottomasti dots}
\kat{Suomen Heimo}{20.12.1931}{P26}{Edellä esitettyjen kihlakuntien hallinnollisten rajojen mukaan Kauko-Karjala, Kuollaa lukuunottamatta, käsitti 136,300 km 2:n pinta-alan, josta Vienan-Karjalan osalle tuli 58,000 km 2 ja Aunuksen Karjalan osalle 78,300 km 2. Petroskoin kihlakunnassa oli 12 kuntaa: Ostretshini, Soutjärvi ( asuu melkein yksinomaan vepsäläisiä ), Ladva, Pyhäjärvi, Suoju, Säämäjärvi, Munjärvi, Kontupohja, Tepenitsa, Jalguba, Suurlahti ja Talvoja. dots}\section{Yksi suku aineistossa}

Tämä luku on eri lajia kuin edelliset. Se ei yleistä mitään; se kertoo, mitä yhdestä
suvusta löytyy, kun neljää miljoonaa sanaa haravoi nimellä.

Sukunimi \emph{Jukola} tuottaa 56 osumaa. Suurin osa niistä ei koske tätä sukua:
mukana on urheilukirjailija Martti Jukola, Nurmeksen Jukolan osuuskauppa ja lukuisia
viittauksia Seitsemään veljekseen. Ne on karsittu pois. Jäljelle jää seuraava.

\subsection{Pekka Jukola, esittely 1935}

\emph{Karjalan Vapaus} julkaisi toukokuussa 1935 sivullaan 19 kuuden karjalaisen
liikemiehen esittelygallerian. Mukana ovat kauppias Vasili Nikitin Aunuksen
Tuuloksesta, kauppias S. Saharoff Rukajärveltä, toiminimi M. Pedaksen johtaja
J. Dorojejeff Mäntyselän Karhumäeltä --- ja Pekka ja Paavo Jukola.

Sivun palstat ovat OCR:ssä sekoittuneet, joten yksittäisen lauseen kohdistaminen
henkilöön ei aina onnistu. Yksi lause on kuitenkin kiistaton, koska se nimeää yhtiön:

\begin{quote}\itshape
Myöhemmin v. 1930 aloitti itsenäisen liikkeen perustamalla Suojärven Lapinjärvelle
sahatavarakaupan sekä harjoittaen puutavaraliikettä, mikä liike v. 1931 muutettiin
Veljekset Jukola O/y-nimiseksi osakeyhtiöksi, jonka toimitusjohtajana on edelleenkin.
\end{quote}

Ja tämä: \emph{``Liikemies Pekka Jukola on synt. 29/1 1898 Mäntyselän pit.''}

Samalla sivulla on lauseita, jotka kuuluvat todennäköisesti hänelle mutta joita ei voi
kohdistaa varmasti: \emph{``Jo 19-vuotiaana hän joutui maailmansotaan''};
\emph{``Sieltä palasi kotiin keväällä v.\ 1918''}; \emph{``Otti osaa karjal.\ vapausl.\
1921--1922''}; \emph{``Samana syksynä muutti Suomeen''}; ja
\emph{``ensimäisiä Suojärvellä toimivan Karjalan Kansalaisseuran perustajia ja kuuluen
sen johtokuntaan sekä toimien jatkuvasti seuran varapuheenjohtajana''}.

Nämä on tarkistettava alkuperäisestä sivusta, jossa palstajako ratkaisee asian
silmämääräisesti.

\subsection{Paavo Jukola järjestömiehenä}

Toinen veljeksistä esiintyy aineistossa toistuvasti ja aivan toisessa roolissa kuin
liikemiehenä.

\begin{center}
\begin{tabularx}{\linewidth}{@{}lX@{}}
\toprule
1934 & \emph{Toukomies} 1.5.: kokouksessa ``pöytäkirjaa piti sihteeri kaupanhoitaja Paavo Jukola'' \\
1934 & \emph{Toukomies} 12., vuosisisällysluettelo: ``Paavo (Jukola): Suojärven kuulumisia'' --- hän kirjoitti lehteen omalla palstallaan \\
1938 & \emph{Karjalaisuuden Ystävä} 7.: liiton sihteeri, ``kauppias Paavo Jukola, os.: Rajasuvanto'' \\
1938 & Sama lehti: toimituskunnan jäsen kokoonpanossa E. V. Ahtia, Ville Helve, Paavo Jukola ja N. Ruotsi \\
\bottomrule
\end{tabularx}
\end{center}

{\footnotesize Taulukko: Paavo Jukolan esiintymiset heimolehdistössä.}

Viimeinen rivi ansaitsee huomion. \textbf{E. V. Ahtia} oli karjalan kielen tutkija,
laajan sanastonkeruun tekijä ja vuonna 1938 ilmestyneen karjalan kieliopin kirjoittaja
--- alansa keskeisin nimi. Suojärveläinen kauppias istui saman lehden toimituskunnassa
hänen kanssaan.

Karjalaisuuden Ystäväin Liitto perustettiin \emph{Karjalaisuuden Ystävän} oman
kertomuksen mukaan Sortavalan Seurahuoneella 20.3.1938; sitä oli edeltänyt
neuvottelukokous kreikkalaiskatolisen kirkkokunnan juhlasalissa 13.2.1938.

\subsection{Veljekset Jukola Oy ilmoittajana}

Yhtiö osti mainostilaa heimolehdistä vähintään vuosina 1935--1939.

\begin{center}
\begin{tabularx}{\linewidth}{@{}lX@{}}
\toprule
\emph{Karjalan Vapaus} 5/1935 & ``WELJEKSET JUKOLA OY.'' ilmoituspalstalla Suojärven liikkeiden joukossa \\
\emph{Karjalan Vapaus} 3/1936 & ``PUUTAVARALIIKE konttori Suojärvi as. VELJEKSET JUKOLA OY.'' \\
\emph{Karjalaisuuden Ystävä} 7/1938 & Puutavaraliike ja sekatavarakauppa; ``Konttori 2, Myymälä 20'' \\
\emph{Karjalaisuuden Ystävä} 12/1938 & Sivuliike Rajasuvanto ja Lapinjärvi; Suojärvi, Jehkilä \\
\emph{Karjalaisuuden Ystävä} 10/1939 & Suojärvi, Välikylä --- viimeinen tavattu ilmoitus \\
\bottomrule
\end{tabularx}
\end{center}

{\footnotesize Taulukko: yhtiön ilmoitukset heimolehdistössä.}

Vuoden 1939 lokakuun ilmoitus on aineiston viimeinen. Kuukautta myöhemmin alkoi
talvisota.

\subsection{Doroiset}

Suvun oma nimi oli \emph{Doroiset}; \emph{Ivanoff} oli rekisterinimi. Aineistossa
esiintyy kaksi kertaa muoto \emph{Dorojejeff}, molemmat samalla vuoden 1935 sivulla:


\kat{Karjalan Vapaus}{05.1935}{P19}{13 ja Oikok. 19 Juho Dorojejeffin poikana. Pekka Jukola Samana syksynä muutti Suomeen.}
\kat{Karjalan Vapaus}{05.1935}{P19}{Jo 1890, siis 22 vuotiaana, perusti hän oman tukkuliikkeen, joka myi kangastavaroita. Kun 1932 hänen toimesta perustettiin T:ni M. Pedas niminen liike, tuli hän sen johtajaksi. Toiminimi M. Pedaksen johtaja J. Dorojejeff on syntynyt Mäntyselän Karhumäellä 24 / 10 -94 mv. Joensuussa kangastavaroilla tukku- ja vähittäiskauppaa harjoittava kauppias S. Saharoff on syntynyt Rukajärvellä helmik. Tämän sahan palon jälkeen v. 1928 dots}

J. Dorojejeff on siis Mäntyselän Karhumäellä 24.10.1894 syntynyt liikemies,
maanviljelijä Juho Dorojejeffin poika. Sukulaisuutta Jukoloihin ei ole tässä
osoitettu, mutta yhdistelmä --- sama pitäjä, sama nimi, sama ammatti, sama
pakolaisuus --- tekee siitä todennäköisen ja tarkistamisen arvoisen.

\subsection{Ympäröivä piiri}

Aineisto näyttää myös sen verkoston, jossa nämä miehet toimivat. Vuoden 1935
gallerian muut nimet ovat karjalaisia pakolaisia, jotka kaikki päätyivät kauppaan:
Nikitin Aunuksen Tuuloksesta, Saharoff Rukajärveltä joka aloitti Pietarissa
kaksitoistavuotiaana ja tuli Suomeen 1917 ``koska tällöin alettiin kauppaliikkeitä
ryöstää ja yksityisyrittelijöitä vainota'', Dorojejeff Mäntyselältä.

Se on sama kuvio, joka näkyy tilastossa: maansa menettäneet talolliset siirtyivät
kauppaan. Galleria on siitä henkilögalleria.

\section{Miten teemat muuttuivat 1925--1939}

Lopuksi kokoan sen, mitä aineisto sanoo muutoksesta. Nämä ovat hypoteeseja, joita
frekvenssit tukevat --- eivät todistettuja tuloksia.

\subsection{Kolme kautta}

\textbf{1925--1930: menetetyn alueen kausi.} Puhe on Tarton rauhasta, pakolaisten
sijoittamisesta, avustuksista ja kielikysymyksestä. Neuvostovaltaa käsitellään
poliittisena vastustajana mutta arjen kuvausta rajan takaa on vähän. Kolhoosisanastoa
ei ole lainkaan.

\textbf{1931--1936: rajan takaisen arjen kausi.} Kolhoosi, karkotus, vankileiri ja
nälkä ilmestyvät sanastoon. Tämä on aineiston tiheintä aikaa: silloin julkaistaan
ensikäden kertomuksia, myös karjalaksi. \emph{Karjalan Vapaus} osuu juuri tähän
ikkunaan.

\textbf{1937--1939: hiljenemisen kausi.} Kolhoosisanasto katoaa, Stalin-maininnat
jatkuvat, mutta rajantakaisen arjen kuvaus ohenee. Samaan aikaan järjestötoiminta
vilkastuu ja uusia lehtiä perustetaan. Puhe kääntyy siitä mitä siellä tapahtuu siihen,
mitä täällä tehdään.

\subsection{Neljä havaintoa, jotka kestävät}

\begin{enumerate}[leftmargin=1.4em,itemsep=0.35em]
\item \textbf{Sanastonmuutos on jyrkkä ja ajoitettavissa.} Kolhoosi tulee sanastoon
vuonna 1934, neljä vuotta itse tapahtuman jälkeen. Se on mitattu viive tiedonkulussa.
\item \textbf{Vakoilusanasto puuttuu lähes kokonaan.} Alle sata osumaa neljässä
miljoonassa sanassa julkaisuissa, jotka muuten kuvasivat rajan takaisia oloja
jatkuvasti.
\item \textbf{Pakolaisuus on identiteetti, kansalaisuus ei ole aihe.}
Pakolainen 1\,160 osumaa, kansalaisuus 15.
\item \textbf{Talouspuhe on järjestötaloutta.} Suuret lukemat syntyvät ilmoituksista
ja jäsenmaksuista, eivät liiketaloudellisesta keskustelusta.
\end{enumerate}

\section{Rajoitteet}

\begin{enumerate}[leftmargin=1.4em,itemsep=0.35em]
\item \textbf{Vuodet 1919--1924 puuttuvat.} Kaikki tämän työn päätelmät koskevat
vuosia 1925--1939.
\item \textbf{\emph{Karjalan Vapautta} on vain 1934--1936.} Alkuperäistä
tutkimuskysymystä ei voi tästä aineistosta vastata.
\item \textbf{Frekvenssi ei mittaa sävyä.} Luvut kertovat mistä puhuttiin, eivät mitä
sanottiin. Liite B on tässä suhteessa tärkeämpi kuin liite A.
\item \textbf{OCR aliarvioi järjestään.} Tunnistamattomat sanamuodot eivät tuota
osumia; todelliset luvut ovat esitettyjä suurempia, eikä virhe jakaudu tasaisesti ---
karjalankielinen ja fraktuurainen teksti kärsii eniten.
\item \textbf{Hakusanasto on tekijän valitsema.} Toisenlaisella sanastolla saisi
toisenlaisen kuvan. Kaikki hakulausekkeet ovat liitteessä A, joten valinta on
tarkistettavissa ja korjattavissa.
\item \textbf{Palstat sekoittuvat.} Liitteen B katkelmissa on kohtia, joissa kaksi eri
juttua tai ilmoitus ja leipäteksti ovat sulautuneet yhteen. Yksittäistä katkelmaa ei
pidä siteerata ilman alkuperäisen tarkistusta.
\item \textbf{Lehtien painotus on Vienan puolella.} Karjalan Sivistysseura syntyi
Vienan Karjalaisten Liitosta, ja aunukselaiset olivat liikkeessä vähemmistö. Aineisto
kertoo siksi vähemmän Aunuksesta kuin Vienasta.
\end{enumerate}

\section{Miten tätä käytetään}

Tämä on työkalu, ei kirja. Kolme käyttötapaa.

\textbf{Ajankuvan hakemisto.} Liitteestä B löytää nopeasti, millä sanoilla jostain
asiasta puhuttiin ja minä vuonna. Romaania kirjoittaessa se on käyttökelpoisempi kuin
tutkimuskirjallisuus, koska se antaa aikalaisen sanamuodon eikä historioitsijan.

\textbf{Hakupohja.} Liitteen A hakulausekkeet ovat sellaisenaan toistettavissa
Kielipankin Korp-palvelussa, ja niitä voi laajentaa. Kaikki tämän raportin luvut
saa uusittua noin puolessa tunnissa.

\textbf{Aukkojen kartta.} Se, mitä aineistossa ei ole, on usein kiinnostavampaa kuin
se mitä on. Tämän työn selvin esimerkki on vakoilusanaston puuttuminen.

\appendix
\section{Liite A: hakulausekkeet ja vuosittaiset frekvenssit}

Jokainen haku on tehty muodossa

\begin{quote}\small\ttfamily
[word="LAUSEKE" \%c] :: match.text\_publ\_title = "Karjalan Vapaus|Toukomies|\\
Viena-Aunus|Suomen Heimo|Karjalaisuuden Ystävä|Itä-Karjala|Heimotyö"
\end{quote}

korpuksissa KLK\_FI\_1925 \dots\ KLK\_FI\_1939.

\begin{center}\scriptsize
\begin{longtable}{@{}lp{52mm}rrrrrrrrrrrrrrrr@{}}
\toprule
\textbf{Teema} & \textbf{Hakulauseke} & \textbf{25} & \textbf{26} & \textbf{27} & \textbf{28} & \textbf{29} & \textbf{30} & \textbf{31} & \textbf{32} & \textbf{33} & \textbf{34} & \textbf{35} & \textbf{36} & \textbf{37} & \textbf{38} & \textbf{39} & \textbf{Yht.} \\
\midrule
\endhead

Venäjä & {\ttfamily Venäjä|Venäjän|Venäjällä|Venäjältä|Venäjää} & -- & -- & -- & -- & -- & -- & -- & -- & -- & -- & -- & -- & -- & -- & -- & \textbf{3684} \\
Neuvostoliitto & {\ttfamily Neuvostoliitto|Neuvostoliiton|Neuvostoliitossa|Neuvostoliittoon} & -- & -- & -- & -- & -- & -- & -- & -- & -- & -- & -- & -- & -- & -- & -- & \textbf{1148} \\
Neuvosto-Karjala & {\ttfamily Neuvosto-Karjala|Neuvosto-Karjalan|Neuvosto-Karjalassa} & -- & -- & -- & -- & -- & -- & -- & -- & -- & -- & -- & -- & -- & -- & -- & \textbf{480} \\
venäläinen & {\ttfamily venäläinen|venäläiset|venäläisiä|venäläisten|venäläisyys|venäläisyyden} & -- & -- & -- & -- & -- & -- & -- & -- & -- & -- & -- & -- & -- & -- & -- & \textbf{1629} \\
ryssä & {\ttfamily ryssä|ryssät|ryssän|ryssiä|ryssien} & -- & -- & -- & -- & -- & -- & -- & -- & -- & -- & -- & -- & -- & -- & -- & \textbf{817} \\
bolsevikki & {\ttfamily bolshevikki|bolshevikit|bolshevikkien|bolshevikkeja|bolsevikki|bolsevikit|bolsevikkien} & -- & -- & -- & -- & -- & -- & -- & -- & -- & -- & -- & -- & -- & -- & -- & \textbf{335} \\
kommunisti & {\ttfamily kommunisti|kommunistit|kommunistien|kommunismi|kommunismin} & -- & -- & -- & -- & -- & -- & -- & -- & -- & -- & -- & -- & -- & -- & -- & \textbf{827} \\
punainen valta & {\ttfamily punavalta|punavallan|punaiset|punaisten|punakaarti|punakaartin} & -- & -- & -- & -- & -- & -- & -- & -- & -- & -- & -- & -- & -- & -- & -- & \textbf{478} \\
Stalin & {\ttfamily Stalin|Stalinin|Stalinia} & -- & -- & -- & -- & -- & -- & -- & -- & -- & -- & -- & -- & -- & -- & -- & \textbf{402} \\
Lenin & {\ttfamily Lenin|Leninin|Leniniä|Leninets} & -- & -- & -- & -- & -- & -- & -- & -- & -- & -- & -- & -- & -- & -- & -- & \textbf{158} \\
GPU/tsheka & {\ttfamily GPU|G.P.U|tsheka|tshekan|tsheKa|Gepeu|gepeu} & -- & -- & -- & -- & -- & -- & -- & -- & -- & -- & -- & -- & -- & -- & -- & \textbf{99} \\
kolhoosi & {\ttfamily kolhoosi|kolhoosin|kolhoosiin|kolhoosit|kolhoosien|kolhosi|kolhozah} & -- & -- & -- & -- & -- & -- & -- & -- & -- & -- & -- & -- & -- & -- & -- & \textbf{133} \\
pakkotyö & {\ttfamily pakkotyö|pakkotyön|pakkotyöhön|pakkotyöleiri|pakkotyöleirin} & -- & -- & -- & -- & -- & -- & -- & -- & -- & -- & -- & -- & -- & -- & -- & \textbf{35} \\
Solovetsk & {\ttfamily Solovetsk|Solovetskin|Solovetski|Solofkoih|Solovetsissa} & -- & -- & -- & -- & -- & -- & -- & -- & -- & -- & -- & -- & -- & -- & -- & \textbf{64} \\
karkotus & {\ttfamily karkotettu|karkotettiin|karkotus|karkotuksen|karkoitettu|karkoitettiin} & -- & -- & -- & -- & -- & -- & -- & -- & -- & -- & -- & -- & -- & -- & -- & \textbf{168} \\
vangitseminen & {\ttfamily vangittu|vangittiin|vangitseminen|pidätetty|pidätettiin|vankila|vankilaan} & -- & -- & -- & -- & -- & -- & -- & -- & -- & -- & -- & -- & -- & -- & -- & \textbf{199} \\
teloitus & {\ttfamily teloitettu|teloitettiin|teloitus|ammuttiin|surmattiin} & -- & -- & -- & -- & -- & -- & -- & -- & -- & -- & -- & -- & -- & -- & -- & \textbf{56} \\
nälkä & {\ttfamily nälkä|nälän|nälkää|nälänhätä|nälänhädän} & -- & -- & -- & -- & -- & -- & -- & -- & -- & -- & -- & -- & -- & -- & -- & \textbf{176} \\
sorto & {\ttfamily sorto|sorron|sortoa|sortotoimet|hirmuvalta|hirmuvallan} & -- & -- & -- & -- & -- & -- & -- & -- & -- & -- & -- & -- & -- & -- & -- & \textbf{274} \\
vakoilu & {\ttfamily vakoilu|vakoilun|vakoilua|vakooja|vakoojat|vakoojan|urkkija|urkkijat} & -- & -- & -- & -- & -- & -- & -- & -- & -- & -- & -- & -- & -- & -- & -- & \textbf{36} \\
tiedustelu & {\ttfamily tiedustelu|tiedustelun|tiedustelija|tiedustelijat|tiedustelijoita} & -- & -- & -- & -- & -- & -- & -- & -- & -- & -- & -- & -- & -- & -- & -- & \textbf{37} \\
rajanylitys & {\ttfamily rajanylitys|rajanylityksen|rajaseutu|rajaseudun|rajavartio|rajavartijat|rajavartiosto} & -- & -- & -- & -- & -- & -- & -- & -- & -- & -- & -- & -- & -- & -- & -- & \textbf{162} \\
salakuljetus & {\ttfamily salakuljetus|salakuljetuksen|salakuljettaja|salakuljettajat|pirtu|pirtun} & -- & -- & -- & -- & -- & -- & -- & -- & -- & -- & -- & -- & -- & -- & -- & \textbf{12} \\
loikkari & {\ttfamily loikkari|loikkarit|loikkarien|karkuri|karkurit|luvaton} & -- & -- & -- & -- & -- & -- & -- & -- & -- & -- & -- & -- & -- & -- & -- & \textbf{20} \\
provokaattori & {\ttfamily provokaattori|provokaattorit|kavaltaja|kavaltajat|ilmiantaja|ilmiantajat} & -- & -- & -- & -- & -- & -- & -- & -- & -- & -- & -- & -- & -- & -- & -- & \textbf{9} \\
pakolainen & {\ttfamily pakolainen|pakolaiset|pakolaisten|pakolaisia|pakolaisuus} & -- & -- & -- & -- & -- & -- & -- & -- & -- & -- & -- & -- & -- & -- & -- & \textbf{1160} \\
kansalaisuus & {\ttfamily kansalaisuus|kansalaisuuden|kansalaisuutta|kansalaistaminen} & -- & -- & -- & -- & -- & -- & -- & -- & -- & -- & -- & -- & -- & -- & -- & \textbf{15} \\
avustus & {\ttfamily avustus|avustuksen|avustusta|huolto|huollon|avustuskomitea} & -- & -- & -- & -- & -- & -- & -- & -- & -- & -- & -- & -- & -- & -- & -- & \textbf{225} \\
asutus & {\ttfamily asutus|asutuksen|asuttaminen|maansaanti|tilaton|tilattomat} & -- & -- & -- & -- & -- & -- & -- & -- & -- & -- & -- & -- & -- & -- & -- & \textbf{239} \\
heimo & {\ttfamily heimo|heimon|heimoa|heimoveljet|heimoaate|heimoaatteen} & -- & -- & -- & -- & -- & -- & -- & -- & -- & -- & -- & -- & -- & -- & -- & \textbf{2603} \\
heimosoturi & {\ttfamily heimosoturi|heimosoturit|heimosoturien|heimosota|heimosodat|heimosotien} & -- & -- & -- & -- & -- & -- & -- & -- & -- & -- & -- & -- & -- & -- & -- & \textbf{192} \\
Suur-Suomi & {\ttfamily Suur-Suomi|Suur-Suomen|Suursuomi|Suursuomen} & -- & -- & -- & -- & -- & -- & -- & -- & -- & -- & -- & -- & -- & -- & -- & \textbf{233} \\
Aunus & {\ttfamily Aunus|Aunuksen|Aunuksessa|Aunukseen|aunukselainen|aunukselaiset} & -- & -- & -- & -- & -- & -- & -- & -- & -- & -- & -- & -- & -- & -- & -- & \textbf{1725} \\
Viena & {\ttfamily Viena|Vienan|Vienassa|vienalainen|vienalaiset} & -- & -- & -- & -- & -- & -- & -- & -- & -- & -- & -- & -- & -- & -- & -- & \textbf{2049} \\
karjalan kieli & {\ttfamily karjalankieli|karjalankielinen|karjalankielellä|karjalaksi|äidinkieli|äidinkielen} & -- & -- & -- & -- & -- & -- & -- & -- & -- & -- & -- & -- & -- & -- & -- & \textbf{225} \\
Tarton rauha & {\ttfamily Tartto|Tarton|rauhansopimus|rauhansopimuksen} & -- & -- & -- & -- & -- & -- & -- & -- & -- & -- & -- & -- & -- & -- & -- & \textbf{515} \\
kauppa & {\ttfamily kauppa|kaupan|kauppaa|kauppias|kauppiaat|kauppiaan|liike|liikkeen} & -- & -- & -- & -- & -- & -- & -- & -- & -- & -- & -- & -- & -- & -- & -- & \textbf{1825} \\
osuuskauppa & {\ttfamily osuuskauppa|osuuskaupan|osuusliike|osuusliikkeen|osuustoiminta} & -- & -- & -- & -- & -- & -- & -- & -- & -- & -- & -- & -- & -- & -- & -- & \textbf{76} \\
puutavara & {\ttfamily puutavara|puutavaran|puutavaraliike|saha|sahan|sahat|metsäkauppa} & -- & -- & -- & -- & -- & -- & -- & -- & -- & -- & -- & -- & -- & -- & -- & \textbf{181} \\
metsätyö & {\ttfamily metsätyö|metsätyöt|savotta|savotan|uitto|uiton|hakkuu|hakkuut} & -- & -- & -- & -- & -- & -- & -- & -- & -- & -- & -- & -- & -- & -- & -- & \textbf{30} \\
työttömyys & {\ttfamily työttömyys|työttömyyden|työtön|työttömät|pula-aika|lama} & -- & -- & -- & -- & -- & -- & -- & -- & -- & -- & -- & -- & -- & -- & -- & \textbf{47} \\
laukkukauppa & {\ttfamily laukkukauppa|laukkukaupan|laukkukauppias|laukkuryssä|sumtsa} & -- & -- & -- & -- & -- & -- & -- & -- & -- & -- & -- & -- & -- & -- & -- & \textbf{13} \\
maatalous & {\ttfamily maatalous|maatalouden|maanviljelys|maanviljelijä|karjatalous} & -- & -- & -- & -- & -- & -- & -- & -- & -- & -- & -- & -- & -- & -- & -- & \textbf{247} \\
raha & {\ttfamily markkaa|markan|rupla|ruplaa|hinta|hinnat|verotus|vero} & -- & -- & -- & -- & -- & -- & -- & -- & -- & -- & -- & -- & -- & -- & -- & \textbf{3061} \\
Jukola & {\ttfamily Jukola|Jukolan} & -- & -- & -- & -- & -- & -- & -- & -- & -- & -- & -- & -- & -- & -- & -- & \textbf{56} \\
Aksenoff & {\ttfamily Aksenoff|Aksenoffin|Aksenov} & -- & -- & -- & -- & -- & -- & -- & -- & -- & -- & -- & -- & -- & -- & -- & \textbf{0} \\
Suojärvi & {\ttfamily Suojärvi|Suojärven|Suojärvellä|Suojärveltä} & -- & -- & -- & -- & -- & -- & -- & -- & -- & -- & -- & -- & -- & -- & -- & \textbf{418} \\
Mäntyselkä & {\ttfamily Mäntyselkä|Mäntyselän|Mäntyselässä|Mäntyselästä} & -- & -- & -- & -- & -- & -- & -- & -- & -- & -- & -- & -- & -- & -- & -- & \textbf{13} \\
Säämäjärvi & {\ttfamily Säämäjärvi|Säämäjärven|Seämärvi|Seämärven} & -- & -- & -- & -- & -- & -- & -- & -- & -- & -- & -- & -- & -- & -- & -- & \textbf{42} \\
Repola & {\ttfamily Repola|Repolan|Porajärvi|Porajärven} & -- & -- & -- & -- & -- & -- & -- & -- & -- & -- & -- & -- & -- & -- & -- & \textbf{494} \\
\bottomrule
\end{longtable}
\end{center}

\section{Liite B: konkordanssit}

Alla kaikki tähän työhön poimitut katkelmat teemoittain. Jokaisen edessä on lehti,
numero ja sivu. Teksti on koneellisesti luettua eikä sitä ole korjattu; ainoastaan
välimerkkien edeltä on poistettu ylimääräiset välilyönnit.

\textbf{Katkelmia yhteensä: 707.}

\subsection{Venäläisyys ja neuvostovalta}
\subsubsection*{ryssä \normalfont\small(30 katkelmaa, 817 osumaa aineistossa)}
\kat{Toukomies}{05.02.1925}{P5}{Venäläiset varmasti vainuaisivat näissä puuhissa suomalaista propagandaa, eikä varoja sellaisten oppilaitosten ylläpitoon valtiolta saataisi ja yksityistietäkin, kun näiden koulujen tarkoituksesta ja tarpeellisuudesta ei julkisuudessakaan saattanut mainita, kertyisi vain verrattain pieni, riittämätön kannatusapu. Pääpaino tämäntapaisessa valistustoiminnassa oli kohdistettava siihen, että suomalainen heimo rajan kummallakin puolella pääsisi yhteisymmärrykseen ja tulisi elävään heimoustajuntaan, niin että tähän asti rajarahvaan tajunnassa säilynyt » ryssä » - ja » ruotsi » -käsite toisistaan havaittaisiin luonnottomaksi ja vääräksi ja sellaisena unohdettaisiin. Sallittanee minun tämän dots}
\kat{Toukomies}{01.04.1925}{P6}{Samoin Petsamon kylä ( Petzfby ) ja eteenpäin useampia muita kyliä, jotka ovat kaikki Valkeanmeren rannalla. Mitä suolaan tulee, niin tunnustaa Nousiaryssä, että ryssät ( karjalaiset ), jotka asuvat Valkeanmeren rannalla, lähettävät Ensimmäinen tiedustelu ( käännös on paikotellen tahallisen ruotsinmukaista, — ehkäpä se antaa luonteenomaisemman käsityksen keskustelusta ):}
\kat{Toukomies}{01.04.1925}{P7}{Käännämme tähän melkein sanasta sanaan alkukirjaa: Ensin kun karjalaiset lähtevät Käkisalmesta kohti Pohjanlahtea, kulkevat he niiden järvien kautta, jotka nyt tässä jälkeen mainitaan. Muistettakoon: Nousia Ryssä tunnustaa, että Pielisjärveltä pohjoiseen on joki nimeltään Lieksanjoki ( Liexäniåkj ), joka lähtee Valkeastamerestä; saman joen varsilla asuu 150 karjalaista. 60: —. » Parhaita ja kohottavimpia romaaneja, mitä meillä on viime aikoina kirjoitettu — todellinen suurromaani, joka on asetettava arvokkaim- pien historiallisten romaaniemme joukkoon. »}
\kat{Toukomies}{01.09.1925}{P4}{Pellavan, hampun ja nauriin siemenistä voisi pusertaa öljyä, mitä siihen saakka kalliilla hinnalla oli ostettu Venäjältä. Jos perustettaisiin karvari- ja säämyskäverstaita, niin ei » ryssä » näitäkään tavaroita enää yksin saisi maassa kaupitella. Se ei ole turhaa omillaankaan ollessa, se pakottaa sivistyksensä ainaiseen korottamiseen ja lujittamiseen ja siten vahvistaa meitä.}
\kat{Toukomies}{12.1925}{P16}{Kun parantamiseen on käytetty jotakin ainetta, esimerkiksi maitoa tai suolaa, ei niillä itsellään suinkaan ole katsottu olevanmitään lääkitsevää arvoa, ne ovat olleet ainoastaan välikappale, johon » loitsu » on » puhuttu », saneltu. — Hupaista lukemista ja erinomaisia pikakuvia Venäjän nykyisistä oloista saa se, joka hankkii Werner Söderström Osakeyhtiön suomeksi julkaiseman ruotsalaisen pilapiirtäjän Albert Engström i n teoksen » Ryssiä », jonka tekijä on itse tunnetulla taidollaan kuvittanut. Akseli G allen- Kallela: Helmikuu-kuvitelma}
\kat{Toukomies}{05.1926}{P11}{Ukko aina valveillansa, Seikkaeli ryssän kanssa, Miihkali ei enää laula,}
\kat{Toukomies}{05.1926}{P13}{Korpiselällä, Ilomantsissa, Pielis järvellä ei kansa enää pidä itseänsä venäläisenä, mutta käsitteet suomalaisesta kansallisuudesta ovat näil akin tienoilla vielä varsin sekavat ja heikot. Pari vuotta sitten vaativat esim. Korpiselän kunnan valtuuston sosialistiset jäsenet, että valtuusto karkoittaisi kuntansa alueelta kaikki pakolaiset pois, perustellen, että ne ovat ryssiä. Tämä kanta olisi yleensä omattava Suomen itärajoilla ja koko Suomen kansankin keskuudessa, mieli rakentaa sitä siltaa, joka heimot yhteen tois ja nostaisi niistä heikommat jaloilleen.}
\kat{Toukomies}{06.1926}{P9}{Varhaisina aikoina se on myös laskenut kahdella haaralla, mutta silloin toinen johti, vielä nytkin näkyviä juonteita myöten, Suomenveden pohjaan, Viipurin kaupungin kohdalle. Toinen erinomainen Venäjäteos on ruotsalaisen pilapiirtäjän Albert Engströmin kuvakirja ja matkateos » Ryssiä », joka viime syksynä Werner Söderström Osakeyhtiön kustannuksella ilmestyi suomeksi. On erityinen syy tutustua tähän Venäjäteokseen, joka kokemuksiltaan ja asioiltaan on nuorin, vasta viime syksynä kirjoitettu, ja joka myöskin suomenkielellämme on viimeksi ilmestynyt, tänä keväänä.}
\kat{Toukomies}{12.1926}{P16}{Pienestä pirtistä alkaa kuulua levotonta. liikettä. — » Yksi ryssä livisti, ota hänet pahus! Eivät kaikki ole vallanneet yhtä onnellisesti taloansa.}
\kat{Toukomies}{12.1926}{P17}{Urhoolliset urohot! » virkahtaa sissipäällikkö, — » voitto on tällä kertaa meidän ja kylämme vainolaisista vapaa. Vihdoin muuan heistä virkahtaa: » Älä suutu, hyvä veli, se ryssä, joka talostasi karkasi, Koko tämä kirottu korpi jäykkäniskaisine asukkaineen on vastassamme, mutta vielä me ne nujerramme, niiden jäykät niskat! »}
\kat{Toukomies}{12.1926}{P21}{Kustaa II Aadolfin ja Mikael Romanovin rajankävijät lyövät merkkinsä Variskiveen Laatokan itärannalla ja tämän rauhan on määrä tulla » ikuiseksi ». Hänen oli tapana sanoa: » En kuole, ennenkuin ryssä on maasta karkoitettu », ja hän saikin elää niin kauan! Tällaisen kohtalon varalle oli pappien tapana oppia jotakin käsityötä, sillä maanpakolaisuudessa ei saanut henkisellä työllä itseänsä elättää!}
\kat{Toukomies}{01.02.1927}{P23}{Lienee todella jotain ennenkuulumatonta, että ihmisluonto on toisen ihmisen ojennukseksi saattanut keksiä näin kyynillistä. Jo ennen olemme käsitelleet Neuvosto- Venäjää koskevia viimeaikaisia teoksia, m. m. ruotsalaisen pilapiirtäjän Albert Engströmin mainiota matkakirjaa » Ryssiä ». Jatkettiin yhteislaululla, jonka jälkeen M. Saaristo luennoi Suomen historiasta, selostaen tärkeimmät tapahtumat v. sta 1157 v-een 19 17. Sanarieskaa esitti T. Lahti.}
\kat{Toukomies}{01.04.1927}{P7}{kirotulvassaan ja ilmaa halkova piiska viuhui säestystään. Ainakaan ei ryssä ollut liiaksi sivistysahjoamme lämmittänyt. Ainakin hän hampaittensa välistä sähisi jotain — barbaarista.}
\kat{Toukomies}{01.04.1927}{P7}{Itse korkea esivallanedustaja teroitteli sisunsa kärkeä hypähtelemällä kukonaskelilla keskellä siltaa kädessä heilahtelevan piiskan yhäti kiihtyvän nopeuden osoittaessa hermostuneisuuden helteistä nousua. Olin päässyt näinä aikoina puheitten ja kirjojen välityksellä siihen käsitykseen, että ryssä kiiltonappisten virkaherrojensa muodossa sorti ja orjuutti meitä karjalaisia. Tässä nyt astellessani — hartioillani kasa haisevia ryssänpeitteitä — muistui mieleeni muuan monista haukkumasanoista, joilla äskettäin pristavi oli meitä pommittanut, nimitys " tsuhnat ".}
\kat{Toukomies}{01.04.1927}{P8}{Pysähdyimme kuuntelemaan, sillä miehen kasvojen epätoivoinen ilme vetosi tunteisiimme. — Syöka heinä! tiuskasi ryssä ja käännähtäen päin ojenti pyytäjälle jalkaa. Taas höhähti Miikkula ja kohotti ylös oikean kätensä, jossa keskisor-}
\kat{Toukomies}{01.04.1927}{P9}{Hän itse näet kunnioitti kaikkia hiukankin epämieluisilta tuntuvia lähimmäisiään tuolla nimellä, ja kun » Hunsvotti » oli kylän ristirahvaasta ylen outo ja naurettava sana, niin annettiin se nimeksi sille, joka oli sen kylään tuonutkin. Ymmärsikö » Hunsvotti » itse tiitinkejään, en ole varma, — ■ kylän pojat kuitenkin väittivät, että Hunsvotti oli yhtä hyvä ruotsalainen kuin Simana oli ollut ryssä, inttivätpä vielä, että tuo tiitinki hänen kädessään pakkasi olemaan usein ylösalasin. Kaunista Konnunkylää Salmin pitäjän ja Aunuksen rajalla.}
\kat{Toukomies}{01.06.1927}{P11}{Ja sieltä sankariajoilta on myös periytynyt heidän taipumuksensa runouteen ja lauluun. A ryssä tuli, orja; ksi teki meidät, vapautemme vei, kalleimpamme vei. Kirkko oli 'koristettu kymmenin dekoratiivisin pyhimyskuvin, jotka pappi itse oli maalannut.}
\kat{Toukomies}{01.06.1927}{P12}{Sellaisena oli hän koko karjalaisen vapaustaistelun vertauskuva ja kaunis vertauskuva. Ja sitten, veli, sai Karjalan kansa olla vuosisadan toisensa jälkeen ryssän sorron alaisena. Nyt on pyydys pyytämässä, ansaraudat vaanimassa, miehet niemien nenissä, saktkytät salmen suissa.}
\kat{Toukomies}{01.06.1927}{P12}{Toisin on nyt toimet siellä, nyt on kaikki karsahasti, poikkiteloin porras pantu, aita rautainen rakettu. Paljon meiltä hukkui venäläiseen mereen, paljon, paljon, suurin osa, ja ryssä on vienyt meiltä alueen toisensa jälkeen. Koko tämä juttu on niin uutta ja suurta, suurempaa kuin mitä hän tähän asti on jaksanut koskaan ajatella, käsittää, uneksia.}
\kat{Toukomies}{12.1927}{P27}{Matkan varrella, missä ikinä oli kyläpahanen, pistettiin aina kisat pystyyn, oli yö tahi päivä. Pahin, jota ryssät kauhistuivat, oli suomalainen puukko, jota asetta karj alaisetkin melkein poikkeuksetta kantoivat vyöllään. Uhtuan kylä on koria kylä, keskellä kylää mäki, talvella siellä varpuset laulaa, kesällä kukkuu käki.}
\kat{Suomen Heimo}{20.01.1928}{P5}{Sillit ei ole Suomen kansan elämässä vielä ollut suurempia ja kauniimpia aikoja, kuin ne, jolloin Suomen harmaa talonpoikaisarmeija kuin taikaiskusta syntyi, ja tuntiessaan itsensä tarpeeksi vahvaksi lopulliseen ratkaisuun, vastustamattomasti alkoi vyöryä etelää kohti musertaen kaikki esteet tieltään ja vapauttaen pitäjän toisensa, kaupungin toisensa jälkeen. 10 vuotta on kulunut Mitä, kun Suomen kansa tarttui aseisiin verin hankkiakseen itselleen ja isänmaalleen vapauden, 10 vuotta bn • vierähtänyt niistä riemukkaista päivistä, jolloin kansamme ryhtyi joulukuussa 1917 julistamaansa itsenäisyyttään toteu t t am a an, ja tämän kevään kuluessa tulee 10 vtiotta kuluneeksi siitä, kun viimeinen dots}
\kat{Suomen Heimo}{20.01.1928}{P8}{Avoimin silmin tappaa Karjalan yapausajatuksen alku, jonka lietsomiselle hain oli vannoutunut, oli hänelle „moraalinen rikos ", kuten hänen omat sanansa kuuluvat. Muutamia päiviä kului, ja lähestyi aika, jolloin Suomen oli täytettävä Tarton rauhan ehdot, joita ryssä puolestaan ei koskaan oletäyttänyt, tuskin aikonutkaan. Kaikille Karjalan ystäville ovat velvoittavana perintönä hänen ulkoasiainministerille juuri ennen kuolemaansa osoittaman avoimen kirjeen loppusanat: „Älkää luulko että ryhdyn tähän tekoon harkitsemattomasti.}
\kat{Suomen Heimo}{20.01.1928}{P8}{Ehken juuri talia teollani, joka itseasiassa on sovittava, täytän minulle antamanne määräyksen, sillä voi olla, että Repolan vapausaate nyt hetkeksi minun mukanani väsähtää — onko se sitten onneksi tai ei, sitä en tiedä. Tällöin ulkoministeriömme, jossa kokonaan oli unohdettu, ettei ryssä koskaan ole vielä lupauksiaan pitänyt, luuli olevansa velvollinen tekemään kaikkensa tätä estääkseen ja lähetti tammikuun 12 p:nä 1921 Sivénille radioteitse käskyn tästä. Niiden puolesta, joiden ääni ei pääse kuuluville julkisen puoluelehden palstoilla eikä puolue-elämän parnassolta, niiden puolesta kaikuu H. H. Sivénin teko vapauttavana: on vielä niitäkin, jotka osaavat taittua, mutta eivät taipua. "}
\kat{Suomen Heimo}{20.01.1928}{P13}{Tätä repivää työtä ei ole suoritettu vain sen aineksen taholta, jolle maamme turvallisuus ei ole minkään arvoinen, vaan ikävä kyllä maamme itsenäisyyden puoltajiksi lukeutuvienkin puolelta. t. k. 23 p:nä siitä, kun suomalaisen Etelä- Pohjanmaan suojeluskunnat kokoontuivat riisumaan ryssät aseista, mutta saatuaan ylemmältä taholta jyrkän käskyn odottaa, vastahakoisesti hajaantuivat; t. k. 21 p:nä siitä, kun Karjalan maaseudun suojeluskunnista kokoontui ensimäinen suurempi joukko, joka sitten seuraavana päivänä miehitti Viipurin, ollen kuitenkin jo sitä seuraavana yönä pakoitettu vetäytymään Venäjänsaareen;}
\kat{Suomen Heimo}{20.01.1928}{P17}{Arhippa aikoi iskeä pistimen haavoittuneen-rintaan, mutta hän esti. He suorittivat vimmatun taistelun, ryssät olivat jättäytyneet jälkeen. Hiljainen tuulenhenki kantoi sieraimiimme vastaniitettyjen heinien väkevää tuoksua.}
\kat{Suomen Heimo}{20.01.1928}{P18}{— Arhippa ja Triihyo! Kosto, ryssän kosto. — Petetty " V\textasciicircum{}ena parka}
\kat{Suomen Heimo}{20.01.1928}{P18}{Sitten vierähti kalpeiden huulien ylitse sydäntäsärkevä, katkonainen valitus: ' Kostonhalu antoi heille voimia, ja he kestivät, kestivät siinä, missä ryssä uupui. hyväksi... elämämme sen todistaa... ja kuolemamme...}
\kat{Suomen Heimo}{20.01.1928}{P19}{N:o 1 1928 Ryssän kosto! Vilho Helanen.}
\kat{Suomen Heimo}{20.01.1928}{P19}{Vilho Helanen. Arhippa-veljeni tapasi ryssän luoti. KIHLA-, vihkimä- ja jaloklvisormukset.}
\kat{Suomen Heimo}{15.03.1928}{P7}{Saisitte vain samanlaisia vastauksia, kuin minulle on annettu. Eivät he vihanneet ryssiä edes ampuessaan niitä kuin jäniksiä. Se on ainoa viha, jota he pystyvät sydämessään kantamaan.}
\subsubsection*{bolsevikki \normalfont\small(30 katkelmaa, 335 osumaa aineistossa)}
\kat{Toukomies}{01.05.1925}{P12}{— Johan Bojtr, Viimeinen viikinki, meriromaani. - Heikki Huhtamäki, Seikkailuissa bolshevikkien maassa. Teos on epäilemättä elämäkerrallisista kuvauksista ensimmäisiä.}
\kat{Toukomies}{01.10.1925}{P1}{Siitä lähtien hoitaa Repolaa Suomen hallitus » Itä-Karjalan toimituskunnan » avulla; vuotta myöhemmin, kuten tulemme näkemään, Porajärvi tulee saman hoidon alaiseksi. Bolshevikit olivat keränneet suuria voimia, ja aunukselaiset, jotka taistelivat vain pienin voimin ja vähin varoin, etupäässä vapaaehtoisen avustuksen turvissa, luonnollisesti hävisivät. Silloin karjalaiset yhdessä Suomesta maaliskuulla 1918 tulleen pienen sotilasretkikunnan avulla vapauttivat melkein koko pohjoisen Karjalan, mutta jo syksyyn mennessä, lokakuussa, oli Venäjä uudelleen herrana Vienassa.}
\kat{Toukomies}{01.10.1925}{P2}{Konstantinopoliin tuli nimitetyksi kreikkalaisen hengellisen seminaarin johtajaksi Galatassa ja sitten metropoliitaksi. Ennen kokouksen alkamista hyökkäsivät bolshevikit kuitenkin monelta taholta karjalais-alueelle — huono enne tietenkin. Tammi-, helmi- ja maaliskuu v. 1919 kuluivat jälleen toivon merkeissä, kansalliskokous saatiin toimeen ja pidettiin maalisk.}
\kat{Toukomies}{06.1926}{P17}{Nykyään on siis, kuten sanottu, tämä oppilaitos muuttunut suomenkieliseksi opettajaopistoksi, ja vaikka se toimiikin vieraassa hengessä ja vieraan esivallan alaisena, ei sen merkitys silti voine jäädä aivan olemattomaksi. Viimemainitussa pyydettiin, että liittolaisvallat miehittämällään ja vaikutusvaltansa alaisella ltä-Karj alan alueella takaisivat järjestyksen sekä hengen että omaisuuden turvan, niin etteivät punakaartilaiset ja bolshevikit saisi terroriserata, murhata ja ryöstää maan asukkaita. Tarton rauhansopimuksessa jäi venäj änpuoleinen Karjala kokonaan Neuvostovenäjän valtapiiriin; suomalaisen suojelusväen tuli poistua Repolasta ja Poraj ärveltäkin, jotka Venäjä lupasi liittää » dots}
\kat{Toukomies}{10.1926}{P13}{Kyprianon kylvö on hävitettävä lopullisesti, ettei siitä jää muuta muistoa kuin hautakivi Valamon luostarin kalmistoon, itse Valamonkin muuttuessa suomalaiseksi ja suomenkielistä uskonnollista valistusta palvelevaksi laitokseksi. Elokuun alussa levitettiin ulkomailla uutista, että Parisissa paraikaa käytäisiin neuvotteluja yhdeltä puolen neuvostohallituksen ja toiselta puolelta menshevikkien ja sosiaalivallankumouksellisten välillä tarkoituksella saada aikaan sopimus bolshevikkien ja menshevikkien ynnä sosiaalivallankumouksellisten välillä. silmät vezih täytty kai.}
\kat{Suomen Heimo}{20.02.1928}{P6}{Svinhufvud olisi ilman muuta suostunut tunnustuksen antamaan, mutta Enckell esitti sellaisen painavan seikan, ettei Suomikaan ollut vielä saanut itsenäisyyttöäi\textasciicircum{}unnustetuksijolloin Svinhuf- I vud myöntyi I hänen mielipi- I teeseensä. Viron itsenäisyys — se tuli näinä päivinä ehdottomaksi vaatimukseksi, Marraskuun 15 p:nä 1917 julistivat maapäivät Viron itsenäiseksi valtakunnaksi; mutta bolshevikit hajoittivat maapäivät ja vangitsivat eräitä maapäivien jäseniä, kuten Tönissonin ja Wilmsin. Tammikuus- Isa maapäivät I antoivat vv- I det valtuudet I Tönissonille, Martnalle, I Menningilie, |pu « tnll< Pii- I pii le, Kullille Ija Virgolle; I heidän oli |ni! i; irä Suomen kautta matkustaa dots}
\kat{Suomen Heimo}{20.05.1928}{P4}{Eikä lähimpien heimolaistemme suhde Suomeen ja suomalaisiin ole merkityksellinen ainoastaan näille heimolaisillemme « itselleen, vaan Suomen etu vaatii sekin mitä selvimmin, että suomalaisen sivistyspiirin rajat ulotetaan niin kauas kuin suomalaista asutusaluetta jatkuu. Olipa sitten kysymys alueen ja kansan liittämisestä Suomeen, johon se oikeudenmukaisesti kuuluu, tai Karjalan itsenäisyydestä, mikä meille suomalaisille olisi oikeastaan saman veroista, taikka sitten sen autonomian' toteuttamisesta, joka Tarton rauhassa Karjalalle on luvattu, vaikka se yhä vielä on vain bolshevikkien käsitteiden mukaista » vapautta *, — niin on olojen nykyisellään ollessa sula mahdottomuus, että karjalaiset dots}
\kat{Suomen Heimo}{15.06.1928}{P14}{Anna vielä hieman voita, muutama puuta vain pitäjää kohti, ja maitoa, marjoja, sieniä, leipää, särkiä, turkkia ja takkia, housuja Ja alusvaatteita, saappaita ja sääryksiä ja — ja — kaikkea anna ristiveli. Kun bolshevikki pistinniekka keväällä tuli Karjalaan, oli silla mukanaan laukun täydeltä erästä tavaraa, jota se kaupitteli maailman parhaana ja ainoana autuaaksi tekevänä aineena, eräänä ihmelääkkeenä, jonka oli keksinyt suurtoveri Lenin, vot Venäjällä oli myös kerran tehty keksintö. Siihen aikaan, kun bolshevikit v. 1920 valtasimt Vienan Karjalan, karjalaisten voimatta asettua vastarintaan, kun ryssien puolella oli myöskin joukko harhaanjohdettuja Karjalan miehiä, kirjoitti eräs dots}
\kat{Suomen Heimo}{15.06.1928}{P14}{Kun bolshevikki pistinniekka keväällä tuli Karjalaan, oli silla mukanaan laukun täydeltä erästä tavaraa, jota se kaupitteli maailman parhaana ja ainoana autuaaksi tekevänä aineena, eräänä ihmelääkkeenä, jonka oli keksinyt suurtoveri Lenin, vot Venäjällä oli myös kerran tehty keksintö. Siihen aikaan, kun bolshevikit v. 1920 valtasimt Vienan Karjalan, karjalaisten voimatta asettua vastarintaan, kun ryssien puolella oli myöskin joukko harhaanjohdettuja Karjalan miehiä, kirjoitti eräs karjalainen seuraavan kertomuksen, joka tässä karjalaisten ( kertomatavalle ja ajatukselle kuvaavana sijansa saakoon: Karjalan veli raapasi korvallistaan. ja mietti, mietti: „Hyvät ne olisivat maakuntaipäivät dots}
\kat{Suomen Heimo}{15.06.1928}{P16}{Ohjelmaan nähden viittaamme lehtemme » ♦Tietoja eri aloilta " -osastossa olevaan selostukseen. Ja yhtä selvää on, että bolshevikit juuri armeijallaan, toisin sanoen asein, aikovat I " " * * * — ' — » « ■ * " l Tuo matala pensas voi kasvaa komeaksi puuksi suotuisassa olosuhteissa.}
\kat{Suomen Heimo}{15.06.1928}{P16}{On ilmeistä, että enemmän selventää nykyisen Venäjän todellisia aikomuksia Leninin lause: „ Venäjän vallankumous on kuolemanvaarassa, ellei pian tapahdu vallankumouksia muissa maissa. " Senpä vuoksi onkin punainen armeija ainoa laitos Venäjällä, jonka hyväksi bolshevikit todella tehokkaasti ovat kuluneiden 10 vuoden aikana työskennelleet ja varoja uhranneet, ja juuri puna-armeijan suhteen huolehditaan parhaiten siitä, että henki taatusti on nykyisten vallanpitäjien tahdon mukainen. Menettelytapa, jota bolshevikit nykyään meilläkin noudattavat, on lyhyesti sanottuna se, että kiihoituksella ja rahalla koetetaan kaikkialla maailmassa synnyttää tyytymättömyyttä, lakkoja, levottomuuksia ja dots}
\kat{Suomen Heimo}{15.06.1928}{P16}{Senpä vuoksi onkin punainen armeija ainoa laitos Venäjällä, jonka hyväksi bolshevikit todella tehokkaasti ovat kuluneiden 10 vuoden aikana työskennelleet ja varoja uhranneet, ja juuri puna-armeijan suhteen huolehditaan parhaiten siitä, että henki taatusti on nykyisten vallanpitäjien tahdon mukainen. Menettelytapa, jota bolshevikit nykyään meilläkin noudattavat, on lyhyesti sanottuna se, että kiihoituksella ja rahalla koetetaan kaikkialla maailmassa synnyttää tyytymättömyyttä, lakkoja, levottomuuksia ja kapinoita, joihin sitten sopivan tilaisuuden tullen koko punainen armeija » » sekaantuu " marssien rajan yli » » työläisveljiä auttamaan ". Jonkun verran Idän salaperäisyyteen verhoutuvasta dots}
\kat{Suomen Heimo}{15.08.1928}{P15}{Kylmyys veljeskansaamme kohtaan on liian suuri, ja tämä itsekylläisyytemme on nyt kantamassa hedelmää: virolaiset kääntyvät länteen päin. Heinäkuussa 1923 havaitsivat bolshevikit hyväksi antaa Karjalalle nimellisen itsehallinnon, ja melkeinpä todellisenkin joillakin ajoilla. Sillä, että edellä olemme puhuneet moitteen sanoja virolaisten veljiemme nykyisestä » orientoitumisesta *, emme lainkaan tahdo puhdistaa itseämme.}
\kat{Suomen Heimo}{18.11.1928}{P18}{Ne on kuitenkin jo Tarton rauhassa tehty aseettomiksi, joten ei muuta kuin nousta maihin vain. Bolshevikit ovat jälleen mitä törkeimällä tavalla loukanneet molempien naapurimaiden välillä voimassaolevia rajasopimuksia. Samoihin aikoihin tunkeutui ryssien koululaiva liian lähelle Lavansaaren. kivikkoista rannikkoa, jääden karille.}
\kat{Suomen Heimo}{18.11.1928}{P18}{Ryssien toistuvista rauhansäilyttämisvakuutteluista huolimatta emme me niihin luota, enempää kuin ryssän puheisiin muutenkaan. — 1 p:nä syyskuuta ilmestyi suuri bolshevikkien laivasto-osasto vedenalaisineen aluevesillemme Lavansaaren rannikolle. Tehtäköön ryssien äskeinen luvaton vierailu huvipursilla rannikoillamme miten viattomaksi tahansa, ei se tosiasiassa sitä kuitenkaan ole.}
\kat{Suomen Heimo}{18.11.1928}{P18}{Kuten tiettyä, nousi viime elokuun 29 p:nä 35 henkilöä käsittävä bolshevikkiseurue ilman passia, ilman lupaa röyhkeästi valokuvauskoi>eineen Suursaaren Suurkylän satamassa komeista huvialuksistaan maihin. Bolshevikkien kaikki toimenpiteet, niin meitä suomalaisia, kuin muitakin kansoja kohtaan, tähtäävät, milloin häikäilemättömän avoimesti, milloin salaisemmin sotaan, haaveilemaansa maailmanvallankumoukseen. Pari päivää myöhemmin kaarteli ryssien tähystäjäkone Koiviston rannikolla, jatkaen sitten tiedusteluaan rannikkoa pitkin, hyvin matalalla lentäen Inoon, josta Terijoen kautta poistui Rajajoen. yli.}
\kat{Suomen Heimo}{18.11.1928}{P19}{Tämän ovat niin monasti ja raskaasti kokeneet esi-isämme ja yhtä kalliiksi käy se meillekin, jos emme tahdo oppia tuntemaan vihollistamme sellaisena, mikä se todellisuudessa on. Jos suomalainen kalastaja Suomenlahdella ( sikäli kun siellä suomalaisilla enää kalastusmahdollisuuksia on ) erehdyksessä tai myrskyn ajamana jäälautalla joutuu aluevesirajamme ulkopuolelle, ottavat bolshevikit hänet kiinni, takavarikoivat veneen ja pyydykset, ja kulettavat hänet kuukausimääriksi kitumaan haiseviin vankiloihinsa. Jokainen jätti jälkeensä jotain rakasta,}
\kat{Suomen Heimo}{18.12.1928}{P14}{aivan liian pessimistiseen sävyyn laadittu. tässä kuussa bolshevikkien hyökkäyksestä Viroon. hyvinkin läheinen tulevaisuus osottava.}
\kat{Suomen Heimo}{18.12.1928}{P14}{REMINGTON POiR- dSP\textasciicircum{}B\textasciicircum{}\textasciicircum{}\textasciicircum{}C / / x' TABLE vastaa kaikkia niitä vaatimuksia, mitä voi- IVylfiy < daan asettaa yksityiskäyttöön sopivalle kirjoitusko- Ä\textasciicircum{}r\textasciicircum{}äÉÉ jJfl / i\textasciicircum{}i / / / neelle. Näin pääsivät bolshevikit, harvalukuisten virolaisten joukkojen voimatta sitä estää, joulukuun kuluessa tunkeutumaan syvälle maahan, vallaten puolet Viron ja uhaten jo Tallinnaakin. Pipping näkyy pitävän yliopiston suomalaisia professoreja suurina ruotsalaisuu-, den sortajina, meidän tahollamme taas on totuttu pitämään heitä suomalaisuuden oikeutettujen vaatimusten vastustajina.}
\kat{Toukomies}{01.04.1928}{P15}{Kiannon lennokas ja omalaatuinen tyyli antaa näille vapaussotamuistelmille erikoisen tenhon ja elävyyden. — Bolshevikkien pyrkimyk " sistä saimme verisen kokemuksen kymmenen vuotta sitten täällä Suomessa ja Karjalassa. — Sotilashuumorilla on meillä vanhastaan vanhat perustat, mutta nyt ilmestyy nykypäivän kasarmijuttuja kirjaksi koottuna.}
\kat{Toukomies}{01.04.1928}{P15}{Miltä olot näyttivät Suomenlahden etelärannalla, siitä antaa tunnetun venäläisen kirjailijan teos erittäin luonteenomaisen ja opettavan kuvauksen. Tekijä oli itse mukana eräässä balttilaisessa suurkaupungissa, joka onnettomuudekseen joutui bolshevikkien länttä kohti suunnatun hyökkäyksen uhriksi. — Katri Ingman, jonka viimevuotinen esikoisteos » Virran viemä » sai erinomaisen menestyksen, on nyt julkaissut uuden romaanin » Lahjakas tyttö ».}
\kat{Toukomies}{12.1928}{P16}{RETTY JÄRNEFELT-RAUANHEIMO: » Vierailla verä- jillä. » L. VATJALAINEN: » Vihreä lotta bolshevikkien maassa. » Kuv. lIVO HÄRKÖNEN: » Runon hirveä hiihtämässä. »}
\kat{Suomen Heimo}{31.01.1929}{P22}{Koko Suomi ja huomattava osa muutakin maailmaa seurasi tätä suurimmaksi, osaksi kynällä ja musteella käytyä sotaa. » Aa bolshevikit, he. ovat venäläisiä, elävätkin venäläisten tavalla k huudahti muuan munkki vilpittömän ihailun väre äänessään. Se näyttää kylläkin sangen toivottomalta vanhan ajanlaskun kannattajain näkökulmalta katsottuna, mutta näillä on kyllä sitkeyttä ja toivoa odottaa.}
\kat{Suomen Heimo}{15.04.1929}{P22}{111111111111111 l HUONEKALUJEN, KOULU- KALUSTOJEN JA LELUJEN ERIKOISLIIKE. P O H J OLA Rahastot Smk. Tuntuu siltä, kuin muukalaiset bolshevikit tarkoituksellisesti tahtoisivat siveettömyydellä Karjalan kansan kokonaan hävittää. 52,000,000: — Palo-, meri-, auto- ylm. vakuutuksia i • Tapaturmavakuutuksia KULLERVO Rahastot ''Smk.}
\kat{Suomen Heimo}{25.05.1929}{P16}{Mietteisiini vaipuneena ja tupakan sauhuja imiessäni löi seinäkello seitsemän kertaa, osoittaen seitsemää illalla. Olihan syytäkin iloon, sillä äskettäin olimime juuri vallanneet bolshevikit sekä heidän varastonsa, sikäli kuin niitä jäljellä oli. Talon isännän toimesta saapuikin \{ pian metsänhoitajia y. m. herroja 4 tai 5 miestä, niiden mukana entinen kaupungin päällikkökin.}
\kat{Suomen Heimo}{15.10.1929}{P16}{Ennen keisarihallitus „kutsui ", kutsui " karjalaisia Venäjän armeijaan keisarillisen isänmaan puolustamisen nimessä, nyt heitä kutsutaan jälleen Venäjän armeijaan, punaisen isänmaan palvelukseen, puolustamaan ja lujittamaan Venäjän kommunistista imperialismia, mikä on omiaan jouduttamaan karjalaisten olemassaolon tuhoa. Bolshevikit ryöstivät kaikki ruokatarpeet, takavarikoivat vaatetavaroita, tappoivat ( karhujen laskuun ) lampaita, hävittivät rakennuksia polttopuiksi j. n. e. — Ja siis: „talous, talous meni rappiotilaan ", josta se ei ole vieläkään noussut, vaikka kansalaissodasta on kulunut jo alun toistakymmentä vuotta. Pääkonttori: Fabianinkatu 15. o Haarakonttori: Museokatu 3.}
\kat{Suomen Heimo}{13.03.1930}{P6}{Ammattiliitot ovat toiminnassaan suhteellisen itsenäisiä ja saattavat, jos niin haluavat, noudattaa järjestön keskuselimestä varsin riippumatonta politiikkaa. Bolshevikit ovat senvuoksi tarttuneet yhteisrintaman oljenkorteen ja puuhaavat nykyään Köpenhaminan sopimuksen pohjalla uutta internationalea, johon molem- mat edellämainitut yhtyisivät ja jossa Profintern suuren jäsenmääränsä avulla — Venäjällähän voidaan tarpeen vaatiessa merkitä miltei koko kansa ammatillisten organien jäsenluetteloihin — heti saisi yliotteen.}
\kat{Suomen Heimo}{15.06.1930}{P20}{Tulevaisuus näyttää, miten käy, mutta nyt on elämä kovin sekavaa ja rauhatonta. Bolshevikit ovat hävittäneet kaikki entiset kirjastot, eivätkä kykene uusia hommaamaan tilalle. Kun talonpojat vapaaehtoisesti eivät tahtoneet liittyä kollektiiveihin, ruvettiin käyttämään pakkotoimenpiteitä.}
\kat{Suomen Heimo}{22.10.1930}{P18}{Sillä vierasta oli Suomenmaa, vaikkakin heimoveljien maa; eihän täällä pakolaisrukkaa pidetty vaivoin siedettyä loista parempana. Sanotaan etta bolshevikit. vangitsevat heti, olipa tulija kuka tahansa, ja raahaavat I Ytroskoihin tshekan kuulusteltavaksi. lupaikkojen hankkimista kesäajaksi veljesmaan ylioppilaille, jotka opiskelevat esim. kansantalous-, maatalousja lääketiedettä.}
\kat{Suomen Heimo}{22.10.1930}{P19}{yleisessä kokouksessa 12. 10. 1930 hyväksyttiin yksimielisesti A. K. S:n päämäärää ja ohjelmaa käsittelevä päätöslauselma, jonka julkaisemme tämän numeron otsikkosivulla. Siinä tämän toivioretkeläisen koko varustus. |a järjen kannalta katsoen olikin suotta enempää viedä, bolshevikit kumminkin olisivat " yliteiskunnallistuttaneet " runsaammat tuomiset. Täytyy kuitenkin ottaa huomioon, että poliisit samoinkuin siviilipukuiset etsivät, joita asemalla aivan vilisi, estivät jatkuvat vihellykset pidättämällä useita viheltäjiä, ja että hurraahuutajien enemmistö näytti olevan Siltasaarelta päin kotoisin, joukossa joitakin ruotsalaisia ylioppilaita!}
\subsubsection*{Stalin \normalfont\small(25 katkelmaa, 402 osumaa aineistossa)}
\kat{Toukomies}{01.04.1925}{P11}{Kuvamme esittävät talon päärakennusta ja sen mahtavaa emäntää sekä toimeliasta tytärtä, Tatjanaa, jota viimemainittua noihin aikoihin pidettiin seutukunnan rikkaimpana perijättärenä. Puoluesihteeri Stalin on pitänyt Moskovassa puheen, jossa kehottaa kääntämään » kasvot kylään päin », koska » kaikki muut liittolaiset paitsi talonpojat ovat pettäneet kommunistijohtajat ». Punaisen Karjalan » ilmoituksen mukaan on Itä-Karjalan komisaarineuvosto päättänyt laajentaa metsänhakkuun alkuperäistä suunnitelmaa myöntämällä » Severoles-trustille lisäoikeuden hakata tänä talvena 149,000 puuta.}
\kat{Toukomies}{01.03.1929}{P19}{on alenemassa. » VENÄJÄN HIRMU: STALIN. RUOTSIN SUOMALAISTEN LUKUMÄÄRÄ.}
\kat{Toukomies}{01.03.1929}{P19}{RUOTSIN SUOMALAISTEN LUKUMÄÄRÄ. Perheensä piirissä on Stalin täydellinen tyranni. Hän elää aivan yksin kaksihuoneisessa asunnossaan Kremlissä.}
\kat{Toukomies}{01.03.1929}{P19}{Hänen isänsä, ammatiltaan suutari, tahtoi pojasta tehdä papin. Stalin on keskikokoinen, kasvot ilmeettömät, koko tyyppi itämainen. Varmalta näyttää hänen osallisuutensa Tiflisiu kuuluisaan kaappaukseen v. 1907.}
\kat{Toukomies}{01.03.1929}{P19}{Kukaan ei nykyisellä Venäjällä tiedä hänen ajatuksiaan, suunnitelmiaan ja aikeitaan. Jossip Vissarionovitsh Dshugashvili, myöhemmin venäläistytetty Stalin, syntyi Georgiassa 1879. Tiflisistä lähetettiin kahdessa postivaunussa suuria arvolähetyksiä Pietariin sotilasvartioston saattamina.}
\kat{Toukomies}{01.03.1929}{P19}{Vain yksi ainoa ajatus vallitsee häntä kokonaan: suunnaton vallanhimo, jolle hän on omistanut koko elämänsä. Bajanov antaa samalla Stalinin hahmoa kuvatessaan diktaattorista muutamia mielenkiintoisia elämäkerrallisia tietoja. torioitsijoita ja esiintyi maailmansodan aikana upseeripuvussa Suomessakin koettaen täällä saada kapinaa aikaan.}
\kat{Toukomies}{01.03.1929}{P19}{Tsaarinvallan poliisille aiheutui hänestä paljon huolta, sillä hän osasi aina kätkeä hyvin vallankumouksellisen toimintansa, niin että todisteita oli vaikea saada. Myöhemmin todettiin tämän upseerin olleen Stalinin uskollinen apulainen Ter Petrosian, toiselta nimeltä Kamo, joka nykyisin 011 bolshevismin tunnetuimpia his- V. 1920 oli suomalaisia Vesterbottenissa 27, Jemtlannissa 20 sekä muualla Ruotsissa 1,172. Ylläolevat 29,028 Norrbottenin suomalaista jaetaan asuinpaikkansa mukaan seuraavasti:}
\kat{Toukomies}{01.03.1929}{P19}{V. 1920 oli suomalaisia Vesterbottenissa 27, Jemtlannissa 20 sekä muualla Ruotsissa 1,172. Ylläolevat 29,028 Norrbottenin suomalaista jaetaan asuinpaikkansa mukaan seuraavasti: Samalla on tämän Stalinin ja Trotskin välien loppuselvittelyn yhteydessä Stalinin oma hahmokin siirtynyt enemmän etualalle vissistä salaperäisyydestä, mihin se aina on verhoutunut. Oikealla hetkellä näin ollen on Stalinin entinen yksityissihteeri Bajanov julkaissut venäläisessä monarkistilehdessä » Vozrozdeniessa » Stalinista luonnekuvan, joka jossain määrin selvittää tämän arvoituksellisen miehen olemusta.}
\kat{Toukomies}{01.03.1929}{P19}{V. 1920 oli suomalaisia Vesterbottenissa 27, Jemtlannissa 20 sekä muualla Ruotsissa 1,172. Ylläolevat 29,028 Norrbottenin suomalaista jaetaan asuinpaikkansa mukaan seuraavasti: Samalla on tämän Stalinin ja Trotskin välien loppuselvittelyn yhteydessä Stalinin oma hahmokin siirtynyt enemmän etualalle vissistä salaperäisyydestä, mihin se aina on verhoutunut. Oikealla hetkellä näin ollen on Stalinin entinen yksityissihteeri Bajanov julkaissut venäläisessä monarkistilehdessä » Vozrozdeniessa » Stalinista luonnekuvan, joka jossain määrin selvittää tämän arvoituksellisen miehen olemusta.}
\kat{Toukomies}{01.03.1929}{P19}{Samalla on tämän Stalinin ja Trotskin välien loppuselvittelyn yhteydessä Stalinin oma hahmokin siirtynyt enemmän etualalle vissistä salaperäisyydestä, mihin se aina on verhoutunut. Oikealla hetkellä näin ollen on Stalinin entinen yksityissihteeri Bajanov julkaissut venäläisessä monarkistilehdessä » Vozrozdeniessa » Stalinista luonnekuvan, joka jossain määrin selvittää tämän arvoituksellisen miehen olemusta. Hän tuskin tuntee tärkeimmät poliittiset kysymykset, taloudellisista ja rahallisista kysymyksistä hänellä ei ole juuri aavistustakaan, hän ei osaa mitään vierasta kieltä eikä tunne venäläistä kirjallisuuttakaan liioin.}
\kat{Toukomies}{01.03.1929}{P19}{Hän tuskin tuntee tärkeimmät poliittiset kysymykset, taloudellisista ja rahallisista kysymyksistä hänellä ei ole juuri aavistustakaan, hän ei osaa mitään vierasta kieltä eikä tunne venäläistä kirjallisuuttakaan liioin. Leninin ehdotuksesta nimitettiin hänet 19 13 vallankumouksellisen keskuskomitean jäseneksi ja kun Lenin 191 7 palasi Venäjälle sai Stalin tehtäväkseen kukistaa Kerenskin Pietarissa, mikä tehtävä hänelle täydellisesti onnistuikin. Ylläolevat Ruotsin tilaston viralliset luvut antavat vahvan lisävalaistuksen Nordisk Familjebokin kuvaukselle Norrbottenin läänin kielioloista: » Suomenkieli erinäisillä seuduilla, ainakin aivan viime aikaan saakka, on voittanut alaa ruotsinkielen dots}
\kat{Toukomies}{01.03.1929}{P19}{Ylläolevat Ruotsin tilaston viralliset luvut antavat vahvan lisävalaistuksen Nordisk Familjebokin kuvaukselle Norrbottenin läänin kielioloista: » Suomenkieli erinäisillä seuduilla, ainakin aivan viime aikaan saakka, on voittanut alaa ruotsinkielen kustannuksella. Trotskin kukistumisen ja nyt viimeksi maasta karkoituksen jälkeen on Kremlin nykyinen diktaattori Stalin vapautunut vaarallisesta kilpailijastaan, joka Venäjän kommunistipuolueen vasemmistopiireissä vahvan kannatuksen omaten kukistettunakin oli vielä vakavana uhkana nykyisille vallanpitäjille. Suomalaiset asuvat ' pääasiallisesti Tornion voutikunnassa, missä v. 1900 yli 3 / 4 väestöstä oli suomenkielistä, sekä sen lisäksi dots}
\kat{Toukomies}{01.03.1929}{P19}{Tietenkin on myös ruotsia puhuva väestö kasvanut suurten malmikaivantojen, valtavan rautatieliikenteen, Bodenin linnoituksen y. m. vaikutuksesta, mutta suomalaisen väestön lisäysprosentti v:sta 1870 v:een 1920 on 107. 1, kun vastaava prosenttiluku Suomen suomenkieliseen väestöön nähden ( v. 1870 1,640,000 ja 1920 2,754,000 ) on 68 pros. Poika kyllästyi kumminkin jo 20 vuoden ikäisenä kirkollisen uran vaikeuksiin ja antautui sanomalehtimieheksi työskennellen useammissa sosialidemokraattisissa sanomalehdissä kuten » Borba Proletariata », » Dro » ja » Bakinski Babotshi » Kaukaasiassa hän tutustui muutamiin myöhempiin hallitustovereihinsa, kuten sota-asiainkomissaariin Frunseen, dots}
\kat{Suomen Heimo}{26.04.1930}{P24}{Ken viljelee paremmin peltonsa ja kasvattaa enemmän viljaa, hänet nimitetään kulakiksi ja ruvetaan kaikin tavoin vainoamaan, sillä ne pelkäävät kapitalismin kehitystä maa- Kun Stalin julkaisi tunnetun kirjelmänsä, jossa kielletään käyttämästä pakkotoimenpiteitä kollektivisoimispuuhissa, alkoivat Inkerinkin talonpojat joukottain karata kolhoseista, joista monet kuolivat aivan alkuunsa. Kulakit ja papit luokkana on hävitettävä, silloin sortuu viimeinen vallankumouksen jarru ja avautuu sosialistinen ihanneyhteiskunta, jossa vallitsee rauha, rakkaus ja ihmisillä hyvä tahto, selittävät kommunistiteoreetikot.}
\kat{Suomen Heimo}{15.06.1930}{P20}{Uskonnollisia ynnä muita suomalaisia kirjoja ei ole saatavana, niitä lainaillaan talosta taloon, joten monesta kirjasta ovat enää kannet jäljellä. Kun julkaistiin Stalinin tunnettu kirjelmä, jolla kielletään käyttämästä pakkotoimenpiteitä kollektiiveja perustettaessa, ovat melkein kaikki yhteistaloudet jo hajonneet. Olipa teidän yhteiskuntanne millainen tahansa, ei minulla ole muuta kuin hyvää muistettavana niiltä ajoilta, jolloin olin siellä puolen punaisen riu'un.}
\kat{Toukomies}{01.04.1930}{P18}{35: —. Stalin peräytyy. Ranskankielestä kääntänyt Jalmari Hahl.}
\kat{Toukomies}{01.04.1930}{P18}{Karjala-Seuran edustajat, jotka myöhemmin kutsuivat avukseen Karjalakerhojen keskusliiton edustajan. Senvuoksi Stalin on peräytynyt ja alkanut uuden suunnan t alonpoikaispoli tilkkaan nähden. Kauan aikaa oli Wecksellin tehdas ainoa, kehittäen vähäpätöisen ja yksinkertaisen käsityönäpertelyn suurteollisuudeksi.}
\kat{Toukomies}{01.04.1930}{P18}{— Heimo- ja pakolaisjärjestöinä toimivat nykyisin maassamme Akateeminen Karjala-Seura, Itä-Karjalan komitea, Karjalakerhojen keskusliitto ja Karjalan Sivistysseura, mistä kaksi viimemainittua, nuorin ja vanhin, ovat varsinaisesti karjalaisia, toiset suomalaisten muodostamia. Riiasta on ilmoitettu ( 17 / 3 ) että, Kalinin ja Voroshilov ovat selittäneet Stalinin nykyisen agraaripolitiikan ja uskonvainojen johtavan epäilemättä talonpoikaiskapinoihin, samalla kun punainen armeija ei ole luotettava, koska siilien kuuluu suuri määrä talonpoikien poikia. Kirjassa annetaan käytännöllisiä ohjeita akkunaverhojen valitsemista, sommittelua ja ripustamista varten y.m. Kirjanen on painettu hyvälle dots}
\kat{Suomen Heimo}{24.01.1931}{P24}{— Eivätkä suomalaiset kansakoululapset ole läheskään niitä ainoita uhrilampaita, joita » luonnollinen kehitys " vaatii. Voidan siis täydellä syyllä sanoa, että Stalinin diktatuuri ja kommunistinen huliganismi Venäjällä ovat suurelta osalta Karjalan tukkipuiden varassa. Se ei tietysti mallaa Helsingin herroille, mutta ettehän te suomalaiset talonpojat kysyneet heiltä lupaa, kun vaaditte Aleksanteri II:n patsaankin poistettavaksi.}
\kat{Suomen Heimo}{08.08.1931}{P11}{Ja jos lapset kulkevatkin niukoin yksityisin varoin ylläpidetyssä suomalaisessa koulurähjässä, täytyy heillä alati olla henkivartio matkassaan ja heidän isänsä ja äitinsä saavat käydä ostamassa särpimensä usein matkojen takaa, kun ruotsalaiset eivät heille anna paljolla rahallakaan. kenties sitäkin helpompaa, kun Stalin on kääntynyt kapitalistisille poluille, maksaa työläisille palkan työnsuorituksen mukaan, tunnustaa sivistyneistön merkityksen ja johtavan työn arvon, pyhittää sunnuntain lepopäivänä j. n. e. Johan puoluesihteeri sanoi, " että se on vain tavallista rajaseutupolitiikkaa ". V\textasciicircum{}aOllßiiK' Pikemmin kiitettävistä järkisyistä kuin sen vuoksi, <'tiii neitosiltamme BrtÉßÉ\textasciicircum{}sf\textasciicircum{}ft muka dots}
\kat{Suomen Heimo}{20.12.1931}{P29}{Agapetus: AKKOJEN KAUHUNA eli tuumailuja tuuli- as in takaa. Essad bey: STALIN. Elämäkerta kuin paras jännitys- romaani. Henrik Vaanila: TIE TAIVAAN PORTILLE. Hiljaisten hetkien tarina.}
\kat{Suomen Heimo}{20.12.1931}{P34}{— Hinta 25: —, sid. Coudenhove-Kalergi: Stalin \& Kumpp. Yrjö Rauanheimo: Merimies ja Albatrossi.}
\kat{Suomen Heimo}{20.12.1931}{P34}{Asevelvollisuusikään varttunalle nuorisolle annetaan neuvoja ja ennakko-opetusta niissä kysymyksissä, joiden tunteminen ja tietäminen on palvelukseen astuttaessa hyvin tärkeä. Neuvosto- Venäjä eli " toiminimi Stalin \& Kumpp. " on valtava uhka kaikille Euroopan maille, eikä mikään niistä pysty sitä yksinään vastustamaan. Tämä suuri murhenäytelmä on antanut Johannes Karman nimellä esiintyvälle tekijälle aiheen historialliseen draamaan, jossa on suurpiirteinen henkilökuvaus ja mahtava kohtalon tuntu.}
\kat{Suomen Heimo}{22.02.1932}{P17}{Kirj. Stalin. Nuijasota.}
\kat{Suomen Heimo}{22.02.1932}{P17}{15 Stalin \& Kumpp. 68 s. 12 mk.}
\subsubsection*{GPU/tsheka \normalfont\small(25 katkelmaa, 99 osumaa aineistossa)}
\kat{Toukomies}{01.09.1925}{P22}{Tshekkoslovakia, kuvitettu esitys, kirjoittanut tohtori V. J. Mansikka. Tsheka, kirjoittanut Georg Popoff, bolshevistisen Venäjän tunnetuimpia kirjailijoita. Parembie laulahtuksie kuin tuloo kevätsiäl, jo tshirkun tshirkutuksie}
\kat{Toukomies}{01.09.1925}{P22}{Laajan ja runsaan kirjallisuussadon on suomalainen kirjanpainatustoiminta antanut kevätkauden myöhäisemmälläkin osalla ja kesäkaudella. Tekijän kehoituksesta on eras suomalainen asiantuntija liittänyt teokseen esityksen, jossa kuvataan tshekan työtä Suomessa. Vain helkytykses tyhjäs ei sendäh hergie suu, kuin kylän kylmän ryhjäs ei kuulu laulu muu.}
\kat{Toukomies}{12.1925}{P12}{— Kuka meistä 50 vuotta takaisinpäin olisi uskonut — virkkoi ministeri Lampén puhaltaen valtavan sauhun suoraan minun silmilleni — että tästä Vienan Karjalastakin repaleesta paisuisi itsenäinen valtio — se oli sellaista hämärää hosumista koko Pro Karelia — heillähän oli käynnissä kuin mikä pontikkatehdas ja ryssäläinen yliopistokin Muurmannilla... kaikki narut johtivat Moskovaan ja Trotskigradiin... — Mutta niinkuin me itse vuonna 1977 Europan kulttuurikansojen pistimien ja kuularuiskujen avulla yhtäkkiä muutuimme vaivaisesta tasavallasta korskeaksi kuningaskunnaksi, samoin Vienan Karjala vuonna 1981 Jaappanin ja Kiinan avulla vihdoinkin vapautui Venäjän Tshekan pirullisista kynsistä ja kun dots}
\kat{Toukomies}{06.1926}{P10}{Vihdoin on mainittava näistä kaikista aikaisimmin ilmestynyt, alussa mainittu englantilaisen Paul Dukesin laajahko teos » Punaisen hämärän maa ». Kustannusosakeyhtiö on aikaisemmin julkaissut kaksi ansiokasta venäjäkirjaa: Georg Popoffin » Neuvostotähden alla » ja » Tsheka ». Nishni-Novgorodin metropoliitta Sergei on tuominnut tämän neuvoston toiminnan kokonaan, määräten sen jäsenille katumuksenteon ja eroittanut heidät toistaiseksi kirkosta.}
\kat{Toukomies}{06.1926}{P10}{Mutta ne mietteet, mitä hän antaa koko Venäjän nykyisen vallan perusteista ja olemuksesta, on kenties iskevintä, mitä näistä oloista on sanottu. » Toinen kirja, » Tsheka », on tärisyttelevä, räikeä kuvaus Venäjän nykyisestä salaisesta poliisista, kirjoitustavaltaan ei niin mieltäkiinnittävä kuin edellinen. vaikka tekijä liikkuu vasta vuosina 1918 — 1919; vangitsemisia tapahtuu, vakoilu kiristyy ja kiristyy ja yhä useampia sivistyneistöstä saa elämällään maksaa uuden onnentilan.}
\kat{Toukomies}{08.1926}{P10}{Jalommat lehtipuut eivät muodosta suuria yhtenäisiä piirejään, mutta pysyttelevät yleensä omien levenemisrajojensa sisäpuolella. Heinäkuun 20 päivänä kuoli Moskovassa Neuvostoliiton korkeimman taloudellisen neuvoston puheenjohtaja, entinen tshekan päällikkö. Usein elävänäkin, jota ihmetteleväisten kummastukseksi kuletetaan kaupungin kaduilla, — toki sentään vain pieni penikka.}
\kat{Toukomies}{01.06.1927}{P19}{' Luotettavalta taholta saadun tiedon mukaan ön toukokuulla Neuvosto-Karjalan rajaseutukunnissa, etupäässä Repolassa, Poraj arvella ja Kontokissa toimeenpantu lukuisia pidätyksiä. Heitä odottavasta kohtalosta ei kotipaikkakunnilla tiedetä mitään, mutta yleensä luullaan, että nämä tshekan epäilyksen ja vainon alaisiksi joutuneet miehet karkoitetaan rajaseudulta epäluotettavana aineksena. Häntä on epäilty, samoinkuin toisiakin syytettyjä, » luvattomasta yhteydestä Suomen kanssa ".}
\kat{Suomen Heimo}{20.05.1928}{P15}{tuleminen on hyvän sattuman varassa; vain köyhälistöluokalla on täysi oikeus, sikäli kuin se ei ole mielipiteiltään hallitusjärjestelmää vastaan, tulla valituksi sekä valita. Piirineuvos-to on riippuvainen niistä osviitoista, joita paikalliset kommunistiset puolue järj estot ja valtiollisen poliisin — tshekan — hallintoelimet niille sanelevat. Jokaisen piirin alueella toimii paikkakunnan kommunistisen puoluesolun ( -osaston ) n.k. aktiivi, joka juuri valvoo piirineuvoston toimintaa ja antaa sille ohjeita.}
\kat{Suomen Heimo}{20.05.1928}{P18}{Pietarissa painetaan Inkeriä varten » Vapaus « -nimistä lehteä, Nuorta Kaartia, Punaista kyntäjää ja eräitä muita tilapäisjulkaisuja. Aikaisemmin jo mainitsimme, että maaseudulla toimii erikoinen puolueaktiivi ja tsheka, jotka valvovat toiminnan kulkua kommunismieliseen suuntaan. Tavarain tasajakoa ei tosin enää sotakommunismin aikakauden päätyttyä v. 1922 jälkeen ole juuri tapahtunut, mistä syystä kansan taloudellisessa hyvinvoinnissa on huomattavissa noususuunta.}
\kat{Suomen Heimo}{15.06.1928}{P19}{Puolaa ja Romaniaa vastaan Ukrainassa: 17 divisioonaa ja 2 ratsuväluarmeijakuntaa ( 2 ratsuväkidivisioonaa kummassakin ). Lisäksi on Venäjällä suuret joukko-osastot varsinaisen armeijan ulkopuolella olevaa santarmistoa, nimittäin taatusti luotettavat n.s. tshekan joukot. 1 ) Jalkaväkirykmentti on kokoonpantu 3 pataljoonasta ja yhdestä 2-patterisesta patteristosta 7.6 cm. kenttäkanuunoita sekä tiedustelu-, yhteys- ja kaasuosastoista.}
\kat{Suomen Heimo}{20.12.1929}{P23}{Terveydellisiä näkökohtia niitä rakennettaessa ei ole otettu läheskään riittävästi huomioon, mikä ilmeisesti on johtunut osaksi köyhyys-syistä ja osaksi myöskin rakennustaidon puutteesta johtuvasta saamattomuudesta. Karjalan metsärajoillahan komintern palkkaa kiinoitta jia ulkomaille, toimeenpanee mielenosoituksia, järjestää lakkoja eri maissa, rakentaa punaisia lentokoneita, vahvistaa punaista armeijaa, lisää tshekan valtaa j. n. e. Ensi talveksi on päätetty haalia metsätöitä varten työvoimaa Karjalan ulkopuolelta Tverin, Smolenskin ja Pensan lääneistä, Valko- Venäjältä ja Keskisen Volgan alueelta yhteensä 22,000 hevosta ja 49,000 työläistä, mikä teko osoittaa, että Neuvostoliiton dots}
\kat{Toukomies}{01.12.1929}{P22}{Tänä vuonna vahvistettu 5-Vuotissuunnitelma edellyttää, että Karjalan metsät on hakattava puhtaiksi. Vallan keskitys, tsheka, vastustajien tuhoaminen ja vainot ovat surullisia todistuksia tästä. Se ei voi muunlainen ollakaan, eikä voi paremmaksi kommunistisen järjestelmän Vallitessa tulla.}
\kat{Suomen Heimo}{13.03.1930}{P8}{Yleensä on kommunistien Tshekan ) Suomessa toimiva haaraosasto. B. KOMMUNISTIEN SALAINEN TOIMINTA.}
\kat{Suomen Heimo}{22.10.1930}{P18}{Sillä vierasta oli Suomenmaa, vaikkakin heimoveljien maa; eihän täällä pakolaisrukkaa pidetty vaivoin siedettyä loista parempana. Sanotaan etta bolshevikit. vangitsevat heti, olipa tulija kuka tahansa, ja raahaavat I Ytroskoihin tshekan kuulusteltavaksi. lupaikkojen hankkimista kesäajaksi veljesmaan ylioppilaille, jotka opiskelevat esim. kansantalous-, maatalousja lääketiedettä.}
\kat{Suomen Heimo}{22.10.1930}{P19}{KESKINÄINEN HENKIVAKUUTUSYHTIÖ SUOMI Rajalla olivat sitten tshekan miehet vastassa. Epäröiden vastasi viimein sieltä pieni kätönen.}
\kat{Toukomies}{01.03.1930}{P7}{Lehtien kertoman mukaan Neuvosto-Venäjä suunnittelee 100,000 puna-armeij alaisen ja punaupseerin lähettämistä maaseudulle talonpoikaisen kansan keskuuteen opettamaan sille maanviljelystraktorien käyttöä ja » maan sosialisoimista ». Venäjän piispat ja arkkipiispat eivät tietysti millään muotoa ole valittaneet ja tämänkin julistuksen he tietysti ovat antaneet — omasta vapaasta tahdostaan, ilman tshekan vähintäkään sekaantumista. Kirkkoja ei kaadeta eikä hävitetä, mutta itse kansa sitä haluaa, — se tapahtuu » väestön toivomuksesta, vieläpä monessa tapauksessa uskovaistenkin toivomuksesta » (! ) — Paavi on julkaissut vastalauseensa, mutta sehän kalastelee Venäjän kirkkoa helmaansa. »}
\kat{Toukomies}{01.06.1930}{P20}{Lukuunottamatta tshekaa, s. o. lähinnä kommunistista puoluetta, toimii se yhteistyössä puolueen ammatillisten ja muiden yhteiskunnallisten järjestöjen kanssa. Sen toimivaltaan kuuluu ennenkaikkea 16 — 23. ikäluokkaan kuuluvan nuorison kasvattaminen kommunistisen diktatuurin ehdottomiksi kannattajiksi ja tshekan asiamiehiksi. Karjalan kansaa rajan takana valistetaan nykyisin Marxin ja Leninin opeilla s. o. opetetaan vanhan maailman rikkirepimistä ja uuden tilalle suunnittelemista.}
\kat{Toukomies}{01.06.1930}{P21}{Kieretistä kertoo taru, että siellä asui Varlo-niminen Tämä tshekan asiamiehenä toimiminen muodostuu usein nuorisoliittolaisten päätehtäväksi. Nuorisoa kiihoitetaan tähän työhön maksamalla hyvä palkka jokaisesta ilmiannosta.}
\kat{Toukomies}{01.12.1930}{P12}{Jos Karjalan kieli olisi hävinnyt, silloin ei Karjalasta ja Karjalan kansasta voisi elia puhettakaan, sillä kansan tai heimon olemassaolon edellytyksenä on sen kieli, jonka kadotettuaan sillä ei ole oikeutta lukeutua kansakuntien joukkoon. Ranskan vallankumouksen kauhut, Neron vainot Roomassa, Kristian Tyrannin ajat Tanskassa ja Ruotsissa, ne kaikki kalpenevat sen järjestön hirmutekojen rinnalla, joka kantaa tshekan nimeä, järjestön, jonka kynnet ulottuvat melkein kaikkialle maailmaan, niin ettei ulkomaillakaan oleva henkilö voi olla täysin turvassa, jos se on hänet saaliikseen valinnut. Eikö tämä ole jo liikaa pienen heimomme osalle, eikö se taakka, joka kannettavaksemme on pantu jo dots}
\kat{Suomen Heimo}{30.08.1932}{P22}{Alussa ne olivat täysin itsenäisiä, sitten, PH:jen kultaaikana kaikki niiden asiat ratkaistiin Moskovassa saakka ja nytkin, monista uudistuksista huolimatta, ne kaikissa suhteissa ovat vain trustien osayrityksiä. Liittovaltion valtiollinen poliisi, GPU, valtio valtiossa, pitää huolen siitä, että kaikki viranomaiset, järjestelmän heille myöntämästä harkintavallasta riippumatta, tarkoin noudattavat kulloinkin vallassaolevan ryhmäkunnan poliittisia tarkoituksia. Voidaanko niille antaa riittävää kiihoketta niin kauan kuin teollisuuden omistusoikeus kuuluu kokonaan valtiolle, siinä on nähtävästi Neuvosto- Venäjän teollisuuden ja uskallammepa väittää koko sen valtiokapitalistisen järjestelmän dots}
\kat{Suomen Heimo}{30.08.1932}{P22}{Niinpä tehtaat nykyään tekevät trustien kanssa sopimuksia raaka-aineitten hankinnasta ja tuotteiden myynnistä, mutta koska tehtaalla ei ole oikeutta tehdä sopimusta muuta kuin oman trustinsa kanssa ja trusti omistaa itse tehtaan, on sopimus ilmeisesti aivan yksipuolinen. Ja jos GPU erikoisen ankarasti valvoo kommunistiseen puolueeseen kuulumattomien ja etenkin ulkomaalaisten teollisuuden toimihenkilöiden edesottamisia, niin ovat luotetuimmatkin kommunistit taasen oman puoluepoliisinsa valvonnan alaisia, joka ottaa määräyksensä yksinomaan puolueen pääsihteeriltä. Trustit on velvoitettu pitämään eri kirjaa eri tehtaiden taloudesta ja eräissä kommunistisen puolueen päätöslauselmissa on dots}
\kat{Toukomies}{01.02.1932}{P9}{Leiritulen ympärillä värjöttävä joukko oli kodeistaan karkoitettuja Itä-Karjalan ja Inkerin kansaa. Vapautta kesti vain vuoden, kunnes taas tsheka hänet vangitsi ja karkoitti vSiperiaan. Hän muistaa nyt ainoan poikansa, jonka oli onnistunut aikaisempina vuosina paeta Suomeen.}
\kat{Suomen Heimo}{30.03.1933}{P17}{Mitä tulee väestösuhteisiin, oli l>olshevikien v. 1926 laaditun tilaston mukaan asukkaita koko alueella 13,546,000 henkeä, josta V. 1924 kaikessa hiljaisuudessa Moskovasta saapunut tshekan erikoisosasto pidätti kansankomissarien neuvoston kokonaisuudessaan ja 28 muuta johtoasemassa olevaa henkilöä. Heidän ansionaan on pidettävä m.m. sitä, etta Tataristanin sotajoukot eivät ottaneet osaa Neuvosto- Venäjän ja Puolan väliseen sotaan.}
\kat{Karjalan Vapaus}{12.1934}{P41}{Opettaja, toimeensa pätemätön ja hermosairas vanhapiika, oli jollakin tavoin suututtanut paikallisen väestön ja nyt oli hän joutunut suoranaiseen piiritystilaan. Tshekan agenttien harjoittamasta provokatsionista aiheutui, että hyvin monet, joiden kanssa jouduin kosketuksiin, ilmaisivat kannattavansa venäläisten koulujen perustamista. Jos Petroskoista Repolaan menevä maantie kulkisi Suomen puolella rajaa, muodostaisi se varmasti yhden enimmänkäytettyjä matkailureittejämme, niin runsaasti tarjoaa se luonnonnähtävyyksiä kulkijalleen.}
\kat{Suomen Heimo}{06.04.1934}{P23}{Kuussaari esitelmän. Ohrana ja Tsheka. Tilaisuudessa piti m.m. ev. luutn.}
\subsubsection*{sorto \normalfont\small(25 katkelmaa, 274 osumaa aineistossa)}
\kat{Toukomies}{12.1925}{P4}{Kuten Attila tai Dshingiskhani aikoinansa vyöryttivät raaat laumansa sivistyneen maailman yli, niin aikasi meitä ja kaikkia Länsimaita ryöstö, hävitys; ja kuolema hillittömien bolshevikilaumojen puoletta. Muu Suomi, kiitos Jumalan ja valkoisen kansannousun, vapautui, mutta kaukainen Karjala, riittävän heimoushengen puuttuessa, oli tuomittu jäämään vieraan kansallisuuden sorron alaiseksi. Missä hyinen Pohjola ennen vain ankaralla työllä, uskomattoman suurella itsensä kieltäymyksellä ja kotiseudun rakkauden innoittamana taipui antamaan niukan toimentulon viljelijälle, siellä riistettiin kansalta nyt henkinenkin vapaus.}
\kat{Toukomies}{12.1925}{P6}{sivistys on tullut monia mutkia myöten; myöskin hunnien sotaretkien jäliltä, joiden hevoskavioiden alta ei sanottu enää ruohon nousseen, on levittäytynyt laajoja ja viljavia sivistysmaita, kuten yksin kaukaisen heimokansamme unkarilaisten historia todistaa. Kirkontornien helotus, kellojen soitto, kuvastojen kaunis kimmeltely, hartaat messut — nuo varmaan synnyttivät myös jotain sisäistä, elämää ja toivoa antavaa, niin toivotonta kuin muuten lienee elämä ollutkin, pelkkää sydämetöntä sortoa. Merkillisellä tavalla sivistyksen monimutkaisia kulkuteitä kuvaa myös se elämänhistoria, jonka tässä aiomme lyhyesti esittää ensimmäisenä noista muutamista menneisyyden miehistä, joita historia dots}
\kat{Toukomies}{12.1925}{P23}{Oikea joulukirja on Werner Söderström Oy:n kustantama ja Jalmari Kauppalan kirjoittama suurikokoinen teos Marttyyrien historiaa, jossa erinomaisen henkevästi ja mielenhartaudella esitetään enimpäin kristillisyyden alkuaikain marttyyrit. Toiseksi heimomielessä: jos sellaisissakin oloissa, missä Liettuan kansa on elänyt — Venäjän vanhan sorron pahimmin kuristamia kansoja — on sellaista omaa hengenviljelystä syntynyt, niin mitä sitten voi toivoa — vaikkapa rajan takana elävän karjalaisen heimonkin tulevaan hengenviljelykseen ja kansalliseen sivistykseen nähden. noustah kuin lumivöis noustah kukkazet muan!}
\kat{Toukomies}{01.1926}{P16}{Luettiin pari kiitoskirjettä myönnettyjen avustusten saaneilta opiskelijoilta ja käsiteltiin kolme uutta avustusasiaa. Onhan Suomi, samoinkuin Kreikka, pitkäaikaisen sorron jälkeen sitkeydellään saavuttanut valtiollisen itsenäisyyden. Oikea kirkko on Jumalan antama ja en sen oikeaksi tunnustaneet ei ainoastaan me, vaan koko kristikunta.}
\kat{Toukomies}{12.1926}{P9}{Leikkiä työn lomassa: Kerhonuonsoa karkelolla Sallisessa * koossapysymisessä ja voimistumisessa suurina sorron aikoina. Laulu on antanut voimaa Unkarin taisteluun vuosisataisia sortajia vastaan.}
\kat{Toukomies}{12.1926}{P18}{Kun saksalainen ritarikunta v. 1283 oli lopullisesti kukistanut baltilaissukuiset muinaispreussilaiset, käänsi se aseensa Liettuata vastaan taivuttaaksensa pakanallisten liettualaisten jäykät niskat pyhän Rooman alle, kuten vanha kronikka kertoo. Alamme tässä joulunumerossamme julkaista prof. A. R. Niemen huomattavaa esitystä Liettuan pienen ja kovia kokeneen kansan vaiheista, jotka erittäin suuressa määrässä ovat kuvaavia kaikkien pienten, erittäinkin Venäjän sorron alla olleiden, kansojen kohtaloille ja. samalla innostavia pienten kansallisuuksien nousulle. Kaikkea tätä ja tämän ajan vapautta todistaa m. m. eräs maininta Novgorodin ja Gotlannin sekä toisaalta Saksan kauppiaiden keskisessä dots}
\kat{Toukomies}{12.1926}{P20}{35, iB:lta 130, yht. Sorto kohtasi etupäässä yläluokkia. Venäläistyttämisprosessia alettiin jouduttaa.}
\kat{Toukomies}{01.02.1927}{P18}{Hän syntyi 23 p:nä marrask. Mutta sorto vaikutti päinvastaista. Liettuan entinen pääkaupunki Vilna.}
\kat{Toukomies}{01.02.1927}{P18}{Liettuan kansa vietti viime marraskuun 23:na p:nä aamun juhlaa, silloin tuli kuluneeksi 75 vuotta sen kansallisen uudestisyntymisen luojan, tri Juhana Basanaviciuksen syntymisestä, — kuten jo kirjoitukseni edellisessä osassa olen maininnut. Uusi sivistyksen riemuvoitto todellakin, voi täydellä syyllä huudahtaa se, ken tietää, mitä on ollut ja mihin pyrki venäläinen sorto valtius, joka järjestelmällisesti vuodesta 1864 saakka tahtoi poisjuurruttaa maan päältä koko Rooman-katoolinen papisto, joka oli suurimmaksi osaksi suoraan kansasta noussut, asettui taistelevien eturiviin, ja niin alkoi vuosikymmeniä kestävä herkeämätön taistelu, jossa ulkonaiset voimakeinot ja raaka väkivalta oli dots}
\kat{Toukomies}{01.06.1927}{P3}{Ei tarvitse olla oikeusoppinut voidaksensa niiden sisällyksestä olla selvillä. Venäläistyttämistyö oli täydessä käynnissä, sorto ja vaino jokapäiväisiä ilmiöitä. Jäljelle on jäänyt ainoastaan Karjan kysymys — Karjalan oikeus.}
\kat{Toukomies}{01.06.1927}{P12}{Sellaisena oli hän koko karjalaisen vapaustaistelun vertauskuva ja kaunis vertauskuva. Ja sitten, veli, sai Karjalan kansa olla vuosisadan toisensa jälkeen ryssän sorron alaisena. Nyt on pyydys pyytämässä, ansaraudat vaanimassa, miehet niemien nenissä, saktkytät salmen suissa.}
\kat{Toukomies}{01.06.1927}{P15}{musta, ajankohta osottautui oikein valituksi. Oma maa ja oma raskasta sortoa kärsivä kansa Basanaviciuksen 010 Bulgariassa ei tällä kertaa ollut tuleva pitkä-aikaiseksi.}
\kat{Suomen Heimo}{20.01.1928}{P12}{Lukemattomissa taloudellisissa ja lukemattomissa poliittisissa vaikeuksissa — niin lapsellinen ja pikkumainen on Latvian hallitus, että se näkee separatistisia pyrkimyksiä liiviläisten liikehtimisessä — on liiviläisten herännyt aines saanut työskennellä henkisten ja aineellisten elinehtojen luomiseksi. Ja niinkuin me suomalaiset riensimme virolaisten avuksi silloin, kun ryssää ajettiin Narvan joen taakse Virosta, tai silloin, kun Karjalan miehet nousivat, on välttämätöntä, että me lähdemme nytkin, jolloin pieni Liivin kansa, on nousemassa lättiläistä sortoa vastaan. On ilmestynyt joukko liiviläisiä kirjoja, ön noussut joukko liiviläisiä runoilijoita, liivin kieltä opetetaan kuudessa dots}
\kat{Suomen Heimo}{20.01.1928}{P18}{Se oli suomalainen, joka houraili — Mutta sorron, tulen ja veren te toitte. Veljekset katsoivat isäänsä, jonka kasvot värähtelivät.}
\kat{Suomen Heimo}{20.05.1928}{P3}{— Tietoja eri Aloilta: Akateemisen Karjala-Seuran kokous. ( Aimo Vuorinen ) — Harjoitetaanko Ruotsin Länsi-Pohjassa kielellistä sortoa. Maamme kaikkien maarajojen takana aauu omia heimolaisiamme, omaa Suomen kansaamme.}
\kat{Suomen Heimo}{20.05.1928}{P10}{Leppälän Kuvastintehdas Uudenmaank. ettei mitään sortoa harjoiteta? 11, puhelin 27023.}
\kat{Suomen Heimo}{20.05.1928}{P10}{Ovatko siis ruotsalaiset oikeassa väittäessään, HARJOITETAANKO RUOTSIN LÄNSI-POHJASSA KIELELLISTÄ SORTOA? Suomalaisten puolelta vastataan tähän kysymykseen myönteisesti, ruotsalaisten puolelta kielteisesti.}
\kat{Suomen Heimo}{20.05.1928}{P10}{Tämä on aivan ymmärrettävää, sillä olihan meillä Venäjän vallan aikana kylliksi huolehtimista omista asioistamme, eikä siis heräävältä suomalaisuudelta voi " nut vaatia mitään maan rajojen ulkopuolelle ulottuvaa herättävää toimintaa. Parhaillaan vireillä oleva Länsi-Pohjan kysymys, josta Suomen ja Ruotsin ylioppilaitten kesken käydään jo väittelyä, on <tässä väittelyssä kärjistynyt kysymykseen » Harjoitetaanko Ruotsin Länsi-Pohjassa kielellisiä ja kansallista sortoa vai ei? « Surullinen tosiasia nim. on, että kysymyksessäolevan väestön enemmistö ei tosiaankaan vastusta ruotsinkielen leviämistä keskuuteensa, vaan katselee sitä, ellei aina suopein silmin, niin ainakin välttämättömänä ja dots}
\kat{Suomen Heimo}{20.05.1928}{P18}{Tästä on ollut seurauksena taloudellinen terrori. Poliittinen ja taloudellinen sorto. Kaikki riippuu Teistä ja — me luotamme Teihin.}
\kat{Suomen Heimo}{15.08.1928}{P11}{Sillä Maammelaulun on runoililut Suomen kansan isänmaallisen hengen luoja jia säveltänyt Suomen taidemusiikin isä. Niin kauan emme voi emmekä saa laulaa näin, kun Karjala ja Inkeri huokaavat ryssän sorron alla. 2 ) „Isänmaan kasvoit " on muuten kaikin puolin sopiva kansallislauluksemme, mutta se on edellä ajastaan.}
\kat{Suomen Heimo}{15.09.1928}{P6}{Lukuunottamatta sitä, että ruotsinkieli perustuslain pyhyyden omaavalla säännöksellä on julistettu maakunnan viralliseksi kieleksi ja maakunnassa toimivien valtion viranomaisten ja valtioneuvoston y.m. virastojen väliseksi sisäiseksi virkakieleksikin, löytyy n.s. takuulaissa Kansainliiton painostuksen johdosta seuraavat, pahaa verta herättävät säädökset: Suomenkielisiä valtion oppilaitoksia ei Ahvenanmaalle saa perustaa, eivätkä maakunta ja Ahvenanmaan kunnat ole velvolliset perustamaan tai ylläpitämään muita kuin ruotsinkielisiä kouluja. Sillä sitä periaatetta, ettei vähemmistö saa omakielisiä kouluja, ei edes ole sanellut mikään ruotsalaisuuden säilyttämiseen tähtäävä ajatus, vaan dots}
\kat{Suomen Heimo}{15.09.1928}{P19}{Tämä tietenkin hieman ällistytti suomalaisten dele i atiota, jonka johtaja kuitenkin kohta rauhallisessa äänilajissa selvitti, ettei riidan hakemisen halu tässä asiassa ollut suomalaisten puolella, ja että jos ruotsalaisilla on halua asioista neuvotella, niin se sopii hyvin tehdä yksityisestikin. Siinä vielä lisäksi annettiin erittäin perinpohjainen selvitys Länsipohjassa asuvan suomalaisväestön tilasta, suuruudesta ja kohtelusta, ja todisteltiin erittäin helly t tel e väs sä äänilajissa, ettei Ruotsin taholta mitään erikoista sortoa harjoiteta, vaikka kylläkin j är jestelmällistä ruotsalaistuttamista. Hyvämuistiset lukijamme muistanevat viime keväänä sattuneen Tukholman ylioppilaskokouksen dots}
\kat{Suomen Heimo}{18.12.1928}{P18}{Säästöpankki S ••••••••• i i • aas t ajan pankki Talletuksistanne syntyvä voitto jää säästäjäin omille, yleishyödyllisille laitoksille, jos talletatte rahanne Säästöpankkiin tai niiden keskusrahalaitokseen SÄÄSTÖPANKKIEN KESKUS-OSAKE-PANKKIIN Helsinki, Aleksanterink. Lyhyesti sanoen suomalainen on ruotsalaisseuduilla aivan turvaton kansanopetuksen suhteen, milloin ei tämä sorto mene pitemmällekin, kuten joku aika sitten Maalahdella, jossa suomalaisen yksityisen kansakoulun ikkunat ja ovet vähän päästä särjettiin! -ten koulujen asioita loppumattomiin, sikäli kun ei asianomaisen maaherran lopulta täydy turvautua asettamaan uhkasakkoja kunnalle, joka ne kukaties vielä maksaakin ja odottaa dots}
\kat{Toukomies}{01.02.1928}{P12}{Ei siis Petr, vaan Pietari tai Petrus tai supisuomalainen Pekka, ei loan, vaan Johannes, Juho j. n. e. Tässä suhteessa olisi pappien otettava heti täysi askel ja jo ristimätilaisuudessa ohjattava kansa antamaan suomalais- tai länsieurooppalaismuotoiset nimet ja ne heti virallisiin lähteihin merkittävä. Vieraiden nimien juuri sitäpaitsi on usein satunnainen tai sitten kaukainen, usein jo itse suvultakin unohtunut, ja joka tapauksessa se on sellainen, joka ei tälle maalle ja kansalle valaise mitään, ei tuo mitään kansallisesti kohottavaa tai suurta mieleen, päinvastoin useinkin masentavaa, ikävää sortoa, vierasta valtaa. Loukkaan ehdottomasti suomalaista sukutuntoa ja korvaa sekä on muutenkin dots}
\kat{Toukomies}{01.03.1928}{P9}{Mutta siitä huolimatta julkaistiin sitä kaikki vuoden kaksitoista numeroa, alkunumeroiden jälkeen laajoina kaksoisvihkoina antamalla niille eri nimet ja vain alaotsikossa viitaten yhtenäiseen julkaisu järjestykseen. Alkavan uuden Suomen sortoa j an vuoksi ei kuitenkaan lehdelle saatu Suomen painohallitukselta lupaa monienkaanluvanhakuyritysten perästä ja niin lehti ei voinut virallisesti ilmestyä säännöllisenä aikakausi ehtenä. Siitä johtuvat vihkojen vaihtelevat nimet: » T ouon p a n o », » Kar j alaisten kesälehti », » Elon aika », » Karjalaisten syyslehti », » Karjalaisten Joulu ».}
\subsection{Terrorin sanasto}
\subsubsection*{kolhoosi \normalfont\small(30 katkelmaa, 133 osumaa aineistossa)}
\kat{Toukomies}{01.04.1931}{P35}{3Battl ) an fanfan elämää opitte tuntemaan fau- neista furoateoffistamme — Niitä aikoja, jolloin nykyiset va nhu k- set olivat nuoria, kuvailee tutkijana ja valokuvaa- jana yhtä etevä maisteri AHTI RYTKÖNEN tavattoman viehättävässä kuvateoksessaan SAVUPIRTTIEN KANSAA Laakapainostamme juuri ilmestynyt teos sisältää 144 taiteellista jäljennöstä tekijän ottamista valo- kuvista. Niinpä mainitsemassamme Raja-Konnun kylässäkin on kymmeniä taloja käsittävä kolhoosi ( yhteistalous' eli kommuuna ), jonka hallussa ovat nyt kaikki kylän parhaat viljelysmaat ja nurmet sekä yksityisiltä ryöstetty omaisuus, kuten karja, työvälineet y. m. Työt tässä kommuunassa ovat järjestetyt suurtilan malliin dots}
\kat{Toukomies}{01.05.1932}{P12}{Kolmen vuorokauden kuluttua lähdimme Uhtualle, jossa saimme matkapassit toimikunnalta. Urakkatyösopimuksia on tehty 11 joen uitoista, niistä 8 eri kolhoosin kanssa tehtyjä. Matkalla kotoani Kiisjoelta jäin Röhössä odottelemaan kylämme hevosmiehiä, jotka olivat matkalla Kajaaniin jauhon hakuun.}
\kat{Suomen Heimo}{30.04.1933}{P20}{Uusi Pesula Osakeyhtiö Ab. Poltetut kolhoosit Etelä-p oh jalainen}
\kat{Toukomies}{01.10.1933}{P12}{Hän on suomentanut m. m. P. Johannes Krysostomuksen liturgian ja ollut » Aamun Koitto » -lehden päätoimittajana 1897 — 1907 sekä vastaavana toimittajana lehden uudenkin, v. 1918 alkaneen vaiheen aikana viime vuosiin. Punaisen Karjalan kirjeenvaihtaja mainitsee, että kuluu varmaankin vielä toiset 3 1 / 2 vuotta, elleivät rakennusaineiden hankinnassa Tunkuan Tpk:n kunnallisosasto, Niikkananvaaran kyläneuvosto ja kylien kolhoosit paranna otettaan. — Helsingin yliopiston muuttamisesta täysin suomalaiseksi on maan hallitus tehnyt suunnitelman ja siitä johtuvan ruotsinkielisten yliopisto-opetuksen järjestämisestä kysynyt m. m. Turun ruotsalaisen akatemian mieltä tarkoituksella että se ottaisi dots}
\kat{Karjalan Vapaus}{08.1934}{P2}{Vienan kansa sananlaskujen valossa, N. V. Kolhoosi » Lokakuun Sankarit », Levanpoika. Karjalavalistusta tehostamaan, L. R — na.}
\kat{Karjalan Vapaus}{08.1934}{P12}{Sitä ei voitu välttää. Ainoastaan kolhoosin johtaja. Kokous kesti koko päivän.}
\kat{Karjalan Vapaus}{08.1934}{P12}{Ja vaikeneminenhan on tavallisesti myöntymyksen merkki. Kolhoosi " Lokakuun Sankarit. " Mutta se " vauras elämä ", josta tov.}
\kat{Karjalan Vapaus}{08.1934}{P12}{— Kenties tällä tapauksella ja talonpoikien vaikenemisella oli jotakin yhteistä. Tässä kolhoosin perustavassa kokouksessa kylän työläiset ja talonpojat näet vaikenivat. Ajattelivatpahan vain itsekseen, että jahka tässä....}
\kat{Karjalan Vapaus}{08.1934}{P12}{Se ensi töikseen järjesti kokouksen ja siinä ilmoitti kyläläisille valmiit ponnet, jotka kuuluivat näin: Potjomkin oli kolhoosin perustavassa kokouksessa väläytellyt teoreettisia utukuvia, jäi saavuttamatta Rokkalan kolhoosissa. Talonpojat väsymyksestä jo haukottelivat ja tätä raukeutta taisi vielä erikoisesti lisätä suuri välinpitämättömyys ponsia kohtaan.}
\kat{Karjalan Vapaus}{08.1934}{P12}{" Uusi elämä " oli alkanut voittaa alaa Rokkalankin kylässä, vaikka se, tämä Rokkala, ei ollutkaan aivan lähellä parempiosaisten asutuskeskuksia. Kylään oli nim. perustettu eräänä vallankumouksen vuosipäivänä ja sen kunniaksi yhteinen talous, kolhoosi ja sille oli annettu nimeksi " Lokakuun Sankarit ". Kaksi vangituista oli ammuttu " neuvostovallan vastustamisesta " jo kylän reunassa ja muista ei ollut mitään tietoa, minne heidät oli viety.}
\kat{Karjalan Vapaus}{08.1934}{P12}{Kaksi vangituista oli ammuttu " neuvostovallan vastustamisesta " jo kylän reunassa ja muista ei ollut mitään tietoa, minne heidät oli viety. Talonpoikien oli vietävä tavaransa ja karjansa, mikäli niitä oli aikaisemmilta pakko-otoilta säästynyt, kolhoosin omaisuudeksi ja heidän itsensä oli liityttävä kolhoosin jäseneksi. Toveri oPtjomkin puheli kokouksessa yksin, teki ehdotukset ja saman tien myös hyväksyi ne lyömällä nyrkkinsä pöytään, itse hän kun sattui olemaan kokouksen puheenjohtajanakin.}
\kat{Karjalan Vapaus}{08.1934}{P12}{Kaksi vangituista oli ammuttu " neuvostovallan vastustamisesta " jo kylän reunassa ja muista ei ollut mitään tietoa, minne heidät oli viety. Talonpoikien oli vietävä tavaransa ja karjansa, mikäli niitä oli aikaisemmilta pakko-otoilta säästynyt, kolhoosin omaisuudeksi ja heidän itsensä oli liityttävä kolhoosin jäseneksi. Toveri oPtjomkin puheli kokouksessa yksin, teki ehdotukset ja saman tien myös hyväksyi ne lyömällä nyrkkinsä pöytään, itse hän kun sattui olemaan kokouksen puheenjohtajanakin.}
\kat{Karjalan Vapaus}{08.1934}{P12}{Toveri oPtjomkin puheli kokouksessa yksin, teki ehdotukset ja saman tien myös hyväksyi ne lyömällä nyrkkinsä pöytään, itse hän kun sattui olemaan kokouksen puheenjohtajanakin. 3 ) Mainitun kolhoosin johtajaksi on piirineuvosto määrännyt toveri Huliganovin, joka on samalla kylän fyysis-kulttuurisen solun tuotannollis-taloudellisen koperatiivin valtuutettuna. 2 ) Tämän nojalla piirineuvosto määrää, että Rokkaian työläiset ja talonpojat luovuttavat kaiken omaisuutensa kolhoosille ja itse liittyvät 10O \%:sesti sen jäser-iksi.}
\kat{Karjalan Vapaus}{08.1934}{P12}{lehmiä, samoin vasikoita ja lampaita, kymmenkunta hevosta, maanviljelyskaluja, veneitä, kalan pyydyksiä metsästysvälineitä\textasciicircum{} rakennuksia j. n. e. — Ei oikeastaan mitään puut unut. il ) Piirineuvosto on juhlakonferenssissaan todennut, että Rokkalan työläiset ja talonpojat ovat yksimielisesti perustaneet kolhoosin ja sen nimi on " Lokakuun Sankarit ". Mutta sen lupasi toveri organisaattori toimittaa piirin keskuksesta, niin että " sosialistinen rakennustyö " rupeaa luistamaan Rokkalassa " ennenkuulumattomilla tempoilla " ja siitä on tuloksena " varakas elämä ".}
\kat{Karjalan Vapaus}{08.1934}{P12}{Potjomkin itsekin vihjaili puheessaan, että kyllä me nujerramme " kulakit ", " opportunistit " ja muut neuvostovallan viholliset, jotka eivät noudata dekreettejämme, taikka jotka eivät " aktiivisesti " osallistu " sosialistiseen rakennustyöhön ". Sen perustaja, toveri Stahvei Potjomkin, entiseltä ammatiltaan kadunlakaisija kaupungista, nykyiseltä arvoltaan teknillis-taloudellisen jatsheikan polittis-yhteiskunnallisen pleenumin kolmannen aluepiirin kolhoosi-osuuskunnallisen sektorin organisaattori, oli keskusneuvostolle lähettämässään " raportissa " muistanut luonnollisesti huomauttaa, että Rokkalan kolhoosi, " Lokakuun Sankarit ", syntyi kylän työläisten ja talonpoikien yksimielisellä dots}
\kat{Karjalan Vapaus}{08.1934}{P13}{Tällaisia huolia on punaisilla spetseilla » kulttuurisen rakennustyön alalla » Karjalassa. Huliganov ilmoitti piirin keskukselle, että Rokkala kolhoosin karja kuoli viimeistä päätä myöten punatautiin. Petroskoissa ovat parhaillaan rakenteilla kaupungin vesijohto, leipätehdas, spesialistien talo ja kansallinen kulttuuritalo.}
\kat{Karjalan Vapaus}{10.1934}{P17}{Karjalan uutisista havaitaan, ovat tuotantovajauksen syyt kalastuksessa sangen moninaiset. Niin ollen on selvää, että sanotun kolhoosin kalastusta on haitannut elintarvepula. Uhtuan piiristä ilmoitetaan, että vesiliikenne Kuittijärvellä " on heikossa kunnossa ".}
\kat{Karjalan Vapaus}{10.1934}{P17}{Ontolahden metsätyöpunktin kalastusprikaadi on jo pitemmän ajan vaatinut itselleen pyyntivälineitä, mutta tuloksetta. Krasnyi-Trud " ( Punainen-Työ ) niminen kolhoosi ( Louhen piirissä ) on jäänyt varustamatta elintarpeilla. Polttoaineen puutteesta ovat joutuneet kärsimään Kiestingissä kaikki muutkin järjestöt kuin autobaasa.}
\kat{Karjalan Vapaus}{10.1934}{P17}{Lisäksi on mainittava, että myöskin polttoaineen puuttumisen vuoksi ovat autot joutuneet seisomaan usein viikkomääriä. " Krasnyi Ryibak " ( Punainen Kalastaja ) niminen kolhoosi on muuten valmistunut syyskalastukseen, mainitsee Pun. Syyskalastuksen tärkeimpänä aikana useat kollektivistit istuivat kotonaan ja esim. " syväkalastukseen he eivät halunneet lähteä.}
\kat{Karjalan Vapaus}{10.1934}{P17}{Kolme autoa, jotka ovat nyt olleet liikkeessä, joutuvat myöskin lähipäivinä pois liikenteestä kumien puuttumisen takia. Karjalan uutisessa, mutta kolhoosit yhtä kaikki " lorvailevat ", eivätkä lähetä kalastajia merelle. Ja niistäkin lehdistä, mitä piireille menee, on melkoinen määrä vapaita, tilaani attomia kappaleita.}
\kat{Karjalan Vapaus}{10.1934}{P18}{Mutta tseka on järjestö, joka saa selville vanhat ia salaperäisetkin asiat. KutsuttLn koolle kolhoosin ja kyläneuvoston täysistunto, johon kutsuttiin huomatuimpia luottamusmiehiä. No selvä on ", sanoi tsekata päällikkö ja kehoitti istuntoa hajaantumaan.}
\kat{Karjalan Vapaus}{10.1934}{P18}{Asiamiehiä otamme hoikille paikkakunnille hankkimaan leh " dellemme tilauksia * Maksamme tilausten hankki- joille edullisen palkkion. Seuraavana aamuna saapui " Neuvostonainen " kolhoosiin suuria virkamiehiä, joukossa itse piirin tsekapääilikkö ynnä piirin kansanvalistusasian asiamies. Se tapahtui niin, että Hotokalle myönnettiin pienoinen lahja kolhoosin puolesta ja päätettiin se antaa hänelle lokakuun vallankumouksen suurena merkkipäivänä.}
\kat{Karjalan Vapaus}{10.1934}{P18}{Seuraavana aamuna saapui " Neuvostonainen " kolhoosiin suuria virkamiehiä, joukossa itse piirin tsekapääilikkö ynnä piirin kansanvalistusasian asiamies. Se tapahtui niin, että Hotokalle myönnettiin pienoinen lahja kolhoosin puolesta ja päätettiin se antaa hänelle lokakuun vallankumouksen suurena merkkipäivänä. Istuttiin, tuumattiin ja pohdittiin pykäliä ja puolueen direktiivejä, ja juotiin teeheinien puutteessa " Pyhän Bohrotsan " heinistä keitettyä juomaa. "}
\kat{Karjalan Vapaus}{10.1934}{P18}{Tämä jo suututti Hotokkaa ja hän sanoa tokaisi: " Mie en taho siun nenänkannattimies, tai suat pityä lahjaskhij jos tuo puumerkkinä ei kelpua. " Juhlallisesti asteli Hotokka monien muitten joukossa kolhoosin neuvostoon ja odotteli vuoroaan: saadakseen sen kauan toivotun ahkeruuden tunnustuksen — kolhoosin lahjan. Ne, jotka siis haluavat hankkia itselleen varsinaisen toimensa ja työnsä ohella lisätuloja, voivat sen helposti järjestää pyytämällä asiamiespaperit ja ryhtymällä vapaa- hetkinään keräämään tilauksia.}
\kat{Karjalan Vapaus}{10.1934}{P18}{Tämä jo suututti Hotokkaa ja hän sanoa tokaisi: " Mie en taho siun nenänkannattimies, tai suat pityä lahjaskhij jos tuo puumerkkinä ei kelpua. " Juhlallisesti asteli Hotokka monien muitten joukossa kolhoosin neuvostoon ja odotteli vuoroaan: saadakseen sen kauan toivotun ahkeruuden tunnustuksen — kolhoosin lahjan. Ne, jotka siis haluavat hankkia itselleen varsinaisen toimensa ja työnsä ohella lisätuloja, voivat sen helposti järjestää pyytämällä asiamiespaperit ja ryhtymällä vapaa- hetkinään keräämään tilauksia.}
\kat{Karjalan Vapaus}{10.1934}{P18}{Sen enempää Hotokka ei tahtonut kuunnella tuota tsekan päällikön kiusallista kuulustelua, vaan pyörähti ympäri, kohautti " emussan helmua ", sanoen kiivaassa äänensävyssä: " Kun et prohvosta uso, niin ole uskomatta. Viisi vuotta ja muutamia päiviä on Hotokka nauranut partaansa, kun pääsi kuin " kutti veräjästä ", ettei tarvinnut istua kolhoosin pahahajuisessa ruokalan ummehtuneessa suojassa, kommunistinen aapinen käsissään, oppimassa luku- ja kirjoitustaitoa. Pää alas painuneena, syvissä mietteissään " paikankokkua " pyöritellen astui Hotokka tsekan päällikön eteen, ottaa tavallisen kolhoosilaisnaisen arastelevan asennon ja tokaisee hiem. ua pelkoa peittelevällä äänensävyllä: " Ka miksi sie dots}
\kat{Karjalan Vapaus}{12.1934}{P17}{Sen " soluja " on perustettu kyliin, joissa on urkkivat ja terrorisoivat väestöä. Uskonnollisten juhlien aikana meidän on järjestettävä erikoisia uskonnonvastaisia tilaisuuksia jokaiselle metsätyömaalle ja jokaiseen kolhoosiin. Tärkeintä osaa tässä työssä näyttelee luonnollisesti " Jumalattomien Liitto ", joka varta vasten on sitä varten perustettukin.}
\kat{Karjalan Vapaus}{12.1934}{P30}{Kirkossa kävijät eivät montakaan sanaa ymmärtäneet papin puheesta, joten kirkon taholta tuleva valistus jäi pakostakin hyvin pintapuoliseksi. Tähän lokaan lankeavat jo keskenkasvuiset tytöt ollessaan suorittamassa ylivoimaista pakkotyönormiaan, joko kolhoosin pelloilla tai metsätrustin tukkityömailla. Siveellisyyssuhteessa ja hyvien tapojen palvomisessa olisi Karjalan kansan kyennyt vielä noin 20 vuotta takaperin kilpailemaan minkä kansakunnan kerällä hyvänsä.}
\kat{Karjalan Vapaus}{12.1934}{P33}{— " Jumala, anna anteeksi, jos tuskissani olen kironnut "... hän sanoo taaskin kuuluvasti, " mutta nämä sielun ja ruumiin tuskat ovat raskaat kestää, ei kuitenkaan niin raskaat, kuin Poikasi ristinpuulla. Kyläneuvoston jäsenet ja köyhälistöryhmän akat viettämässä lokakuun vallankumouksen vuosijuhlaa kolhoosin uuden navetan edustalla. Kaiken!}
\kat{Suomen Heimo}{27.10.1934}{P14}{Siellä ensikerran näimme oikean karjalaisen talon, pätseineen ja patsaineen, laskeuduimme kolpitsan kautta karsinaan ja tarkastimme asiaankuuluvasti jauhinkivien paikan. Mutta ne näkevätkin siellä joka päivä rajan ja ryssän kolhoosin ja vanhan rukoushuoneensa, tsasounansa, jonka raja jakaa keskeltä kahtia —. Tuberkuloosiesitelmää ja pitkää isänmaallista puhetta kuunneltiin hievahtamatta, — eikä niiden ajaksi lähdetty " pufettiin ", niinkuin täällä lännessä usein on tapana.}
\subsubsection*{Solovetsk \normalfont\small(30 katkelmaa, 64 osumaa aineistossa)}
\kat{Toukomies}{12.1925}{P17}{Niitä lahkojen johtajat — usein tehovoimaisia henkilöitä — ■ pitivät öillä, ja luonnollisestikaan ei saanut laiminlyödä pääsiäisyötä. Sen juuret nähtävästi juontuivat siltä ajalta, kun Solovetskin luostari taisteli urhokkaasti lahkon puolesta, elleivät juontutaan. Sitä taas varmasti edisti se, että papit Vienan- Karjalassa olivat vieraskielisiä ja kirkko ei avannut oviaan kansan kielelle.}
\kat{Toukomies}{12.1926}{P20}{Tämän jälkeisen historian voimme lyhyesti kertoa. Perustetaan Valamon ja Solovetskin ja muita luostareja. Erään laskelman mukaan v:lta 1900 tuotiin Liettuaan vv.}
\kat{Toukomies}{12.1926}{P26}{20: —. Kertomus viisaasta, monitaitoisesta, inhimillisiä ominaisuuksia osoittavasta Toto simpanssista. Järkyttäviä kuvauksia Solovetskin luostariin sijoitetutta valtiollisesta vankilasta, jonka hirmujen vertaa vain lienee keskiaika tuntenut. — Tässä teoksessaan White kertoo vuoden 1849 suuresta kultakuumeesta, joka tarttui nuoriin ja vanhoihin ympäri koko pohjoisen Amerikan.}
\kat{Toukomies}{01.02.1927}{P23}{Varsin vakuuttavasti tuo tekijä siten esiin palan ihmisluonnon kataluutta — ja ehkä nimenomaan juuri yksilön kataluutta, paljastaa erään sellaisen puolen tässä hyvin tunnettua systeemiä, joka ei helpoimmin tule suuren ihmiskunnantietoon, mutta joka tiedoksi tultuaan ei jaksa kylliksi voimakkaasti saarnata tällaista ihmisten parantamista ja onnellistuttamista vastaan. Teoksen, joka ei ole erittäin laaja, on kirjoittanut entinen venäl. upseeri S. A. Malsagov, joka ensin, bolshevikkivallan alkuaikoina, taisteli n. s. vastavallankumouksellisten armeijoissa Kaukaasiassa, joutui bolsheviikkien käsiin ja vieriskeli heidän monissa vankiloissaan, kunnes lopulta, v. 1924, siirrettiin heidän suureen dots}
\kat{Suomen Heimo}{20.01.1928}{P19}{Mutta talo on rakennettava nopeasti mainitulle palstalle, muuten alue otetaan taas pois, minkä vuoksi liiviläisten yhdistys pyytää taloudellista apua myöskin Suomesta. Lisäksi on talollinen Valerian Reijo, joka on ollut useita kuukausia Jaaman ja Shpalernajan vankiloissa, nyt lähetetty Solovetskin vankileiriin, luultavasti elinkaudeksi. Eikö olisi aika jo siivota uskonnollinenkin sanastomme, olkoon Harjoitus kuinka tieteellinen hyvänsä, ja ennen kaikkea vielä päivälehdestä käyvässä julkaisussa, jota laajoissa piireissä luetaan.}
\kat{Suomen Heimo}{20.05.1928}{P19}{Inkerin nuorisoa ja Inkerin kansaa ei ole saatu kommunismiin taivutetuksi, mutta nuorison sielu on saastutettu, se on myrkytetty kommunismibasilleilla, jotka uhkaavat tuhota tämän iterveen ja kansallisesti sitkeän Suomen sukupuun vankkojen juurien versoavat vesat. Tehtyjen laskelmien mukaan on tällä hetkellä parisataa inkeriläistä karkoiltéttuna poid kotipaikka- * kummiltaan, kuka Siperiaan, kuka Sisä- tai Pohjois- Venäjälle, ja onpa muutamia lähetetty \} kauhuistaan kuuluisaan Solovetskin vankileiriin Pohjoisella Jäämerellä. Useita Pohjois-Inkerin miehiä on hakenut neuvostoviranomaisilta lupaa päästäkseen tänä keväänä käymään Suomessa, Kun mitään pätevää käyntiluvan estettä ei ole keksitty, dots}
\kat{Toukomies}{01.02.1928}{P16}{Jäseneksi ilmoittautumisia ja lahjoituksia seuran tarkoitusten hyväksi ottavat vastaan kaikki seuran jäsenet, jotka tästä asiamiestoimestaan saavat palkkiokseen 15 \% hankintojensa summasta. Hovatta Hovatanpoika kerran muuanna kesänä, suloisimpana suvena lähti viemähän venettä halki kuulun Kuitti ] arven, matkalla manasterihin, Solovetskin saaren suulle. Jäsenmaksut; vuosi jäseniltä 50 mk. vuodessa, tavall. ainaisilta jäseniltä 500 mk. kertakaikkiaan, ainaisilta lahjoittajajäseniltä 2,000 mk. ja ainaisilta säätäjä jäseniltä 5,000 mk. kertakaikkiaan.}
\kat{Toukomies}{01.06.1928}{P14}{- N:o 6 — 7 Solovetskin kuulu saari V ienanmeressä. Voologdan läänin puolelle sekä Aunuksen läänin että vielä Vienan lääninkin kaukaisten rajojen yli.}
\kat{Suomen Heimo}{15.08.1929}{P20}{Ja kun pappi kertoi heille Aikateemisesta Karjala- Seurasta ja siitä, kuinka on Suomessa tuhatlukuinen joukko nuoria miehiä, jotka ovat vannoneet uho-aavansa työnsä ja elämänsä Isänmaan, Karjalan, Inkerin ja Suomen puoleista, välähti noiden miesten silmissä uuden toivon kipinä. He olivat nähneet joukkojensa hupenevan viimeiseen mieheen, he olivat nähneet kyliensä ja kotiensa jäävän ryssän käsiin., jopa sortuvan liekkeihin, heidän omaisiansa oli raahattu Pietarin, Novgorodin ja Solovetskin vankiloihin ja he olivat joutuneet vieraalle maalle maanpakolaisen niukkaa leipää syömään. Myöhäisestä ajasta huolimatta kokosi isäntäni kylässä majailevat pakolaiset kotiinsa kuulemaan Suomesta saapunutta dots}
\kat{Suomen Heimo}{15.06.1930}{P19}{Moni olisi valmis jättämään pienen kotikontunsa ja lähtemään esi-isiensä maahan, Suomeen. Yksistään Solovetskin kauhunsaarilla on parisen sataa Inkerin talonpoikaa odottamassa varmaa kuolemaa. Mutta yhä suuremmaksi on kasvanut rajamuuri, ylitsepääsemättömäksi on käynyt kohtalokas Rajajoki.}
\kat{Toukomies}{01.12.1930}{P10}{Oli nimittäin pääsiäinen meneillään. Solovetskin luostari Vienan meressä. 18 vuosisadalla kävi engelsmanni luostaria pommittamassa.}
\kat{Toukomies}{01.12.1930}{P10}{oli saanut läpikäydä —, havaitsin heti, kun isäntäväen kanssa astuin ruokasaliin. Luostarin sanotaan perustaneen kaksi pyhimystä nimeltään Isossima ja Savaatei, Solovetskin ihmeidentekijät. Kirkko ei saanut olla koskaan tyhjänä, aina piti olla siellä jonkun seisomassa.}
\kat{Toukomies}{01.12.1930}{P12}{Tämän kalliin perinnön ovat esi-isämme meille jättäneet ja sanoneet: » Edustakaa te meitä ja jättäkää muistomme jälkimaailmalle yhtä kunniakkaana ja tahrattomana kuin sen itse perinnöksi saitte. Alituiseen saamme kuulla jonkun omaisistamme joutuneen punaisten verihurttien uhriksi; kuka on joutunut Solovetskin vankilahelvettiin, kuka Siperian vaalenneiden ihmisluiden ja pääkallojen kammottavaan valtakuntaan. Toiselta puolen Ruotsi ja toiselta Venäjä ovat Karjalassa törmänneet yhteen hävittäen maata, välittämättä siitä, että heihin kumpaankaan kansallisesti ja valtiollisesti kuulumaton maa saa kärsiä seuraukset.}
\kat{Suomen Heimo}{24.01.1931}{P19}{Tämän johdosta mo-B nen talonpojan oli pakosta liityt- B tävä yhteistalouteen pelastaak-B seen itsensä ja perheensä nälkä- B kuolemasta. Puhumattakaan tunnetuista Pohjois- Venäjän ja Solovetskin karkoitusleireistä, on viime aikoina ruvettu käyttämään pakkotöitä kaikkialla valtion työmailla. Jokaisen henkilön tehtäväksi on annettu sahata ioo m 8 halkoja ja hevos- I miesten on ajettava 200 m 8 hal- I koja rautatieasemille.}
\kat{Suomen Heimo}{16.06.1931}{P28}{Yksityisen yritteliäisyyden pohjalla syntynyt taloudellinen nousu merkitsi bolshevikien mielestä entisten yhteiskuntapoliittisten voimien keräämistä, jolle oli ajoissa asetettava este. Useat Inkerin parhaimmista miehistä ovatkin saaneet kuulan kalloonsa tai viruvat nälkäisinä ja kauhean rääkkäyksen alaisina Siperian ja Solovetskin vankileireillä. Tästä onkin ollut seurauksena yleinen tavarapula, jota lisää vielä se, että vähistäkin tavaroistansa koettavat bolshevikit viedä ulkomaille " huolimatta siitä, että oma maa näkee nälkää.}
\kat{Toukomies}{01.02.1931}{P8}{Moni huokaisee epätietoisena tuntemattoman taistelun tuloksesta, mutta samalla monen rinnassa syttyy palamaan uskon tuli siitä, että tämäkin taistelu päättyy voittoon ja vie kappaleen eteenpäin asiaamme. Sanomat tammik. lopulla, kertoen m. m. Kyminlinnan pakolaisleirille karjalaisten keskuuteen sijoitetuista Solovetskin metsätyöorjuudecta paenneista umpivenäläisistä y. m. karkoitusvangeista, joista osa on tavallisia rikollisia. Kalastusluvan saannin edellytyksenä on, että valtiolle suoritetaan veroa 5 verkolta tai rysältä 1 rpl., 10 verkolta 2 rpl. j. n. e. Koukuista, uistimista ja muista pyydyksistä on suoritettava myös veroa.}
\kat{Toukomies}{01.04.1931}{P17}{Karjalaisten Joulu- Sanomat 1919 n:o 23 — 24. ) ja että sieltä päin kulki tie Vie'in kautta Juukaan yli Pielisen ( Heikki Kokkonen, » Maantapa ». 2 ) Tässä on huomautettava Paalasmaan saarien välillä olevan salmen nimestä » Sirnitsansalmi », joka viittaa Solovetskin luosta- rin valtakauteen ( Heikki Kokkonen, » Sirnitsa ». Lieksalaiset ja viekiläiset ovat myöhemmissäkin muistoissa tunnetut, niinkuin yleensä Pielisen itäpuolen asukkaat, suurista monihankaisista veneistään ja näillä etelään päin tekemistään pitkistä retkistään.}
\kat{Toukomies}{01.04.1931}{P35}{3Battl ) an fanfan elämää opitte tuntemaan fau- neista furoateoffistamme — Niitä aikoja, jolloin nykyiset va nhu k- set olivat nuoria, kuvailee tutkijana ja valokuvaa- jana yhtä etevä maisteri AHTI RYTKÖNEN tavattoman viehättävässä kuvateoksessaan SAVUPIRTTIEN KANSAA Laakapainostamme juuri ilmestynyt teos sisältää 144 taiteellista jäljennöstä tekijän ottamista valo- kuvista. Niinpä mainitsemassamme Raja-Konnun kylässäkin on kymmeniä taloja käsittävä kolhoosi ( yhteistalous' eli kommuuna ), jonka hallussa ovat nyt kaikki kylän parhaat viljelysmaat ja nurmet sekä yksityisiltä ryöstetty omaisuus, kuten karja, työvälineet y. m. Työt tässä kommuunassa ovat järjestetyt suurtilan malliin dots}
\kat{Toukomies}{01.02.1932}{P9}{Tähdet syttyivät kirkkaina taivaalle ja kuu näkyi kelmeänä nousevan. Hän jäi paikoilleen, kunnes hänet laahattiin Solovetskin vankilaan. Hän oli pitänyt kansankielellä joskus saarnoja ja raamatunselityksiä kansalle.}
\kat{Toukomies}{01.02.1933}{P7}{N:o I — 2 - Niin syntyi sinne kuulu Solovetskin luostari. Pitäjissä vallitsi näihin aikoihin täydellinen käskyvalta.}
\kat{Toukomies}{01.02.1933}{P7}{Saarella oli hän nähnyt rakennettavan asuntoja sekä kuullut puhuttavan outoa kieltä; ja takaisin palattuaan oli hän kerännyt miehiä hyökkäykseen, mutta ei ollut onnistunut karkoittamaan saaren asukkaita saareltaan. Kaukonäkemisen avulla oli hän nähnyt, että Solovetskin saarelle oli tullut outoja mustia miehiä ja niinpä hän saari saarelta, luoto luodolta ui neljäkymmentä virstaisen vesitaipaleen saadakseen varmuuden oudoista eläjistä saarella. Tämä melkein ijäisyydenpituinen sotapalvelus oli erittäin raskas karjalaisille, kun he olivat tottuneet metsien vapaihin oloihin ja palvelemaan vain omia jumaliaan ja tottelemaan vanhoja heimopäälliköitään: Soti-Trohkimaa, Ahmaa- Rokatsua y. m.}
\kat{Toukomies}{01.12.1933}{P7}{Siinä perustettiin Rurikin luoma Venäjä, vaikka sen vallanpidon keskus pian siirtyi siitä etelämmäksi, Kieviin. H Hän esiintyy yhtenäH Karjalan orjuutta-H jista, hänelle luovu-H tetaan suuria maitaH pohjoisessa Vienassa, H jopa Solovetskin luos-H iarikin. Novgorodin Ilja-! munkki kastoi ensimmäisinä Karjalan kansaa, Novgorodista tai sen tahdosta tulivat Valamon, Solovetskin ja enimpäin muiden Karjalan luostarien perustajat.}
\kat{Toukomies}{01.12.1933}{P7}{H Hän esiintyy yhtenäH Karjalan orjuutta-H jista, hänelle luovu-H tetaan suuria maitaH pohjoisessa Vienassa, H jopa Solovetskin luos-H iarikin. Novgorodin Ilja-! munkki kastoi ensimmäisinä Karjalan kansaa, Novgorodista tai sen tahdosta tulivat Valamon, Solovetskin ja enimpäin muiden Karjalan luostarien perustajat. Novgorodin — tai niinkuin venäläiset sanovat, » suuren Novgorodin » — kukistuminen on kaikkein mielenkiintoisimpia itäisen Euroopan, myöhemmin Venäjän valtakunnan kansapiirin historian kerrottavista.}
\kat{Karjalan Vapaus}{09.1934}{P7}{Hän on kerran painuva maata vasten! Solovetskin uhrien kuolinhätä Ja silloin kavahtakaa!}
\kat{Karjalan Vapaus}{09.1934}{P7}{On kerran eessänne Kain! Solovetski, Karjalan hirsipuu, Solovetski, pyöveliparatiisi, sinä koston istutit uumeniisi.}
\kat{Karjalan Vapaus}{09.1934}{P7}{Solovetski, Karjalan hirsipuu, Solovetski, pyöveliparatiisi, sinä koston istutit uumeniisi. Pian väkevänä on vastassanne se, jolla kahlittu, kova ranne.}
\kat{Karjalan Vapaus}{09.1934}{P11}{Vuosikymmen, ylikin on kulunut jo siitä sana tuo kun miesten mielet vältas. » Ei Solovetskin tuntematon, synkkä tie pois miesten mielist' ole sanaa tuota saanut. Laulu lakkas, sammui Sammon taika, kotiliess' kylmäksi nyt jäi.}
\kat{Karjalan Vapaus}{12.1934}{P6}{Se innoittaa mieltämme ja kohottaa katseemme tulevaisuuteen sekä käskee paimenien tavoin julistamaan päivästä päivään joulun suurta ilosanomaa » joka on tuleva kaikille kansoille ». Solovetskin " monasteri " ennen bolshevikivallankumousta Oi ystävä!}
\kat{Karjalan Vapaus}{12.1934}{P30}{Ja apu useimmissa tapauksissa tulikin, niin ihmeelliseltä kuin se tuntuneekin. ■ — Lapsi parani. Solovetskin luostari, joka nykyään kelluu " helvettinä " Vienan merellä, oli hyvin suosittu käyntipaikka. Kirkossa kävijät eivät montakaan sanaa ymmärtäneet papin puheesta, joten kirkon taholta tuleva valistus jäi pakostakin hyvin pintapuoliseksi.}
\kat{Toukomies}{01.03.1934}{P9}{Sen jälkeen hän painui laajoihin ja. osaksi eriskummallisiin paikannimi- ja muinaistutkimuksiinsa, julkaisten ne kahdessa Suomi-kirjassa nimellä » Tietoja suomalais-ungarilaisten kansain muinaisista olopaikoista » ( 1868 — -1870 ), joissa kosketellaan myös karjalaisia alueita. Vuosien 1842 — 1844 » Suo mi » -kir joissa julkaisee M. A. Castren esityksensä Solovetskin luostarikronikasta ja » Savolotscheskaja Tcshudeista », jotka sittemmin julkaistiin myöskin hänen kootuissa teoksissaan » Nordiska Resor och Forskningar » ( 1852 j. n. e. ). ( Kaikki kuvatut sivut esiintyvät noin puolen kokoisina. )}
\subsubsection*{pakkotyö \normalfont\small(25 katkelmaa, 35 osumaa aineistossa)}
\kat{Toukomies}{01.03.1931}{P15}{— » Mikael 11, mustalaisten kuninnas. » ■ — Venäjäkin muka tutkimaan pakkotyön käyttöä. Vashihgtonista ilmoitetaan ( U. S:n mukaan ) helmik.}
\kat{Toukomies}{01.03.1931}{P15}{Tammikuun 28 p:nä Moskovassa varustelukomitean istunnossa annetussa selostuksessa on todettu, että erikoisneljänneksen suunnitelmasta koko mahtavan valtioliiton alueella on täytetty vain 72,4 ° / 0. San. kirjeenvaihtaja sähköttää Riiasta 24 / 4 Moskovasta ilmoitetun, että on asetettu erikoiskomitea Sinovjevin puheenjohdolla tutkimaan pakkotyön esiintymistä kapitalistisissa maissa ja laatimaan asiasta raportin. Tästä huolestuneena varustelukomitea on katsonut välttämättömäksi » ratkaisevan käänteen » aikaansaamisen metsätöissä, jotta huhtikuun 1 päivään mennessä olisi suunnitelmasta yleensä toteutettu tuo häämöttävä 70 prosenttia.}
\kat{Toukomies}{01.09.1931}{P11}{Kärsimysten malja oi jo tästä kaikesta kukkuroillaan; kuitenkin kärsii tällä hetkellä kokonainen kansakunta, Karjalan heleä heimo, sydänjuuriinsa saakka enimmän siitä, ettei tiedä, ei aavista nykyistä kovaa kohtalonhetkeään kauemmaksi, sanalleen kuten erään runon sanat sattuvat: Oi synnyinmaa, mik' osakses' lie luoto, iloko lie vai murhe sulle suotu — salattu kunut yhtämittaa tuskainen hätähuuto korviimme Aunuksen ja Vienan mailta, mihin on liittynyt äskettäin sydäntäsärkevä Inkerinkin itku, tuhansiin nousevin joukoin olemme nähneet heidän pakenevan rajan tälle puolen hakemaan hengen turvaa läheisen heimokansan hoivissa, ja tietoomme on tullut, että ne herkkäuskoiset, jotka sortajan dots}
\kat{Suomen Heimo}{28.07.1932}{P7}{johtuvaa, että suurin osa tästä kirjallisuudesta käsittelee teollisuutta, onhan teollisuuden kehittäminen jo vallankumouksesta saakka ja etenkin 5-vuotissuunnitelman aikana ollut Venäjän yhteiskuntauudistuksen keskus ja epäilemättä vaikein ja ulkomaidenkin etuja lähimmin koskeva tehtävä. Suorittamatta jääneestä verosta määrätään työvelvollisuutta, esim. halkojen valmistusta 230 *. Kun näin rangaistu ei kykene kotoa käsin lätä työtä määräaikana suorittamaan, niin myydään omaisuus ja " velallinen " lähetetään toiselle paikkakunnalle pakkotyöhön. Niin kauan kuin kukaan ei ole julkaissut bolshevistis-porvarillista sanakirjaa, joka selittäisi kaikki salasanat tavalliselle porvarille dots}
\kat{Toukomies}{01.10.1932}{P11}{Volgan ja Marian kanavan välisen vesitien uudestijärj estäminen edellyttää linjan, joka yhdistää Volgan ja Pietarin. Työssä on noin 250,000 pakkotyö! äistä, jotka on jaettu tämän 250 kilometriä pitkän kanavan eri jaksoihin. Mainittu pato on myöskin iso ja järjestetään sillä Telokan joen ja muutamien pienempien järvien vedenpinta Matkaj arvelle saakka.}
\kat{Suomen Heimo}{24.02.1933}{P7}{Kun itse olimme venäläisten sorron alla sanoimme me: " VieV uusi päivä kaikki muuttaa voi. " V. 1927 lähtien on Karjalaan karjalaisväestön keskuuteen tuotu satojatuhansia venäläisiä pakkotyöhön metsiin ja kanavoimistöihin. Tuossa joukossa on professoreita, insinöörejä, piispoja, pappeja, opettajia ja muuta sivistyneistöä sekä rahvasta.}
\kat{Karjalan Vapaus}{12.1934}{P33}{Hän ei kironnut maailmaa ja siksi oli Hänen ylösnousemus, maailman vapahdus ja ijänkaikkinen elämä. Pakkotyöleirin ylin työnvalvoja ilmestyi ärtyisänä alimpien työnvalvojien eteen, jotka odottivat hänen määräyksiään. Sitten vanhusta tärisytti outo vavahdus yli koko ruumiin ja oikea käsi kohoutuu hitaasti ja tekee ristinmerkin.}
\kat{Karjalan Vapaus}{12.1934}{P34}{— Karkoitetun näännytetty henki kuoleman lähestyessä ei tunnu uhrilta, vaan pääoman säästöönpanolta. perillisille sillä eikö vanhempi aina säästä lapsilleen. Tämä on jouluyön evankeeliumi — minulle pakkotyöleirin sairaalle vanhukselle — se on vapahduksen sanoma kiduttetun vanhuksen ansaittu ijänkaikkisen rauhan ja levon lupaus. — Makarielan Kirlän sairastusvuoteelta, jonka jo lumi on melkein kokonaan peittänyt, lähtee ilmoille heikko korahdus, joka häipyy heti tuulen hornamaiseen vihellykseen ja lumimyrskyn pauhuun.}
\kat{Suomen Heimo}{04.06.1934}{P8}{Sen pää kävi Venäjällä, ja katso, kun hän sieltä tuli, niin hän oitis kirjoitti lehteen, että meidän pitäisi kiiruhtaa tekemään enemmän kauppaa Venäjän kanssa, koska suurvallatkin niin tekevät. Ja toiseksi siinä, mitä Venäjältä meille ostetaan, puutavarassa, viljassa ja Siperian voissa, on sellaista, mitä oma maa meille antaa ja antaisi ainakin yhtä halvalla kuin pyhä Venäjä, jos meilläkin olisi pakkotyö systeemissä. " Etelä-Suomen Sanomat " sietäisivät kiihoituslakia, mutta kun se 011 säädetty toisia tahoja varten ja sen käyttö, kuten puserolainkin, riippuu siitä ( lue Suomen Sosialidemokraatista, että on kysyttävä, mitä edustaja Erkko tarkoitti puserolailla, ennenkuin sitä sovelletaan ), dots}
\kat{Itä-Karjala}{11.1935}{P7}{Mutta turhaan neuvostolehdistö koettaa peitellä sitä, mistä tässä puhdistuksessa todella on kysymys. Löivät he orjaksi Karjalan kansan, Heittivät kaulaan kavalan ansan, Veivät he veljemme pakkotyöhön. Senpä takia ei myöskään nyt viimeksi tapahtunut voi olla herättämättä mielipiteitten vaihtoa rajan tällä puolen.}
\kat{Karjalan Vapaus}{03.1935}{P5}{Valokuvahan tuleekin parempi liikkumattomista, esineistä, varsinkin huonolla koneella talvisena pimeänä aikana näpätessä. Pakkotyö voiman käyttämisestä metsätöissä on muodollisesti luovuttu, mutta tosiasiallisesti ovat kaikki kansalaiset neuvostoteissa Harjoitetaanpa suoranaista metsien ryöstöä, koska esimerkiksi viime vuoden hakkuumäärä nousi 18 \% suuremmaksi vuotuista metsien lisäkasvua.}
\kat{Karjalan Vapaus}{04.1935}{P16}{Vaimo laski hellävaroin kantamuksensa maahan mättäälle ja aukaili sen suojana olleita vaateryysyjä. Se tiesi matkaa Pietariin vangittuna ja sitten Siperiaan pakkotyöhön siksi kunnes siellä menehtyi. Siitä olivat urkkijat saaneet vihiä, ja jonkun ajan kuluttua tultiin vanha isäntä pidättämään.}
\kat{Karjalan Vapaus}{04.1935}{P16}{He olivat kumpikin laihoja, repalaisia ja nälkiintyneitä, ja heidän silmissään asui takaa-ajetun eläimen arka, pelokas katse. Mutta kaikki olivat kuitenkin varmoja siitä, että hänet oli viety jonnekin pakkotyöhön, luultavasti Siperiaan, kuten niin yleisesti tehtiin. Siihen mättäälle vaipui äiti huokaisten syvään ja elonmerkkiä enää osoittamatta, puristaen kuolleen lapsensa elotonta ruumista suonenvedontapaisesti rintaansa vasten.}
\kat{Karjalan Vapaus}{05.1935}{P21}{Viimeksimainittuna vuonna perusti oman kangas- ja sekatavarakaupan, jota hän edelleenkin harjoittaa. Kun 1920 joutui punaisten pakkotyöhön, pakeni Suomeen, sai liikeapulaisen paikan kauppias J. Peilamolla Juuvassa. Tultuaan Suomeen 1919 syyskesällä pääsi Suomalainen puuliike Oy:n palvelukseen ollen sillä konttoristina ja työnjohtajana vv.}
\kat{Karjalan Vapaus}{07.1935}{P8}{Mutta sitten vähitellen rupesin selkenevässä tapunnassani ymmärtämään asemani, meidän kaikkien aseman, toivottomuuden ja synkkyyden. Heimolaistemme karjalaisten kohtalosta lukiessa ei meidän, joille heidän asiansa on rakas ja kallis, sydämemme voi olla tuntematta mitä syvintä sääliä noita onnettomia veljiämme ja sisariamme kohtaan, jotka on viety luultavasti pakkotyöhön kauas tuntemattomille seuduille — nääntymään ja kuolemaan. Toivomme, että lehtemme saisi osaltaan auttaa vanhojen tietojen keruutyötä, tietojen, jotka sisältyvät taikoihin, sananparsiin, sananlaskuihin, tapakuvauksiin, elämäkertoihin, lauluihin, itkuvirsiin, paikannimiin, kielimurteeseen jne. Olem me vakuutetut, että dots}
\kat{Suomen Heimo}{07.02.1935}{P15}{S. H. 2 / 35 = Suomen Heimo N:o 2 vuonna 1935. ) Koska kirja, johon lisätietoja liimataan, vähitellen voi tulla kovin paksuksi ja lisätietoja voi vuosien varrella kertyä päällekkäinkin, on AKS varannut halukkaille mahdollisuuden. tilata itselleen " Itä- Karjala " -kirjan välilehdillä varustettuna ja tavallista löyhemmin sidottuna. — Vaazeni 5 Revselga — länteen ( pakkotyöleiri halonhankintaa varten Pietariin ) 2 Kontupohja — tehdas ( sekä kapea- että leveäraide ) 5 Karhumäki — satama '. 2 Karhumäki — pohjoiseen ( pakkotyöleiri ) 2 Maaselkä — Merimaaselkä ( kanavalle ) „ 32 Maaselkä- — Liistepohja ( Seesjärvi ) < 10 Parantova — luoteiseen ( pakkotyökin ) 15 Sorokka — ulkosatama Rasnavolok 10 dots}
\kat{Suomen Heimo}{07.02.1935}{P15}{S. H. 2 / 35 = Suomen Heimo N:o 2 vuonna 1935. ) Koska kirja, johon lisätietoja liimataan, vähitellen voi tulla kovin paksuksi ja lisätietoja voi vuosien varrella kertyä päällekkäinkin, on AKS varannut halukkaille mahdollisuuden. tilata itselleen " Itä- Karjala " -kirjan välilehdillä varustettuna ja tavallista löyhemmin sidottuna. — Vaazeni 5 Revselga — länteen ( pakkotyöleiri halonhankintaa varten Pietariin ) 2 Kontupohja — tehdas ( sekä kapea- että leveäraide ) 5 Karhumäki — satama '. 2 Karhumäki — pohjoiseen ( pakkotyöleiri ) 2 Maaselkä — Merimaaselkä ( kanavalle ) „ 32 Maaselkä- — Liistepohja ( Seesjärvi ) < 10 Parantova — luoteiseen ( pakkotyökin ) 15 Sorokka — ulkosatama Rasnavolok 10 dots}
\kat{Suomen Heimo}{07.02.1935}{P17}{178. Eräs pakkotyömuoto on se, että henkilöllä, jota ei muuten voida tuomita, teetetään puolipakolla työsitoumus, urakka. Ellei hän sitä pysty täyttämään, tuomitaan sakkoa 300 — 400 ruplaa, karkoitettavaksi tai I — 21 — 2 vuoden pakkotyöhön. Kysymyksen tätä vaihetta selostetaan tarkemmin kahdessa kirjoituksessa " Tartto Rediviva ", Suomen Heimon numeroissa 9 ja n v. 1934.}
\kat{Suomen Heimo}{26.02.1935}{P30}{Kotonani sain olla noin 4 vuotta, kun taas eräät vihamieheni kantelivat päälleni, että minä muka olin eturivin miehiä Aunuksen vapaustaistelijani riveissä ja että vieläkin olen mielipiteeltäni " valkoinen "... Kun olin ollut tällä työmaalla 2 vuotta ja 2 kuukautta, käännyin pakkotyöleirin johtohenkilöiden puoleen ja ilmoitin heille, että vankeusaikani on jo ajat sitten päättynyt ja eikö minua voitaisi vapauttaa ja päästää Karjalaan. Jos vanki sairastui, sai hän lojua ahtaassa parakissa niin kauan kunnes kuoli, sillä paranemisesta. ei ollut koskaan toivoa, kun joutui hoidon- puutteessa virumaan nälkäisenä ja viluisena kylmissä ja kosteissa suojissa.}
\kat{Suomen Heimo}{16.09.1935}{P14}{Kuluvan vuoden toukokuussa Suomen hallitus tiedusteli Neuvostohallitukselta, miksi Karjalan ja Inkerin suomalaista väestöä suurin joukoin kuljetettiin Siperian tuntemattomiin osiin. Helmikuussa 1933 yhdessä ainoassa Aunuksen piirikunnassa lähes sata henkeä tuomittiin vaihteleviin rangaistuksiin — aina kymmenen vuoden pakkotyöhön saakka.,. Joku saattaa ihmetellä, miksi Neuvostohallitus on niin innokas pitämään tämän kapean alueen, joka on osaksi vettä, osaksi suota ja osaksi karua ja hedelmätöntä maata.}
\kat{Suomen Heimo}{16.09.1935}{P16}{Suomalainen kommunisti Antikainen jättää oikeudelle valituskirjelmänsä huomenna Helsingissä. Antikainen todettiin syylliseksi ja tuomittiin elinkautiseen pakkotyöhön. Kommunismin vastustaminen on sentähden Suomessa mitä läheisimmin kytketty puolustuskysymykseen.}
\kat{Suomen Heimo}{16.09.1935}{P16}{Väitettiin myös, että hän oli saanut venäläisen kunniamerkkinsä juuri tästä teosta. Hänet pidätettiin ja tuomittiin kahdeksaksi vuodeksi pakkotyöhön luvattomasta maahan saapumisesta. Ehkäpä Venäjän hallitus torjuu tämän innoittavan syytteen ja täten auttaa vankia vapauteen.}
\kat{Suomen Heimo}{16.09.1935}{P16}{Karjalaiset — jotka ovat alkuperältään suomalaisia ja joiden onnettomuudekseen on täytynyt jäädä Venäjän puolelle rajaa — ovat tuskin kärsineet vähemmän Neuvostohallituksen alaisina, jotenka suomalaiselle sana 'ryssä " merkitsee " sortajaa " nyt kuten aina. Toisessa paikassa lehteä julkaisemme yksityiskohtia kommunisti Antikaisen oikeudenkäynnistä ja piakkoin jätettävästä valituksesta, Antikaisen, joka tuomittiin elinkautiseen pakkotyöhön mielikuvituksellisen, keksityn syytteen nojalla ainoastaan, jotta saataisiin aineistoa ankaraa kommunisminvastaista kiihoitusta varten. KESKUS-PAPERIKAUPPA KONTTORI- KIRJOITUS- KOULU-TARPEITA PAPERIA. PAPERITAVAROITA ', \textasciicircum{} HELSINKI, ALEKSANTERINKATU 44 K. dots}
\kat{Itä-Karjala}{12.1936}{P7}{Toimittajat kyllä tulevat tekemään velvollisuutensa — kunhan vain me sen teemme. Se on ajettu metsään pakkotyöhön, ja kotoiset pellot ovat kylminä ja talot on poltettu. Erittäinkin vuoden lopulla ja juuri nyt joulujuhlan aikana Karjala-kysymyksessä astuu esiin valoisampia puolia.}
\kat{Karjalan Vapaus}{03.1936}{P37}{— V. 1919 otin osaa Aunuksen vapaustaisteluihin n.s. suomalaisen retkikunnan riveissä. Tämän tuomiopäätöksen mukaisesti sain elokuussa 1925 lähteä Solovetskin kauhunsaarille pakkotyöhön. Solovetskissa oli tullut tavaksi kaivaa haudat kuolleita varten jo syksyllä sulan aikana.}
\subsubsection*{karkotus \normalfont\small(25 katkelmaa, 168 osumaa aineistossa)}
\kat{Toukomies}{01.06.1925}{P4}{Vuosien 1911 — 17 välisestä ajasta antaa » Karjalan Sivistysseuran » perustavan kokouksen pöytäkirja seuraavan oloja ja aikoja kuvaavan karakteristiikan: » On tuntunut siltä, että huomattavampi toiminta vahingoittaisi sekä itse rajantakaisen Karjalan että vielä Suomenkin asiaa. Kirjoittajan viittaus tarkoittaa epäilemättä syyskuussa 1907 tapahtunutta kolmen karjalaisen kansanvalistusmiehen laitonta vangitsemista, joista kaksi Vasili Jeremejeff-Räihä ja H. Remsujeff sittemmin karkoitettiin Temirin kaupunkiin kirgiisien aroille, missä, kuten muistetaan, Räihä sittemmin, 2. XI. 1908, kuoli keuhkotautiin. Karjalan Sivistysseura on päättänyt lähettää kerhon jäsenille » Toukomiehen » nyt dots}
\kat{Toukomies}{01.06.1925}{P12}{Paikalliskerhot nimittävät itseään sen ja sen paikkakunnan Karjalan-kerhoksi ja keskustoimikunnan nimenä on Karjalan Kerhojen Keskustoimikunta. Ja edelleen hän kirjan aiheesta lausuu: » Vaikka aikaa ja älyä on paljon tuhlarttu, ei Kristusta vieläkään " ole karkoitettu maan päältä. Vielä päätettiin, että kerhot voivat, missä se mahdolliseksi osottautuu, kantaa pientä jäsenmaksua jäseniltään, josta on mainittava kerhon säännöissä.}
\kat{Toukomies}{01.1926}{P11}{Heidän anomuksensa on Viipurin läänin maaherra palauttanut Sortavalan maalaiskunnalle lausunnon antamista varten ja tämä on yksimielisesti päättänyt evätä puoltolauseen antamisen. Sitten kohta sen jälkeen karkoitettiin yksi pahimpia viime syksyisiä kiihottajia, jokin munkki „Polykarpus ", suoraan omaan maahansa, Venäjälle. Asian saadessa tämän käänteen, ovat muutamat munkit, s. o. osa niistä, jotka ovat vielä Venäjän kansalaisia, ruvennut pyytämään pääsyä Suomen kansalaisiksi.}
\kat{Toukomies}{01.1926}{P11}{maanpinta rupeaa vähitellen laskemaan etelään ja kaakkoon päin loivarinteisinä laaksoina ja yhä pehmeäpiirteisemmäksi ja matalammaksi käyvinä mäkinä, kunnes Suomenlahden rannalla tapaa se melkoisen mataluuden ja samalla yhä suuremman vehmauden, lähinnä Itä-Uudenmaan maisemia muistuttavan. Ensinnä karkoitettiin maasta jo aikaisemmalla karkoituksella maasta poistettaviksi määrätyt viisi pahinta, munkkikiihottajaa — joillakin munkkinimillä Joannik, Teodosij, Iraklij, Juvenalij ja Teofil varustettuja, mutta tosiasiassa tavallisia venäläisiä mushikoita —, Ulotiuupa tämä kaarre — jos Sortavalan jyrkästi vuoriset, vaikka savikkolaaksoiset seudut jätämme pois — taas sen taakse aina Salmiinkin dots}
\kat{Toukomies}{05.1926}{P7}{Seuran jäsenluku, vuotuinen kirjeenvaihto, etupäässä vapaaehtoisina lahjoituksina saadut tulot — m. m. Kontron ja Kuosmasen 100,000 markan suuruinen lahjoitus — ja moni muu seikka osoittavat melkoista uhrimieltä ja toimeliaisuutta. Viitattu teko, jota ei tahdottu todeksi uskoa, oli kolmen Vienan karjalaisten liiton ylläpitämän kansakoulun sulkeminen ja kolmen kansallismielisen karjalaisen vangitseminen, joista kaksi: nuori kunnallis- ja kansanvalistusmies Vasili Jeremejeff Räihä ja eräs toinen karjalainen H. Remsujeff karkoitettiin kauas kirgiisien ahoille, Aasiaan. Odottamattoman hyvällä menestyksellä karjalaiset tähän saakka ovat kestäneet pitkällisessä taistelussaan kansallisuutensa dots}
\kat{Toukomies}{06.1926}{P10}{Htenkin ensinmainittu on ansiokas kuvaus siitä, mitä korkeissa nykyisvenäläisissä diplomaattisissa piireissä liikuskellut henkilö on saanut nähdä. Kun Tihonin seuraaja Pietari Moskovasta karkoitettiin, niin muodostettiin kirkollinen neuvosto, joka otti hoitoonsa tihonilaisen kirkon johdon. Pienemmät Tuleman-, Kirkko- ja Miinalanjoki ovat tyypillisiä Salmin savikkorantojen jokia, osittain tuovia Aunuksen puolelta.}
\kat{Toukomies}{12.1926}{P21}{Kustaa II Aadolfin ja Mikael Romanovin rajankävijät lyövät merkkinsä Variskiveen Laatokan itärannalla ja tämän rauhan on määrä tulla » ikuiseksi ». Hänen oli tapana sanoa: » En kuole, ennenkuin ryssä on maasta karkoitettu », ja hän saikin elää niin kauan! Tällaisen kohtalon varalle oli pappien tapana oppia jotakin käsityötä, sillä maanpakolaisuudessa ei saanut henkisellä työllä itseänsä elättää!}
\kat{Toukomies}{01.03.1927}{P7}{Siks iloni mun, minä syleilen teitä kuin vaivasta-vapahtajaani. On viety aina Siperiaan asti, karkoitettu Pohjoiselle Jäämerelle. 'Itse Inkeri on vielä vuoden ollut täydellisesti eristetty Suomesta.}
\kat{Suomen Heimo}{20.01.1928}{P12}{On ilmestynyt joukko liiviläisiä kirjoja, ön noussut joukko liiviläisiä runoilijoita, liivin kieltä opetetaan kuudessa koulussa, Margrete Stalte johtaa laulukuoroa, joka rohkenee tehdä laulumatkoja ulkomaillekin, kaksi valittua liivi-' laista opiskelee Suomen kouluissa ja Liivin liitto toimii innokkaasti. Erään toisen kylän opettaja Polman käänsi sen itäliiviksi — ja tämä Polman tuli sitten liiviläisen kansallisuusaatteen ensimmäiseksi marttyyriksi; hänet karkoitettiin kauas Venäjälle Novgorodin soille ja siellä hän- kuoli kolkossa yksinäisyydessä; ja hänen kaikki runonsa hajoitettiin likaisen ryssäläisen kylän tielokaan. Kun sen valmistus tähän mennessä on erinäisistä syistä ollut ra- dots}
\kat{Suomen Heimo}{20.01.1928}{P19}{Suurta surua on Raumalla herättänyt radiossa kuulutettu uutinen, että » Voima " on Rauman edustalla turhaan pyrkinyt Mäntyluotoon. Neuvosto- Venäjänpuolisesta Kullankylästä on hallinnollisina vankeina karkoitettu Siperiaan neidit Anna Olavi ja Hilda Reijo, edellinen Omskiin, jälkimäinen Novo-Simbirskiin. Nyt on koulun Osoaviahim-jatseikka ( = ilma- ja kemiallisen sodan yhdistyksen. solu ) lahjoittanut rahat » Vastauksemme Ghamberlainille " -rahastoon, siis sotatarkoituksiin!}
\kat{Suomen Heimo}{20.02.1928}{P5}{Tähän liittoon, jonka puheenjohtajana oli puolalainen Lédnicki, kuului myöskin virolainen duuman jäsen Jaan Tönisson. Tultuaan otetuksi vangiksi illegaalisen toiminnan vuoksi hänet karkoitettiin Kasaniin, mutta hän pakeni ja saapui Geneveen. Venäjän vähemmistökansallisuuksiin kuuluvat duuman jäsenet perustivat keskuuteensa n.s. federatistien liiton, joka ajoi Venäjän vähemmistökansallisuuksien autonomiaa.}
\kat{Suomen Heimo}{20.05.1928}{P15}{Kansanopetus on se astinlauta, jolle bolshevistinen hallitusjärjestelmä on lujimmin tallannut. Entiset opettajavoimat karkotettiin maasta vastavallankumouksellisina tai muka kansan sielun turmeli joina. Kommunistien lukumäärä Inkerissä on sangen vähäinen, niin vähäinen, että neuvostotilastot ovat vaienneet siitä kymmenen vuoden ajan.}
\kat{Toukomies}{01.10.1928}{P8}{Kun seura erosi, antoi Falkman Vornaselle alfeniidikannun ja vaikka nimenomaan huomautettiin, että kannu ei ollut hopeaa, herätti se sekä Vornasessa että hänen kotiseudullaan suurempaa huomiota kuin hopeakannu. Ostettiin saarrettu karhu pesästään ja kun sitten talviuninen otus oli karkoitettu ilmoille ja nopeasti ja hupaisesti passitettu parempiin maailmoihin, alkoivat monipäiväiset ilonpidot ja kemut sen kunniaksi. Tiedon puussa on sekä huono että hyvä hedelmä ja me voimme ymmärtää Vornasta, kun hän ajatteli, että ne virkamiehet ja liikkeenharjoittajat, jotka esiintyivät hänen seudullaan, olivat saaneet maistaa enemmän edellisestä, kuin jälkimäisestä hedelmästä.}
\kat{Suomen Heimo}{31.01.1929}{P22}{( Me, rajantakaiset karjalaiset, niin valistumattomia kun olimmekin, ymmärsimme, ettei ole niinkään tärkeä, milloin pyhä vietetään, kuin se, ni it e n se vietetään, ja siirryimme ilman muuta uuteen, joka on oikeampi ja käytännöllisempi ). Joitakin vuosia sitten sai harva se päivä nähdä sanomalehtien palstoilla seuraavanlaisia otsakkeita: > Valamoa rettelöivät munkitc > Joukko Valamon munkkeja karkoitettu ulkomaille< >Munkit tervanneet johtajansa asunnon ovenc » Mathilda Wrede valittanut munkkien puolesta kansainliittoon * j. n. e. Sattuipa tässä päivänä moniahana kouraamme paksuhko, noin kolme ja puoli sataa sivua käsittävä teos, jonka julkisivulla, komeili seuraavanlainen teksti: » Uuden dots}
\kat{Suomen Heimo}{25.05.1929}{P20}{Heimopakolaisten asema ja kohtelu maassamme ei millään seudulla ole erikoisen loistavaa, mutta sellainen menettely, joka viranomaisten harjoittamana nyt on saatu päivän valoon, ei ole sukuakaan oikeusvaltion virkamiesten toiminnalle, vaan korkeimmassa määrässä rikollista, ei ainoastaan sen uhriksi joutuneisiin pakolaisiin nähden, vaan myöskin koko heimoystävyysaatetta kohtaan. Oikeudenkäynnin kestäessä ilmieni lisäksi, ettei nimismies Karlsson eivätkä poliisikonstaapelit voineet edes väittää minkään rikollisen teon jäävän Busareffin syyksi, vaan että siis tätä ilman laillista syytä oli kauan pidetty vangittuna, ja hänet sitten karkoitettu miltei varmaan kuolemaan, josta hän kuiteinkin oli dots}
\kat{Suomen Heimo}{15.09.1929}{P9}{Edellisen päivän sateet olivat saaneet isäntämme eivät olleet aikoinaan paenneet rajan tälle puolen tai joita ei vielä oltu karkoitettu Solovetskiin ja Siperiaan, kuten esimerkiksi mielen synkäksi. Ja kun juhlä-aika läheni, täyttyi juhlakenttä, korkea paikka aivan rajan vieressä, hartaasta Inkerin kansasta.}
\kat{Toukomies}{01.12.1929}{P15}{Miksi Suomessa pitää olla kaikkialla kaksikieliset kyltit, kun » sivistyneessä » Ruotsissa ei ole umpisuomalaisellakaan alueella muuta kuin yksikieliset, ruotsalaiset? Hän oli nuorempana ollut upseeri, mutta juopottelun ja muun hurjan elämän vuoksi oli hänet upseerin toimesta erotettu, karkoitettu Karjalaan ja pantu papiksi. M. Terho.}
\kat{Toukomies}{01.12.1930}{P21}{niin voimat vehmaat valkoalttarille antoi ja valkoparvet paisui, kävi kaikki mukaan, ne kalleimpansa eestä kansan kantoi, ei henkeänsä kartellunna heistä kukaan. jo vapahaksi pettureista vallan pääsi ja himot herjain meni hukkaan mahdotonna ja vallan vieraan ikeen yltään maa se haasi ja keskellämme karkotettiin konna. Saat olla pojistasi noista ylvään ylväs ja historiaasi luoda heistä lehden oivan: ne valkomiehet ovat eelleen suojapylväs ja suovat sulle, kuten silloin, turvan, hoivan.}
\kat{Suomen Heimo}{20.03.1931}{P17}{Turaanilaisen Seuran toiminta koskee suomalaisia siinä määrässä, että sen yhteydessä on erikoinen suomalaisvirolainen osasto johtajana tri Aladår Bå n ja näin suomalais-ugrilainen heimoliikekin saa juuri Turaanilaisen Seuran kautta laajempaa kannatusta Unkarissa. Siellä, missä talonpoikien pienet maatilat joutuivat osaksi rajan toiselle puolelle, pienen unkarilaisen kylän kirkontornista katselee eräs koulustaan karkoitettu opettaja joka päivä syntymäkyläänsä, rakasta kotitaloansa, missä hänen äitinsä odottaa yhä vielä poikaansa... Raivaustyöt vielä jatkuvat, mutta tähänastisten tulosten perusteella voidaan jo nytkin nähdä linnan koko sisäinen rakennuskaava jättiläismittoineen, ja se dots}
\kat{Suomen Heimo}{25.04.1931}{P3}{Me allekirjoittaneet pyydämme kunnioittavimmin esittää Suomen Tasavallan Hallitukselle seuraavaa: Monet Inkerin kirkoista ovat jääneet kylmille, kun seurakuntien hoitajat on karkoitettu paikkakunnalta. Lisäksi harjoitetaan viranomaisten toimesta hillitöntä uskonnonvastaista kiihoitusta sekä yksityisenkin kotihartauden vainoamista.}
\kat{Suomen Heimo}{25.04.1931}{P3}{Kohta rauhansopimuksen jälkeen ovat neuvostoviranomaiset sekä Karjalassa että Inkerissä ryhtyneet toimenpiteisiin, jotka ovat ilmeisessä ristiriidassa edellämainittujen vakuutusten kanssa. Kansallisen kokonaisuuden hävittämiseksi on lisäksi tämän vuoden aikana karkoitettu yksistään Pohjois-Inkeristä 350 nimeltä tiedettyä perhettä ja muista osista Inkeriä myöskin useita satoja perheitä. Neuvosto-Karjalaan on liitetty puhtaasti venäläisiä alueita sekä siirtolaisina tuotu useita kymmeniä tuhansia venäläisiä, joten karjalainen väestö on joutunut vähemmistöksi eikä siis voi nauttia sille vakuutettua kansallista itsemääräämisoikeutta.}
\kat{Suomen Heimo}{25.04.1931}{P8}{Miten voimaperäiseksi ja innokkaaksi, laajoja heimoystäväpiirejä mukaan tempaavaksi tämä liike tuli, voidaan nähdä Vienan Karjalaisten Liiton jäsenluettelosta, joka ensimmäisiä nimiä myöten sisältää sivistyneistöä ympäri maata, karjalaisten kera yhteensä puolikahdeksattasataa nimeä. Kaikki perustetut kansankieliset koulut Vienan Karjalassa lakkautettiin, kaksi kansallisessa työssä mukana ollutta Vienan miestä karkoitettiin maanpakoon, kirjastot ja lukutuvat suljettiin armotta, Liiton toimesta aloitettu tientekokin lopetettiin ja vienankarjalaiset " Pakinat " oli ehdottomasti jätettävä rajan tälle puolen. Vuodesta 191 1 vallankumouskevääseen asti oli Liiton toiminta samoin muodollista, dots}
\kat{Suomen Heimo}{16.06.1931}{P20}{Kun tästä edesottamisesta annettiin tieto lehtiin, katsottiin erikoisen tarpeelliseksi huomauttaa, ettei ylioppilaiden toivomuksilla, mielipiteenilmauksilla, uhkauksenluontoisilla vetoomuksilla y. m. s. ole ollut mitään vaikutusta konsistorion päätökseen. Vaikka toisaalta, tuskinpa Inkerin köyhiä talonpoikia olisi karkoitettu konnuiltaan, jos he olisivat ajatelleet puoluepaavi Wiikin lailla, mutta kaikilla ei, onnekseen tai onnettomuudekseen, ole yhtä joustavat periaatteet ja yhtä venyvä omatunto kuin hänellä. Jos tätä muuten joku muiluttaisi, niin eihän se edellyttäen logiikan paikkansapitävyyden, mitään erikoista olisikaan: 1 ) tavallinen muilutus vain, 2 ) oikeistoporvarien provosoima, 3 dots}
\kat{Suomen Heimo}{20.12.1931}{P5}{Suomen kansaa ei ole ollut syytä varoittaa liiasta kansallistunnon hehkumisesta, tätä kansaa, joka lienee maailman välinpitämättömin kansa heimolaisistaan. öinä useimmasta Inkerin kylästä ja eräs karkoitettu kirjoittaa: " Näinä päivinä on moni inkeriläisnainen kadottanut järkensä valon. " Karjalan ja Inkerin vapauden aate on liian pyhä, jotta sitä voisivat menestyksellä ajaa yksilöt, jotka ovat omassa siveellisessä taistelussaan sortuneet.}
\kat{Suomen Heimo}{20.12.1931}{P10}{Siellä asuu myös tuo lapsellisen herkkä ja hetkellisiin mielialoihin taipuisa rahvas, joka vain elävän kosketuksen kautta suomalaiseen sivistykseen voi kohota muitten tasolle. Kaikki sellaiset, joiden suinkin epäillään imevän jollakin tavoin vaikutteita rajan tältä puolen, on armottomasti karkoitettu ja onpa alettu kuljetella rajakyliin taatuksi siiviläksi ryssäläisiä perheitä jostakin sisämaasta. Ja suosituimmat ovat juuri sel- I laiset tulijat, jotka eivät tee mitään I suurinäänisiä valmistuksia ja luota I pelkkiin virallisiin elimiin tai tilai- I suuksiin, vaan jotka tulevat tänne I omalaatuiseen rajanurkkaan vain et- I simaan kosketusta rahvaan kanssa.}
\subsubsection*{vangitseminen \normalfont\small(25 katkelmaa, 199 osumaa aineistossa)}
\kat{Toukomies}{05.02.1925}{P4}{Toista lukuvuottansa alottaneet Vienan Karjalaisten Liiton perustamat ja ylläpitämät suomenkieliset koulut suljettiin, kirjastot ja lukutuvat hävitettiin, suoranainen postiyhteys Suomen ja Vienan välillä lakkautettiin. Niinikään lakkautettiin » Karjalaisten Pakinoita » -niminen Suomessa julkaistu vienankarjalaisten äänenkannattaja, ja päälle päätteeksi kaikki kotisallaan silloin tavatut valistustyön tekijät vangittiin ja tyrmään teljettiin sekä poissa olevilta toimihenkilöiltä kiellettiin määräämättömäksi ajaksi oleskeleminen Vienan Karjalassa. Tapahtumista tiedon saatuansa ryhtyivät silloin Suomessa olleet kansansa parasta harrastavat vienankarjalaiset ja Karjalan asian ystävät toimiin dots}
\kat{Toukomies}{01.05.1925}{P12}{Vielä vähemmin on perustetta sillä valtiovallalle osotetulla rohkealla neuvolla, jonka valtuusto lausunnossaan suvaitsi julkituoda. Bolsheviikit sulkivat hänet pari vuotta sitten vankilaan, mutta vapauttivat sitten, hänen annettuaan tunnetun julistuksensa. Äsken kuolleen suuren ranskalaisen kirjailijan kirja hänen nuoruutensa ja lapsuutensa muistoista ja Ranskassa tekemistään matkoista.}
\kat{Toukomies}{01.10.1925}{P4}{Vuottojärvellä asuessa olivat hänen kolme poikaansa käyneet hänen luonaan, pyytäneet häntä lähtemään takaisin Venäjän puolelle ja taipumaan uuteen uskoon. Hänet oli vangittu Salmissa ja syytettiin, että hän vietteli Ruotsin valtakunnan uskollisia alamaisia » uuteen venäläiseen lahkoon, polttamaan itsensä. » Kun ukko tätä kertoi tutkijoilleen, kihlakunnanoikeudelle, painoi hän nöyrästi päänsä alas ja sanoi mieluummin henkensä menettävänsä kuin luopuvansa uskostaan.}
\kat{Toukomies}{03.1926}{P4}{Teille ovat tuoneet kärsimystä, mutta myös mainetta ne suuret teot ja taistelut, joita te olette saaneet suorittaa suurten maailmantapausten yhteydessä, osittain vierailla mailla ( esim. kolmekymmenvuotisessa sodassa ). 30 vuotta sitten, kun eräs karjalainen toivoessaan kansalleen valistusta toi Suomesta kuormallisen aapisia lukemisen alkuopastukseksi, sai hän luovuttaa kuormansa venäläisen valtiomahdin vartijalle ja itse lähteä 3 kuukaudeksi vankilaan Kemissä. Monesta voi tuntua, että pimeys on johtunut kansamme tarmottomasta luonteesta ja moni voi hyvälläkin syyllä ihmetellä, että meillä ei ole ollut miehiä, jotka olisivat pystyneet astumaan etualalle ja tuomaan valoa kansalleen.}
\kat{Toukomies}{04.1926}{P2}{Vielä virkeänä oleva vanhus liikkui nuorempina päivinään paljon rajantakaisessa Karjalassa runon keruussa. Iguumeni ja johtokunta on pidätetty virantoimituksesta ja sijalle asetettu erityinen hoitokunta. Wossinoi, Jalaja ja Mökerikkö jäävät Kurkijoen kunnan maakirjaan ja käsittävät 239,87 ha.}
\kat{Toukomies}{05.1926}{P7}{Seuran jäsenluku, vuotuinen kirjeenvaihto, etupäässä vapaaehtoisina lahjoituksina saadut tulot — m. m. Kontron ja Kuosmasen 100,000 markan suuruinen lahjoitus — ja moni muu seikka osoittavat melkoista uhrimieltä ja toimeliaisuutta. Viitattu teko, jota ei tahdottu todeksi uskoa, oli kolmen Vienan karjalaisten liiton ylläpitämän kansakoulun sulkeminen ja kolmen kansallismielisen karjalaisen vangitseminen, joista kaksi: nuori kunnallis- ja kansanvalistusmies Vasili Jeremejeff Räihä ja eräs toinen karjalainen H. Remsujeff karkoitettiin kauas kirgiisien ahoille, Aasiaan. Odottamattoman hyvällä menestyksellä karjalaiset tähän saakka ovat kestäneet pitkällisessä taistelussaan kansallisuutensa dots}
\kat{Toukomies}{10.1926}{P13}{16 p:nä erotettiin kauppakomissaari Kamenev toimestaan ja hänen sijalleen nimitettiin neuvostoliiton toimeenpanevan keskuskomitean jäsen Mikojan. Samoina päivinä tiedotettiin ulkomaille, että Moskovassa oli heinäkuun lopulla vangittu lukuisia punaarmeijan upseereja ja sadottain muita henkilöitä kommunistien luotetuimmistakin joukoista. Jos tästä sopimuksesta tulee tosi, tulisivat menshe vikit puheenaolevan tietotoimiston mukaan saamaan yhden kolmanneksen kansankamissaarien ja muiden korkeimpien neuvosto virkailijoiden paikoista.}
\kat{Toukomies}{12.1926}{P15}{— Yhtäkkiä hän hätkähtää, tarttuu nopeasti kivääriinsä ja kyyristyy kylyn nurkkaan. Eikö hän kuullut jotakin, eikö kajahtanut pidätetty aivastus jossakin aivan lähellä? Kuuluu ainoastaan suksen suhahtelua korven pimeydessä ja toisinaan vilahtaa epämääräinen varjo.}
\kat{Toukomies}{12.1926}{P20}{Liettuan roomankatoolinen papisto ei voinut jäädä kansansa kohtalolle kylmäksi. Yksinkertaisten hartauskirjani levittäjät tuomittiin säälittä vankilaan tai maanpakoon. Tämän viidenneksen verokirja v:lta 1500 on yksi parhaita Karjalan historian lähteitä.}
\kat{Suomen Heimo}{20.02.1928}{P22}{Voima-lii- I ton puuhiin Herman Stenberg laittautui koko sydämellään ja I sielullaan, mutta vasta jääkäri- I liikkeessä hän löysi ihanteensa, Ija päämääränsä toteuttamisen I välikappaleen. Tien päässä heitä itseään odotti vankila, maanpako -taikka hirttonuora, mutta niiden takaa kangasti heidän silmiinsä Isänmaan vapaus — Herman Stenbergille vielä enemmän: Suomen heimon vapaus. Kun aktivisti I toisensa jälkeen santarmien toi- I mesta joutui ilmi ja joko vankeu- I teen taikka maanpakolaisuuteen, I oli Herman Stenberg viimeisiä, |joka Venäjän maaliskuuhvaliankumouksen aattona suostui jättämään Isänmaan.}
\kat{Suomen Heimo}{15.04.1928}{P19}{Vaikka me emme heti ensi otteissamme pystyisikään hienoudessa ja laadun jaloudessa kilpailemaan Ruotsin kanssa, niin antaisi meille arvoa kuitenkin se, että me olemme tuoneet esiin omaamme, jotakin ennen löytämätöntä suomalaisen sielunelämän aarnioista ja että me ajattelemme omilla aivoillamme. t. k. 28 ja 29 p:nä Viipurin valloituksesta monipäiväisten katkerien taistelujen kautta, jotka olivat enimmäkseen asevelvollisten joukkojen suoritettavana. \_ Viipurin valloituksen jälkeen oli pääasiallisesti ainoastaan hajallelyötyjen vihollis joukkioiden saartaminen ja vangitseminen enää jälellä, mikä kuitenkin vielä vaati vakavia taisteluja; Palkittuja \&111 1:118 palkinnolla.}
\kat{Suomen Heimo}{20.05.1928}{P16}{On lisäksi vielä huomattava eräs seikka, joka parhaiten kuvastaa opettajajoukon puutteellista kokoonpanoa: opettajista on 70 \% naisia; tämä seikka parhaiten vakuuttaa nykyisen opettajiston tilapäisluonnetta ja epäpätevyyttä. Luettelossa mainituista muista oppilaitoksista on ennen kaikkea huomattava Rääpyvän maatalous- ja karjanhoitoteknikumi, joka onkin ainoa todellinen ammattioppilaitos koko Inkerissä ja jonka johdossa on toiminut inkeriläinen agronoomi, nyttemmin kyllä vangittu. Kansanvalistusasiain johto on keskitetty Pietarissa olevaan Lännen vähemmistökansalli suuksi en sivisty » valistu « jaostoon, jonka johdossa on suomalaisia punaisia, virolaisia ja lättiläisiä, ja sieltä sanellaan dots}
\kat{Suomen Heimo}{15.09.1928}{P5}{— Huomioonotettava näkökohta on myöskin se, että ahvenanmaalaiset asevelvolliset täyttämällä luotsi- ja majakkalaitoksemme kohtuuttomasti ruotsalaistuttavat tärkeätä haUintohaaraa ja sen laitoksia. Kieltämättä on myönnettävä, että tärkeimmät lainsäädännön alat on pidätetty Suomen eduskunnalle, mutta turhaa nöyryytystä Suomelle tämä » Oolannin lakitehtailu « kuitenkin kaikitenkin on. Tälläkin tietysti yritetään estää uusien asukkaiden tuloa mannermaalta käsin, ja on tämä säännös varmastikin paljon tehokkaampi ase suomalais vastaisessa taistelussa kuin edellämainittu kansallisen ja maakunnallisen äänioikeuden menetys.}
\kat{Toukomies}{01.02.1928}{P4}{Toinen mies olisi kai siinä leikissä jo mennyt tajuttomaksikin, mutta herra Henrikki oli vielä vanhoillaankin lujatekoinen äijän karilas. Välistä vain pääsi pidätetty ähkäisy hänen viinakurkustaan ja tämä ähkäisy palkittiin joka kerran äänekkäällä naurunremahduksella. Ja entäs kun nyt tuli vielä kaupan päällisiksi tämä vanki, tämä verraton vanki — Pielisen pokosten ryöväritiuni.}
\kat{Suomen Heimo}{15.09.1929}{P24}{Muuten näkyy hra P:lle Karjalan kysymys olevan yksinomaan hyötykysymys, eikä hänelle mitään merkitse se, että heimo laisemme pitäisi pelastaa länsimaiselle sivistykselle. Norjalaiset, englantilaiset ja saksalaiset valtaisivat metsiä, ruotsalaiset perustaisivat tulitikkutehtaita, karjalaiset pitäisi panna vankilaan, koska ovat kommunisteja 1 Mikä nai visuus! Kun tällainen persoonallinen kanssakäyminen on enemmän kuin mikään muu omiaan lujentamaan Eesti — Suomi-siltaa, on toivottava, että Suomen puolelta osanotto nyt tulisi olemaan erittäin runsas.}
\kat{Suomen Heimo}{15.11.1929}{P19}{Jälkisatoa A I:n rintakuvan tervau ksesta. Nykyisen ylioppilaspolven isät ja veljet raastoi ryssä vankilaan. Voimme hra Maliselle vakuuttaa, että se ei tule pysymään.}
\kat{Toukomies}{01.05.1929}{P15}{TOUKOMIES Se olikin — vankila. Oli jo myöhäinen ilta.}
\kat{Toukomies}{01.09.1929}{P18}{Maataloudesta puhuen juohtuu mieleeni kirkkoherra Nuutin puheet rajantakaisen Karjalan maataloudesta. — • Mistä syystä olette joutuneet vankilaan ja noin kurjiin oloihin? Olin jo saanut kyllikseni Petroskoin oloista, ihmisistä ja kadulla mustanaan vilisevistä naakkaparvista.}
\kat{Suomen Heimo}{13.03.1930}{P20}{— Mitä oli tällaisesta julkilausumasta seurauksena: se, että maamme sanomalehdistö harvinaisen yksimielisesti ja isällisesti, harvoja poikkeuksia lukuunottamatta, hyökkäsi ylioppilaitten kimppuun — ja kielsi " lapsia politikoimasta ". Ja epäilemättä hän karkeakouraisesti vaatisi vankilaan toimitettavaksi ne idän valtakunnan rahalla ostetut kätyrit, jotka lupauksin, valhein ja kaunomaalauksin saavat monet rehellisetkin suomalaiset seuraamaan narrinkulkusin varustettua lippuaan. Ja minä tiedän, että hänen sydämensä olisi avoinna Suur-Suomen aatteelle, tuolle huimaavalle aatteelle, joka vanhojen ja viisaiden leirissä tuomitaan utopiaksi, mutta joka on nuorisolle reaalinen, dots}
\kat{Suomen Heimo}{26.04.1930}{P24}{Jos täällä kenen on leipä ostettava, on kovin vaikea tulla toimeen, sillä leipää eikä jauhoja saa ostaa mistään yhtään kiloa. Useita I kymmeniä talonpoikia ja I kirkonpalvelijoita on vangittu, lähetetty pakkotöihin aina Solovetskiin saakka ja heidän koko omaisuutensa takavarikoitu. Hyvinvoinnin asemasta I kansa on saanut jo toista- I kymmentä vuotta kärsiä mi- I ta suurinta puutetta eikä ok- I olemassa mitään toivoa 010- I jen muutokseen.}
\kat{Suomen Heimo}{26.04.1930}{P25}{Huoneuston suuruudesta saa käsityksen mainit- J taessa, että siinu' asetetaan n. 35 näytepöytää, kukin * • esittäen oman kirja-alansa. Nyt jälleen ovat kulakkivainot enemmän kiihtyneet, vangitsemisia on tapahtunut paljon, vangitut istuvat kuukausimääriä vankiloissa eivätkä tiedä, miksi vangittiin. " Seurakuntain taloudenhoidolle asetetaan esteitä suurien verojen muodossa, vaikeutetaan seurakuntamaksujen keräämistä ja asetetaan kirkolliset yhteisöt lain edessä niille epäedulliseen poikkeukselliseen asemaan. "}
\kat{Suomen Heimo}{15.06.1930}{P19}{Jos luoksesi saapuu A. K. -S:n jäsen, pyytäen Sinua kannatusyhdistyksen jäseneksi, liity siihen epäilemättä. Monesta perheestä on raahattu isä vankilaan, kodin kaikki omaisuus ryöstetty ja äiti lapsineen häädetty tuuliajolle. Ainaiset vangitsemiset, ryöstöt, murhat, nälkä ja puute ovat saattaneet suomalaisen väestön ennenkuulumattomaan hätään.}
\kat{Suomen Heimo}{15.06.1930}{P19}{Se ei kuitenkaan tahdo vielä menettää viimeistä toivoaan, vaan uskoo edelleen, että neuvosto-orjuudenkin päivät loppuvat. Kaikkia niitä henkilöitä, jotka käyvät asioillaan Suomen konsulaatissa valvotaan tarkoin ja useita heistä on sen tähden vangittu. A. K. S-. tarvitsee kannatusta ja rahaa, ilman rahaa se ei voi tarpeellisella voimalla toimia korkeiden päämääriensä hyväksi.}
\kat{Suomen Heimo}{18.11.1930}{P13}{— Sulo-Weikko Pekkolan muistelmateos — Lähettä kää minulle allamerkits emäni Werner Söderström on tyhjentävä kertomus sen lähes 100- O. Y:n kustantamat teokset ( mitä en halua olen pyyhkinyt ): jéÉfc * \textasciicircum{}åaå miehisen suomalaisjoukon vaiheista, jot- Huntuvuori: Kokemäenmaan kuningas. 180: SPI ovat asianomaisten itsensä kertomat, sev- nästesk o - H e p o rauta: Ursula Keivaara 1 UI. mm raavat toisiaan siinä järjestyksessä kuin Ni ( l 135 - — / sid 180: — heidät vangittiin v Juuri i\textasciicircum{}stvnyt I osa Kait eri jääkärit I. Nid. Nyt tiedetään, että valtiollinen taju ei ole koskaan ollut täällä sammuksissa, vaan on usein puhjen- nut oikeaksi kansallistunnoksi. iv\textasciicircum{} n n rji n imi / P a g l\&iyjlKiulNl<y> / dots}
\kat{Toukomies}{01.02.1930}{P18}{Suomi se oli, joka monena huonona vuotena rajantakaisille elatuksen antoi. 550,000 henkilöä on pidätetty lain rikkomisen johdosta ja 230,000 teljetty vankiloihin. Alkoholismin tai myrkyllisen alkoholin seurauksiin on tänä aikana kuollut 34,000 henkeä.}
\subsubsection*{teloitus \normalfont\small(25 katkelmaa, 56 osumaa aineistossa)}
\kat{Suomen Heimo}{15.08.1929}{P23}{Vaadimme täyden ja selvittävän vastauksen näihin kysymyksiin. ta ammuttiin arvan perusteella paristakymmenestä siihen osallistuneesta yksi — viaton. Tuntui tcsiaankin koomilliselta lukea kansallismielisen lehden palstoilta tuollaista törkyä.}
\kat{Toukomies}{01.09.1929}{P25}{Maahan tuli uusia pappeja, ja vaikka Ansgarius taas matkusti Saksaan, jäi hänen työnsä jatkajia ja opille jalansijaa maassa, ja sitä mukaa tämä levisi ja juurtui edelleen. Pianpa Ansgariuksen lähdettyä tämä usko sieltä tyrehtyikin, kansa nousi Vastaan, yksi uusista lähetyssaarnaajista surmattiin j. n. e. 850-luvulla tuli Ansgarius kuitenkin takaisin Ruotsiin ja silloin hänen työllään oli parempaa menestystä. Iltapäi väliä jatkuivat juhlallisuudet juhlijan kauniilla kesäasunnolla, jossa lukuisat puhujat toivat esiin kunnioituksen ja kiitoksen tunteitaan sekä tunnustuksensa juhlijan vaatimattomasta, hartaalle työlle omistetusta luonteesta ja miellyttävästä personallisuudesta.}
\kat{Suomen Heimo}{13.03.1930}{P17}{Mutta aamupuoleen vaikeni suomalainen patteri ja konekivääri tulen suojassa lähetti vihollinen joukkoja maihin. Kehittyvässä tykistötaistelussa ammuttiin Viteleen kylä tuleen, jolloin jääkärikapteeni Kuisma kuolettavasti haavoittui. Weilojan ja kapt. Hämäläisen joukkojen kanssa muodostaa Viel järvestä etelään rintama majuri Talvelan suojaksi.}
\kat{Suomen Heimo}{28.10.1931}{P18}{46484, Suosittelee arv. yleisölle Aamiaisia / Päivällisiä / Illallisia / Annoksia Iltakonsertti päivittäin kello 20 — 24 kapellimestari Ahrenbergin johdolla. Niinpä lokakuun 28 päivänä 1929 ammuttiin tuntemattomasta syystä 65 entistä tsaariarmeijan upseeria, joista osa sittemmin oli palvellut Koltsakin johtamissa joukoissa. Huone, johon kasakkaupseeri joutui, oli suuruudeltaan 17 X 25. Siinä oli kolmessa kerroksessa laudoista tehdyt makuusi j at, keskellä lattiaa välillä käytävä.}
\kat{Suomen Heimo}{20.12.1931}{P5}{Tuo yksinäiselle vartiopaikallensa kaatunut Inkerin mies on myöskin AKS:n jäsenille velvoittava esimerkki varsinkin sellaisina aikoina, jolloin heidän rivinsä horjuvat ja jolloin he ovat vaarassa unohtaa sen minkä vuoksi ovat vannoneet elävänsä ja kuolevansa. Ja samaan aikaan kun kauhunviestit kuuluivat Inkeristä, päätti eduskunta kieltää avustuksen Eestin vapaussodassa haavoittuneilta suomalaisilta sotilailta, jolloin se samalla sen kielsi myös Inkerin vapaajoukkojen suomalaisilta invaliideilta, sillä on totta niinkuin on sanottu: " Tälläkin puolella Suomenlahtea elää vielä muutamia raajarikkoja, joilta ammuttiin käsi tai jalka silloin kuin Inkerin vapaajoukot rynnistivät kohti Krasnaja dots}
\kat{Suomen Heimo}{17.12.1932}{P30}{Se oli kaunein joulu elämässäni. Everstin määräyksestä ammuttiin ilmaan satoja valoraketteja. Eräs mies sytytteli parastaikaa kynttilöitä pitkällä kepillä.}
\kat{Suomen Heimo}{30.03.1933}{P11}{Pian olisi hetki tullut, jol- loin korvaamattomat tiedot olisivat lopullisesti hävin- neet unholaan. » Pojan veli otti osaa murhayritykseen Aleksanteri 111:ta vastaan ja teloitettiin nuoremman veljen ollessa 17-vuotias. Pätevän Venäjän-tuntijan tunnollinen, laajaan lähdekirjallisuuteen pohjautuva elämäkerta tulee varmasti herättämään paljon huomiota.}
\kat{Toukomies}{01.02.1933}{P4}{Kerrottiin, jotta illalla pitäisi kädenisku tapahtua... jos se nyt totta lie... Ammuttiin muutamia laukauksia pirtin kattoon tuli jäisiksi ja sulhaskansa kävi sisään, puhemies edellään. Ja muista sanuokin, kun kumarrat, jotta tuattosen, olet syöttänyt juottanut, tiijä tuumakin antaa..}
\kat{Karjalan Vapaus}{12.1934}{P33}{Sitten vanhusta tärisytti outo vavahdus yli koko ruumiin ja oikea käsi kohoutuu hitaasti ja tekee ristinmerkin. Poika ammuttiin, mutta hänet muka armahdettiin vanhuutensa tähden, mutta karkoitettiin Siperiaan. Erikoisena määräyksenä hän ilmoitti, että joulujuhlan lähestymisen johdosta työpäivä pidennetään kolmella tunnilla entisestään.}
\kat{Toukomies}{01.10.1934}{P10}{Puheiden sisältöä en muista, tuskinpa niitä niin hyvin ymmärsinkään, mutta nyt kyllä hyvin helposti sen arvaan. He nähtävästi aikoivat ottaa tepluskat sen perään, mutta jostakin ammuttiin ja veturi iski melkoisella vauhdilla vaunuihin. Tuntikausia me istuimme heidän parissaan, silloin kun he päästelivät silakoita verkoistaan, ja kuuntelimme heidän puhettaan.}
\kat{Suomen Heimo}{26.02.1935}{P30}{Myöskin. vankien joukkoteloitukset olivat niihin aikoihin miltei jokapäiväisiä. Ammuttuja syyttivät karkulaisen avustamisesta, mutta teloitus tapahtui tuomiotta. Janoisena ja nälkäisenä join keiton, vaikka se haisi tavattoman pahalle.}
\kat{Suomen Heimo}{26.02.1935}{P30}{Kymmenniekka ja 2 sotilasta tulivat yöllä riiheen, pidättivät minut ja veivät Keskuskylään, johon oli matkaa 16 km. Satuin kerran näkemään omin silmin, kun 62 vankia ammuttiin sen vuoksi, että yksi heidän joukostaan oli päässyt pakenemaan. Neuvostoliiton yhteiskunnallisen kokeilun ja kansallisen hävittämis politiikan pyörteissä on yksilö kuin avuton lastu, jolla kohtalo leikittelee.}
\kat{Suomen Heimo}{15.03.1935}{P14}{Silloin kertantti Helge Halminen ja säämäjärveläinen vapaustaistelija joutuivat aivan Harjun kylän rinnassa bolshevikien käsiin. Heitä ammuttiin, mutta ei kuolettavasti, kidutettiin niin, että kidutettavain hätähuudot kuuluivat Harjun kylään. Kersantti Helge Halminen oli ollut mukana Viron retkellä, oli pelkäämätön, uljas nuorukainen ja paloi sankari-intoa.}
\kat{Suomen Heimo}{01.04.1935}{P8}{Lönnrot, kansanomaisin ja vaatimattomin suurmiehemme, jonka suurtyötä Kalevalaa koskevan, ennen v. 1927 painetun kirjallisuuden pelkkä luettelo sisältää 261 suurikokoista sivua, on tähän asti ollut arvonsa mukaista elämäkertaa vailla, siis todellakin vaatimattomassa asemassa. Niinpä kaikki nämä kyläkapitalistit, maanomistajat ym. omistavan luokan ainekset teloitettiin summittaisesti, vanhat maanjaon rajapyykit hävitettiin ja uusi jako toimitettiin, verot julistettiin anteeksi jne. Sitten ryhdyttiin rakentamaan uutta pakkokomentoa, kovempaa kuin tuskin minkään sotilastyranninkaan aikana. Kalevalan ilmestymisestä lähtien hän oli itseoikeutettu auktoriteetti suomen kielen}
\kat{Suomen Heimo}{15.05.1935}{P8}{Mutta koska kolmas rivistö minua läheni, niin minä tunsin hakkapeliitat, suomalaiset ratsumiehet, jotka vietiin kauas etelään Ruotsin ja ruotsalaisten kunnian koroittajiksi, silloin kuin Suomessa talot autioituivat ja idässä Novgorodin mahti kasvamistaan kasvoi. Olivat saatossa myös ne Rautalammin harteikkaat uroot, jotka kuninkaan käskystä raivasivat Värmlannin erämaat viljelykselle, mutta joita susikoirien kera vainottiin, kun he eivät suosiolla luovuttaneet perkaamiaan korpia vieraan maan miehille; olivat myös Länsi-Pohjan talonpojat, joilta kansallistunto kehtoon surmattiin. " MIKÄ ON NIMI teidän, jotka vielä tänäpänä rikostanne jatkatte, eristyen ja eristäen, hajoittaen ja halliten, dots}
\kat{Suomen Heimo}{28.06.1935}{P15}{Aino-uimapuvuista pidetään yhtä paljon kuin i on valmistamiemme miesten alus- ja urheilu- sekä voi misteluvaatteiston rekisteröity laatumerkki. Jollei henkeään menettänyt, niin ainakin häneltä silvottiin nenä tai korvat, mutta toisia ammuttiin suoraan omaisten silmien edessä. Kun joutuivat omaisen tai ystävän taloon, oli heillä hyvä, mutta kun majapaikka oli kommuunaan liittyneen kodissa, oli tuo yksi vuorokausi hirvittävä.}
\kat{Viena-Aunus}{01.02.1935}{P10}{Uusiksi varajäseniksi päätettiin kutsua kirjailija Eero Eerola, kansak. opettaja Mikko Karvonen, konttoristi Lasse Kuntijärvi ja työnjohtaja Pekka Tiilikainen. Teon suorittaja Nikolajev ja lukuisia muita — kerrotaan yli sadasta hengestä — sen jälkeen kohta teloitettiin, ja sen jälkeen on syyllisiä etsitty yhä laajemmilta aloilta. Ja sillä teollaan he myös osoittivat ja näyttivät kieltämättömän todeksi, että Suomen kansan kaikilla heimoilla oli suuri osuutensa ja tehtävänsä tämän kansan ja suomalaisen kansallisuutemme henkiin herättämistyössä.}
\kat{Itä-Karjala}{02.1936}{P15}{Mutta sitten tuli Moskovasta Tsekan pomo Ivan Ivanovits Pyövelinski, ja kuten tiedämme, ei Tsekan " lintukoiralta " pysy mikään salassa — ei edes tuollainen mitätön tapahtuma kuin parin kaivoksen katoaminen. Edelleen on kerrottu, että kaikki itäkarjalaiset sissit olisi Antikaisen toimesta muitta mutkitta heti teloitettu, mutta Lassy ja hänen vääpelinsä viety Moskovaan asti muka todistuskappaleiksi siitän että Vienan sodassa oli Suomen armeijan upseereita johtajina. 35: —. W. S. O. Y. AAAAAAAAAAAAAAAAAAAAAAAAAAAA}
\kat{Itä-Karjala}{09.1936}{P10}{Seuraavat Pekka Hirsso, Lempessalmi ja Nehvonen. Kaikki salaliiton jäsenet, myös GPU:n agentit-provokaattorit, ammuttiin. Tähän naisten kolmiotteluun ei ilmaantunut enempää osanottajia.}
\kat{Karjalan Vapaus}{03.1936}{P14}{Laivasta kuului suomenkielinen kehoitus meille: „tulkaa, tulkaa rantaan, he antavat tupakkaa. " Sillä Tuulosjoen suussa oli ammuntaa, vaikka emme saaneet selvää, mistä ja minne päin ammuttiin. Mutta siitä huolimatta avattiin tykkituli bolsheviikkeja vastaan ja tervehdittiin heitä siinä muodossa merkkipäivän johdosta.}
\kat{Karjalan Vapaus}{03.1936}{P20}{Pienempää vartio joukkoa varten ottivat karjalaiset haltuunsa lisäksi Enonsuun kylän, joka sijaitsee Uhtualta ja Luusalmelta Vuokkiniemelle johtavan vesiväylän varrella. Ihmisiä vangittiin, toisia raakamaisella tavalla teloitettiin, kuten Pappinen Jyväälahden kylästä, muita mainitsematta, toisia lähetettiin Venäjälle ja toisia vainottiin. V. 1920 juuri siihen aikaan, jolloin vallassaolijat vaativat veroja ja tekivät ryöstöjä, oli Karjalassa ankara nälkä. Éi tiedetty, millä yli päivän eletään.}
\kat{Karjalan Vapaus}{03.1936}{P37}{kohtalo oli edessämme. Sellaiset ammuttiin näet heti. Kokemuksia neuvostomaasta}
\kat{Karjalan Vapaus}{03.1936}{P37}{Vangitsemisen jälkeen pidettiin minua 3 päivää Vitelen komendatuurissa, josta sitten siirrettiin 5 päiväksi Aunuksen vankilaan ja sieltä Petroskoin vankilaan, jossa kuulustelunikin vasta alkoivat. Niinpä esim. saimme kokea ensimmäisen näyn tästä jo kohta mentyämme Solovetskiin, sillä marraskuussa v. 1925 ammuttiin 3 miestä, joilla karkaamisaikeen sanottiin olleen mielessä. Tässä yhteydessä sanottakoon, että Solovetskiin matkalla en ollut taaskaan yksin, vaan paljon oli seurassani toisiakin onnettomia ja jo Popovan saarella meidän annettiin aavistaa, minkälainen}
\kat{Suomen Heimo}{15.09.1936}{P16}{Varsinaisena alkuna käsilläolevaan suurpuhdistukseen voidaan kuitenkin pitää " Pravdan " pääkirjoitusta elokuun 7 päivänä " Osattava tuntea vihollinen ", jonka johdosta puolueen paikalliset elimet järjestivät kansalle valistustilaisuuksia jo seuraavana päivänä. Kominternin entinen puheenjohtaja Sinovjev ja kansankomissaarienneuvoston entinen pää Kamenev sekä 14 muuta tunnettua kommunistia ammuttiin lyhyen oikeudenkäynnin jälkeen syyllisinä terrorihankkeisiin, jotka kohdistuivat Stalinin, Vorosilovin ym. neuvostojohtajien henkeä vastaan. Mistä sitten on lähdettävä etsimään syitä, jotka ovat saaneet Stalinin ryhtymään juuri nykyisenä ajankohtana näin ujostelemattomaan välienselvittelyyn, dots}
\kat{Suomen Heimo}{30.11.1936}{P20}{Vain yksi pikakäynti tutustuttaa meidät ukko Klercker iin, joka Viaporissa yöllä vastaanottaa yllättävän sotasanoman. Tuomionsa saivat tällä kerralla Kullervo Manner, joka todennäköisesti ammuttiin, Villiam Tanner sekä Hanna Malin. Mitähän palkintoja meille vielä sataakaan, kun skandinaviset suhteet vielä edelleen lujittuvat!!!}
\subsubsection*{nälkä \normalfont\small(25 katkelmaa, 176 osumaa aineistossa)}
\kat{Toukomies}{05.02.1925}{P12}{poroksi polttivat huonehiamme, Surma ja nälkä ja uskonvaino laulaa Arvi Jännes, tai:}
\kat{Toukomies}{01.06.1925}{P7}{Sellaisella sivistyksellä on voima vaikuttaa kaikkiin ja jalostaa kaikkia. vei nälkää, vaaroja voittamaan, vei valoa voittamaan meille! » Monen pitkän viikon taivalluksen jälkeen joutuivat valitut miehet vihdoin kullankiiltävään Moskovan kaupunkiin.}
\kat{Toukomies}{01.09.1925}{P20}{Mutta kumppanini Pekka jäi vielä siihen katsomaan, saavatko pulkat olla rauhassa. Nälkä käski näille töille vanhana varastamahan, ikälopun ilkatöihin. sen antavinaan, mutta äkkiä rajusti tempasi sen käsiinsä ja ojensi miehiä kohti.}
\kat{Toukomies}{01.09.1925}{P20}{— Elkää, veikkoset, ne ampuvat! huusivat miehet ja lähtivät pois. Nälkä tuli näille maille, Tuoppajärven jäätä myöten, pieni kassara käessä, — Setä, lue sinä » Isämeidän » minun kanssani, jatkoi hän sitten.}
\kat{Toukomies}{08.1926}{P16}{Yhä nopeammin ja nopeammin suhahtelivat sukset hohtavaa hankimetsää. Eivät he huomanneet väsymystä, eivät nälkää, eivät pakkasta. Toverinaan kalma kauhistava, kymmenen hyisessä hyhmässä makaavaa vainolaista.}
\kat{Toukomies}{01.02.1927}{P7}{Nuori Lokki halkaisee lauman kahtia, kiitää, kiitää silmäinsä edessä lumivalkoinen johtajapeura. Häntä ahdistaa kylmyys, majan kosteus, nälkä, mutta ei niidenkään vuoksi pyhä mies kärsi, ei! Häntä vastaan tuijottaa kaksi suurta, mustaa, tuskasta hehkuvaa silmää, ja hän kuulee heikon huokauksen.}
\kat{Toukomies}{01.02.1927}{P7}{ovat pöhöttyneet paksuiksi, muodottomiksi, niitä pistelee, mutfei niiden vuoksi pimeän metsän erakko voihki, ei! — Tuska, nälkä, vilu, Kalma ovat hänen ystäviänsä, vanhoja tuttaviansa, mitäpä hän niiden vuoksi huokailisi? Luonnonvoimat riitelevät ja toraavat pimeän metsän syvyyksissä toinen toisensa kimppuun syöksyen ja täyttäen pimeyden jylhällä pauhinalla.}
\kat{Toukomies}{01.03.1927}{P10}{Pysähdyspaikoissa kuolleet ehkä huomataan, ellei, niin hekin saavat jatkaa matkaansa. Ne vaikeroivat: — Maammo, maammoseni, on niin vilu, on niin nälkä! Pois, pois alta, pois surman suusta, kalman kidasta — pois!}
\kat{Toukomies}{01.09.1927}{P5}{Iltaman ohjelmassa oli maisteri vSetälän, Keynään ja Ky öttisen puheet, viulunsoittoa, yksinlaulua, lausuntoa y. m. Miten yksinkertaisia tavoissaan, miten vähääntyytyväisiä, vilua ja nälkää kestäviä, miten lapsenmielisiä olivatkaan nuo vanhukset tuolla saloilla! Toisena päivänä kaupungintalolla piti toim. Kotonen puheen, Arvi Mansikka lausui Eero Eerolan kirjoittaman juhlarunon ja juettiin saapuneet sähkösanomat.}
\kat{Suomen Heimo}{20.01.1928}{P17}{Pitkä ja vankka, mies, kolmenkymmepenviiden maissa. Nälkää, kylmää, rasituksia, epätoivoista kamppailua. Silloin hän puhkesi äkkiä puhumaan, kertomaan elämästään.}
\kat{Suomen Heimo}{15.04.1928}{P10}{Me tiedämme vain, että vuosituhansia sitten suomensukuiset kansat pitivät hallussaan suunnattomia Itä-Euroopan alueita sekä että ne hajaantuivat useihin heimohaaroihin, jotka vaelsivat pois alkuperäisiltä asuinsijoiltaan. Erämaiden elämä ei ollut kuitenkaan rauhaisaa, sillä lukuunottamatta keskinäisiä kovia kamppailuja kalavesistä ja metsästysmaista, suomensukuisten kansojen perivihollinen ilmestyi tuon tuostakin idästä jättäen jälkeensä nälkää ja ruttoa, kuolemaa ja savuavia raunioita. Rohkenen epäillä, että tuo hajaannus yhteisestä alkukodista, jonka yhtenäisyyden säilyttäminen olisi varmaankin merkinnyt Suomen heimon suuruutta, johtui eräistä luonteemme perusominaisuuksista, vaikka dots}
\kat{Suomen Heimo}{15.04.1928}{P11}{Suomessa ei vielä osata täysin ymmärtää nälkää ja vihollista vastaan, taisteluihin kotiemme puolesta. Ja nyt kun tulevaisuus on edessämme kaikkine mahdollisuuksineen, on tehtävämme suunnata katse eteenpäin.}
\kat{Suomen Heimo}{15.06.1928}{P17}{Osa sijoitetaan vakinaiseen seisovaan armeijaan, jossa on 2 vuoden rauhanaikainen palvelusaika, osa ter r i tori a a 1 i armeijaan, jossa palvelevilla on 8 — 11 kuukauden harjoitusaika, ja loput muodostavat ylijäämän, jota harjoitetaan vain 4 — 6 kuukautta. suuresti väheni lännessä, mutta tästä huolimatta, ja vaikka maailmansota, vallankumous, nälkä ja teloitukset ovat tuottaneet suurta mieshukkaa, on Venäjän väkiluku nykyään arvioitu noin 145 miljoonaksi, mikä on yhtä paljon kuin ennen maailmansotaa silloisella alueella. Punaisen Itämerenlaivaston toimintaa sodan sattuessa vaikeuttaa suuresti se, että Suomenlahden rannat ovat Suomen ja Viron käsissä, mutta kun se on Itämeren rantamaiden dots}
\kat{Suomen Heimo}{15.09.1928}{P13}{Jos jätämme poikkeusasemassa olevat unkarilaiset huomioon ottamatta, näemme, •että kaikista Neuvosto- Venäjän suomensukuisista kansoista ainoastaan suomalaiset ja tsheremissit ovat pystyneet säilymään epäkansallistumiselta nykyään elävän sukupolven aikana. Parhaiten säilyvät suomensukuisista kansoista pystyvät aivan ilmeisesti pitämään suhteellisen hyvin puolensa elämän kamppailussa Neuvosto- Venäjänkin sodan ja valr lankumousten, nälän ja kulkutautien jo kohta 15 vuotta myllertämässä noidankattilassa. „Uuden Suomen " 194 numerossa on julaistuna henkivakuutustarkastaja V. Keynään haastattelu, jossa hän m.m. lausuu, että „eräältä suomalaiselta taholta on koetettu saada eripuraisuutta ja dots}
\kat{Suomen Heimo}{15.10.1928}{P17}{Monesti Snellmanin ja Kiven aikojen jälkeenkin on julkisuudessa huomautettu siitä, että » Maamme " - laulua voisi melkein mikä muu kansa tahansa, joka elää samantapaisissa maantieteellisissä olosuhteissa kuin me, käyttää omana laulunaan. Ajatellen » Maamme " -laulun ajatussisältöä hän jatkaa: „Turve et ole isänmaa, orjan ottelu nälkää vastaan ei ole kansalliselämää, jäljet vihollisen hevosen kaviosta ei historiaa, halu kotikontuun ei isänmaallisuutta ". Kiven ainoaan isänmaallis-aatteelliseen tuotteeseen, runoon » Suo. m imi maa ", jonka otsakkeena alkuaan on ollut „Maamme ", sisältyy tavallaan Runebergin runon arvostelua, kuten kirjallisuudentutkijamme toht. J. V. Lehtonen on huomauttanut.}
\kat{Suomen Heimo}{18.11.1928}{P17}{Kauhea löyhkä täytti ilman ja kaikki eloon jääneet alkoivat pyrkiä pois kaupungista. Kelirikon vuoksi ei voitu tuoda ruokatavaroita kaupunkeihin, joten nälkä ja puute oli varsin suuri. Samoin tapahtui vuonna 1158, Jolloin kronikoitsijan sanojen mukaan Novgorodissa kadut olivat täynnä}
\kat{Toukomies}{01.02.1928}{P18}{Onttoni Miihkali. Nälkä ja vilu alkavat ahdistella. Totuus itsestämme.}
\kat{Toukomies}{01.03.1928}{P14}{Karpau oli työläs, suorastaan ylivoimaisen vaikea, selittää edes itselleen tätä nurjaa, sysimustaa tilannetta, missä hän Ontro-poikansa tyttösen kanssa yritti elää ja minkä tulirautaisessa kädessä oli jokaisen läheisen heimolaisen elämä. Hänen oma elämänsä oli arpinen, runneltunut, kuuruinen ja kärsivä vastaus: Kuinka monelle sodan tantereelle oli jo hänen verensä vuotanut vieraan veren ja vieraan heimon etujen puolustamiseksi, kotijoukon nähdessä sillä välin sekä ruumiillista että henkistä nälkää! Kaiken varalta oli Karppa kuitenkin pari päivää ennen mellastajain tuloa vienyt kelkalla kaiken ruokatavaran — leivät, jauhot, voit, suolat ja lihat — ja osan vaatteitakin muutamain virstain dots}
\kat{Toukomies}{01.10.1928}{P10}{\_ yi maamme Suomi, synnyinmaa, soi sana korkea! soi sota päällä laaksojen, toi halla nälän, puutiehen. n köyhä maa, ei köyhempää, jos kullan kiihkoon saat.}
\kat{Toukomies}{12.1928}{P12}{Kylän naiset kantoivat kaikki parhaat herkkunsa Litin eteen. Olihan litillä jo nälkä, vaan eipä uskaltanut pyytää. Nauraen se Litin varpaan nenille nousi ja alkoi painaa.}
\kat{Toukomies}{12.1928}{P12}{Pirtissä asuvat kokko Järveläiset käydessään jauhattamassa. Silloin alkoi hähelle nälän näivertäjän, iki-ikävän ja puutteen aika. Muurahaiset, jääjalkaiset pikku eläjät, juoksentelivat selkänikamiani myöten.}
\kat{Toukomies}{12.1928}{P13}{Ja takaa meren kymmenen, kedoilta Onnellisten Maan toi hunajaa se nokassaan ja mettä kukkien. Jok' ainut kerta, päällä maan kun janosin ja nälkää näin, se ruokki minun sydäntäin kuin emo poikastaan. — Paennut: kerran retkellään se palas luokse sydämen ja — liudat haaskalintujen se näki pesällään.}
\kat{Suomen Heimo}{25.05.1929}{P8}{Vaikein kuitenkin oli rinnassa, syömmessä haava: hylkäsi heidät, poikansa,. syntymämaa! A. V. A., kylmää kärsiä saivat, nälkää ja janoakin. kestivät heltehen tuskat ja marssien vaivat mielin riemuavin. Toisien vieläkin ammotti vanhat arvet polttaen, kirveltäin, kun kotikonnultaan taistojen tuoksinaan > lähtivät puolesta heimon ja syntymä / maan.}
\kat{Suomen Heimo}{15.08.1929}{P20}{Myöhäisestä ajasta huolimatta kokosi isäntäni kylässä majailevat pakolaiset kotiinsa kuulemaan Suomesta saapunutta vierasta, — Kuulijaksi siinä kuitenkin itse lopulta jäin, sillä puhuttuani rupesivat pakolaiset vilkkaasti kertoilemaan merkillisiä hakemuksia an ennen kaikkea Inkerin onnettomasti päättyneestä vapaussodasta. He olivat valloittaneet Kaprion linnoituksen lyöden ryissäläiset rykmentit hajalle, heidän joukkonsa olivat lisääntyneet 500 mieheksi, he olivat huonosti aseistettuina, nälkää kärsien kuukausimääriä taistelleet vihollisjoukkojen saartamina, valloittaneet YiLenmäen patterit, -nählneet jo Kronstadtinkin yllä valkoisen lipun, samonneet valkoryssäläisten emigranttijoukkojen dots}
\kat{Suomen Heimo}{15.09.1929}{P14}{EDISTYTTE taloudellisesti käyttämällä aina paikkakun- nan teollisuuden tuotteita. Nälkä siellä on esiintynyt kroonillisena ilmiönä koko kommunistisen hallituskauden ajan. Karjalan rahvas saa olla syrjässä ja odottaa mitä heidän yhteiskunnassaan tapahtuu.}
\subsection{Raja ja tiedonhankinta}
\subsubsection*{vakoilu \normalfont\small(30 katkelmaa, 36 osumaa aineistossa)}
\kat{Toukomies}{06.1926}{P10}{Toinen kirja, » Tsheka », on tärisyttelevä, räikeä kuvaus Venäjän nykyisestä salaisesta poliisista, kirjoitustavaltaan ei niin mieltäkiinnittävä kuin edellinen. vaikka tekijä liikkuu vasta vuosina 1918 — 1919; vangitsemisia tapahtuu, vakoilu kiristyy ja kiristyy ja yhä useampia sivistyneistöstä saa elämällään maksaa uuden onnentilan. Henkeä pidätellen lukee tekijän uhkarohkeita seikkailuja, matkoja Suomesta Venäjälle ja takaisin ja hänen toimiskelujaan tuossa uuden vapauden maassa oman maansa ( Knglannin ) salaisen tiedustelulaitoksen asiamiehenä.}
\kat{Suomen Heimo}{11.08.1930}{P28}{Hinta 40 mk., sid. Fenimore Cooperin " Vakooja ". Jatkoa sivulta 178.}
\kat{Toukomies}{01.06.1930}{P22}{Kukin osanottaja kussakin lajissa esittää vain yhden kappaleen. Nyt on » Vakooja » ilmestynyt I. K. Inhan suomentamana WSOY:n Koululaiskirjastossa. ja omintakeiset tuotteet on lähetettävä jäljennöskappaleina samalle henkilölle.}
\kat{Toukomies}{01.03.1931}{P11}{Smirnoff oli vetänyt taskustaan paperin. Kätyri — urkkija! ärjäsi Miihkali. Karjalais äidin kuolema.}
\kat{Suomen Heimo}{14.10.1933}{P7}{Emme väitä, että viranomaisten pitäisi vastaanottaa ehdottomasti kaikki Karjalan pakolaiset. Ei kukaan luonnollisesti asetu karkoitusta vastustamaan, jos pakolainen todella on kommunistien urkkija. Olen tietoinen, että täältä Kainuusta on tänä kesänä karkoitettu muitakin kuin nämä mainitsemani pakolaiset.}
\kat{Toukomies}{01.02.1934}{P8}{Viime syksynä paljastui täällä Suomessa suuri vakoilujuttu, joka sittemmin on osoittautunut ulottuneen moniin suurvaltoihinkin. Venäjältä Suomeen kohdistunut vakoilu tapahtuu luonnollisesti enimmäkseen Venäjällä olevien suomalaisten pakolaisten kautta ja välityksellä. Siitä on meillä valitettavasti tuoreitakin esimerkkejä, mitkä osaltaan ovat aiheuttaneet näiden rivien kirjoittamisen.}
\kat{Toukomies}{01.02.1934}{P8}{Kun karjalaiset joutuivat viimeksi vuoden 1922 alussa siirtymään Karjalassa puhjenneen vapausyrityksen aikana ja johdosta suurin joukoin Suomen puolelle, oli niiden joukossa vain vähän sellaisia, joihin täällä ei voitu poliittisesti luottaa. Tämä vakoilu siihen liittyvine likaisine ja häikäilemättömine keinoineen on erikoisesti vaarallinen Venäjää rajoittavissa valtakunnissa, joten niiden ennen kaikkea on oltava valppaita ja varuillaan sekä tietoisia tällaisen toiminnan vaarallisuudesta. Eikä suinkaan silloin tänne lähetetä sellaisia karjalaisia, jotka täällä entuudestaan tiedetään kommunisteiksi, vaan miltei yksinomaan joko täällä aikaisemmin olleita tai sitten Karjalassa pysyneitä dots}
\kat{Toukomies}{01.02.1934}{P8}{Jos me täällä Suomessa olevat karjalaiset voimmekin varmemmin olla vakuutettuja bolshevikivastaisesta mielialastamme ja luottaa itseemme ja toisiimme — sekalaista kuitenkin lienee seurakunta tässäkin suhteessa • — •, niin tätä ei voida varmuudella enää sanoa rajan takaa Karjalasta nyt tulleista tai tulevista pakolaisista. Tällöin luonnollisesti saattaa sattua erehdyksiäkin ja väärin on varmaan ajatella, että kaikki raj antakaa tulleet olisivat urkkijoita, mutta Suomen viranomaiset kyllä ovat oikeassa siinä, että kaikkia rajan yli tulleita on syytä epäillä, sillä eihän ole olemassa mitään varmuutta, kuka on todellinen, luottamusta ansaitseva pakolainen, kuka taas lähetetty vakooja. On siis}
\kat{Karjalan Vapaus}{04.1935}{P16}{Se tiesi matkaa Pietariin vangittuna ja sitten Siperiaan pakkotyöhön siksi kunnes siellä menehtyi. Siitä olivat urkkijat saaneet vihiä, ja jonkun ajan kuluttua tultiin vanha isäntä pidättämään. Venäläiseltä taholta oli tullut määräys kirkon sulkemisesta ja kaikenlaisten uskonnollisten menojen kieltämisestä.}
\kat{Suomen Heimo}{16.12.1935}{P19}{Ensimmäisenä ja polttavimpana taistelutehtävänä täytyy tällä hetkellä olla Suomen nuorison joukkomielenilmaisu Italian rosvousretkeä vastaan, Suomen taantumuksellisen porvariston hitleriläisfasistista taantumusta vastaan, Suomen kansan rauhan, vapauden ja riippumattomuuden puolesta. Ylioppilasnuorisossa, jonka keskuudessa fasismin vaikutus on ollut ja on vieläkin kaikista voimakkain, on viime aikoina ruvennut ilmenemään suurta vakoilua ja syviä ristiriitoja, jotka alkavat järkyttää natsihenkisen Akateemisen Karjala-Seuran yli kymmenvuotista ylivoimaista johtoasemaa. son velvollisuus on tukea Suomen Nuorison Liiton itsenäisyyttä ja loukkaamattomuutta fasististen kiihkoilijoin hyökkäyksiä dots}
\kat{Suomen Heimo}{23.06.1936}{P14}{( Ja me kun luulimme, että se oli opettamista varten ). Suomalaisilla tiedemiehillä onkin tutkimustyö pelkästään vakoilun ja propagandan naamioimista. Selostus perustuu nähtävästi Seuran omiin toimintakertomuksiin ja on yleensä asiallinen.}
\kat{Suomen Heimo}{23.06.1936}{P18}{Olen aina nähnyt ryssien kunnioittavan ulkonaista loistoa ja röyhkeyttä, jos se vain esiintyy sopivissa puitteissa. Eihän kenenkään päähän voinut pälkähtää, että englantilaisten vakooja ajelisi pitkin valtateitä troikalla Pietaria kohti. Saman henkilön elämänvaiheet ovat niin merkilliset, että niistä sietäisi saada perinpohjainen selvyys.}
\kat{Suomen Heimo}{15.09.1936}{P10}{Moskova on pilannut paljon osakkeittensa kurssia kansainvälisten kannattajiensa, kaikkien intoilevien " antifascistien " silmissä tuolla " kerettiläisvainolla ", joka on aukaissut monen silmät näkemään bolsevismin koko suvaitsemattomuuden, ahtauden ja julmuuden. Upseereja valvotaan tarkoin, poliisin keskuudessa harjoitetaan vakoilua, sotatarveteollisuuden salaisuuksia kavalletaan jne. Lisäksi Ranskassa kiinnitetään aivan erikoista huomiota katoliseen kirkkoon, jota pidetään taantumuksen ja fascismin todellisena tukilinnana. NE 8 OSASTOA, JOTKA AMERICAN Olla alituisesti ovat palveluksessanne A 1 ) Kemiallista pesua 5 ) Taideparsintaa. j\textasciicircum{}\textasciicircum{}P|i|g A|jt M 2 ) Prässäystä 6 ) Värjäystä Yrjönkatu io dots}
\kat{Viena-Aunus}{01.09.1936}{P18}{Suurin osa Petroskoin suomalaisia punapakolaisia ja osa paikallista väestöäkin eli melkeinpä yksinomaan » Hooverin komitean » jakaman amerikkalaisen avustuksen varassa. Eihän kukaan voinut varmuudella tietää, etteikö puhetoveri ollut tshekan salainen urkkija, tai valkoisten kanssa yhteydessä oleva vakoilija. Torilla näkyi korkeintaan joitakin epätoivoisia optimisteja, jotka kuvittelivat siellä saavansa vaihtaa syötävää vaateriekaleillaan tai vallankumouksessa säilyneillä arvoesineillään.}
\kat{Suomen Heimo}{20.01.1937}{P15}{Tämä yhteis- ja kansanrintamaryhmä on tullut jo niin rohkeaksi, että se esiintyy varsin röyhkeänä siellä, missä se on päässyt valtaan, esim. Helsingin sos. dem. kunnallisjärjestössä, jossa " vasemmistososialistinen " ryhmä on hyvin pinnalla. On varmaa, että näillä on suunnitelmat ja ohjeet, joiden mukaan toimien he hyökkäyksen meitä vastaan alkaessa aloittavat — ellei vallankumousta, niin ainakin tuhoisan puolustustamme turmelevan toiminnan, rintamantakaisten yhteyksien häiritsemisen, vakoilun ja terroriteot. Voimme luetella ainakin seuraavat toimintamuodot, jotka voidaan viedä SKP:n edellä selostettujen taitavasti toisiinsa kytkettyjen legaalisten ( = puuhailu " viattomissa " ja dots}
\kat{Suomen Heimo}{15.05.1937}{P10}{Kk:aan kuuluu seuraavat osastot: maaosasto. kansanvalistusosasto, finanssiosasto, sisämaankaupan osasto, terveysasiani osasto, yhteiskunnallisen huollon osasto, yleinen osasto, tieosasto, suunnittelukomissi. Synnyinmaan pettäminen: valan rikkominen, siirtyminen vihollisen puolelle, vahingon aiheuttaminen valtion sotilaalliselle mahdille, vakoilu — rangaistaan lain koko ankaruudella kaikkein raskaimpana rikoksena. " — Lopuksi huomautamme perustuslakiehdotuksen 17 §:n määräyksestä, jonka mukaan Neuvosto-Venäjän ja kaikkien, muiden SNTL:n liittotasavaltojen kansalaiset nauttivat Neuvosto-Karjalan alueella samoja oikeuksia kuin Karjalan ASNT:n kansalaisetkin!}
\kat{Suomen Heimo}{25.11.1937}{P11}{Takaisin pyörtäminen raskaine matkalaukkuineni olisi saattanut herättää Imatrallakin parveilevissa urkkijoissa epäilyksiä, ja niinpä päätin pistää pääni " leijonan kitaan " luottaen viattomaan naamaani ja kivuloiseen ulkonäkööni ja hiivin kuin hiivinkin junaan. M. m. jo iäkäs isoäitini oli neuvokkaasti kuljettanut niitä kainalossaan pahvirullaan kätkettyinä, eikä Viipurin asemalla kyttäilevät urkkijat osanneet aavistaakaan, että' pyylevä, puuskuttava, vanha vaimoihminen muodosti sillä hetkellä täydellisen liikkuvan " arsenaalin ". Olin Ensoon poikkeamatta keinotellut itseni Imatralle, missä m.m. alustavasti neuvottelimme ratkaisun hetkellä Vuoksenniskalla majailevien matruusien aseiden dots}
\kat{Suomen Heimo}{31.01.1938}{P3}{/. L.: Espanjan tragedia V 1 / 2 — 12, VI 3 — 39, VII 4 - 55, VIII 6 — 100. O. P.: Vakoilua ja vastavakoilua Venäjän vallan viime päivinä 18 — 142, II 9 — 157, 111 10 — 180, IV 11 — 195. Finno: AKS Jäisen ajatuksia 1 / 2 — 1, 3 — 29, 4 — 45, 5 — 65, 6 — 85, 7 — 109, 8 — 133, 9 \textasciitilde{} 150, 10 — 169, n — 185, 12 / 13 — 201, 14 — 233, 15 — 253, 16 — 269, 17 — 285, 18 — 301, 20 — 341.}
\kat{Suomen Heimo}{15.03.1938}{P9}{Takana väikkyy tietysti " Juudas " Trotskin hirmuhahmo, oikeastaan yhtä todellisena kuin manalan haamut Macbethissa. — Tämä yhtymä on järjestetty ulkovaltain toimeksiannosta ja sen tehtävänä on ollut harjoittaa vakoilua Neuvostoliitossa Juuri sitä kaikessa: syytettyjen entisen aseman ja henkilön, syytekirjelmän laajuuden, kaikkien tehokeinojen moninaisuuden kannalta.}
\kat{Suomen Heimo}{16.04.1938}{P5}{Niinpä ovat suomalaiset " kieltäneet " vierailta miltei kaiken yhteyden ruotsinkielisiin ylioppilaspiireihin. " Taitaapa toveri Välläri olla trotskilainen diversantti, ellei suorastaan tuholainen ja fascisti-Saksan vakooja? ' Täten lukija jää siihen käsitykseen, että venäläiset olisivat halunneet kaapata koko johdon käsiinsä.}
\kat{Suomen Heimo}{30.04.1938}{P12}{Kansan vapaus, veljeys ja tasa-arvoisuus näytti olevan vaarassa. VAKOILUA JA VASTAVAKOILUA VENÄJÄN VALLAN VIIME PÄIVINÄ Pakkomääräyksiä paitsi sensuurin suhteen annettiin myöskin matkustamisen suhteen.}
\kat{Suomen Heimo}{30.04.1938}{P13}{Torniossa on. kuvernööri määrännyt luovuttamaan siellä takavarikoidut tinar ja kumitavarat omistajilleen. Niin vakoilu kuin vastavakoilukin saatiin alulle kesän kuluessa ja se näyttikin toimivan aika ripeästi. Ennen vallankumousta toimivat rajavartiostopunktit santarmien apuna pidätyksissä ja kotitarkastuksissa sekä vartioinnissa.}
\kat{Suomen Heimo}{16.05.1938}{P11}{71 806 4; Hatunprässäystä ja korjausta fäSRStmSZZ 8 Puh. VAKOILUA JA VASTAVAKOILUA VENÄJÄN VALLAN VIIME PÄIVINÄ Harjoituksia oli kaarti pitänyt kerran viikossa, tavallisesti sunnuntaisin.}
\kat{Suomen Heimo}{31.05.1938}{P11}{Kun eräs kirjeenvaihtaja protestoi sen johdosta, että hänen sanomastaan oli pyyhitty pois punaisten tappiota koskeva kohta, sensori vastasi vain: " Hallitus ei koskaan kärsi tappioita! " iMiliisimiehet ja kaikenlaiset urkkijat kiertelivät autoissa, joiden kylkeen oli maalattu sana " Cheka ", tunkeutuen kahviloihin ja yksityisasuntoihin löytääkseen " Quinta Columnan " jäseniä, Francon salaisia kannattajia. Kansallisten armeijan ehdittyä pureutua kiinni Madridin laitaan hallitus, joka oli aikaisemmin kieltänyt kaiken poistumisen pääkaupungista, koetti saada naiset ja lapset pois, mutta ainoastaan harvat lähtivät, koska toivoivat kaiken pian olevan ohi.}
\kat{Suomen Heimo}{31.05.1938}{P14}{Koko vakoilutoimintaa johtaa Ruotsista käsin konsuli Edvard Stjenivall. VAKOILUA JA VASTAVAKOILUA VENÄJÄN VALLAN VIIME PÄIVINÄ VÄRVÄTTYJEN TIEDOITTAJIEN ILMIANTOJA. Kirje Sodankylästä 1 p. maalisk.}
\kat{Suomen Heimo}{30.06.1938}{P13}{Päätöksestä on annettu tieto kuvernööreille. VAKOILUA JA VASTAVAKOILUA VENÄJÄN VALLAN VIIME PÄIVINÄ Eri piioliiery liniien katsantokantoja Ruotsin avunannosta}
\kat{Suomen Heimo}{15.08.1938}{P29}{Aikaisemmin Venäjä, kun tyytymättömyys maassa oli yleinen, ryhtyi sotaan jotakin naapuriaan vastaan ja sai siten joukkojen huomion kiintymään muuhun. Milloin vain neuvostojärjestelmän rappiotila käy liian ilmeiseksi, vangitaan ja teloitetaan umpimähkään alempia virkailijoita, jotka muka ovat harjoittaneet vakoilua ja sabotaasia. Tiedustelijat, jotka ovat kfiyneet aina Kemissä ja Sorokassa asti, eivät ole voineet kertoa mitään sellaista, joka olisi antanut aihetta huolestumiseen ja suurempaan varovaisuuteen.}
\kat{Suomen Heimo}{15.12.1938}{P5}{Joo — jatkuvasti on tämä " Arvid Luhtakanta " ollut kotimaassa aktiivisessa vallankumoustoiminnassa, ja se kai " tulenkantajien " lukijoille on paras suositus " syyskauden ja joulun huomattavimmalle suomalaiselle alkuperäisteokselle ". Moskovassa ilmestyvä " The -new Russia " näet kertoo teloitetun saksalaisen vakoojan Fahrburgin tunnustaneen, että ns. kummituslentäjät ovat saksalaisia koneita, jotka haluavat tutustua Skandinavian pohjoisosiin, erittäinkin Haaparannan — Luulajan — Narvikin alueen liikenneyhteyksiin. Fahrburgin tunnustuksen johdosta ei ole syytä kiinnittää huomiota sellaisiin pikkuseikkoihin kuin ettei kukaan koskaan ole nähnyt saksalaisia lentokoneita matkailla dots}
\kat{Suomen Heimo}{15.12.1938}{P23}{Suomen kansan kaskut kuuluvat osana pilasatujen laajaan ja elinvoimaiseen ryhmään. Hän on " vieraan vallan vakooja, ja sellaisena häntä pitää myöskin kohdella. " Pistäytykää meille ja löydätte varmasti esineitä, joita voitte ostaa lahjoiksi omaisillenne.}
\kat{Viena-Aunus}{02.06.1938}{P13}{Tämä julma loppuselvitys sai armeijan kirjaimellisesti järkytyksiin. Siihen on viskattu vihaa, epäluottamusta ja vakoilua... Yksityiset ryhmät ja henkilöt alkoivat pyrkiä esittämään määrättyä osaa.}
\subsubsection*{tiedustelu \normalfont\small(30 katkelmaa, 37 osumaa aineistossa)}
\kat{Toukomies}{01.04.1925}{P6}{Mitä suolaan tulee, niin tunnustaa Nousiaryssä, että ryssät ( karjalaiset ), jotka asuvat Valkeanmeren rannalla, lähettävät Ensimmäinen tiedustelu ( käännös on paikotellen tahallisen ruotsinmukaista, — ehkäpä se antaa luonteenomaisemman käsityksen keskustelusta ): Sitävastoin kuinka pitkä matka on Petsistä Kuolansuuhun, ei hän tiedä, sen vain, että koilliseen Kuolansuusta luostari on.}
\kat{Toukomies}{01.04.1925}{P9}{Mikä vielä antoi nuorelle kotiseutujen runoperuja keräilevälle miehelle miettimisen aihetta, oli se, että runoja ei tähän asti kuljetulla matkallani ollut paljoa ilmestynyt. Sitten aamulla vielä turha tiedustelu kylän olemattomista runoista ja lauluniekoista, joista yksikantaan lausuttiin: » Ne veikkonen, on jo aikoja paremmille laulumaille menty. » Monissa on vanha kansanomainen kunto yhtynyt tarmokkaaseen pyrkimykseen, vakavaan taloudellisten oloauhteitten vaarinottoon sekä vihdoin yleisempäänkin kansallisten ja yhteiskunnallisten pyrkimysten kannattamiseen.}
\kat{Toukomies}{01.06.1925}{P16}{Jos rajaseutukysymys on hoitoonsa nähden yksipuolinen, on heimopakolaisten hoito tähän asti ollut liian hajanainen. Lyhyt puhe, joku tiedustelu ja joku tutustuminen, siinä kaikki, mikä on mahdollista sellaisella käynnillä. Tosinhan keskuksessa muina jäseninä lienee joku suomenpuolinen karjalainen, mutta ei yhtään rajantakais- ja pakolaiskarjalaista.}
\kat{Toukomies}{01.1926}{P9}{Tiedusteluja on lähetetty m.m. Ilomantsin suurten runolaulajain hautain selvillesaamiseksi. Karjalan Sivistysseura on pannut toimeen tiedustelun Suomen Karjalan runolaulajain hautojen tilasta. joka on Loimolan kalmistossa, mutta jo tuntemattomassa tilassa, kuten ilmoitus kuuluu.}
\kat{Suomen Heimo}{20.12.1929}{P28}{Jos me oletetun vaaran — sillä sen rohkenemme sanoa, että kommunistinen liike on meillä sittenkin asiallisesti niin vaaraton, että me jaksamme sen hyvin ilman ruotsalaistenkin apua että piakkoin annetaan eduskunnalle ehdotus tasavallan suojeluslaiksi, viskasi Antti Juutilainen esille tiedustelun, kuinkahan ruotsalaiset siihen suhtautuvat, kun se ilmeisesti kohdistuisi heidän äänestysystäviään vastaan. Jos nyt joku kansallismielinen suomalainen asettuu puoltamaan suomalaisuusliikkeen syrjäännakkaamista, jotta ruotsalaiset saataisiin mukaan kommunismin lopettamiseen, niin tuntuu siltä kuin hän liian myöhään astuisi ingmanilaisjin asemiin.}
\kat{Suomen Heimo}{16.12.1930}{P34}{Sen sai aikaan päivän juhlien arvokas ohjelmisto ja edelliseen vuoteen verrattuna paljon valoisampi sisäpoliittinen tilanne. Siitä ovat todistuksena ne lukuisat vastaukset, jotka on lähetetty Hallituksen Seuramme jäsenille tekemän tiedustelun johdosta. Lukukauden alussa lähetettiin kaikille oppikouluille kiertokirje, jossa toverikuntia kehoitettun perustamaan keskuuteensa Karjala- Seuroja.}
\kat{Suomen Heimo}{16.06.1931}{P26}{Tällaisella pienellä pahanteolla sai retkikunta aikaan, etteivät viholliset, jotka olivat pakottaneet Kuisman joukot perääntymään Uhtualta ja Vuokkiniemeltä, enää halunneet edetä Repolan alueelle. Loput sijoitettiin Repolan kirkolle, josta tarpeen mukaan lähti liikkeelle pitkän matkan tiedustelijoita, jolloin koetettiin saada selkoa tilanteesta Muurmanin radalla ja mahdollisimman syvältä vihollisen selän takaa. Samoin Alpo Ikonen, nykyinen Rautavaaran lukkari, joka tuli kerran yllätetyksi maatessaan omin lupinsa eräässä saunassa Akonlahdessa toverinsa Veli Heiskasen kanssa, mutta selviytyi kuitenkin tällä kertaa ehjin nahoin.}
\kat{Suomen Heimo}{30.04.1933}{P13}{Lainaamme tähän pari lausetta eräästä lehdestä tämän mielenkiinnon laatua mitä parhaiten kuvaavina: " Muutama vuosi takaperin oli Inkerissä hyvin vankka suomalainen propaganda. Ja aina riigikogua myöten on tämä mielenkiinto levinnyt, jota todistaa eräässä sen valiokunnassa hiljakkoin tehty, suomalaista vaikutusta Inkerissä koskeva tiedustelu. Niinpä kommunistiset aatteet ovatkin löytäneet lujan jalansijan varsinkin venäläiskylissä mutta osaksi myöskin inkerikkojen keskuudessa, minkä merkitys valtakunnan turvallisuuden kannalta ei ole suinkaan vähäpätöinen.}
\kat{Suomen Heimo}{10.06.1933}{P20}{Varsinkin kokoomuksen pää-äänenkannattajaan nähden täytyy sanoa: " vestigia ferrent ", kun muistelee sen päätoimittajan aikaisempia esiintymisiä yliopiston suomalaistamisasiassa. Tähän päätelmään antoi aiheen m.m. eräs tiedustelu, joka sanomalehden yleisön osastossa kohdistettiin IKL:n edustajaehdokkaille ja jossa kyseltiin heidän kantaansa yliopiston kieliasiassa. Koko kevään on Helsingin Yliopiston suomalaistaniidestä puhuttu, jos on kirjoitettukin, kokouksia on pidetty, komiteoita asetettu ja komiteat ovat istuneet ja antaneet lausuntonsa.}
\kat{Suomen Heimo}{29.08.1933}{P20}{Ja kuitenkin pyydämme vaivata arv. lukijaa palauttamaan mieleensä tuona ajankohtana tapahtuneen tasavallan presidentin vierailun Ahvenanmaalle. Mikäli meille on kerrottu, joukko AKS:läisiä lähetti taannoin Ajan Suunnassa julkaistavaksi seuraava IKL:n johdolle osoitetun tiedustelun: Teemme seuraavassa erinäisiä pistokokeita pakinanaihe varastoomme ja yhdymme niihin valitteluihin\textasciicircum{} jotka lukija kohdistaa käsittelemiemme asioitten tuoreuteen.}
\kat{Suomen Heimo}{29.08.1933}{P21}{Mikäli asiantuntevista lähteistä ilmoitetaan, on Selityksiä voidaan antaa, mutta ne eivät tyydytä kaikkia tiedustelijoita. O> SILMIENNE TYÖKAUSI ALKAA! y Hankkikaa itsellenne oikein sovitetut silmälasit.}
\kat{Toukomies}{01.05.1933}{P7}{Pääomat kasvavat kohdakkoin määrättyihin summiin. Asia on ollut keskustelun ja tiedustelun kannalla. — Johtokunnalla on toimivuoden ajalla ollut 6 kokousta.}
\kat{Karjalan Vapaus}{10.1934}{P5}{Komppanian oli määrä olla Matrosan itäpuolella hyhökkäysvalmiina toukokuun 2 p:nä klo 18,00 ja aloittaa hyökkäys heti kun tulesta päättäen oli todettavissa 3. K:n toiminnan alkaminen rintamasta käsin. Heti Prjäshan valloituksen jälkeen oli eteenpäin lähetetty tiedustelu todennut, että 16 — 17 km:n päässä Prjäshasta Petroskoihin päin oleva pieni Matrosan taloryhmä oli parilla konekiväärillä vahvistetun viholliskomppanian miehittämä. Tehtävä annettiin kahdelle reippaasta hyökkäyshengestään ja nopeasta toiminnastaan tunnetulle, uutta järjestelyä vastaavan I patl:n komppanialle, vänrikkien Heikkisen ja Juseliuksen komppanioille, vahvistuksenaan sekä raskaita että keveitä kk:ejä.}
\kat{Suomen Heimo}{28.04.1934}{P13}{Tällaisten maastoleikkien merkitys olisi varmasti maanpuolustuksen kannalta katsottuna mitä suurin. Tiedustelu juoksut ja hiihdot sopisivat sensijaan erinomaisesti edellämainittujen retkeilyjen yhteyteen. Ter vey sopin taas tulisi sisältää myös kaasusairaille annettava ensiapu.}
\kat{Suomen Heimo}{28.04.1934}{P13}{Fysiikassa taas Voidaan ainakin pitemmällä linjalla perusteellisemmin läpikäydä vino heittoliike ja muutkin sellaiset kohdat, jotka ovat tarpeen. Oppikoulukomitean mietinnössä on voimistelun ja urheilun oppiennätyksiin laitettu viimeistä edelliselle luokalle tiedustelu juoksua ja hiihtoa sekä viimeiselle luokalle lisäksi mahdollisesti ampumista salonkikiväärillä. Nykyään kun tykeillä ja heittäjillä ammutaan melkein yksinomaan epäsuorasti ja konekivääreillä taas ollaan kaikkialla maailmassa siirtymässä epäsuoraan ammuntaan, asettaa tällainen ammunta alemmille upseereille, jotka joutuvat sitä useinkin johtamaan, erinomaisen suuria vaatimuksia.}
\kat{Suomen Heimo}{28.04.1934}{P21}{Ei auta enää repäistä remseästi, sillä se on se kiihoituslaki sitä paljon " tulkintakysymys ", että saattavat vaikka jotkut tykätä pahaa näistäkin meikäläisen sivalteluista ja lykätä pian pään telkien taakse meidän näpertelymme ansiosta. Ja mutistuamme itseksemme vastaukseksi loppulauseen jälkimmäiseen momenttiin tiedustelun, että " Koska tässä ennen sitten juttu on jäänyt tekemättä ennen lehden ilmestymistä? ", aloimme kävellä kiltisti murjuumme ryhtyäksemme sanoista työhön. Samaan paikkaan missä Pohjois-Pohjalaisen osakunnan miehet ovat viime aikoina mainetta niittäneet jopa siinä määrin, että täytyy taaskin puhjeta puhumaan runokieltä ja siteerata kansallisrunoilijaa tällä kertaa sanasta dots}
\kat{Toukomies}{01.02.1934}{P22}{— Älä, veikkonen, pahastu, jos teen sinulle äkkinäisen ja suoran kysymyksen, mikä koskee sinun omia periaatteitasi. ■ — - Ennenkuin vastaan kysymykseesi, olisin halukas tietämään, mikä aiheuttaa moisen tiedustelun puoleltasi? ■ — Aika on ankara ja monen on pakko turvautua sellaiseenkin, johon ehkä ei haluaisi, puolusteli isäntä.}
\kat{Toukomies}{01.09.1934}{P19}{Myöntäen viimeksimainitun ainoan myönteellisen kohdan asialliseksi ja huomioonotettavaksi voimme sanoa, että muu lausunto ei vastaa sitä asemaa, jonka Keskusvaliokunta heimotyön eri järjestöjä edustavana elimenä on ottanut tai sanonut itsellään olevan, eikä edistä sen tärkeän asian kehittämistä, jommoinen heiinotyössä on pakolaisheimolaistemme kouluuttamis- ja valistamistyö varsinkin valtion yleisillä varoilla. Heimojärj. keskusvaliokunta on viimein elok. n p:nä antanut lausuntonsa Viipurin pakolaisavustuskeskukselle tämän sille viime keväänä lähettämän tiedustelun johdosta, olisiko sillä jotain muistutettavaa pakolaisavustuskeskuksen johdolla suoritetun heimolaispakolaisten valistu dots}
\kat{Suomen Heimo}{15.01.1935}{P21}{Kokous hyväksyi ponnet, joissa osoitettiin monia lakiehdotuksen puutteita. Saman tiedustelun teki Suomalaisuuden Liitto muille kansanedustajille. Konsistori hyväksyi 26. V. 1932 komiteansa laatiman ehdotuksen..}
\kat{Suomen Heimo}{15.03.1935}{P14}{Näitten kahden keskuksen väliset neuvottelutkaan eivät ole johtaneet tuloksiin, minkä vuoksi uusi sisällissota on taas ovella. Tiedustelijat poikkesivat Suomen puolella olevassa Ignoilan kylässä Kononovissa pääsiäislauantain aamuna ja jatkoivat matkaa Hyrsylän kautta Harjun Suitan Bek Bahtijaar oli Kashgarin riippumattoman hallituksen puoiustusasiainkomissaarin ja ylipäällikön apulaisena, Kashgarin komendanttina ja itsenäisyyskomitean varapuheenjohtajana.}
\kat{Viena-Aunus}{01.02.1935}{P13}{ole kokonaan ruumiillisesti tuhottu tai asuinsijoiltaan karkoitettu siihen aikaan päästessä, jolloin taas sinnepäinkin voidaan laillista yhteyttä pitää. Kalevala on alkuperältään venäläinen opus ja tuo teidän Lönnrotinne oli — suomalaisen imperialismin ensimmäisiä tiedustelijoita Karjalassa, muistakaa se! Kenties kaikki olisikin mennyt hyvin, jos olisin aikaisemmin tutustunut kommunistien ajatuksiin kaiken kirjallisuuden luokkaluonteesta ja muihin bolshevistisen » ideologian » kummallisuuksiin.}
\kat{Karjalan Vapaus}{03.1936}{P12}{Vähitellen tuntui koko taistelulinja olevan äänessä. Mies mieheltä saapuivat tiedustelijat ja kuiskasivat tietonsa johtajalle. Ryömien ja hiipien lähestyivät miehet ryssien nuotiota jokainen taholtaan.}
\kat{Karjalan Vapaus}{03.1936}{P12}{Ja viimeisenä juoksi Alikersantti, vilkaisten tuon tuostakin taaksensa, kaksi kivääriä olkapäillä remmistä ja kolmas laukaisuvalmiina kädessä. — Saivatpa „kotona " olijatkin jotain tekemistä! " tuumivat tiedustelijat, latasivat uudelleen kiväärinsä ja lähtivät juoksemaan poispäin. — Aikovatko nyt elävältä ottaa, pirut! kirosi Alikersantti, kääntyi ympäri ja laukaisi pari kertaa umpimähkään hieman viivyttääkseen takaa-ajajia.}
\kat{Karjalan Vapaus}{03.1936}{P16}{Mistä mahtoi olla kysymys, ei kukaan osannut sitä aavistaa. Arvasin, että tiedustelijat olivat liikkuneet tarkastelemassa parakkia. Hänen katseestaan huomasin, että hänellä oli myös sama ajatus.}
\kat{Karjalan Vapaus}{03.1936}{P31}{Nämä jätettiin sivuun ja kahlattiin taasen. Tiedustelu olisi saavuttanut siten tarkoituksensa. Ainoastaan pieni peltoaukea oli metsän välillä.}
\kat{Suomen Heimo}{15.04.1936}{P21}{Jokunen opettaja on tosin katsonut, että ovat ne vain " sinänsäkin " poliittisia, koska ne vievät joukkonsa vappukulkueisiin punaisten lippujen perässä kävelemään y.m. Opetusministeriöstä on nimenomaisen tiedustelun johdosta asiaa " tulkittu " niin, että nuo osastot " eivät sinänsä ole poliittisia yhdistyksiä " eikä kielto näin ollen tarkoita niitä. Kirjoittajan muisti poimii ihmeteltävällä tarkkuudella tapauksia ensimmäisten askelten ajoilta, johdattaa mieleen lapsen maailmankuvan hiljaisen ja ihmeitä täynnä olevan avartumisen ja viipyy onnellisena lapsen elämän suurissa tapauksissa.}
\kat{Suomen Heimo}{30.09.1936}{P8}{Voi kauhistusta, miten sotilaallinen merkitys sillä sanalla onkin! Kaiken sotilaallisen toiminnan edellytyksenä on erittäin tarkka tiedustelu. 2. Partiolainen täyttää aina velvollisuutensa kotia ja isänmaata kohtaan.}
\kat{Viena-Aunus}{01.12.1936}{P5}{Keväällä 191 7 virkosi vienan- ja aunuksenkarjalaisten kansallinen pyrkimys uuteen eloon, uusi kansallinen sivistysjärjestö muodostettiin entisen pohjalta, ja samalla alkoi toinen vaihe itäisten karjalaisten lehtien historiassa. Ne olivat kuin mainituille heimomaille lähetetyt tiedustelijat, vietinviejät täältä kansallisuutemme päämailta, ja samalla myös lähetetyt tiedustamaan, olisiko siellä vielä kansallista henkeä, vanhaa kalevalaissuomalaista tuntoa, — olisiko siellä vastaanottoa omille, omakielisille sanomille. Lämmin hengähdys tuulahti sieltä kansallisuutemme itäisiltä ääriltä, rajan takaa, Karjalan heimon vuosisatojen ajat eristetyiltä, unohdetuilta, vähän tietyiltä, sorron ja dots}
\kat{Suomen Heimo}{20.01.1937}{P18}{Ryssät yrittivät tammikuun alussa kaksi kertaa Kuusenvaarasta edeten valloittaa Kellovaaran kylää, mutta I / MSR. 11 / MSR ja IV / MSR löivät heidät verissä päin takaisin, ja vihollinen peräytyi Rukajärvelle poistuen säikäyksissään sieltäkin pariksi vuorokaudeksi. Täältä oli tosin tiedustelijoita lähetetty etelään, mutta he eivät melle.}
\kat{Suomen Heimo}{25.11.1937}{P27}{Tämän puheen johdosta huomauttaa Uusi Suometar, että " Kaikitenkin on kielikysymys elävältä haudattu, koska se, vaikka onkin julistettu polttopisteen ohittaneeksi, pääkohdilta perille ajetuksi, kuitenkin nousee ylös ". Pari päivää myöhemmin rehtori kylmästi ilmoitti molempia kieliä käytetyn, ja samana päivänä, kun tiedustelu Uudessa Suomettaressa julkaistiin, todella ilmestyikin taululle suomenkielinen ilmoitus — vieläpä samalle päivälle päivättynä kuin ruotsinkielinenkin. Mutta toiselta taholta, vieläpä yliopiston kateederista, julistaa heti vuoden alussa rehtori Hjelt avajaispuheessaan, että " kielitaistelu on sivuuttanut polttopisteensä ", ja tämä käsitys tuntuu sitten olevankin dots}
\subsubsection*{rajanylitys \normalfont\small(25 katkelmaa, 162 osumaa aineistossa)}
\kat{Toukomies}{05.02.1925}{P6}{Ja niin Novgorod lähetti nuoren ruhtinaansa Aleksanteri Jaroslavinpojan Nevalaisen sotajoukkoineen ja kaikkine apuvoimineen venein soutavia, ristein tenhoavia ja välkkyvin keihäin uhkaavia joukkoja vastaan. Se veljesviha, jota suomalaisen heimon leppymättömät periviholliset ikimuistoisista ajoista meidän päiviimme asti ovat rajaseudun kansaan kaikin keinoin lietsoneet, oli jo siinä määrin himmentänyt yhteenkuuluvaisuuden Ja siinä he olivat yhtä oikeassa kuin Ruotsin ja Suomen vallan — sekä sen ohella oman mahtinsa — laajentamista unelmoiva Tuomaspiispa, joka jäntevänä ja suuripiirteisenä miehenä näki, että niinpian kuin Nevan suu, Inkerin lukko ja Karjalan avain, oli valloitettu, oli idän dots}
\kat{Toukomies}{01.03.1925}{P5}{Siihen puoleen kuuluvat niiden laajat sanakirja- ja kielioppiosat. — Rajaseudun Kansankorkeakoulu, Joensuuhun perustettava kansansivistysahjo. Kielitieteilijöille kuuluu näiden teostenkielitieteellisen arvon lähempi määritteleminen.}
\kat{Toukomies}{01.09.1925}{P23}{Vain kuin se päivä lämmennyh kuttshuu, pilvilöiks ilmah hyö höyrytäh tuas. Zane Grey n uusi teos Rajaseudun henki, raikas ja tervettä voimaelämää uhkuva teos. Kiändi suuren angliskoin runonkirjuttajan Shakespearen ( Shekspiirin ) kai suuret näytelmöt suomekse.}
\kat{Toukomies}{12.1925}{P10}{— Eino Leino on alkanut julkaista muistelmia elämästään. tulokkaat melkein jokaisesta rajaseudun vanhasta karjalaiskalmistosta. On täällä Timovaaran laella kerran taisteltu oikea taistelukin.}
\kat{Toukomies}{01.1926}{P8}{Nyt 20 vuotta viimemainitusta lukien, ensi kesänä vietetään samassa kaupungissa kolmannet suuret laulujuhlat, tällä kertaa taas yleiset, Kansanvalistusseuran järjestämät. kylvöjään ja päästäneet ensimmäisiä opettajasarjojaan maailmalle, rajaseudun kansakoulut olivat alkaneet työnsä ja kaikinpuolinen uudempi elämä oli päässyt viriämään jo maanäärien saloillakin. Juuri näillä kohdin kulkee muun Suomen ja Savon vanha, nykyisissä käsityksissä vakiintunut raja Suomen Karjalaa vastaan, viistosti lounaasta koillista kohti, melkeinpä Salpausselän mukaisesti.}
\kat{Toukomies}{01.1926}{P19}{Ja kun vuorovaikutus loppuisi, silloin » hedelmöittäminenkin ». — Suomen Poh jois-Kar jalan ja rajaseudun siMri kansankorkeakoulu. Senlisäksi ovat pakolaisten lapset saaneet opetusta suomalaisissa kouluissa.}
\kat{Toukomies}{01.1926}{P24}{20 mk. ja Viimeinen ajo, nid. 28 mk., Rajaseudun henki, nid. — Konrad Lehtimäki: Yhdennellätoista hetkellä.}
\kat{Toukomies}{01.1926}{P24}{Norjalaisen saarnamiehen puheita, jotka ovat herättäneet vastakaikua kautta kaikkien pohjoismaiden. Viehättävä rakkaus- ja seikkailuromaani, jatkoa saman tekijän romaanille " Rajaseudun henki ". Kibunoi, suoraal a i z i e runoloi aunukse-laz i 1 I. Kääntänyt livo Härkönen.}
\kat{Toukomies}{01.1926}{P24}{Keskisuomalainen säveltäjä lähettää julkisuuteen jo viidennen kokoelman hengellisiä kaksinlaulujaan, jotka ovat saavuttaneet laulunharrastajain taholla suurta suosiota. Romaani, joka kertoo samojen henkilöiden vaiheista, joita kirjailija on kuvaillut romaaneissaan " Rajaseudun henki " ja " Viimeinen ajo ". Ylipainoksena tästä lehdestä on ilmestynyt yllämainittu ensimmäinen kokoelma, jossa on pantu alulle aunukseksi käännettyjen suomalaisten runojen suurempi valikoima.}
\kat{Toukomies}{10.1926}{P14}{— Liisa Rautamaa: Lolu ja hänen ystävänsä. — Zane Greyn uusi romaani: Rajaseudun sissit. Päivöi, paistele hellimbäh korbeh Suomeni kylmimbäh!}
\kat{Toukomies}{01.06.1927}{P9}{Maan akateeminen nuoriso ja muu sivistyneistö ovat muodostaneet voimakkaita seuroja ja yhdistyksiä, ajaakseen suomalaisuuden etuja ja oikeuksia. Suomen Rajaseutuyhdistyksen aikakauskirja " Rajaseutu " om, tehdäkseen maan rajaseutuja tunnetuksi, ryhitynyit julkaisemaan erifcoisnutaeroita rajaseutualueista. Joku aika sitten hallituksessa esillä ollut asia, heimoopetuksen ottaminen oppikoulujen ja seminaarien opetusohjelmaan, on saanut hylkäävän päätöksen.}
\kat{Toukomies}{01.06.1927}{P24}{Kustakin näistä alueista, niiden merkkiprkoista, oloista ja elinkeinoista, kulkuneuvoista ja matkareiteistä kilometrimäärineen tehdään lyhyesti selkoa ja samalla mainitaan se kirjallisuus, mikä kutakin matkaa varten olisi hankittava ja siihen tutustuttava. Oppaassa tehdään ensin selkoa siitä, mikä on rajaseutua ja ryhmitetään rajaseutualueemme seitsemään ryhmään, jotka ovat: Länsi-Pohja ja Länsi- Lappi, Itä-Lappi ja Petsamo, Koillinen rajaseutu, Kainuun rajaseudut, Pielisen ja Ilomantsin Karjala, Raja-Karjala sekä Karjalan kannas. Antti Vanninen: Karjalainen maisema, kansikuva; esittelykuva; uutisia; Kesän tullessa, kirj. prof. A. R. Niemi; Miten terva tulee alas Vienan rajoilta, kuva; dots}
\kat{Toukomies}{01.09.1927}{P22}{Koko neuvostot asav altain liitossa on yhteinen armeija, mutta Valko- Venäjällä on kuitenkin olemassa valkovenäläinen divisiona, johon kuuluu yksinomaan valkovenäläisiä. Samoinkuin tämän majan on koko erämaan ja rajaseudun sairasmajaketjun syntysanat lausunut maaherra Gustaf Ignatius, tehden alotteen sairasmajojen rakentamisesta pitkin maamme rajaa. NUIJAMIEHIEN MARSSHI. Meil on hangi ja jiä, meil on halla ja yö, meil on murdajat käskyt kohtalon: käzi meijän kui kenen muan tazah lyö, se muas jo on!}
\kat{Suomen Heimo}{20.05.1928}{P11}{Ja surullista on, että Länsi-Pohjan suomalaiset ovat langenneet ruotsalaisen kielipolitiikan hienosti asettamaan ansaan ja ihmettelevät vain, miksi Sitäpaitsi alkaa Suomen puoleinen rajaseutu määrätietoisen rajaseututyön tuloksena olla taloudellisesitö yhtä pitkällä kuin Ruotsin puoleinenkin, meneepä joissakin suhteissa edellekin. Siitä ovat todistuksena Tärännön ja Karungin kuuluisat papinvaalit, samoin useitten pitäjien kansakoulunopettajien ehdotus suomen kielen ottamisesta osittaiseksi opetuskieleksi kansakouluissa.}
\kat{Suomen Heimo}{15.07.1928}{P16}{A-kirjaimesta yksistään löysimme miltei hakematta seuraavat: Aarto, Erik, ingeniör; Ahola, Olga, fröken; Airio, Paavo, sekr. i Poststyrelsen; Arokallio, Gust., prost; Arvonen, Maja, fil. mag. ynnä muita. Kuvaavana piirteenä rajaseutu jemme vielä omituisille oloille mainittakoon, että A. K. -S:n 1. 7. 28 Terijoella vietetyssä muuten erittäin onnistuneessa kesäjuhlassa eräs ryssäläinen nainen 1 esiintyi humalaisena häiriten ohjelman suoritusta, niin että oli pakko. poistaa hänet juhlasta. Helsingin Lehden " puolue näyttää ajavcn puhtaasti kielipoliittista päämäärää, jonka toteuttamisesta ei myöskään kai liene varsin paljoa pahaa sanottavana, vaikkakin sikäläisen leirin aitosuomalaisuudella on dots}
\kat{Toukomies}{01.02.1928}{P22}{Teos on oikeata » kulttuurilukemista », tietorikas, pirteä, eurooppalaiseen henkeen kirjoitettu älykäs keskustelukirja, joka on merkittävä meidän oloissamme erikoislaatuiseksi uutuudeksi. Romaanin juonena on suomalaisen nuorukaisen ristiriita kahden rakkauden välissä: turmeltuneen venäläisen yläluokkaneitosen ja suomalaisen valkolakkitytön, joka viimein vapauttaa hänet rajaseudun kirouksesta. Maailmansodan aikainen elämä rintaman takana kuvataan tavalla, joka johtaa mieleen sellaiset kertomataiteen mestarit kuin Tolstoin ja Balzacin, joiden jättiläissuku jo näyttää kuolleen sukupuuttoon.}
\kat{Toukomies}{01.04.1928}{P2}{Tervehdysten jälkeen paavi Pius XI kasvoissaan hiukan iroonisesti hymyilevä ilme odottaa kysymyksiä, joita tulkki hänelle tekee. Suomen Heimori, Ylioppilaslehden, Aitosuomalaisen, Virittäjän ja Rajaseudun tämän vuoden numerot ja vihkot ovat ilmestyneet. Tuuli ensinnä puhalsi talon alas ja sieppasi sitte miehen, kantoi hänet ilman läpi pari sataa kyynärää ja asetti hänet omenapuun oksalle.}
\kat{Toukomies}{01.04.1928}{P6}{Ne ovat täten näkyvillä hengen siteillä sidottavat muuhun Suomeen. Rajaseudun väestö on edellä muiden kansallistutettava, saatava isänmaalliseksi. Sivistysjärjestöillemme ja suurille kustannusyhtiöillemme on tässä suuri ja}
\kat{Toukomies}{01.05.1928}{P4}{Mitä ne sitten sisältävät, mitä heimoasia ja heimotyö tietävät? Rajaseudun poikana hän tunsi tuon kansan puutteet, joita hän monella tapaa koetti poistaa. Karjala on meistä kuin tuo runon iso tammi, sen oksat ulottuvat pilviin, juuret syvälle maaemoon.}
\kat{Toukomies}{01.06.1928}{P28}{Koko Suomi tuntee hänet, mursukasvoisen — Napoleon Bonabarte Kenosen sana — halkopatruunan, jolle jo on ennättänyt kasautua vuosia hartioille, mutta jonka sydän on yhä nuori ja jonka edesottamukset ja teot todistavat, että vanhuudenheikkous on vielä kaukana hänestä. — Rajaseutu N:o 2 011 ilmestynyt laajana, runsaasti kuvitettuna Laatokan karjalaisten erikoisnumerona sisältäen m. m. seuraavat kirjoitukset: Laatokan Karjala; Piirtonen Karjalan asutushistoriaa, U. Karttunen; Kansanopetuslaiton Raja-Karjalassa, I. H. Holma: Maatalous ja sen mahdollisuudet Raja-Karjalassa, U. Bränder; Kreikkalaiskatolinen ortodoksinen kirkkokunta Karjalassa, 5. Solntseiv; Raja-Karjalan luonnonrikkauksista ja dots}
\kat{Suomen Heimo}{15.08.1929}{P16}{KIHLA-, vihkimä- ja jalokivisormukset, tarkkakäyntiset tasku-, ranne-, " seinä- ja herttysKELLOT, kulta-, hopea- ja alpaccateokset ostatte edullisimmin J. H. Solnin Kultaseppä- ja Kelloliikkeistä HELSINKI, I. Robertink. Mitä tämä vaikuttaa rajaseudun puolustusmahdollisuuksiin ja missä määrin se on omiaan kohottamaan rauhan aikana turvallisuuden tunnetta näillä seuduilla, joissa Suomella ei saa edes olla linnoituksia, on ilman muuta selvää. karia koskee, että se sijaitsee aivan Pietariin johtavan suuren laivareitin varrella, niin että sopimusta kyllä voidaan puolustaa tältä kohti sillä, että Venäjän kauppa ja merenkulku, välttämättä vaativat oikeutta saada kulkea vapaasti läheltä sanottua dots}
\kat{Suomen Heimo}{15.08.1929}{P21}{Emme puutu puheenaolevaan kirjoitukseen tällä kertaa sen enempää, sillä kuvaavaa on, että asiallisesti ei siinä ole mitään sellaista, joka kumoaisi ainoatakaan suomalaiselta taholta tehtyä väitettä olojen nurjuudesta Länsipohjassa. - Matarengin kuuluisan kansanopiston, Länsipohjan ruotsalajustuttamisen pääahjon rehtori, nimeltään Oarlström\textasciicircum{} lausuu kirjoituksen mukaan, että » ruotsinkielen opet-\textasciicircum{}---lainisen tarkoituksena on estää se, että kielierjoavaisuudet nostattaisivat muurin rajaseudun ja muun maan välille — ei enempää eikä vähempää ". Esimerkiksi Suomalaisuuden Liiton julkaisema Valde Näsin teos » Länsipohjan kysymys ", josta jo on otettu toinen painos, on reposteltu läpikotaisin, onpa dots}
\kat{Toukomies}{01.06.1929}{P10}{- N:o 6-' 1? RAJASEUDUN MATKAILUOLOT. Muutamat olivat jo viimeisellä luokalla.}
\kat{Toukomies}{01.10.1929}{P8}{— Kävin kaupungin varakkaimpia kauppiaita — Kuttua ja Tuhkaa kansansivistystyötä tekevä itäisen rajaseudun tutkimuslaitos, jonka kohottava vaikutus tuntuisi kauas. Pohjois-Kar jalan erikoislaatuisen sivistyslaitoksen luonne, saavutukset ja päämäärät.}
\kat{Toukomies}{01.10.1929}{P8}{kauppala, Nurmeksen maalaiskunta, Pielisensuu, Pielisjärvi, Polvijärvi, Pälkjärvi, Rautavaara, Rääkkylä, Tohmajärvi, Tuupovaara, Valtimo. Myöskin rajaseudun kirjastoliikkeen johto ajateltiin keskitettäväksi tähän laitokseen, josta vuosien vieriessä saattaisi ja pitäisikin muodostua laajaa ja syvää paan sivistystyön alalla ensiksikin tarjota valtiovallalle valmiiksimietitty rajaseutujen kansankohottamisohjelma ja sitten myöskin viipymättä — vaikkapa vähäisinkin voimin — ryhtyä sitä toteuttamaan.}
\subsubsection*{salakuljetus \normalfont\small(12 katkelmaa, 12 osumaa aineistossa)}
\kat{Toukomies}{01.10.1929}{P14}{9 p:nä 1864, ollen kuollessaan siis vasta 65 vuoden ikäinen. Onhan se Vävy-Eerokin, joka on kapteeni ollakseen ja salakuljettaja! Vähän jälkeen päivän nousun tuli Adelaide herättämään miestään}
\kat{Suomen Heimo}{16.12.1930}{P46}{40: —. H VAL. KATAJEV: Hummaajat, romaani, 25: —, sid. 38: —. | I EINO NYYSSÖLÄ: Salakuljettaja, kertomus, 25: —, sid. 38: —. | TEPPO KORVENKELO: Keisarin kaveri, kuvaus, 22: —, sid.}
\kat{Toukomies}{01.03.1930}{P16}{Puujalka nimittäin heti, kun tullivanhin oli valehdellut Matsedän Juonaan pakoyrityksestä, oli kohottanut ilmaan keppinsä, jota hän käytti aina siitä ajasta asti, kun puujalan sai, ja löi sen paukkeella pöytään tullivanhimman nokan edessä ja huusi: Minähän hyvässä uskossa sytytin tulen jaalaani, enkä osannut aavistaa, että yhtäkkiä sukeltaa pimeydestä kuin syliini salakuljettaja, jota minä olen etsinyt koko ikäni enkä ole löytänyt ja niin ollen jäänyt ilman viran- ja palkankorotusta — minä suuriperheinen mies, köyhä mies! Kepillään oli Puujalka löytänyt kujan pohjan, ja oli alkanut kävellä kotiaan kohti, mutta ei muistanut siltaa ja kujan mutkaa siinä, ja otti niinsanotun harha-askeleen dots}
\kat{Toukomies}{01.12.1930}{P26}{— Erkki Kivijärvi: Painomuste ja punaviini. — Eino Nyyssölä: Salakuljettaja, kertomus. ■ — /. Sinnemäki: Yli virran...}
\kat{Toukomies}{01.03.1932}{P11}{Tämän rajantakaisen papin usein uudistuvat matkat suurten tavarakollien kera saivat jatkua pitkän aikaa, sillä tullivartijat eivät voineet aavistaakaan, että pappi oli semmoisissa puuhissa. Vaikka tullivartijat olisivat olleet kuinka valppaita tahansa, niin salakuljettaja osasi melkein aina johtaa tullivartij an harhaan, ja viinakuormat saapuivat Suomen puolelle. Pietarilaisen » Vapaus » -lehden mukaan pitäisi tänä vuonna ryhdyttämän rakentamaan Ustj-Ishoran kylän lähettyville Nevan rannalle Inkerejoen suulle uutta » sosialistista kaupunkia » Kolppinan uuden metallurgisen tehtaan työväestöä varten.}
\kat{Toukomies}{01.05.1933}{P13}{Pidettiin m. m. ' suotavana, että ryhdyttäisiin aikaansaamaan sellaista toimitusten sarjaa, jossa julkaistaisiin huomattavain Karjalanystäväin elämäkertoja. Hilippa siirtyy nuoren vaimonsa kanssa kotikyläänsä Pankaj arvelle, hylkää salakuljetuksen ja alkaa elää maanviljelijänä onnellista ja rauhallista elämää. Niihin aikoihin on Hilippa joutunut elämään vainotun ja pakolaisen alamaa — hänen kotitalonsa on poltettu; herra Henrikki kiroaa tyttärensä ja ajaa hänet pois kotoa.}
\kat{Toukomies}{01.12.1934}{P15}{Olen kaiken parhaan ikäni jalkapatikassa tarponut 'ja kaahlinut pitkin pahoja, mutta aina kiinnostavia kinthipolkuja noilla Taka- ynnä kolatsut — sekä lisäksi » hauvottu maito » ■ — ■ kuuma kerma ( ei pirtu ) ) kuuluu ehdottomasti vienankarjalaiseen vier asp öy tään. — Tämä nuori slantilainen kirjailija kohosi äkkiä maineeseen romaanillaan » Morsiuspuku », joka viime jouluksi ilmestyi suomeksi ja saavutti meilläkin suurenmoiset arvostelut.}
\kat{Toukomies}{01.12.1934}{P16}{Basilier kirjoittaa v. 1886 Virittäjän ensi jakson II vihkoon innostuneen kirjoituksen » Shemeikäisistä » ( Shemeikoista ), joita hän pitää Salmin kihlakunnan ( Raja- Karjalan ) parhaana laulajasukuna, julkaisee saman vuoden lopulla U. Suomettaressa neliosaisen kirjoitussarjan » Kansakouluja Itä-Suomen kreikkalaiskatolisiin kuntiin » ja v. 1890 Suomal. Samana vuonna kun Koidon kertomukset ilmestyivät julkaisee Antero Kanerva ( A. Kor jonen ) Viipurissa kertomuksen » Salakuljettaja » Rajajoen tienoilta, ja Raja-Karjalassa uutisviljelijänä asuva Antti Kaksonen nimimerkillä Juuso Sortavalassa yhtenä vihkosena kaksi kertomusta: » Rämeen Paavo » ja » Antti Heimosen talouden vaiheita », joiden dots}
\kat{Viena-Aunus}{05.04.1937}{P22}{Johtokunta valitsi seuran varapuheenjohtajaksi prokuristi Juho Tiihosen, rahastonhoitajaksi johtaja Pekka Mattisen ja sihteeriksi ja toimistonjohtajaksi allekirjoittaneen livo Härkösen. Kuin haaskalinnut kerääntyivät tänne kaikenkarvaiset onnenonkijat ja salakuljettajat ympäri Afrikan, ja alkoivat kaikkia keinoja ja nykyajan keksintöjä ja välineitä käyttäen suurisuuntaisen timanttien salakuljetuksen. Mutta tästä toimenpiteestä eli seurauksena kova taistelu, joka kaikissa piirteissään suuresti muistutti Suomessa ja Amerikassa kieltolain aikana käytyä kamppailua laitonta spriin salakuljetusta ja myyntiä vastaan.}
\kat{Viena-Aunus}{05.04.1937}{P22}{Johtokunta valitsi seuran varapuheenjohtajaksi prokuristi Juho Tiihosen, rahastonhoitajaksi johtaja Pekka Mattisen ja sihteeriksi ja toimistonjohtajaksi allekirjoittaneen livo Härkösen. Kuin haaskalinnut kerääntyivät tänne kaikenkarvaiset onnenonkijat ja salakuljettajat ympäri Afrikan, ja alkoivat kaikkia keinoja ja nykyajan keksintöjä ja välineitä käyttäen suurisuuntaisen timanttien salakuljetuksen. Mutta tästä toimenpiteestä eli seurauksena kova taistelu, joka kaikissa piirteissään suuresti muistutti Suomessa ja Amerikassa kieltolain aikana käytyä kamppailua laitonta spriin salakuljetusta ja myyntiä vastaan.}
\kat{Viena-Aunus}{05.04.1937}{P23}{Kap-kaupungista pohjoiseen, maanviljelijänä muutamia vuosia sitten, liikkui sekä valkoisten että mustien kesken Namaqua-timantteja, jotka epäilyttävällä tavalla oli hankittu, ja tarjosi minullekin useasti joku kuljeksiva neekeri timanttia muutaman shillingin hinnasta, jolloin katsoin velvollisuudekseni tehdä siitä ilmoituksen piirin Kaikilla näillä toimenpiteillä, joihin vielä kuului lentokoneiden tarkastusmatkat erämaan teillä, saatiin muutamassa vuodessa timanttien salakuljetus rajoitetuksi siihen määrään, että se käytännöllisesti katsottuna voidaan pitää voitettuna, mutta jatkuvasti on vielä tänäkin päivänä pakko aina yhtä valppaasti jatkaa tätä » maavartiolaitosta ». 18. p:nä se kertoi dots}
\kat{Viena-Aunus}{05.04.1937}{P26}{Yksityishuonettakin oli esiintyjälle tarjottu tehtaan puolesta, koska hän oli lyönyt jonkun ennätyksen kutomisessa, mutta hän ei suostunut siihen vaihtamaan asumista yhteisessä asunnossa toverittarien ja toverien kanssa. — • Suomi ja Ruotsi ovat vaatineet virolaisen spriin salakuljetuksen lopettamista, ja sen ehkäisemiseksi on nyttemmin tehty näiden kolmen maan välinen sopimus, jossa Suomi ja Ruotsi sitoutuvat ostamaan Virolta määrätyn summan arvosta vuosittain spriitä, ehdolla, että senj aikeen on Viron spriin salatuonti täysin lopetettava. 6. Heitnovirsi.}
\subsubsection*{loikkari \normalfont\small(20 katkelmaa, 20 osumaa aineistossa)}
\kat{Suomen Heimo}{18.11.1928}{P18}{— 1 p:nä syyskuuta ilmestyi suuri bolshevikkien laivasto-osasto vedenalaisineen aluevesillemme Lavansaaren rannikolle. Tehtäköön ryssien äskeinen luvaton vierailu huvipursilla rannikoillamme miten viattomaksi tahansa, ei se tosiasiassa sitä kuitenkaan ole. Suursaaren maihinnousu on koetettu selittää viattomaksi huvimatkaksi, ovatpa " ryssät — tapansa mukaan valehdellen — sanoneet saaneensa jostakin suullisen luvan.}
\kat{Toukomies}{12.1928}{P26}{— Ensikertaa esiintyvän kirjailijan kertomusteos, jonka aiheena ovat ihmisten salaperäiset, mystilliset kyvyt, elämän järkkymätön syyn ja seurausten laki ja ihmisen paremman ja huonomman minän kamppailu. Nyt juikaistussa kokoelmassa on siis paitsi laajaa runokokoelmaa seuraavat Kiven teokset: Kullervo, Nummisuutarit, Kihlaus, Karkurit, Alma, Yö ja päivä, Lea, Koto ja kahleet sekä Seitsemän veljestä. Werner Söderström Osakeyhtiö on ryhtynyt julkaisemaan sarjaa, johon aikaa myöten tulevat kaik- Men kulttuurikansojen sankarisadut. ■ — Ensimmäisena osana on julkaistu Claes Lindskogin kokoamana Kreikkalaisia Jumalaistaruja ja satuja.}
\kat{Karjalan Vapaus}{09.1934}{P23}{'Kulissit " ja " esirippu " kiinnitettiin seinästä seinään vedettyyn nuoraan " pyykkipojilla ". Loikkari, eräs neitonen, joka oli näytelmän nähnyt moneen kertaan, oli kertonut koko historian juurta jaksain. Kuitenkin kaikitenkin, voin kehumatta sanoa, että näytelmää esitettiin yhä uudestaan ja uudestaan ja jopa joltisella menestykselläkin.}
\kat{Suomen Heimo}{15.11.1934}{P12}{Merkki on erittäin kaunis, AKS:n merkin muotoi- nen, hieman sitä pienempi ja on siinä samat ( Karjalan ) värit kuin AKS:n merkissäkin, kuitenkin ilman miekkaa. Eliel Aspelin-Haapkylä SEITSEMÄN VELJESTÄ KOTO JA KAHLEET I 1 os a näytelmiä: KULLERVO NUMMISUUTARIT KARKURIT KIHLAUS, Yö JA PÄIVÄ LEA JA MARGARETA Suomalaisen Kirjallisuuden Seura 207}
\kat{Karjalan Vapaus}{03.1936}{P11}{Ojia ja aitovarsia painuttiin joen rantaan. Kuiskaten selitti Karkuri miehilleen suunnitelman. Ensin kuitenkin saappaat tyhjiksi}
\kat{Karjalan Vapaus}{03.1936}{P12}{Kosken kohina, takaa-ajajien kiväärien paukkina ja kuulien vinkuna korvissa sulautui helvetilliseksi musiikiksi. Karkuri oli saanut joukkonsa jo omalle rannalle ennen, kuin huomasi kolmen miehen puuttuvan joukosta. Karkuri viittoi miehet kokoon, antoi niinikään viittomalla jokaiselle määräyksensä tiedustelusuunnasta ja paikasta.}
\kat{Karjalan Vapaus}{03.1936}{P12}{Karkuri oli saanut joukkonsa jo omalle rannalle ennen, kuin huomasi kolmen miehen puuttuvan joukosta. Karkuri viittoi miehet kokoon, antoi niinikään viittomalla jokaiselle määräyksensä tiedustelusuunnasta ja paikasta. Huudot ja komennukset kaikuivat, laukaukset pamahtelivat ja maa tömisi takaa-ajajien askelista.}
\kat{Karjalan Vapaus}{03.1936}{P31}{Tuvan ovea koetettiin ulkoa päin avata, mutta se ei auennutkaan. Karkurit ryömivät edelleen niin kauan, kunnes kuului selvästi puhetta. Kahdesti putosi kiväärin tukki miehen pääkalloon, jolloin vastus loppui.}
\kat{Karjalan Vapaus}{03.1936}{P31}{Kylässä nousi hirmuinen elämä, kun siellä huomattiin molempien vankien karanneen. Karkurit olivat ilman suksia, ja maasto oli aivan tuntematon yön pimeydessä. Vuorinen, käsittäen matkan jatkamisen turhaksi, kehoitti toveriaan menemään yksin.}
\kat{Suomen Heimo}{25.11.1937}{P33}{Pohjanpäänkin kirjassa on rakkautta, mutta hän kuvaa puhdasta ja pyyteetöntä nuoruudenrakkautta. Niitten useimpien probleemana on joku, tavallisesti luvaton, rakkaussuhde ja sen käsittely. " Päivä laskee länteen " on Väinö Piiraisen teoksen nimenä, mutta ei kuvaa mielestäni sisältöä.}
\kat{Suomen Heimo}{31.05.1938}{P14}{Kansankomissaarien neuvoston puheenjohtaja Uljanov on antanut Valkeasaaren komissaarille käskyn pidättää kaikki Suomeen osoitetut rautatielastit, paitsi sotilastavarat. Mikäli minä tiedän, ovat karkurit käyttäneet useampia läpi Sodankylän kulkevia salopolkuja, mutta suurin osa heistä on kulkenut Sodankylän läpi. Hän tietää neuvoa, missä Jenaro tai Jetsharo järvi on ( tarkoittaa Inaria, mutta kun on ryssä niin ei tiedä missä Inari on ).}
\kat{Suomen Heimo}{20.06.1939}{P25}{Partiot Kemppilä ja Syvänne valmistuvat lähtöön! määräsi kersantti. Oli taasen sattunut Venäjältä vieraita, " loikkari " ryssästä. Muuallakin on sattunut verta vaatineita yhteenottoja saksalaisten ja tsekkien välillä.}
\kat{Suomen Heimo}{20.06.1939}{P25}{Aseistus ja patruunat lähtömuonan päälle ensiksi, sitten ruokatavarat, teltta ja leirikalut, määräys ja lähtö. Hautavaaran torppari on syöksähtänyt vartiotuvalle ja huutanut: " Loikkari — ryssästä loikkari, ja kaamea kappale! Sekä saksalaiset että protektoraatin omat viranomaiset ovat antaneet varoituksia väestölle pysyä erillään saksalaisvastaisesta kiihoituksesta.}
\kat{Suomen Heimo}{20.06.1939}{P25}{Aseistus ja patruunat lähtömuonan päälle ensiksi, sitten ruokatavarat, teltta ja leirikalut, määräys ja lähtö. Hautavaaran torppari on syöksähtänyt vartiotuvalle ja huutanut: " Loikkari — ryssästä loikkari, ja kaamea kappale! Sekä saksalaiset että protektoraatin omat viranomaiset ovat antaneet varoituksia väestölle pysyä erillään saksalaisvastaisesta kiihoituksesta.}
\kat{Suomen Heimo}{20.06.1939}{P25}{Toinen askel kantaa, toinen askel tuuskautti nokalleen ja rajamies kastuu: askel otti ryypyn ja seuraava toisen ja selkäpiitä karmivat kevätkelin tuntumat, vilustua ja rasitukset. Nuokkuu ja horjuu ja huutelee leipää — leipää —.. < —! Se on meille tuttua: Saariselän vartiossa ollaan jo totuttu näihin tämänlaisiin uutisiin: loikkari ryssästä —! Niinpä protektori käytti ensimmäisen kerran veto-oikeuttaan kumoamalla hallituksen ehdokkaan nimittämisen " hintadiktaattoriksi ", koska tämä ehdokas oli liian iähellä työnantajapiirejä, eikä siis voisi puolueettomasti valvoa myös työläisten etuja.}
\kat{Suomen Heimo}{20.06.1939}{P26}{Etelästä joutunut kurkiparvi nokki rajasuolla karpaloita. Tämä on nyt se loikkari —! Osan hinta 45: —, yhteensid. kangask.}
\kat{Suomen Heimo}{20.06.1939}{P26}{272 sivua 75: —. Sigrid Boon tämänkeväinen uutuus on " MUTTA MUUTEN MEILLÄ OLI HAUSKAA ". Mutta samassa on velipekankin aivoissa helähtänyt tilanteen todellinen ja vakava merkitys: loikkari ja rajamies. Hän huoahti ja heitti nyt reppunsa ja aseensa emäkuusen juurelle ja viritti sitten nuotion.}
\kat{Suomen Heimo}{20.06.1939}{P26}{VIHREÄ VERÄJÄ " on " Mat- kan pään " tekijän valoisa englantilainen perheromaani, 320-sivuinen raikashenkinen romaani on Suomen Kuvaleh- den kuponkikirja. Siellä oli loikkari: se kuuhki tuolta korvesta hoippuen ja harhaillen — miehenkuva lähestyi Kurkijärven perukkaa Ryssästä käsin ja etsiskeli suolta syötävää. Ja siinä se nyt oli jälleen: tämä Kurkijärven suunta oli tuonut hänelle raskaita, onnettomia muistoja, velimies ja sitten hänen omat kärsimyksensä.}
\kat{Viena-Aunus}{20.12.1939}{P18}{Mummo valvoi ja käänteli itseään tilallaan. Miksi minun piti sekin luvaton työ tehdä. Siinä ei ollut tulta ollut moniin vuosiin.}
\kat{Heimotyö}{1939}{P98}{Kurjet. Karkuri. Romaani.}
\subsubsection*{provokaattori \normalfont\small(9 katkelmaa, 9 osumaa aineistossa)}
\kat{Toukomies}{01.12.1930}{P26}{Sen arvoa lisäävät eri tiedemiesten ja erikoisasiantuntijain laatimat, retken kulkua valaisevat katsaukset, piirrokset, kartat ja lukuisat hyvät kuvat. Kirpeä, peittelemätön kuvaus Neuvosto- Venäjän oloista, joissa n. s. hummaajat — valtion varojen kavaltajat — ovat muodostaneet melkein oman ammattiluokkansa. Ruotsin Antropologis-Maantieteellinen Seura sai tehtäväkseen hallituksen puolesta huolehtia löydöstä sekä — vähän myöhemmin — retkikunnan vaiheita kuvailevan julkaisun toimittamisesta.}
\kat{Viena-Aunus}{01.12.1935}{P17}{Japani on ilmoittanut valtaavansa Ulko-Mongolian, jos se neuvostohallituksen vaikutuksen alaisena uhkaa Mandshuriaa ja Pohjois-Kiinaa, jotka Japani aikoo suojella punaiselta valtaukselta. Sodan kestäessä on Abessinian puolella jo ilmennyt yksi kavaltaja, maansa pettäjä Gugsa, joka on jonkun satakunnan miehen kera siirtynyt italialaisten puolelle. Abessinian sisämaiden vaikeapääsyisissä rotkoissa ja huimaavan korkeissa vuoristoissa odotetaan Italian kohtaavan todellista sodanleikkiä tähänastisen melkein vastarinnattoman etenemisen sijasta.}
\kat{Suomen Heimo}{31.08.1937}{P16}{Johtavissa asemissa työskentelevät edelleen mm. Kivekäs, S. Kokko ja I. Heikka, jotka ignoreeraavat Leninin- Stalinin kieltä, eivät edes kykene selvittämään kilpailijoille alkeellisia ohjeita ja sääntöjä. " Metalliverstaalla työskentelee murhaoppositiolainen, puolueesta eroitetu Puhakka ja tunnettu Suomen työväenliikkeen kavaltaja, Juudas-Trotskin suomalaisen apurin Vällärin käsikassara Kalle Söderström, joka tekee salaista hajoitustyötä työläisten keskuudessa ja kerää ympärilleen natsionalistisia aineksia. " Tällä pyritään tietenkin tavallaan kohottamaan karjalaisten itsetuntoa ainakin suomalaisiin nähden, mutta toisaalta on sivutarkoituksena. kuitenkin todennäköisesti se, että Tverin dots}
\kat{Suomen Heimo}{16.05.1938}{P21}{Saksalaiset eivät näet suinkaan ole ainoa tyytymätön väestöryhmä, vaan muutkin vähemmistöt samoinkuin suuri osa virallisesti " tsekkoslovakkilaiseen " isäntäkansallisuuteen luetuista slovakeistakin on kiivaasti oppositiossa. Vaikka Henleinin miehet kaikesta päättäen ovat kuriin tottunutta väkeä, kieliryhmien keskeinen ärtymys sekä nähtävästi myös provokaattorit ovat aikaansaaneet jo monia välikohtauksia, jotka tietysti suuresti vähentävät ristiriidan selvittämismahdollisuuksia. 4422.}
\kat{Suomen Heimo}{15.12.1938}{P3}{Yhtenä syynä siihen, että kansa ei suhtaudu kyllin tuomitsevasti henkirikoksiin, lienee juuri tämä " reiluus ". Miehcntappajat ovat usein luonteeltaan suorasukaisia ja raakuudessaankin rehellisempiä kuin murhaajat, kavaltajat, varkaat tai siveellisyysrikoksien tekijät. Kun rehtori Kauppinen ei suositellut ilmiantamista, on hänet otettu ministerin mustiin kirjoihin malliesimerkkinä sopimattomasti käyttäytyvästä opettajasta.}
\kat{Viena-Aunus}{17.01.1938}{P19}{Silloin yritettiin ryövätä tunnettu kommunisti Friedrich Miintzenberg, joka kuitenkin tarpeeksi ajoissa sai vihiä suunnitelluista aikeista. Hänhän, kuten kaikki provokaattorit oli paljastetuksi tultuaan vaarana, ja sentähden piti hänet saada katoamaan. Burtzeff sanoo huomanneensa heti kenraali Kutjepoffin katoamisen jälkeen, että eräs kenraalin piireissä liikkunut naishenkilö on GPU:n agentti.}
\kat{Suomen Heimo}{28.02.1939}{P9}{Mahtaa nyt tuntua pääministeri Cajanderista ja kunnallisneuvos Mandiata hauskalta ajatella näin: Osa kommunisteja äänesti Lassilaa, mutta äänet eivät riittäneet. Nimittäin, että se mies, joka sitten hänen plakaattejaan liimaili, se kommunisti, olikin edistyspuolueen oikeiston kätyri ja provokaattori. Ellei tämä tai muu seikka muuta kehityksen suuntaa, olemme jonkin vuosikymmenen kuluttua siinä asemassa, että meidän on todettava saavuttaneemme väkilukumme ylärajan.}
\kat{Suomen Heimo}{08.07.1939}{P16}{Tämän tunnustavat ei vain ystävämme vaan myöskin vihollisemme. Vakoojat, kavaltajat, diversantit ja murhaajat juoksentelivat kaupungissa. Roistomaisin kaikista kavaltajista Juudas-Trotski tietoisesti saattoi epäjärjestykseen rintamaa.}
\kat{Suomen Heimo}{08.07.1939}{P16}{Joka tapauksessa voin sanoa, että Pietari on valmis kaikkiin mahdollisiin vihollisen taholta tuleviin yllätyksiin. Vallankumouksen kavaltajat, jotka olivat pesiytyneet 7. armeijakunnan esikuntaan, tekivät parhaansa luodakseen viholliselle otollisen tilanteen. Toiseksi tilanne itse Suomessa ja Eestissä ei ole edullinen valkokaartin muodostamiseen, sillä siellä on vallankumouksellinen liike laajenemassa.}
\subsection{Pakolaisuus}
\subsubsection*{kansalaisuus \normalfont\small(15 katkelmaa, 15 osumaa aineistossa)}
\kat{Suomen Heimo}{15.04.1929}{P12}{Suomi on yhtäkkiä hyljännyt suurimman osan kansalaisistaan itäisessä naapurimaassa, ja nämä menettävät senkin vähäisen turvan, mikä Neuvosto-Venäjällä ulkomaan kansalaisilla on verrattuna omiin kansalaisiin. Kesäkuun 17 päivänä 1927 annettu laki Suomen kansalaisuuden menettämisestä on saanut — varmasti ilman, että nyt tapahtumassa oleva on tarkoitettu, — sellaisen muodon, että ne Suomen rajojen ulkopuolella asuvat henkilöt, jotka eivät ole käyneet vähintään kahta vuotta koulua Suomessa tai suorittaneet täällä asevelvollisuuttaan, menettävät Suomen kansalaisuutensa, ellei heille anomuksesta anneta oikeutta edelleen Suomen kansalaisena pysyä. Tätä puolta ei varmaankaan ole tarpeellisella dots}
\kat{Toukomies}{01.05.1933}{P9}{Vihaa syventää vielä se seikka, että suomalaiset suhtautuvat virolaisiin ylemmyyden tunteella lausuen julkisesti, että virolaiset rypevät niin paljon venäläisten hännässä, ettei heistä ole mitään hyötyä vähemmistökansallisuuksien etujen puolustamisessa. Mutta mielenkiintoisempaa ehkä on tutustua siihen, mitä entinen virolainen kommunistijohtaja kertoo n. s. » suomalais-virolaisista riidoista ja tulevaisuuden haaveista » sekä pienen vepsäläisen kansalaisuuden » löydöstä » Neuvosto-Karjalan eteläpuolelta ja osittain sen sisäpuolelta. » Uuden Suomen » julkaisemassa toisessa sarjassa näitä muistelmia ( Sunnuntailiite n:o 17, 23 / 4 1933 ) Kasemets antaa ensinnäkin katsauksen suomalaisten dots}
\kat{Toukomies}{01.06.1933}{P5}{myös sydämeeni, joka synkkä on, kun olen orpo, maaton, osaton. Tosin sanotaan, että uusi kansalaisuus ehkä heikentää aatteellisesti pakolaisuuden voimaa. Tuo toteutuvan näy ei haavehein, ei auringossa sula kahlehein.}
\kat{Suomen Heimo}{03.03.1934}{P11}{Mutta näihin kaunista kotiseututunnetta ilmaiseviin alotteisiin enempää kuin moniin muihinkaan emme tahdo tällä paikalla puuttua, kiinnitämmepä vain huomiomme erinäisiin suomalaisuusasiaa eteenpäinvieviin esityksiin ja lausumme julki tyytyväisyytemme niiden johdosta. Lakialotteen perusteluissa huomautetaan m. m., että koska Suomen kansalaisuuden saavuttamiselle on tähän saakka kielellisenä vaatimuksena ollut joko suomen tai ruotsin kielen taito, on ylivoimaisesti suurin osa Suomen kansalaisiksi pyrjiöistä hankkinut itselleen vain ruotsin kielen taidon. Tässä kohden on siis kaikin puolin hyödyllistä ja ruotsinkielisen opiskelevan nuorisonkin edun mukaista, että sen on dots}
\kat{Suomen Heimo}{20.01.1936}{P15}{Erikoisen huono on niiden pakolaisten asema, jotka elävät Kainuun ja Pohjois-Karjalan rajaseuduilla. Toivottavasti asia järjestyy, kun kansalaisuutta koskeva lainsäädäntö lähiaikoina uusitaan. Kun rajaseutuja muutenkin tuetaan, olisi samalla muistettava myös rajaseudun pakolaisia.}
\kat{Suomen Heimo}{29.02.1936}{P18}{Bolsevikkiteoreetikot pitävät federalismia periaatteena, jonka pitäisi olla omiaan edistämään muiden kansojen liittymistä Neuvostoliittoon. Huomattava on sekin, että valtiosääntö tuntee vain Neuvostoliiton kansalaisuuden eikä osavaltioilla siis ole omia kansalaisia. Näissäkin on ylimpänä elimenä oma neuvostokongressi, ja sen ohella muut elimet aiva samaan tapaan kuin jäsentasavalloissa.}
\kat{Suomen Heimo}{25.11.1937}{P12}{Pakolainen on sekä henkisesti erillään ympäristöstään että taloudellisestikin näytään toisinaan häneen suhtauduttavan liian kylmästi. Heimotyössä luonnollisesti ei Suomen kansalaisuuden saavuttamisella ole suurtakaan merkitystä eikä siihen ole kiinnitettävä huomiota. Ja sehän on jokseenkin jokaisen itäkarj alaisen ja inkeriläisen osana ollut, jotka ovat turvapaikan maassamme löytäneet.}
\kat{Suomen Heimo}{25.11.1937}{P12}{Vertauksen vuoksi mainittakoon, että muita Venäjän pakolaisia oli maassamme vuoden 1930 lopussa 7,982 ja viisi vuotta myöhemmin 6,912. Siinä on siis tapahtunut 770 hengen vähennys ja samaan aikaan on hyväksytty Suomen kansalaisiksi 902 henkeä. Niissä ei nimittäin ollenkaan ole lueteltu niitä syntyperältään heimopakolaisia, jotka ovat saaneet Suomen kansalaisoikeudet, yhtä vähän kuin niitä naisia, jotka ovat menettäneet ulkomaalaisasemansa solmimalla avioliiton Suomen kansalaisuutta olevan miehen kanssa. Että tällaisten tapausten lukumäärä ei ole mikään pieni, selviää siitä, että yksinomaan Inkerin pakolaisseurakunnan vuosikertomusten mukaan kuulutettiin viisivuotiskautena 1931 — 35 dots}
\kat{Suomen Heimo}{25.11.1937}{P19}{Eräs englantilainen on kuvannut maansa sivistyneistöä seuraavasti: Kaikille oli yhteistä pyrkimys siveellisen luonteen ja oikean kansalaisuuden kasvattamiseen. Tämän päämäärän ovat englantilaiset sattuvasti kiteyttäneet tuohon vanhaan sanaan gentleman.}
\kat{Suomen Heimo}{31.01.1938}{P2}{Életnér Bako: Paljonko unkarilaisia asuu Tsekkoslovakiassa 16 272 H. J:nen: Itä-Karjalan pakolaiset ja Suomen kansalaisuus 14 244 Olavi Lähteenmäki: Vapautuksen kevään muistot ja velvoitukset 9 154}
\kat{Suomen Heimo}{15.09.1938}{P14}{Se on kuitenkin huomattu käytännössä vaikeaksi toteuttaa. ITÄ-KARJALAN PAKOLAISET JA SUOMEN KANSALAISUUS Viimeinen amnestia Karjalan pakolaisille luvattiin vuonna 1925.}
\kat{Viena-Aunus}{07.05.1938}{P21}{Tämän seurakunnan keskus, Vaasan kaupunki, sanoo hän, on kovin kaukana ja itse seurakunta hyvin laaja. Luonnollisesti on otettava myös huomioon kansalaisuuden tuomat yleiset oikeudet, jotka ovat suuremmat kuin ei-kansalaisilla. Tshekkoslovakian unkarilaisten järjestö, joka aikaisemmin on ollut kielletty, on pitänyt kokouksensa monen vuoden jälkeen.}
\kat{Viena-Aunus}{07.05.1938}{P21}{Tämä seurakunta taas, laajuudestaan huolimatta, on niin vähäjäseninen, ettei useampia pappeja ja " keskuksia siihen voine ajatella mahdolliseksi asettaa; sitäpaitsi se korottaisi verot vielä suuremmiksi. Pakolaisseurakunnalle maksettavat verot, kuten kirjoittaja mainitsee, saattavat olla vähäisempiä tai olemattomia, mutta Suomen kansalaisuus asettaa samanlaiset velvollisuudet kaikille kansalaisille, niin kirkollisessa kuin muussakin suhteessa. Näin ollen emme voi antaa kirjoittajallemme muuta neuvoa kuin että jos hän huomaa saaneensa aivan kohtuuttomat verot, hän siitä asianmukaisella tavalla valittaisi seurakuntansa neuvostoon tai sen verotuslautakuntaan.}
\kat{Viena-Aunus}{15.09.1938}{P13}{Sen sanotaan johtuvan siitä, että häntä pidetään syyllisenä Neuvostoliiton ja Englannin välisten suhteiden kylmenemiseen. Harbinin kaupungissa asuvista Neuvostoliiton kansalaisista on 2.000 jo tullut Mandshukuon kansalaisiksi ja 1.800 on hakenut samaa kansalaisuutta. Yksistään Pitkänrannan tehtaat ovat siten menettäneet tänä vuonna 4.5 milj. mk, jonka tähden niille onkin nyt myönnetty rahtialennuksia.}
\kat{Itä-Karjala}{03.07.1939}{P6}{5000 m. K. E. N. Mattisella 15.57 ( SV UL ) 1930. Helpotuksia heimolaisille Suomen kansalaisuutta haettaessa. Siinä on varmasti aukkoja kun kaikkien tuloksia ei ole käsilläni.}
\subsubsection*{heimosoturi \normalfont\small(25 katkelmaa, 192 osumaa aineistossa)}
\kat{Toukomies}{01.10.1925}{P11}{Se on suppea, kotiopintoihin soveltuva » Oma Maa », joka antaa kaikilta puolin valaistun kuvan isänmaastamme. Ensimmäisenä niteenä on ilmestynyt Santeri Ivalon mainio historiallinen romaani Suomen heimosotien ajoilta » Juho Vesainen » 6:na painoksena. Meillä on jo kauan kaivattu kunnollista valikoimaa nuorisoromaaneja, kirjalliselta arvoltaan täysipätöisiä kirjoja, jotka vastaavat nuorekkaan mielikuvituksen vaatimuksia.}
\kat{Toukomies}{12.1926}{P18}{Voitto ei kuitenkaan jäänyt Saksalaiselle ritarikunnalle; päinvastoin se kärsi musertavan lopullisen tap- V. 1143 mainitaan karjalaisten nimi ensikertaa venäläisten kronikoissa — karjalaisten ja hämäläisten välisten heimosotien yhteydessä. Tämä viittaa paitsi sotaisiin, myös muunlaisiin kanssakäymisiin skandinaavien kera, jota lienee ollut jo Novgorodin}
\kat{Toukomies}{12.1926}{P19}{Novgorodin nuoren ruhtinaan Aleksanteri Newskin ( Jaroslavinpojan ) johtamat voimat joka tapauksessa perinpohjin lyövät Tuomaspiispan joukot heinäk. Ruotsalainen ja vähitellen aina Roomasta asti tuleva vaikutus tulee hämäläisiä painostamaan ja siitä saavat heimosodat vauhtia. V:na 1227 kastaa Novgorod jo karjalaisia ja kun v. 1240 Tuomas piispa tekee suuren ristiretken Nevajoen suulle novgorodilaisia}
\kat{Suomen Heimo}{28.07.1932}{P3}{Silloin on itä astunut askeleen meitä lähemmäksi, astunut siihen rajaan asti, jonka molemmin puolin vielä nyt asuu samaa heimoa, mutta joka silloin on jo kansallisuusraja. Kymmenen vuotta sitten käytyjen heimosotien tantereilla tehdään nyt loppuselvityksiä l vastaisten " heimosotien mahdollisuuskin tahdotaan poistaa hävittämällä niiden vankin pohja ja aiheuttaja, kansallinen vähemmistö. Toivottavasti hallitus mitä pikimmin puuttuu asiaan, varsinkin kun karkoitettujen joukossa on Suomen kansalaisia, joitten suhteen meidän on täytettävä tarkoin määrätyt velvollisuudet ja joitten oikeuksia meidän on ehdottomasti puolustettava.}
\kat{Suomen Heimo}{28.07.1932}{P3}{Silloin on itä astunut askeleen meitä lähemmäksi, astunut siihen rajaan asti, jonka molemmin puolin vielä nyt asuu samaa heimoa, mutta joka silloin on jo kansallisuusraja. Kymmenen vuotta sitten käytyjen heimosotien tantereilla tehdään nyt loppuselvityksiä l vastaisten " heimosotien mahdollisuuskin tahdotaan poistaa hävittämällä niiden vankin pohja ja aiheuttaja, kansallinen vähemmistö. Toivottavasti hallitus mitä pikimmin puuttuu asiaan, varsinkin kun karkoitettujen joukossa on Suomen kansalaisia, joitten suhteen meidän on täytettävä tarkoin määrätyt velvollisuudet ja joitten oikeuksia meidän on ehdottomasti puolustettava.}
\kat{Suomen Heimo}{14.10.1933}{P32}{AKS:ssä on alusta alkaen vallinnut lämmin myötätunto vähäosaisia kohtaan.,. Koska, kuitenkin seurassa toimivat jäsenet eivät opiskelijoina ole tulleet paljoakaan tekemisiin yhteiskunnallisten epäkohtien kanssa, on toiminta ennen Jos AKS:n " alkuperäisenä ohjelmana " olisi ollut vastustajien parjaus, niin seura ei olisi koskaan noussut muutamien heimosoturien veljespiiristä mahtavaksi järjestöksi, joka vuosikymmenen kuluessa on näyttänyt suuntaa suomalaiselle ylioppilasnuorisolle. Muistettakoon esim., miten nämä eheyttäjät keväällä 1932 ja tänä syksynä aiheuttivat riitavaalit Ylioppilaskunnan puheenjohtajistoa valittaessa, vaikka kummallakin kerralla oli uudelleen puheenjohtajaksi lupautunut dots}
\kat{Toukomies}{01.02.1933}{P11}{Nimitykset » Bjarmalanti », » Perma » kansallisina yhteisökäsitteinä sekä Karjalan ranta, Taipaleentaus y. m., kuin myös lukuisat muut seikat ja käsitteet saavat näistä Tallgrenin Karjalan rautakautta ja » Bjarmalantia » koskevista kirjoituksista valaistuksensa. Paavi näet v. 1256 karjalais-suomalaisten heimosotien aikana saarnautti kirkoissaan ristiretkeä karjalaisia vastaan, valittaen, että » Kristuksen viholliset, joita tavallisesti Carialeiksi sanotaan, ynnä muut pakanat, olivat tehneet julman hävitysretken Suomeen.... » Viimeksimainittua esittävässä kirjoituksessa selvitellään ensin karjalaisten itsenäisyysajat ennen Novgorodin vallan vakiintumista, sittemmin » varastussotien », dots}
\kat{Karjalan Vapaus}{08.1934}{P22}{Pankinjohtaja Pekka Vartiainen soitti kanteleella. Aamulla pidettiin Heimosoturien liiton juhlakokous. Pastori Elias Simojoki piti puheen vainajille.}
\kat{Karjalan Vapaus}{08.1934}{P22}{Ohjelma näissä juhlissa oli erittäin monipuolinen ja arvokas. Heimosoturien liiton seppeleen laskivat everstiluutnantti Kuussaari ja kapteeni Aarnio. Heimosoturien liiton tehtävänä on näitten taistelumuistojen vaaliminen.}
\kat{Karjalan Vapaus}{08.1934}{P22}{Heimosoturien liiton seppeleen laskivat everstiluutnantti Kuussaari ja kapteeni Aarnio. Heimosoturien liiton tehtävänä on näitten taistelumuistojen vaaliminen. Lopuksi Sortavalan sk:n soittokunta soitti Karjalan marssin, johon yleisö yhtyi.}
\kat{Karjalan Vapaus}{08.1934}{P22}{Vielä laskettiin seppeleitä heimojärjestöjen keskusvaliokunnan, Akateemisen Karjala-Seuran ja Naisylioppilaiden Karjala-Seuran puolesta. Illalla pidettiin niinikään Heimosoturien liiton illalliset seurahuoneella, jossa pidettiin useita puheita. Herra Lasse Kuntijärvi piti puheen, Karjalakerhojen kuorot lauloivat ja vapaan sanan aikana pidettiin useita puheita.}
\kat{Karjalan Vapaus}{08.1934}{P22}{Väliajan jälkeen Kemin seudun Karjalakerhojen naisjoukkue esitti kaksi laululeikkiä » Me Karjalan lapsia » ja » Kalastajatytön ». Heimoveljiemme sorto rajan takana on herättänyt heimotunnetta ja pitää vireillä heimosotien aseveljissä todellista heimorakkautta. Perjantaina klo 14 lähti Sortavalan kauppatorilla heimolippujoukkue marssimaan Kauniiseen Vakkosalmen puistoon, jossa klo 14,30 alkoi avajaisjuhla.}
\kat{Karjalan Vapaus}{08.1934}{P22}{Juhlayleisön asetuttua paikoilleen ja kaupungin suojeluskunnan torvisoittokunnan esitettyä » Reppurin laulun » ja » Karjalan marssin », toimitti ms. Jumalanpalveluksen jälkeen kokoontuivat heimosoturit ja juhlayleisö luterilaisen kirkon luo, josta rivistöinä torvisoittokunta ja liput etunenässä lähdettiin sankaripatsaalle. Ne 16 erillistä sotaretkeä, jotka on rajojemme taakse tehty, eivät ole antaneet suuria tuloksia, mutta kuitenkin voidaan positiivisia saavutuksia havaita.}
\kat{Karjalan Vapaus}{08.1934}{P22}{Aamupäivällä oli myöskin Karjalakerhojen Keskusliiton kokous, jossa puheenjohtajina toimivat P. Rahikainen ja pastori Lauri Homanen sekä sihteereinä I. levala ja V. Paavilainen. Maisteri P. Supinen piti paikkakuntalaisten puolesta tervehdyspuheen toivoen, että Heimosoturien liitto muodostuisi vahvaksi ja että se vielä kerran voisi olla toteuttamassa heimoaaletta. Kokouksen avasi liiton puheenjohtaja, everstiluutnantti Kuussaari kosketellen puheessaan tapahtumia 15 vuotta sitten ja huomauttaen ajasta, jolloin oltiin Sortavalassa koolla Aunuksen sotapoluilla.}
\kat{Karjalan Vapaus}{08.1934}{P22}{Puhuja. toi esille niitä mielialoja, jotka Aunuksen soturien keskuudessa vaiherikkaana aikana liikkuivat ja toivoi, etteivät kaikki ne uhraukset ja kärsimykset, jotka silloin koettiin, olisi hukkaan menneitä. Näiden juhlien toimeenpanosta huolehti toimikunta, jossa olivat edustettuina seuraavat järjestöt: Akateeminen Karjala-Seura, Naisylioppilaiden Karjala- Seura, Heimojärjeslöjen Keskusvaliokunta, Karjalan Kansalaisliiton Itä-Karjalan komitea, Heimosoturien Liitto. Tämän jälkeen kapteeni Aarnio selosti liiton toimintaa ja nitä syitä, jotka ovat pakottaneet tällaisen liiton perustamiseen mainiten, ettei heimosotien invaliideista mikään järjestö tai valtio pidä huolta, minkä vuoksi dots}
\kat{Karjalan Vapaus}{08.1934}{P22}{Näiden juhlien toimeenpanosta huolehti toimikunta, jossa olivat edustettuina seuraavat järjestöt: Akateeminen Karjala-Seura, Naisylioppilaiden Karjala- Seura, Heimojärjeslöjen Keskusvaliokunta, Karjalan Kansalaisliiton Itä-Karjalan komitea, Heimosoturien Liitto. Tämän jälkeen kapteeni Aarnio selosti liiton toimintaa ja nitä syitä, jotka ovat pakottaneet tällaisen liiton perustamiseen mainiten, ettei heimosotien invaliideista mikään järjestö tai valtio pidä huolta, minkä vuoksi heimosoturien itsensä on ryhdyttävä asiaa hoitamaan. Tämän jälkeen Vienan sissijoukkojen mainehikas päällikkö, J. Takkinen piti puheen, jossa hän muisteli niitä mielialoja, jotka rajantakaisessa kansassa heimosotien dots}
\kat{Karjalan Vapaus}{08.1934}{P22}{Näiden juhlien toimeenpanosta huolehti toimikunta, jossa olivat edustettuina seuraavat järjestöt: Akateeminen Karjala-Seura, Naisylioppilaiden Karjala- Seura, Heimojärjeslöjen Keskusvaliokunta, Karjalan Kansalaisliiton Itä-Karjalan komitea, Heimosoturien Liitto. Tämän jälkeen kapteeni Aarnio selosti liiton toimintaa ja nitä syitä, jotka ovat pakottaneet tällaisen liiton perustamiseen mainiten, ettei heimosotien invaliideista mikään järjestö tai valtio pidä huolta, minkä vuoksi heimosoturien itsensä on ryhdyttävä asiaa hoitamaan. Tämän jälkeen Vienan sissijoukkojen mainehikas päällikkö, J. Takkinen piti puheen, jossa hän muisteli niitä mielialoja, jotka rajantakaisessa kansassa heimosotien dots}
\kat{Karjalan Vapaus}{08.1934}{P22}{Tämän jälkeen kapteeni Aarnio selosti liiton toimintaa ja nitä syitä, jotka ovat pakottaneet tällaisen liiton perustamiseen mainiten, ettei heimosotien invaliideista mikään järjestö tai valtio pidä huolta, minkä vuoksi heimosoturien itsensä on ryhdyttävä asiaa hoitamaan. Tämän jälkeen Vienan sissijoukkojen mainehikas päällikkö, J. Takkinen piti puheen, jossa hän muisteli niitä mielialoja, jotka rajantakaisessa kansassa heimosotien aikana vallitsivat sekä toivoi lopuksi, että Suomen kansa kohtelisi heimopakolaisia niinkuin oman kansan jäseniä. Jouhiorkesterin esitettyä tuomari Härkösen johdolla » Sikermän suomalaisia kansanlauluja » piti maisteri Anna-Liisa Heikinheimo puheen, jossa hän dots}
\kat{Karjalan Vapaus}{08.1934}{P23}{Asiakkaillemme, avustajillemme ja asiamiehillemme täten tiedoltamme, että kaikki rahalähetykset ja laskut on lähetettävä suoraan lehtemme taloudenhoitajalle, kanslisti Pekka Harjulle, os. Kyminlinna. Juhlien päätyttyä tehtiin retki Salmiin, jossa heimosoturit laskivat seppeleen Aunuksen vapaussankarien patlurit ja muiden järjestöjen edustajat laskivat seppeleitä Aunuksen vapaussankarien patsaalle. Luettiin sähkösanomat, joita olivat lähettäneet kenraali Sihvo, jääkärimajuri von Hertzen, eversti Talvela, Karjalan ulkomaisen Valtuuskunnan puheenjohtaja Keynäs, arkkipiispa Herman sekä soturijärjestöt, monet yhdistykset ja yksityiset.}
\kat{Suomen Heimo}{03.03.1934}{P7}{SUOMALAISEN KIRJALLISUUDEN SEURA Suomen Heimosoturien Liitto. Veljesseuraansa tervehtii NYKS.}
\kat{Suomen Heimo}{06.04.1934}{P4}{Pakosta he sen tekisivät eikä vakaumuksesta, ja sentähden heidän tekonsa jäävät puolinaisiksi. Puhujina oli- * * \textasciicircum{} ° H täynnä Juhla \textasciitilde{} vat heimosoturit past. yleisöä. Varastomme ka- i sittää pääasiallisesti maailmankuu- luja Starmount ja Starling paito- ja, N:o 35 alh.}
\kat{Suomen Heimo}{04.06.1934}{P23}{Tänä Aunuksen vapaustaistelun merkkivuonna karjalaiset, Karjalan auttajat ja ystävät kokoontuvat Karjalan heimojuhlille näyttämään ja näkemään, etteivä 15 vuoden takaiset uhrit ole olleet turhia, vaan että näistä on ollut siunauksellisia tuloksia, jotka lupaavat onnellisempaa tulevaisuutta koko Suomen suvulle. AKATEEMINEN KARJALA-SEURA KARJALAISTEN KANSALLISLIITTO HEIMOJÄRJESTÖJEN KESKUSVALIOKUNTA KARJALAKERHOJEN KESKUSLIITTO KARJALAN KANSALAISLIITON ITÄ- NAISYLIOPPILAIDEN KARJALA-SEURA KARJALAN KOMITEA SUOMEN HEIMOSOTURIEN LIITTO y " " l|| | jf * * Näin saattaa sanoa maamme vanhimman verka- V:': l i: ->\textasciicircum{}\textasciicircum{}' 'i;: 'l;:\textasciicircum{}; \textasciicircum{}V: -:\textasciicircum{};:V<''-; -; \textasciicircum{} V * |f| m ||ip| * | tehtaan laatuvalmisteista, dots}
\kat{Suomen Heimo}{20.07.1934}{P16}{Hämeenlinnan Paperiaitta O. Y. Hämeenlinna. Suotuen Heimosoturien Liitto. Paluumatkalla käydään Valamossa.}
\kat{Suomen Heimo}{20.07.1934}{P16}{Torvisoittoa Sortavalan sk:n soittokunta Puhe Maisteri Eino Parikka Lausuntaa Metsänvartija 'Santeri Kontu Kuorolaulua Sortazulan Mieskuoro Klo 8,00 Karjalaisten Kansallisliiton vuosikokous. „ 8,00 Karjalakerhojen Keskusliiton vuosikokous. „ 8,00 Suomen Heimosoturien Liiton liittokokous. Näissä yhteistyön ja eheyttämisen merkeissä kutsumme kaikkia Suomen heimojen jäseniä Karjalan heimojuhlille ja Aunuksen vapaustaistelun 15-vuotismuistojuhlaan Sortavalaan elokuun 3 — 5 päiviksi 1934.}
\kat{Suomen Heimo}{20.07.1934}{P17}{Puhe Pastori Elias Simojoki Suomen Heimosoturien Liitto. on kärsimyskiehk v ras I}
\subsection{Kieli}
\subsubsection*{karjalan kieli \normalfont\small(25 katkelmaa, 225 osumaa aineistossa)}
\kat{Toukomies}{05.02.1925}{P8}{Matteuksen evankeliumi Tverin Karjalan murteella. V. iByo ilmestyi kolmas karjalankielinen kirja. Ylläolevaan tapaan maisteri Ahtia.}
\kat{Toukomies}{05.02.1925}{P8}{Sankt Peterburg pri Svjätejshem Praviteljsvujushtshem Sinode goda 1804. Siina luemme venäjän kirjaimin karjalaksi ( tverinkarjalan mukaisella kielellä ): Herran mijän Shjundju-Ruohtinan Svjätoj Jovangeli Matveista, Karjalan kielellä, petshatoidu...}
\kat{Toukomies}{05.02.1925}{P8}{Herran mijän Shjundju-Ruohtinan Svjätoj Jovangeli Matveista, Karjalan kielellä, petshatoidu... Oliko näiden toimenpiteiden ansiota tai ei, mutta v. IQ2O ilmestyi toinen karjalankielinen kirjanen, nim. — Tämän kirjeen johdosta ruhtinas Golitsyn kirjotti kirjeen Novgorodin metropoliitalle Ambrosiukselle ( 14 p. helmik.}
\kat{Toukomies}{05.02.1925}{P8}{Taatto meiän, sinä olet taivahal, i pyhitäh nimi sinun, i tulov tsarstva sinun, i lienöv valdu sinun, kui taivahal i maal, leiby meiän heitelematoi anna meile-nygöi, i jätä meile velgat meiän, kui i myö heitämmyö velguniekoin meiän, i elä vie meidä pahah, i päästä meidy ounahas. Tätä kirjaa tarkotetaan myös eräässä kirjeessä, jossa Aunuksen koulujen esimies v. 1868 käskee sitä lähettää 200 kpl. karjalaisille ja lausuu sen ohella: » Keskustellessaan hänen ylhäisyytensä kansanvalistusministeri Dmitrij Andreijevitsh Tolstoin kera oppilaitosten tarkastuksesta, hänen jaloutensa Pietarin koulupiirin tarkastaja ruhtinas Liewen lausui, että olisi tarpeellista kääntää karjalaksi muutamia raamatun dots}
\kat{Toukomies}{01.04.1925}{P13}{Kibunoi tuas uusi tuhjo Otteita evankeliumeista karjalankielellä vain syväinpohjaspa vain hengävys tuas ylenöö.}
\kat{Toukomies}{01.04.1925}{P13}{2. Kuin on kirjutettu Jesai-prorokan ( profietan ) kirjas: » Katsho, minä työnnän miun anhelit siun silmies icl, ja hain valmistaa siun ties ( dorogas ) »; Koska karjalankielinen kirjallisuus pääasiassa sisältää evankeliumikirjain, rukousten ja muiden uskonnollisten lukemisten käännöksiä, julkaisemme silloin tällöin otteita ensinmainituista näytteiksi nykyisemmältä tavalla suoritetuista käännösyrityksistä. Kergeinä kiitelgäh suvenmaid' igikiildäjie toizet, kaunehis kallehin mua on minul syndymämua!}
\kat{Toukomies}{01.04.1925}{P16}{45. Vain mendyö pois mies rubei leviel kuulovoittamah ja azhies tieduo levittämäh, jotta lisus ei eniä suattannuh julgi mennä linnoih, vain oletteli niijen ulgopuolel tyhjis paikois; ja joga paikas tuldih hänen luo. 1. H.; Otteita evankeliuuneista karjalankielellä; Kibunoi, suomalaisia runoja aunukselaisille, sivut 9-16. — Kuvia: Kalevalan heimoa kansallisilla kesäjuhlillaan; Kaksi karjalaisia esittävää kuvaa Olaus Magnin teoksesta v. 1567; Myöhemmän ajan karjalaisia kauppamatkoilla; Kauppias Ivan Mitropanof f: Wanhalan emäntä ja tytär ja Wanhalan talo Wuokkiniemellä. — Karjalan Sivistysseuran apurahat Kontro \& Kuosmasen lahjoitusrahaston viimevuotisten korkojen Suomen puolelle jaettavista dots}
\kat{Toukomies}{01.09.1925}{P6}{Seuraava venäläisten toiminnan aika tässä suhteessa alkaa v. 1894; mutta sitä ennen jo professori E. N. Setälä laajassa osaksi karjalankielinen nimikin: » Ljuhjukjaini ( lyhykäinen ) svjäzhennoi jstorija ( pyhä historia ); kratkaja svjäzhennaja istorija ». Seuraavana vuonna, 1895, jo julkaistaan samalla Vienankarjalan murteella ja heidän käytäntöään varten laajempi, 57-sivuinen biblian historia, jolla jo on}
\kat{Toukomies}{01.09.1925}{P8}{Kun Aunuksessa tähän asti on kirjakielenä ollut venäjä ja sen avulla jo yksinkertaisemmankin kansanmiehen näkee yksinpä matkoillaankin — kuten olen havainnut — ahmivan Gogolia, Pushkinia y.m. venäläisiä kirjailijoita, voi heidän yksinkertaisemmissaan herätä kenties sekin ajatus, että kauniita kirjaanpantavia asioita ei olekaan muilla kuin venäläisillä. ■ Mitähän sanoisivatkaan aunukselaiset, jos. he Gogolin » Reviisorin » jälkeen kohta saisivat lukea esim. Aleksis Kiven » Nummisuutareita »? Karjalankielisten kirjoitelmien, runokäännösten ja muiden pitempienkin kirjoitustuotteiden saattamisella tähän murrekieleen tai useammasta murteesta kehitettyyn uuteen » kirjakieleen » saatettiin katsoa dots}
\kat{Toukomies}{01.09.1925}{P16}{Midä sanot itshestäs? » Otteita evankeliumeista karjalankielellä Hai vastai: » En. »}
\kat{Toukomies}{01.09.1925}{P25}{Mukana ollut; Karjalan kuohukolla. puhe. pit. I. H.; Otteita evankeliumista karjalankielellä, käänt. I. H.; Karjalan kirkon toinen lakimääräinen kirkolliskokous, kirj.}
\kat{Toukomies}{12.1925}{P26}{Rajaseutukysymys ja heimopakolaisten hoito I, kirj. livo Härkönen. Otteita evankeliumista karjalankielellä I: 1 luku Markuksen evankeliumista. Karjalaisia I: Pelgui, kirj. livo Härkönen ( 1 kuva ).}
\kat{Toukomies}{12.1925}{P26}{Karjalaisia I: Pelgui, kirj. livo Härkönen ( 1 kuva ). Otteita evankeliumeista karjalankielellä II: 1 luku Johanneksen evankeliumista. Karjalankielisestä kirjallisuudesta 111, kirj. livo Härkönen ( 2 kuvaa )}
\kat{Toukomies}{01.1926}{P9}{Relanderia ja kirjailija Artturi Leinosta; juhlarunoilijoiksi Turun yliopiston rehtoria V. A. Koskennientä ja kirjailija livo Härköstä sekä juhlasaarnan pitäjiksi piispa Erkki Kailaa ja arkkipiispa Hermania. Historiallisten olojensa vaikutuksesta tunnustetaan siinä kahta eri uskontoa, kansantavat vaihtelevat ja kielimurteita puhutaan siksi eroavaisia, että kielimiehet ovat puhuneet eri kielistäkin ( esim. varsinainen karjalankieli ). katselkaammepa mitä alueen osaa tahansa: Reima, merenhenkinen virolahtelainen eroaa jo melkein täydelleen karjalankannakselaisesta ajuritoimiin taipuvaisesta maamiehestä ja niin myös Muolaanpuolen muinaisin rekikauppaa harjoittaneesta talonpojasta.}
\kat{Toukomies}{01.1926}{P18}{Sillä jokaisella kansalla oma kieli ei ole ainoastaan sen henkisen omituisuuden uskollisin ja tarkin ilmaus, vaan myös sen välttämätön suojelija ( missä eivät rotu tai erikoislaatuinen maantieteellinen asema eroita sitä naapurikansoistaan ). Samoin jos karjalaisuuden on mieli säilyä henkisenä, vaikuttavana tekijänä ja omalaatuisuutena, niin se voi tapahtua vain siten, että karjalankieli saapi arvonsa ja tulee välineeksi Kalevalan kansan omintakeisuuden ylläpitämisessä ja tehostamisessa. Jääpi lopuksi selvitettäväksi viimeinen epäilys viitotun tehtävän suhteen: ansaitseeko karjalaisuus yleensä säilymistä, kannattaako ponnistella ja vaivaa nähdä tämän tarkoituksen hyväksi, — eikö se liene dots}
\kat{Toukomies}{03.1926}{P12}{Valistus- ja kansallistuttamistyöstä puhuja lausui: » Kansallisuuskysymyksen alalfa on jatkettu toimintaa tsaarin politiikan jätteiden tuhoamiseksi. Gavrilov Aunuksesta puhui karjalaksi huomauttaen m. m. että talonpoikain tarvitsemia tuotteita olisi enemmän tuotava Suomesta. Lainsäädännölliset ja hallinnolliset asiakirjat samoin kuin federatiivisen ja liittovallan Karjalaa koskevat määräykset ovat julaistavat molemmilla kielillä.}
\kat{Toukomies}{03.1926}{P13}{Karjalan tasavallan alueellakin käytetään hyvin monia murteita, jotka suuresti eroavat toisistaan. Mitä karjalankieleen yleensä tulee, lausui siitä selostaja: » Karjalankieli sinänsä on murretta. Päätöksisssään kokous yleensä tunnusti ja hyväksyi hallituksen toiminnan ja toimenpiteet.}
\kat{Toukomies}{03.1926}{P13}{Ja lapsien joukos siivil kaibavon hain ettshikkiä, ken muuttumatoin on, ja hänel tshirkutelgua erikseh Venäjän- ja suomenkielinen opetus on verrattain helposti järjestettävissä, kirjallisuutta kun on, mutta karjalankielinen on vaikeata. Syndy Orihmattilas Uv velmu ai tammikuun 16 p:nä 1849, kävi Jyväskylän kanzakouhmopettajaseminaarin ja toimi koulunopettajanna eri paikois.}
\kat{Toukomies}{03.1926}{P13}{Valistustyön tilasta esitti selostuksen Gabbojev ( Karjalan valistusasiain komisaari ) ja lausui siinä m. m.: » Jyrkkä eroavaisuus eri seutujen kielenkäytössä pakottaa opetuskielenä käyttämään eri alueilla Tämän takia voidaankin karjalaisten opetus ratkaista ainoastaan siten, että opetus suoritetaan karjalankielellä, mutta kielen ja kirjoituksen opetuksessa käytetään suomenkieltä, joka on läheisin ja jonka karjalaiset helpoimmin omaksuvat. Karjalan kansantaloudesta, Karjalan asevelvollisuuskutsunnasta omiin joukko-osastoihin, Karjalan maataloudesta, Karjalan valistustyön tilasta ja tehtävistä, terveydenhoidon tilasta maaseudulla ja vastaisista tehtävistä tällä alalla, Karjalan valtion dots}
\kat{Toukomies}{06.1926}{P13}{Kunpahan vielä ilmestyisi karjalaistenkin omasta keskuudesta oma » Nortamonsa », joka luonnehtisi karjalaisten omalla murteella heimonsa luonnetta ja ajatustapoja. Se ei tosin ole aivan uusi kysymys, onhan meillä eräältä taholta hankkeiltu jo karjalankielistä sanakirjaakin, siis tahdottu tehdä erillinen karjalankieli kiistämättömäksi tosiasiaksi. Tällainen toiminta tietysti pohjautuu siihen olettamukseen, että tästä kielestä on luotava itsenäinen sivistyskieli, kieli, josta tulisi itsenäisen Karjalan valtion kieli omine kirjallisuuksineen.}
\kat{Toukomies}{06.1926}{P15}{evankeliumikappaleet, jotka n. s. » Karelskoje Bratsvo » ( Karjalainen veljeskunta, vaikka karjalaisia tässä veljeskunnassa tuskin lienee ollut nimeksikään ) oli enemmän venäläistyttämis- kuin karjalaistuttamistarkoituksessa julkaissut. Ja jos missä, niin juuri tässä aitokarjalaisessa pikkukaupungissa pysyi karjalankieli sitkeästi elossa, sitä kuuli puhuttavan ei vain torilla ja kaupoksa, mutta myöskin kaikissa virastoissa ja yhteiskunnallisissa laitoksissa. Panfinskaja ( suursuomalainen ) propaganda » oli ennen sellainen hirviö, joka pani kaikkien virkamiesten — kuvernööristä uratniekkoihin asti — ja myös pappien sydämet pahasti lyömään ja tätä » panslavismia » uhkaavaa vaaraa oli dots}
\kat{Toukomies}{06.1926}{P16}{Nyt on siellä toisin toimet, vieras on siellä viestintuoja, outo on opastajana, häneltä kysyttiin, onko hänellä mielitietty, jota myöskin karjalankielellä sanotaan toisinaan dushaksi. Antoi Ahtola osansa, talvella Tapio antoi, ehätti ruskean repoisen, miehen etsivän etehen.}
\kat{Toukomies}{10.1926}{P10}{Muhamettilaisilla mullilla ja roomalaiskatolisilla papeilla tosin on hiukan toisenlainen puku, pitkä maahan asti yltävä musta viitta, mutta paljoa käytännöllisempi ja eurooppalaisempi lienee sekin kuin vanha venäläisen papin viitta. Mutta irvikuva se on sellaisesta jo siksikin, että se on napitettu ja että se halaattimaisuudestaan huolimatta pyrkii kuitenkin jäljittelemään ihmisruumiin muotoja ja siten vihdoin lähentelee venäläisaasialaista talonpoikaisnuttua ( karjalaksi podjeuhka ). Puvuston välikysymyksessä on luonnollisesti tultava kansallisiin tai yleiseurooppalaisiin väreihin ja päähineistä poistettava tuo tyypillinen venäläis-itämainen ylöspäin laajeneva muoto, jonka sijaan lienee dots}
\kat{Toukomies}{12.1926}{P8}{kaskenpolttoa, kalastusta ja pyytää metsänriistaa. So i, äidinkieli, soi! On ilmeistä että kumpasillakin on yhteinen alkuperä}
\kat{Toukomies}{12.1926}{P8}{Hainarin vSuomen Keskiajan historiassa mainitaan Suomen Itä- Karjalasta Käkisalmen linnaläänissä Kurkijoelta pohjoiseen olevissa pogostoissa olevan myös perevaaroja eli piirikuntia, joista jokainen käsitti useampia kyliä ( s. 515 ). Soi, kallis äidinkieli, soi soi esi-isäin virsi, am, Suomen kannel, huminoi, pois sula, syvä kirsil Se kirsi, joka kahlita sun tahtoi korven yöhön, kuin valtavoimaa vangita vois muka orjan työhön. Soi, Suomen kannel, helise, sun kaikus kauas kantaa, se kautta maiden, meritse käy maailman jo rantaa, se astunut on arvossaan rinnalle vapahitten, se tuntee jo nyt tunnossaan: oon juurta sankaritten.}
\subsection{Talous}
\subsubsection*{puutavara \normalfont\small(25 katkelmaa, 181 osumaa aineistossa)}
\kat{Toukomies}{01.09.1925}{P17}{41. Ondrei, Simon Pedrin velli, oli töine niis kahtes, kudamat oldih kuuldu midä livan sanoi, da käveldy Isusan jälles. Tämän valtion sahan johtaja on monta kertaa myöntänyt karjalaisten ansion työssä, mutta näyttää hänellekin olevan vaikeata saada esim. kerhomme tarkoituksiin jotain pientä huonetta, jota olemme pyytäneet. 52. Da Hai sanoi hänele: » Tovesti, tovesti sanon teile: 'täs ielleh suatto nähtä taivahan olevan augi, da Jumalan anheliloin nouzevan yläh da laskeivan alah Inehmizen Poijan piäle.}
\kat{Toukomies}{08.1926}{P18}{Huolimatta matkain pituudesta oli suurin osa jäseniä ansiopaikoiltaan tilaisuuteen saapunut. Kokouksessa päätettiin pitää samalla sahan vanhalla pirtillä illanvietto saman kuun 30 p:nä. Kerhomme taloudellinen asema ei anna myöten maksaa pienintäkään palkkaa yhdellekään toimihenkilölle.}
\kat{Toukomies}{01.02.1927}{P23}{Tässä onkin eräs sen ansioita. 9 p:ksi sahan vanhalle pirtille. N:ot 1 — 2 -}
\kat{Toukomies}{01.05.1927}{P13}{Runsaasta ohjelmasta, joka alkoi pianonsoitolla ja kerhon toimihenkilön Mikko Saariston tervehdyspuheella, on mainittava neiti Klaudia Karvosen karjalanmurteiset lausuntanumerot, tervehdysten lukeminen, nti Viitasen yksinlaulut, lehtori Tikanojan juhlapuhe ja kerholaisten 'esittämä nelikuvaelmainen esitys " Karjalan häät ". 20 p:nä Toppilan sahan vanhalla pirtillä ohjelmallisen illanvieton niinikään runsaalla ohjelmalla ja äskettäin pääsiäisen aikana on se myös pitänyt keskuudessaan illanvieton. vankaan.}
\kat{Toukomies}{01.09.1927}{P2}{1930 — 31: Povenetsiin rakennetaan puuseppätehdas, kust. 1929 — 31: Kemiin rakennetaan 6-raaminen saha, kust. 1931 — 33: Kemiin rakennetaan 4-raaminen saha, kust.}
\kat{Toukomies}{01.09.1927}{P2}{1929 — 31: Kemiin rakennetaan 6-raaminen saha, kust. 1931 — 33: Kemiin rakennetaan 4-raaminen saha, kust. Metsän kulutus vuodeksi 1927 — 28 lasketaan 650,000 kuutiosyleksi.}
\kat{Toukomies}{01.09.1927}{P2}{Neuvosto- Venäjällä on ennenkin laadittu suurisuuntaisia rakennusohjelmia, samaten kuin tsaari- Venäjällä. 1 50,000 r. ( n. 2,85 milj. Smk. ) XTnitsan 2-raaminen saha suljetaan Tilaushinnat: 30 ' mk. vuosik. muilta paitsi karjalaispakolaisilta 2CTmk.}
\kat{Toukomies}{01.09.1927}{P2}{Metsätöissä oleville maksettiin viime vuonna palkkaa 8 milj. ruplaa ( noin 160 milj. Smk. ) Mainitun lehden ilmoituksen mukaan oli puuteollisvmden nettovoitto (? ) 20 milj. ruplaa ( 40 milj. Smk. ) 1927 — 28: rakennetaan Petroskoihin tärpättitehdas, kustannukset 30,000 ruplaa ( n. 570,000 Smk. ), hiomokiviverstas Ylä-Rutsheviin, kustannukset 50,000 r. ( n. 900,000 Smk. ) Iljmskin 3-raaminen saha suljetaan. 1928 — 29: Kontupohjaan rakennetaan 4-raaminen saha, kustannukset 800,000 r. ( n. 15 milj. Smk. ), Kontupohjaan niinikään sulfaatti-selluloosatehdas, kustannukset 1,5 milj. t. ( n. 28,5 milj. Smk. ), Petroskoihin rakennetaan hartsitehdas, kustannukset 30,000 r. ( n. 570,000 Smk. ), Olenin dots}
\kat{Toukomies}{01.09.1927}{P2}{1927 — 28: rakennetaan Petroskoihin tärpättitehdas, kustannukset 30,000 ruplaa ( n. 570,000 Smk. ), hiomokiviverstas Ylä-Rutsheviin, kustannukset 50,000 r. ( n. 900,000 Smk. ) Iljmskin 3-raaminen saha suljetaan. 1928 — 29: Kontupohjaan rakennetaan 4-raaminen saha, kustannukset 800,000 r. ( n. 15 milj. Smk. ), Kontupohjaan niinikään sulfaatti-selluloosatehdas, kustannukset 1,5 milj. t. ( n. 28,5 milj. Smk. ), Petroskoihin rakennetaan hartsitehdas, kustannukset 30,000 r. ( n. 570,000 Smk. ), Olenin saarelle kalkkitehdas, kustannukset iqo, oqo r. ( n. 1,9 milj. Smk. ). 50p. palstamil- joulunumeronamarras-jouluk. vaihteessa Karjalan Sivistysseura, os. Helsinki lim., pysyvistä ilm. sopimuksen dots}
\kat{Suomen Heimo}{20.01.1928}{P9}{1 Vallila Puh. Vallilan Puutavara Oy. SUOMEN HEIMO}
\kat{Suomen Heimo}{20.01.1928}{P26}{2600 \& 12 31 unnwPK Al 1 1 ift m nni t i it Suorittaa kaikenlaatuisia asfaltti-, massaIat- nuUWCKALUJLW, KUULU-, ti sementti-, betoni- ja rautabetonitöitä. Kirjoittakaa tänään, myötallittakäl tämä ilmoitus; Oy- TwaMvaraate POHJOLA. Helsinki Suomalainen Puuliike Oy. i «. •- „ -, „ nm un m Unh.. iut Kiviniemen saha. Kirjallinen 2 vuo- KspyS\textasciicircum{}f' den takuu seuraa. ■sfo\textasciicircum{}\&AJ Naisten rannekello, kultadubleets, 5 vuoden * takuulla, sveitsiläistä tekoa, 6 kivellä ranne- KUOPIOSSA, KAJAANISSA nauba musasta silkistä, ainoastaan Smk.}
\kat{Suomen Heimo}{20.02.1928}{P2}{Aistikkaita Kulta- ja Hopea; esiMK to huokein Puuseppätehdas: hi finnitt - ' ' ninnvirl PiiK' fi7 QK KI / \textasciicircum{}allSam 1 / II • * » Maan vanhin ja suurin rakennuspuu- j • tValllOn tVCllO " " Ja seppätehdas. ' ■ ■ ' ■ \textasciicircum{} ' I; Paperi- ja Kirjoitusvälineitä IB\textasciicircum{}\textasciicircum{}\textasciitilde{} 4\textasciicircum{}\textasciicircum{}B PIIRUSTUSTARPEIDEN W / \textasciicircum{}jjff\textasciicircum{}Sjk ERIKOISOSASTO IATrBJJM JBb I y VV in pr fck\textasciicircum{}s\textasciicircum{} s\textasciicircum{}S\textasciicircum{} « pv, I n. f f ULrr f\&Qjwfij\textasciicircum{} OSAKEYHTIÖ ifcJjuUÄ\textasciicircum{}j HELSINKI j ■ j|\textasciicircum{}S2p\textasciicircum{} P. Esplanadinkatu 43 I m\textasciicircum{}M; Siilen Puutavara Oi. » <fö t\textasciicircum{}\textasciicircum{}r-W ZEHITH HELSINKI Puh. 32 544. Oma piirtämä.}
\kat{Suomen Heimo}{20.02.1928}{P23}{1 Vallila Puh. Vallilan Puutavara Oy. Pianon- korjauksia tehdään.}
\kat{Suomen Heimo}{20.02.1928}{P28}{T « Ml I V TOIVA\textasciicircum{}FN rv,.. n. Tähän kuvatun herrain rannekellon, joka • ITII J. ¥• 1VI » riJLrll PS3JP9 on oikeata kultadubleeta ja varustettu hie- --\textasciicircum{} - = -HHH-i * \$\&ipii nojnta sveitsiläistä valmistetta olevalla ank- PySZ\textasciicircum{}y\textasciicircum{} 1 kurikoneistoHa. O y- Tavaravarasto POHJOLA, Helsinki Suomalainen Puuliike Oy. i n; nn mn nniumn Uvhirilnf Kiviniemen saha LNPII P R I M ULA UMU Sortavala, puh. 140, 112. Laivurinkatu 10. Puhelin 2©3 \& 68 OS L. Heikinkatu 14. „ 5'5 24 \& 21 679 Myy kaikenlaisia SAHATUOTTEITA, sekä RIMA- ja METSÄHALKOJA. Ä.... \_. w. v Insinööritoi mis t o O. Y. HELY LA, \textasciicircum{} os. HelylÄ. ± \_ Jl-V jTjL V\textasciicircum{} Oimist. 29 04\&9Ö5 PYRAMID SUORITTAA KAIKKIA RAKENNUSALAAN KUULUVIA SUUNNITTELU- JA dots}
\kat{Suomen Heimo}{20.02.1928}{P29}{— Laivausp.: Koivisto. Rautakauppa ttoH Ahjo inkilän saha oy. Kyselkää hintojamme! pohjola ] « X8£57l Rahastot Smk.}
\kat{Suomen Heimo}{20.02.1928}{P29}{41, puh. 89 00. Palo-, meri-, auto- y.m. vakuutuksia Tapaturmavakuutuksia KULLERVO Ensiluokkaisia liinavaatteita, lääk iiri- ja Rahastot Smk. Sähkösanasto: erittäin edullisesti kaikkia rauta- Zebra Code 3:rd Edition, kauppa-alaan kuuluvia tavaroita. „....,.., Saha: Inkilän asema. Helsingin Osuuskauppa r. L 15 Siirtomaatavaramyymälää •a Hankinnan lisäys vuonna 1927 suu- T., °....... * Lmamyymälaa rempi kuin kolmella muulla maamme vanhimjmalla yhtiöllä yhteensä.}
\kat{Suomen Heimo}{15.03.1928}{P6}{Siinä olette oikeassa. Vallilan Puutavara Oy. Liiankin selvää kaikki.}
\kat{Suomen Heimo}{15.03.1928}{P22}{Suomalainen Puuliike Oy. Saha: Inkilän asema. — Laivausp.: Koivisto.}
\kat{Suomen Heimo}{15.03.1928}{P22}{T:MI J. V. TOIVASEN Osakeyhfiö Valistus « shbmhsh — = = — a — SH-SSS- — ess-H-sHH — Perustettu 1901 Kulta- Ja KcllosepänlHkkcitä Helsinki " Annankatu 29 * V MAAMME SUURIN JA VANHIN Suositellaan. Leipomo primula Raivilat Laivurinkatu 10. Puhelin 268 \& 68 08 Myy kaikenlaisia L. Heikinkatu 14. „ 66 24 \& 21 5T9 SAHATUOTTEITA, sekä RIMA. ja METSÄHALKOJA. O. Y. HELYLÄ INKILÄN SAHA OY. os. Helylä. 29 04\&966 • PYRAMID SUORITTAA KAIKKIA RAKENNUSALAAN KUULUVIA SUUNNITTELU- JA RAKENNUSTEHTÄVIÄ}
\kat{Suomen Heimo}{15.04.1928}{P9}{Eivätkä Vallilan Puutavara Oy. Nio 4 IÖ2Ö}
\kat{Suomen Heimo}{15.04.1928}{P24}{Suomalainen perheenemäntä käyttää / MP\textasciicircum{}T\textasciicircum{}U. ainoastaan kotimaisia pesuaineita! Sähkösanasto: Koetelkaa ja arvostelkaa / Zebra Code 3:rd Editio " Saha: lnkilän asema. RAKENNUSOSAKEYHTIÖ PYRAMID Viipuri, Pellervonkatu 10. Puhelimet: Sähköosoite: joht. 29 24 kontt.}
\kat{Suomen Heimo}{15.04.1928}{P24}{insinööritoimisto Osakeyhi iö Valistus A\textasciitilde{}r\textasciicircum{} T\textasciicircum{}7V T \_ T f Perustettu 1901 X ly ±jL \textasciicircum{}S Helsinki - Annankatu 29 -.. INKILAN SAHA OY. Konttori: Viipuri, Torikatu 2. Puhelimet: 30ö ja 916. I / ennätin-os.: Inkoy, Viipuri. 26 00 \&. 12 31 p '. Suorittaa kaikenlaatuisia asfaltti-, massalat- s tia-, sementti-, betoni- ja rautabetonitöitä.}
\kat{Suomen Heimo}{20.05.1928}{P4}{1 Vallila Puh. Vallilan Puutavara Oy. SUOMEN HEIMO}
\kat{Suomen Heimo}{20.05.1928}{P23}{INKILÄN SA * * A OY - lepo PRIMULA Katit; £2£ Viipuri, Torikatu 2. Puhelimet: 306 ja 916. Lennätin-os.: Inkoy, Viipuri. Sähkösanasto: Laivurinkatu 10. Puhelin 2©3 \& 08 03 Zebra Code 3:rd Edition L. Heikmkatu 14. „ 56 24 \&21 579 Saha: Inkilän asema. Puhelimet: ■ I I lAMCI / AI I ITT Konttori 25 088, Varasto 32 754, nUAJINtIXALU I Johtajalle 25373. ostatte edullisimmin meiltä niiiniiiiiiiii HELSINKI Aleksanterinkatu 17, 2. kerros.}
\kat{Suomen Heimo}{15.06.1928}{P15}{115 Vallilan Puutavara Oy. SUOMEN HEIMO}
\subsubsection*{osuuskauppa \normalfont\small(25 katkelmaa, 76 osumaa aineistossa)}
\kat{Suomen Heimo}{20.02.1928}{P29}{Sähkösanasto: erittäin edullisesti kaikkia rauta- Zebra Code 3:rd Edition, kauppa-alaan kuuluvia tavaroita. „....,.., Saha: Inkilän asema. Helsingin Osuuskauppa r. L 15 Siirtomaatavaramyymälää •a Hankinnan lisäys vuonna 1927 suu- T., °....... * Lmamyymälaa rempi kuin kolmella muulla maamme vanhimjmalla yhtiöllä yhteensä. 24,000,000 pilrustustakkeja y. m. Keskinäinen Vakuutusyhtiö Ét\textasciicircum{}\textasciicircum{}\textasciicircum{} KALEVA SUf\&ä ) Hyvä ja varma kotimainen palovakuutusyhtiö Ensimmäinen kotimainen henkivakuu- Antaa myöskin tuslaitos. auto-, murto- ja lastvakuutuksia Helsinki, Imatran talo, Kasarmink.}
\kat{Suomen Heimo}{20.12.1929}{P7}{243 OSUUSLIIKE ARINA rl. Tänään on itsenäisyysläivä.}
\kat{Suomen Heimo}{20.12.1929}{P7}{Jäsenmäärä yli 2,000. Vuosimyynti n. 15,000,000 mk. OULU Puolueeton, aina kuluttajien etuja valvova osuusliike. Siirtmnfi, a, -, ruoka- ja sekatavara-, liha- ja leipämyymälöitä.}
\kat{Toukomies}{01.02.1929}{P18}{Kalastus, joka on tullut käytäntöön verrattain myöhään, on nykyään melkein kokonaan lamassa eikä sen uudelleen käytäntöön tulosta ole suuria toiveita. ( Vehkaojan — Vehrutshei ) osuuskauppa on ryhtynyt myymään tahkoja ja kovasinkiviä niin lähiseuduille kuin kaukaisempiinkin seutuihin, jotka tarvitsevat tahkokiveä. Harvoissa kylissä enää harjoitetaan kalastusta ja se vähä, mitä saadaan, menee melkein kokonaan omaan talouteen; hyvin harvoilta joutaa kalaa myötäväksi.}
\kat{Toukomies}{01.12.1929}{P23}{Mutta kuitenkin se ansaitsee huomiotamme yhtenä ilmauksena siitä mielenkiinnosta, mitä venäläiset tutkijat viime aikoina ovat osoittaneet nimenomaan suomalaisia kansoja kohtaan, joiden pariin meikäläisillä tiedemiehillä pitkiin aikoihin ei ole ollut pääsyä. Selitetään petun valmistusta ja kerrotaan, että petun käyttö, joka vallankumouksen jälkeisinä nälkävuosina oli hyvin yleistä, nyttemmin on supistunut melkein olemattomiin, kun asukkailla on tilaisuus hankkia riittämiin tuloja metsä- ja uittotöillä ja osuuskauppa hankkii viljaa paikkakunnalle. Todetaan, että lähellä rajaa sijaitsevissa kylissä leivotaan kesällä, kun on mentävä ulkotöihin kauas}
\kat{Suomen Heimo}{15.01.1930}{P23}{Miesten ja naisten kasarmit ovat erikseen. OSUUSTOIMINTA NEUVOSTO-KARJALASSA. „Pistäysin tässä kerran Kierretin yhteisasunnoissa.}
\kat{Suomen Heimo}{15.01.1930}{P24}{Ja kirjoittaja ihmettelee, että „mistä, mistä tämä johtuu? " Kylässä toimii Kiimasjärven osuuskaupan haaramyymälä, palvellen koko aluetta. Kiimasjärven osuuskaupan toisesta haaramyymälästä valitetaan näin:}
\kat{Suomen Heimo}{15.01.1930}{P24}{Kylässä toimii Kiimasjärven osuuskaupan haaramyymälä, palvellen koko aluetta. Kiimasjärven osuuskaupan toisesta haaramyymälästä valitetaan näin: Niitä annettiin jäsenelle 1 kg. ja jäsenen perheenjäsenille " Yz kg.}
\kat{Suomen Heimo}{15.01.1930}{P24}{Kaupassa ei ole edes paperosseja, joita aina voisi hankkia riittämään asti. „Mikkilän ( Säämäjärvellä ) osuuskaupan eräät asiat eivät ole ihan oikealla tolalla. Sitä todistaa se seikka, että kaupasta puuttuvat aina välttämättömimmätkin kulutustavarat.}
\kat{Suomen Heimo}{15.01.1930}{P24}{Seuraavana päivänä kun kuntikyläläiset tulivat hakemaan normejaan, olivat ne jo loppuneet. „Kiimasjärven osuuskaupan hallinto ei ollenkaan huolehdi Luvajärvcn haaramyymälän tavaroilla varustamisesta. Hallinto varmaan pakoittaa osuuskaupan jäseniä kesselillä hankkimaan tavaransa 25 ja 50 kilometrin päästä.}
\kat{Suomen Heimo}{15.01.1930}{P24}{„Kiimasjärven osuuskaupan hallinto ei ollenkaan huolehdi Luvajärvcn haaramyymälän tavaroilla varustamisesta. Hallinto varmaan pakoittaa osuuskaupan jäseniä kesselillä hankkimaan tavaransa 25 ja 50 kilometrin päästä. Tosiaankin! „Ettei, Ettei maksa vaivaa lähteä ( tavaraa hakemaan ) ennenkuin Ujanpäiväksi. "}
\kat{Suomen Heimo}{15.01.1930}{P24}{Sieltä puuttuu nykyään kaikki välttämättömät tarvikkeet, kuten tulitikut, sokeri y. m. Tavaroiden kuljetuksesta hallinto 25 kilometrin taipaleelta maksaa 3 kop. kilolta. Mistäpä se Kiimasjärven osuuskaupan kommunistinen hallinto voisi „ välttämättömiä tarvikkeita " tilata, kun näitä tarvikkeita ei ole koko Neuvostoliitossakaan olemassa kuin silloin tällöin. Jos olisi tavaroita: sokeria, tulitikkuja ja leipää, niin silloinhan se olisi,, porvarillista meininkiä ", jonka johdosta maailmanvallankumousihanne kadottaisi todellisen proletaariluonteensa.}
\kat{Suomen Heimo}{15.01.1930}{P24}{Jos olisi tavaroita: sokeria, tulitikkuja ja leipää, niin silloinhan se olisi,, porvarillista meininkiä ", jonka johdosta maailmanvallankumousihanne kadottaisi todellisen proletaariluonteensa. Mitä tåas siihen tulee, että,, hallinto varmaan pakoittaa osuuskaupan jäseniä kesselillä hankkimaan tavaransa 25 ja 50 kilometrin päästä ", niin se on verrattain helppo tehtävä kantamiseen tottuneille karjalaisille. Varmaankin lukemattomat „m atkalaiset, tilapäiset asukkaat ja paikalliset palveluskuntalaiset " ovat kommunismin olemassaolon aikana saaneet „kärsiä ", kärsiä " enemmän kuin kerran hivuttavaa nälkää ja sen johdosta monasti katuneetkin kohtaloansa.}
\kat{Suomen Heimo}{15.01.1930}{P24}{Ruokatarpeista mainitussa osuuskaupassa on alinomaan huutava puute — paitsi niistä tavaroista, jotka ovat normilla, kuten jauhot, sokeri, joita on aina saatavissa — vaikka tavaran saantiin kesänkin aikana ovat melkoisen hyvät mahdollisuudet, sillä eihän rautatie ole täältä kuin 30 km. päässä, jonne maantie on aivan tyydyttävässä kunnossa. " „Moni matkalainen, tilapäinen asukas, sesonkityöläinen, paikallinen palveluskuntalainen Leholla ollessaan saa kärsiä nälkää, sillä kylä on piirin keskusta, joten liikenne on kokolailla vilkasta ja kylän taloudet eivät voi tyydyttää yllä esitettyjen tarpeita, sillä ne itsekin tarvitsisivat osuuskaupan apua. „No, on täällä kooperatiivi ", arvelee moni ja dots}
\kat{Suomen Heimo}{22.10.1930}{P20}{KkCSUUSI\&\textasciicircum{}-J Musiikkikeskus Oy. yÉßt£\textasciicircum{} Täydellinen soitin- ja nuotti- kauppa sekä kustannusliike. Osuuskauppa muuten voisi varustaa itsensä tavaralla paljon paremmin, kuin mitä se on tehnyt. Tässähän se, hyvät toverit, on juuri se kohta, joka aiheuttaa maito- ja lihapulan.}
\kat{Suomen Heimo}{22.10.1930}{P20}{Talonpojat näet sanovat, että kauppaasiain kansankomisariaatti saa itse hoitaa ja kehittää karjatalouttaan, koska se kontrolloi hintojakin. Osuuskaupan hallinnon olisi tehtävä sopimus paikallisten kalastajien kanssa, että ne luovuttaisivat saaliinsa osuuskaupalle paikallisen väestön tarpeiksi. " Yhtä lohduton uutinen on tämäkin kuin edelliset. „Vitelen keskikylässäkin " näkyy olevan vaikea elä ». varsinkin maattomien työläisten ja toimitsijain.}
\kat{Toukomies}{01.12.1930}{P18}{Siellä on ilmeisesti kaksi kovaonnista kalastajaa, jotka ovat joutuneet oikullisen meren kuljeteltavaksi. Kansaa alkaa kokoontua osuuskaupan pihamaalle neuvottelemaan ja säälittelemään, mutta sinä päivänä ei ole mitään tehtävissä. Hän on erottavinaan kaukaisella jääkentällä kaksi tummaa pistettä, jotka eivät näytä kärpästä suuremmilta.}
\kat{Suomen Heimo}{16.06.1931}{P28}{Ei siis olekaan ihme, jos lapset kouluista pääsevät sangen vaillinaisilla tiedoilla, Useasti lukutaidottominakin. Niiden hankkimisen helpoittamiseksi on jokaiseen kouluun järjestetty koululaisten osuuskauppa, mutta sillä on etupäässä poliittinen tarkoitus. Vaikka opettajista onkin suuri osa kommunistisen puolueen ja nuorisoliiton jäseniä, niin ei opettajisto sentään ole oloihinsa tyytyväinen.}
\kat{Toukomies}{01.02.1931}{P8}{Luonnollisesti on tarkoitus saada ne " vapaaehtoisesti " liittymään kollektiivitalouksiin. Lisäksi on pakko liittyä osuuskaupan jäseneksi ja suorittaa vuosittain jäsenmaksuna 30 ruplaa. Vihollinen perääntyy takaisin, taistelu on päättynyt — heimoveljiemme avustamana ovat karjalaiset saaneet voiton.}
\kat{Suomen Heimo}{29.04.1932}{P7}{Seuraavassa numerossa lehti ilmoittaa saaneensa selville, että sen lukijakunta saartokommunistien keskuudessa on siksi pieni, että se voi ottaa vastaan seuraukset. Törkeämpää terrorisoimista harjoitti Morajärven punakaarti, joka tammikuussa tunkeutui valtuuston kokoussaliin ja uhaten ryöstää läheisen osuuskaupan såi kiristetyksi valtuustolta rahaerän. Mies- ten pukimiemme leikkaaja on työstään näyttelyssä palkittu, ja ulkomailla opis- kellut naisten pukimien leikkajamme on erinomaisen etevä m. m. erikoismallien sommittelussa. «}
\kat{Toukomies}{01.06.1933}{P18}{käsin, ja näin syntyi useampaan otteeseen kii vasta laukausten vaihtoa molemmin puolin, kunnes viimein bolshevikit polttivat tykkitulellaan suuremman osan mainitusta Raja-Konnun kylästä ja valkoisten toimista taas hävisi melkein kokonaan mainittu Kavainon kylä, palaen sen laidassa oleva kappelikirkkokin sitä ympäröivine vanhoine kuusikkoineen. Näissä kahakoissa paloivat poroksi myös Raja-Konnussa sijaitsevat entisen venäläisen tullivartiopalvelusta suorittaneen ratsuosaston kasarmi- ja virastorakennukset, joissa Aunuksen väliaikainen hallitus ja hoitokunta olivat saman vuoden kesäkuussa sijainneet; samoin paloi myös rajan Suomenpuolella Virtelän ja Manssilan kylien osuuskauppa sekä jokunen dots}
\kat{Suomen Heimo}{28.04.1934}{P19}{VIIPURI I V 12 myymälää, kahvila, ■ f säästökassaliike. Haminan ja Ympäristön Osuuskauppa r. l. RAKENNUSOSAKEYHTIÖ Konttori ja keskusvarasto Haminassa. Sivuliikkeet: Laajalti ympäristössä tunnettu KOUVOLAN RAUTAKAUPPA, Kouvola W. -...}
\kat{Suomen Heimo}{04.06.1934}{P15}{Kangas- ja valmis vaatekauppa LEO PAUKKONEN ÄÄNEKOSKI — PUHELIN Paikkakunnan suurin kangas- ja pukuvarasto. 362. Suur-Savon Osuusliike r. l. Mikkeli ja ympäristö O Ostaa ja myy kaikkea maalais-, teollisuus- ja siirtomaatavaraa. Edellä on ollut mahdollisuus vetää esiin vain muutamia piirteitä tapahtuneesta kehityksestä sekä mainita vain muutamia kansainvälisen politiikan perusongelmia.}
\kat{Suomen Heimo}{04.06.1934}{P28}{2. Maan jakautuminen eri omistajaryhmille Itä- Karjalassa v. 1905 ( kihlakunnittain ). OSUUSLIIKE KULTASEPPÄ T Aé Myymälät: HELSINKI, Mikonkatu 5. TURKU, Kauppiaskatu 5. Itä-Karjalan kylät ovat yleensä suuria, mikä seikka muodostaa tavallaan hyvän edellytyksen kollektiivitalouden kehittymiselle.}
\kat{Suomen Heimo}{24.09.1934}{P21}{Tämän voi siis myös Neuvostoliitto sen ollessa syytettynä ilman varauksiakin tehdä. JALKINELIIKKEIDEN OSUUSLIIKE R. L. UNIONINKATU 41. PUH. 21 556. Edelleen hän väitti ylioppilasnuorison edustajien lähteneen liikkeelle ottamatta ensin tarkempaa selkoa asioista.}
\subsubsection*{laukkukauppa \normalfont\small(13 katkelmaa, 13 osumaa aineistossa)}
\kat{Toukomies}{08.1926}{P5}{Suomeenko aikojen kuluessa asettuneiden ja täällä oleskelevien karjalaisten keskuudessa ja toimesta, vai Vienassa kotosalla olevan kansan havahtumisesta? Toisaalta tuli sitä yhtä itsetiedottomasti karjalaismiesten omain matkojen kautta Suomeen — ikivanhoista ajoista tehtyjen laukkukauppa 3a muiden matkojen kautta. Sen, mikä syttyi täällä, voi sanoa syttyneen sielläkin, mitä siellä ajateltiin, se sai paremman toteutumisen ja näkyväisemmän muodon täällä.}
\kat{Toukomies}{08.1926}{P6}{oli laukut täysinä nyt. Laukkuryssä ei ollut vielä häipynyt mielestä. Mut laukussa ei katinkulta,}
\kat{Toukomies}{01.03.1934}{P3}{Näiden kesken kaikilla aloilla ja lukuisissa piireissä on kehittynyt ja noussut myöskin yleisemmän edistys- ja sivistyselämän harrastajia, yhteiskunnallisia toimihenkilöitä, lahjoittajia, kunnallisiin ja yleisiin sivisty spyrkimyksiin innolla osallistuneita. Sitkeydessään melkein vertoja vetämätöntä on ollut Vienan miesten ikiajoista harjoittama laukkukauppa, joka toiselta puolen on tavaran hakuun ja ostoon houkutellut aina Volgan rannoille asti ja toisaalta saattanut lähtemään kaikissa oloissa ■ — lukuunottamatta Lehti ilmestyy Helsingissä: nrot I — 2 * / >> 3 i / 8, 4 i / 4, 5 i / 6, 6 — 7 i / 6 > B — 9 x / 9, 10 x / 10, 11 — 12 x / i 2 ( joulunumero ).}
\kat{Toukomies}{01.06.1934}{P13}{Samana vuonna kuin Genetz- Jännes alkaa toinenkin Karjalan uuttera kuvaaja ja tutkija, valkjärveläinen /. M. Salenius laajan karjala- aiheisen tuotantonsa. * ) Lisäyksenä kirjoituksemme edell. osaan mainittakoon, että Aleksanteri Rahkonen v. 1869 julkaisi pienen näytelmämukailun » Laukkuryssä ». 1881 ) ja » Tutkimus Aunuksen kielestä » ( Suomi-kirja II 17, 1884, erikseen 1885 ), jotka aina näihin asti ovat olleet perusteoksina karjalankielen alalta.}
\kat{Karjalan Vapaus}{01.1935}{P8}{1 ) Kirjoitettava vain toiselle puolelle paperia; » Kuljeksin ta kuljeksin, a sumtsa selköä painau... » Idästä pilvi taas nousi, tuli raesade, oraan alut löi, ruhjoi kuin ruttotauti.}
\kat{Itä-Karjala}{03.1936}{P8}{Kun maan entiset vartijat, syrjäänit ja lappalaiset oli tungettu tieltä tai alistettu veronalaisiksi, oli käytävä alituisia taisteluja rajoilla, hämäläisiä vastaan lännessä, norjalaisia pohjoisessa ja ruotsalaisia vastaan milloin Suomenlahdella, milloin Laatokalla. Mainehikkaitten soturikauppiaitten, jotka miekka vyöllä olivat verottaneet Lappia tai ryöstäneet Ruotsia, jälkeläiset saivat lopuksi " sumtsa " selässä " laukkuryssinä " kiertää heimomaata poliisit ja vallesmannit kintereillään tai rekikauppiaina myyskennellä ryssänparaskeja ja savipotteja perheensä hengenpitimiksi. Ei olekaan mikään pelkkä sattuma se historiallinen totuus, että Venäjän ensimmäinen litämeren-laivasto oli itäkarja dots}
\kat{Viena-Aunus}{01.10.1936}{P5}{/ /. Seitin kauppa. / / /. Laukkukauppa. Kalan hinta ostopaikalla vaihteli sääsuhteiden mukaan.}
\kat{Viena-Aunus}{01.07.1938}{P12}{Viisi, kuusi vuosikymmentä sitten erikoisesti Karjalan rajantakaista osaa pidettiin Suomen suvun suurten muistojen maana, jossa Kalevalan sankarit olivat sen omia sankareita, Laatokan ja Äänisjärven rannat suomalaisten heimojen kehto, seutu, josta käsin nykyisen Suomen valtaus oli aikoinaan tapahtunut. Kirjoittaja toteaa, miten karjalaisten harjoittama laukkukauppa Suomessa on ollut tärkeä kahdessa suhteessa: se oli aikanaan todellisen tarpeen vaatima ja » yksitoikkoiseen elämään hän ( laukkukauppias ), niin korpien takaa kuin tulikin, toi tuulahduksen suuresta maailmasta, loi runoutta arkeen, houkutteli ostamaan sellaista, jota ehkä ei aina tarvittu, mutta joka sittenkin jollakin tavalla dots}
\kat{Viena-Aunus}{01.07.1938}{P12}{Viisi, kuusi vuosikymmentä sitten erikoisesti Karjalan rajantakaista osaa pidettiin Suomen suvun suurten muistojen maana, jossa Kalevalan sankarit olivat sen omia sankareita, Laatokan ja Äänisjärven rannat suomalaisten heimojen kehto, seutu, josta käsin nykyisen Suomen valtaus oli aikoinaan tapahtunut. Kirjoittaja toteaa, miten karjalaisten harjoittama laukkukauppa Suomessa on ollut tärkeä kahdessa suhteessa: se oli aikanaan todellisen tarpeen vaatima ja » yksitoikkoiseen elämään hän ( laukkukauppias ), niin korpien takaa kuin tulikin, toi tuulahduksen suuresta maailmasta, loi runoutta arkeen, houkutteli ostamaan sellaista, jota ehkä ei aina tarvittu, mutta joka sittenkin jollakin tavalla dots}
\kat{Viena-Aunus}{18.10.1938}{P13}{Nyt sen sijaan voidaan antaa vähäinen rahallinen korvaus niille, jotka lähettävät seuralle omia tai muilta muistiin merkitsemiään tietoja kirjallisessa muodossa kerättäviltä aloilta. Sellaisina historiallisina muistitietoina, joita vanhimmalla karjalaispolvella itsellään edelleen on, mainittakoon esim. » varassotien » aikuiset tapahtumat rajan molemmin puolin sekä » laukkukaupan » Saadakseen historiallisten muistitietojen ja muinaismuistojen keräysharrastuksen heräämään ja virkenemään Suomessa olevien karjalaisten keskuudessa, on Karjalan Sivistysseuran johtokunta varannut huomattavan rahaerän tätä tarkoitusta varten.}
\kat{Viena-Aunus}{12.12.1938}{P15}{Ufotta askelen alla, ah, i / stävä, on muisto / a kauniita niin. Laukkukauppa, jota viimeksi on harjoitettu maamme syrjäseuduilla, loppui tämän vuosisadan alkuvuosina eli Marraskuun alkupuolella kerrottiin puhjenneen vakavia levottomuuksia Siperiassa, missä koko Kaukoidän puna-armeija on ollut kuohumistilassa.}
\kat{Viena-Aunus}{12.12.1938}{P21}{Rahdin ja tukin ajo oli myöskin tärkeä tulolähde kylän asukkaille. Sen sijaan Suomessa laukkukaupan harjoittaminen oli juuri uhtualaisille ominaista. Tällä seikalla oli erittäin suuri merkitys paikkakunnalle kansallisessakin mielessä.}
\kat{Viena-Aunus}{12.12.1938}{P23}{Jyvälahdessa vieras on huomaavinaan vielä läheisemmän Suomen rajan tunnun kuin kunnan muissa kylissä. Kylän monet miehet, samoin kuin luusalmelaisetkin, olivat ahkeria laukkukaupan harjoittajia Suomessa. Kylän elämässä emme havaitse mitään erikoisleimaa, jolla se poikkeaisi muista selostamistamme kylistä.}
\subsubsection*{työttömyys \normalfont\small(25 katkelmaa, 47 osumaa aineistossa)}
\kat{Suomen Heimo}{18.11.1928}{P16}{Tällaisia talonpoikaismaita kutsuttiin nimellä obza, jopka suuruus arvioitiin siihen kylvetyn viljamäärän mukaan. Työttömyys ja köyhyys lienee tullut Nevan seutujen asukkaille, jotka alkoivat muuttaa työansioihin muualle! Mutta yleensä obzien suuruus vaihteli suuresti riippuen maanlaadusta, paikallisista olosuhteista y.m.}
\kat{Suomen Heimo}{15.02.1930}{P28}{Sokerista Sokeri ei ole vain mauste. Ennen oli sokeri harvinainen mauste, ja pula-aika opetti taas säästämään sitä. Sokeri on nimittäin parhaita ravintoainei- tamme, ja myöskin hintansa suhteen edullisimpia.}
\kat{Suomen Heimo}{11.08.1930}{P36}{42 haarakonttoria; \_ ",... Ennen oli sokeri harvinainen mauste, ja pula-aika opetti taas säästämään sitä. Sokeri on nimit- täin parhainta ravintoainettamme ja myöskin hintansa suhteen v edullisimpia.}
\kat{Suomen Heimo}{16.12.1930}{P38}{Sen I osa, " I s a ja t y t a r ", kuvastaa sisämaan pienen sahayhteiskunnan oloja venäläissorron aiheuttaman passiivisen vastarinnan vuosina ja kansanelämän sielunliikkeitä ennen sosialidemokratian tunkeutumista maahamme. Mielikuvituksellinen aihe: eräs huomattava suurliikemies vetäytyy syrjään tavallisesta liike-elämästä, rakennuttaa itselleen loistoasunnon pohjoiseen, erämaiden keskelle, ja luo siellä jättiläissuunnitelman, minkä tarkoituksena on Pohjolan teollisoiminen, työttömyyden ja köyhyyden karkoittaminen maasta. " Kun itämyrsky vinkui " on isänmaallinen romaani ja tekijä — " Jääkärimarssin " runoilija — haluaa siinä osoittaa, minkä arvoisen uhrin Suomen nuoriso kantoi isänmaalleen dots}
\kat{Suomen Heimo}{24.09.1931}{P18}{On muuten omituista, että ruotsalaiset kykenevät toisiaankin, vieläpä vanhoja aatelismiehiä ja kartanonomistajia pitämään niin lujassa kurissa, että näiden on vastoin tahtoaan ajettava suomalaiset työläiset maantielle, kuten äskettäin Koskenkylässä. Onko syytä heittää rahaa kuin " Kankkulan kaivoon " pelkistä koristeellisista syistä, samalla kun maanviljelijäin tilat joutuvat vasaran alle, työttömyys painaa työläisiä, virkamiesten palkat alenevat ja suomalaisen koululaitoksen menoja supistetaan? Sellaiseksi ei kai voitane leimata niitä Oulun koulun määrärahoja, jotka on muilutettu vastoin hallituksen enemmistön mielipidettä tulo- ja menoarvioon, enempää kuin sitäkään Turun teollisuuskouluun dots}
\kat{Suomen Heimo}{20.12.1931}{P22}{Herrojen Lavi ja Aminoff tietojen mukaan ovat 28 lapsen vanhemmat ilmoittaneet haluavansa panna lapsensa suomalaiseen kouluun ja kielitaitonsa perusteella olisi heidän mukaan 22:n oppilaan saatava opetusta suomenkielellä. Kun pula-aika on voitettu ja elämä taas palautuu normaalioloihin, saavuttaa maanviljelys entisen laajuutensa, tehtaitten ja sahojen pyörät alkavat pyöriä entiseen tapaan, jolloin uutta suomalaista työväestöä virtaa paikkakunnalle. Suomenkielisten kansakouluolot olisi ratkaistava ehdottamaamme suuntaan aivan lähiaikoina ja viimeistään normaalioloihin saavuttaessa, sillä ei toki ole kohtuudenmukaista, että suomenkielisestä lapsesta kansakouluun mennessä tulee ruotsinkielinen.}
\kat{Toukomies}{01.04.1931}{P61}{C 925 -193? 7 J\textasciicircum{}\textasciicircum{}M htä LUJAA kuin Sw T\textasciicircum{}' - kaunistakin: MB\& lama en rf / f\textasciicircum{}\textasciicircum{}W ] forssan kangas jSM Banoiiava vrjSlMm vm — kansallis-osake-pankki gtanlututautfankktf ij Perusi cl Ui v. 1889.}
\kat{Toukomies}{01.12.1931}{P2}{TAMPELLA pöytäliinoja — pyyhinliinoia — liinoja, Vuodehuopza, Pöytä- ja Lautasliinoja... [ lakanakankaita 'neljän sukupolven kokemus Pesusta pellava paranee kutomateosten valmistajana' \textasciitilde{} ■ ERÄS = = = PÄTEVÄ, TUNTOMERKKI, KI\textasciicircum{}JH kaitten puolesta: BR'Ä » YHDISTYNEITTEN » miesten- ja naistenkankaat on K? \% ' " \&JM palkittu korkeimmilla palkinnoilla kotimaisissa näytte- * |? |\textasciicircum{}\textasciicircum{}g lyissä. Finlaysonin valmisteiden erinomainen arkioloissakin. kestävyys ja ehdottoman varma värinpitävyys osoitta- Vanhoihin damasteihin, samoinkuin Tampella-,....., j liinoihin ei koskaan » totu », ne eivät koskaan vat, että Finlayson, maamme vanhin puuvillatehdas, on... * J / \textasciicircum{} tule arkipäiväisiksi ja kuluneiksi, lama dots}
\kat{Suomen Heimo}{22.02.1932}{P31}{Niinpä eräs perheellinen Karjalan pakolainen ansaitsee siellä noin 7 markkaa päivässä paperipuun valmistuksessa, sillä työtä tehdään harvennushakkauksena ja työmaa on 10 kilometrin päässä asunnosta. Kun nykyisenä pula-aikana työttömät pakolaiset eivät jaksa hankkia lapsilleen oppikirjoja ja kuntakaan ei heitä tässä kohden avusta, niin heidän on pakko ottaa lapsensa pois koulusta. Toinen perheellinen pakolainen, joka oli hankkinut hyvänä aikana pienen torpan, mutta nyt menetti sen, kun " ulkomaalaisena " ei voinut saada ilman erikoista toimenpidettä siihen omistusoikeutta sekä tarpeellista lainaa.}
\kat{Suomen Heimo}{28.07.1932}{P8}{Omistajat eivät aluksi, niin kauan kuin nämä elimet rajoittuivat huolehtimaan raaka- ja polttoaineitten sekä elintarpeitten hankkimisesta ja kulutusosuuskuntien perustamisesta, älynneet niitä vastustaa, pikemminkin he katsoivat niiden helpottavan heidän työtään. Yleisen tuotannon lamautumisen ja työttömyyden pelko yllytti työväkeä jatkamaan vallankumousta tälläkin alalla, äärimmäiset ainekset saivat vallan ja jo keväällä 191 7 oltiin niin pitkällä, että tehdaskomiteain pääasiallisena ohjelmana oli, ei tuotannon uudelleen järjestäminen, vaan työläisten työvelvollisuuden vähentäminen ja etujen lisääminen. Jos sattuneet poikkeukset olisivat vain tilapäistä laatua, olisivat ne varsin dots}
\kat{Suomen Heimo}{06.11.1932}{P15}{Nyksitär. Koskeeko nykyinen pula-aika Nyksiä? Naisen sydän on tunnetusti hellä.}
\kat{Suomen Heimo}{06.11.1932}{P15}{— Kesäkiertuilu on oikeastaan aika mukava kansanvalistustapa. Sillähän kyllä on ainainen pula-aika, kroonillinen rahapula. Se näet saa kaikki hopealusikat rahassa.}
\kat{Suomen Heimo}{17.12.1932}{P12}{Hän on palannut isiensä lakeuksille — ei tosin Kurikkaan, josta kerran lähti — vaan häjyjen, puukon, paksuna virtaavan veren ja körttivästin Härmään, jota me tavallinen kansa olemme tottuneet pitämään eteläpohjalaasten napapitäjänä. Nyt on kirjojen maailman piiri suunnattomasti laajentunut, ja pula-aika — pula rahasta ja ajasta — ( " Suomen Heimon " kirjakronikasta puheenollen: pula myöskin tilasta ) — sallii tutustua vain harvoihin uutuuksiin. Miten paljon kansatieteellisen aineiston rikkautta, henkilöhistoriallisia tietoja ja kansanpsykologiaa ( viimeksimainittua etenkin häjyjen merkillisen ilmiön kuvauksessa ) Paulaharju on saanut mahtumaan hyvien valokuvien, omintakeisten piirustusten dots}
\kat{Toukomies}{01.02.1932}{P15}{Lumi kinostuu, kattaa kaikki On oikea työttömyyden aika. Kuka sellaisella ilmalla metsään!}
\kat{Suomen Heimo}{10.06.1933}{P6}{Mikä toisin sanoen on se ohjelma, joka pystyy lyömään itsensä läpi nykyisten vallanpitäjien tuiman vastustuksen, ohjelma, joka sisäiseltä voimaltaan on, niin suuri, että se kykenee vetämään yhteistyöhön kaikki kansankerrokset ja siten poistamaan yhteiskunnastamme sen luokkien välisen kuilun, jota turhaan yritetään kauniilla puheilla täyttää. Mitä on muuta kuin suunnitelmallisuuden puutetta se, että kymmenettuhannet työttömät pannaan rakentamaan kaupungeissa aivan tarpeettomia katuja, satamalaitureja j. n. e., maalla tuottamattomia autoteitä, sen sijaan että kunnat vapautettaisiin työttömyysasioista — ja samalla demoralisoivista soppajonoistaan — ja työttömät pantaisiin maanparannus- ja dots}
\kat{Suomen Heimo}{10.06.1933}{P6}{Mikä toisin sanoen on se ohjelma, joka pystyy lyömään itsensä läpi nykyisten vallanpitäjien tuiman vastustuksen, ohjelma, joka sisäiseltä voimaltaan on, niin suuri, että se kykenee vetämään yhteistyöhön kaikki kansankerrokset ja siten poistamaan yhteiskunnastamme sen luokkien välisen kuilun, jota turhaan yritetään kauniilla puheilla täyttää. Mitä on muuta kuin suunnitelmallisuuden puutetta se, että kymmenettuhannet työttömät pannaan rakentamaan kaupungeissa aivan tarpeettomia katuja, satamalaitureja j. n. e., maalla tuottamattomia autoteitä, sen sijaan että kunnat vapautettaisiin työttömyysasioista — ja samalla demoralisoivista soppajonoistaan — ja työttömät pantaisiin maanparannus- ja dots}
\kat{Suomen Heimo}{10.06.1933}{P6}{Eihän toki sovi väittää, että meillä rappeutuneisuus poliittisessa enemmän kuin taloudellisessakaan suhteessa olisi niin suunnattomasti levinnyt, kuin esim. Saksassa vielä joku aika sitten, erittäinkin kun puhtaasta vapaamielisyydestä on pala palalta tingitty, mutta nykytilanne riittää kuitenkin synnyttämään meilläkin yleisen tyytymättömyyden oleviin oloihin ja uudistusten odotuksen. Senjälkeén kun ylläoleva jo oli kirjoitettu, olemmekin saaneet nähdä I. K. L:n " taloudelliset suuntaviivat " julkisuudessa, ja on suoraan sanottava, että ne vastaavat, niin ylimalkaisia kuin ne ovatkin, edelläsanotussa lausuttuja toivomuksia,. Erittäinkin juuri työttömyyden torjunta ja asutustoiminnan dots}
\kat{Suomen Heimo}{22.12.1933}{P4}{Toisena mahdollisuutena on vain asian silleen jättäminen — ja siihen ei Suomen kansalla ole yksinkertaisesti varaa. Kunnat on vapautettava työttömyyden niille aiheuttamasta sietämättömästä taakasta, joka valtion on otettava kantaakseen, ja annettakoon sitä vastaan valtiolle rajaton oikeus määrätä Xi tarvitse nykypäivinä olla erikoisen tarkkakatseinen havaitakseen, että valtiollisen ja yhteiskunnallisen totuuden etsintä on käynnissä, ja vielä varsin kiihkeästi, kautta koko sivistyneen maailman.}
\kat{Toukomies}{01.04.1933}{P12}{Mutta mitä on sitten niistä sanottava, jotka pitävät rahansa kotonaan tai taskussaan eivätkä tahdo käyttää niitä tarpeisiinsa eikä myöskään niillä mitään teettää? Ensiksikin sitä on aiheuttanut tämä tunnettu taloudellinen pula-aika ja toiseksi henkisten voimiemme puute, sekä lisäksi meidän karjalaisten heikko yhteisymmärryksellinen toiminta. Moni voi huomauttaa, pitääkö jokaisen, jolla on rahaa paljon tai vähän, pyrkiä siitä vapautumaan ja ostamaan ensimmäistä parasta tavaraa, mikä eteen sattuu.}
\kat{Toukomies}{01.05.1933}{P2}{Tampella- pellavaa kannattaa käuttää aina. / 1 • 111 ivlingendahl, JNUt\textasciicircum{} * Wm. rw T M P E R E a Lfrijanii:\textasciicircum{} ju jttw w i Jim\textasciicircum{}r Kaivattuun kesälomaan li » VIIDEIM TEHTAAN » ihanine aurinkoisine päivineen, retkineen, iloineen I I |- — | Sn\textasciicircum{}mmAlsolia\textasciicircum{}iuaTla * on enää vain lyhyt aika jäljellä. -....,....,....,, IIH aina alhaisi m n i n a. sista tärkeimmät kumipohjaiset kesäkengät! liVVf\textasciicircum{}ri T-T VVT. IV-l I „.. „,.. „.... „ 1 •' h |jMÉÉiÉiMiJÉjMÉliÉ>iiii|iij Lama toimenpide koituu reidänkin Hankkikaa Tekin sentähden ajoissa itsellenne eduksenne. L\textasciicircum{}S\textasciicircum{}L Ai " if * * KOTIVILLAI NEN miesten kangas C\textasciicircum{} \textasciicircum{}; " HÄMÄLÄINEN " ja samasta raa- JFfj| / / x\textasciicircum{} / ■ ka-aineesta, mutta halvempi valmiste i V\textasciicircum{}y dots}
\kat{Toukomies}{01.10.1933}{P11}{Mutta niiden keskeltä kohosi pitkänlainen uusi kolhoositalouden talo ja toisaalla näyttivät olevan vielä telineet uuden » neuvostokansakoulun » seinillä. Komeanlaisia isonpuolisia taloja, suomalaisen kansakoulun uusi talo korkealla mäellä; ■ — kulunut vaikea pula-aika vain oli vaikuttanut tännekin asti talonpojan elämään. Kun tulimme Suojun lautalle ja tapasimme joen varrella kauniit raja- ja tullivartioston asunnot sekä kuulimme niiden ikkunoista iloista gramotooninsoittoa, emmepä uskoneet}
\kat{Suomen Heimo}{03.02.1934}{P23}{Ylläesitetylle painostukselle päinvastainen ilmiö on se, että eräillä paikkakunnilla, varsinkin suuremmissa asutuskeskuksissa kaksikielisten, vieläpä ruotsinkielisten kotien lapsia käy suomalaisessa koulussa. Monet ovat työttömyyden uhallakin säilyttäneet äidinkielensä ja mieluummin valinneet hädän ja puutteen kuin sellaisen alistumisen, joka äidinkielen kustannuksella takaisi paremmat elinehdot. Tämä johtuu ensinnäkin siitä, että niistä huolehtivat järjestöt ja henkilöt eivät ole voineet läheskään riittävästi avustaa niiden toimintaa, ja toiseksi se, että valtioapu on pienentynyt.}
\kat{Suomen Heimo}{03.03.1934}{P19}{Kuluvana lukuvuonna on niissä ollut oppilaita alakouluasteella 125 ja yläkouluastcella 219 eli yhteensä 344 oppilasta. tulee kysymykseen uuden työvoiman hankkiminen joten tämä teollisuus voi myös suorittaa oman osuutensa työttömyyden poistamisessa. Siinä voi jokainen sekä oman, että " omavaraistuvan Suonien " edun vuoksi tehdä vain päätöksen:}
\kat{Suomen Heimo}{06.04.1934}{P7}{Juutalaiskysymys on Unkarissa ainakin toistaiseksi ratkaistu siten, että heihin nähden on\textasciicircum{}käytännössä numerus clausus, s. o. vain prosenttimääräänsä vastaava osa pääsee yliopistoihin ja saa suorittaa tutkintonsa. Työttömät ylioppilaat ovat keksineet omalaatuisia keinoja olemassaoloaan helpoittaakseen, esim. joukko juristeja vuokraa 2 — 3 huoneen huoneiston, sijoittaa sinne joukon sotilassänkyjä, keittää itse ruokansa tullen siten hyvin halvalla toimeen. Olemme alkaneet pitää entistä run- saampaa valikoimaa myöskin valmiina KÄYKÄÄ KATSOMASSA! P. ESPLANAADINKATU39}
\kat{Suomen Heimo}{30.04.1935}{P11}{Ei ainakaan pitkiin aikoihin. Pelkkä pula-aika ei riitä yleisön välinpitämättömyyden puolustukseksi. Eikä se lakkaa niihin uniinsa myös uskomasta.}
\subsection{Paikat ja nimet}
\subsubsection*{Suojärvi \normalfont\small(30 katkelmaa, 418 osumaa aineistossa)}
\kat{Toukomies}{05.02.1925}{P10}{Ken siltojen kera elon kauniimman tois. ruman häätäisi pois? ( Suojärven ensimmäisen nuorisoseuratalon vihkiäisiin Leppäniemellä 21 / g ig24. ) rajat maan, sulut heimon ja hengen yö, sun suurin ja voittoisin vartiotyö:}
\kat{Toukomies}{05.02.1925}{P10}{rajat maan, sulut heimon ja hengen yö, sun suurin ja voittoisin vartiotyö: Suojärvi, suurien salojen maa, Suojärvi, suurien vetten vyö, Suojärvi, min rantoja rajoittaa satu hohdollaan salot peittänyt maan, runon kulta ja kunnia välkkynyt, miss' oisi se hohto nyt hohtelemaan, joka uudestaan unisuudestaan}
\kat{Toukomies}{05.02.1925}{P10}{rajat maan, sulut heimon ja hengen yö, sun suurin ja voittoisin vartiotyö: Suojärvi, suurien salojen maa, Suojärvi, suurien vetten vyö, Suojärvi, min rantoja rajoittaa satu hohdollaan salot peittänyt maan, runon kulta ja kunnia välkkynyt, miss' oisi se hohto nyt hohtelemaan, joka uudestaan unisuudestaan}
\kat{Toukomies}{05.02.1925}{P10}{rajat maan, sulut heimon ja hengen yö, sun suurin ja voittoisin vartiotyö: Suojärvi, suurien salojen maa, Suojärvi, suurien vetten vyö, Suojärvi, min rantoja rajoittaa satu hohdollaan salot peittänyt maan, runon kulta ja kunnia välkkynyt, miss' oisi se hohto nyt hohtelemaan, joka uudestaan unisuudestaan}
\kat{Toukomies}{05.02.1925}{P10}{satu hohdollaan salot peittänyt maan, runon kulta ja kunnia välkkynyt, miss' oisi se hohto nyt hohtelemaan, joka uudestaan unisuudestaan Suojärvi, suurene kummuillas, sun palttees ja rantasi kultahan luo suur maailma olkohon suuntanas, mut omista juuristas voimasi juo, kuin malmit järvies pohjasta tuo, Maan äärille nouskaa te talot elämän, talot valkeuden nouskaa te paistelemaan, isot ikkunat illassa välkkymähän ja aamun kullassa kuin palamaan, kuin palamaan vain sydän korpisen maan rusohohteissaan!}
\kat{Toukomies}{05.02.1925}{P17}{Karjalan Sivistysseura, Vienan Karjalaisten Liiton perijä, laskee täten syvimmällä kaipauksen, kiitollisuuden ja kunnioituksen tunteella tämän seppeleensä tohtori Benjamin Mitron haudalle. G. Gulin, os. Pitkäranta; Salmin jakopaikka, hoitajana vääpeli V. Jakobson, os. Salmi ja Suojärven jakopaikka, hoitajana vääpeli I. Halonen, os. Suojärvi, Suvilahti. Kajaanin jakopaikka, hoitajana V. Rasi, os. Kajaani; Kuhmoniemen jakopaikan hoitajana V. Vaara, os. Kuhmoniemi; Kiehimän jakopaikan hoitajana P. Kanerva, os. Kiehimä; Ristijärven jakopaikan hoitajana F. Semenoff, os.}
\kat{Toukomies}{05.02.1925}{P17}{Karjalan Sivistysseura, Vienan Karjalaisten Liiton perijä, laskee täten syvimmällä kaipauksen, kiitollisuuden ja kunnioituksen tunteella tämän seppeleensä tohtori Benjamin Mitron haudalle. G. Gulin, os. Pitkäranta; Salmin jakopaikka, hoitajana vääpeli V. Jakobson, os. Salmi ja Suojärven jakopaikka, hoitajana vääpeli I. Halonen, os. Suojärvi, Suvilahti. Kajaanin jakopaikka, hoitajana V. Rasi, os. Kajaani; Kuhmoniemen jakopaikan hoitajana V. Vaara, os. Kuhmoniemi; Kiehimän jakopaikan hoitajana P. Kanerva, os. Kiehimä; Ristijärven jakopaikan hoitajana F. Semenoff, os.}
\kat{Toukomies}{05.02.1925}{P18}{1 inkeril., Lepaan puutarhaopistossa 2 inkeril., Kangasalan kotitalousopistossa 1 inkeril., Viipurin maanviljelyslyseossa 1 inkeril., Turun kauppakoulussa 1 inkeril., Hämeenlinnan kauppakoulussa 1 inkeril., Tarvaalan maanviljelyskoulussa 1 inkeril., Elisenvaaran maanviljelyskoulussa 1 inkeril., Otavan maanviljelyskoulussa 1 inkeril., Vakka-Suomen maamieskoulussa 1 inkeril., Ylä-Vuoksen maamieskoulussa 1 inkeril., Mustialan karjanhoitokoulussa 1 inkeril., Etelä-Karjalan karjanhoitokoulussa 1 inkeril., Reitkallin puutarhakoulussa 1 inkeril., Käkisalmen käsityökoulussa 1.7 inkeril., Oulun kauppakoulussa 1 karjal., Suomen saha- ja teollisuuskoulussa 1 karjal., Orismalan maatalouskoulussa 1 dots}
\kat{Toukomies}{01.03.1925}{P5}{— Kannaksen Vartion n rot I ja 2 ovat ilmestyneet. Lukemistopuolen sisältönä ovat sadut, kootut Suojärven pitäjän eri kylistä. Lokakuussa 1922 alkoi Helsingissä ilmestyä tämänniminen lehti kirjailija Antti Tiittasen toimittamana.}
\kat{Toukomies}{01.03.1925}{P5}{( Tekijän aikaisemmissa julkaisuissa kielennäyttein esitettyä suomenpuolista karjala-aunuksen aluetta ei ole otettu näissä kirjoissa mukaan ). Toisessa puolessa tutkimustaan esittää hän Suojärven karjalan kieliopin ja toisessa yhtä laajassa puolessa kielinäytteitä samalla kielimurteella. Huhtikuun 1 p:stä 1923 ilmestyy Helsingissä toimivan Akateemisen Karjala-Seuran julkaisemana heimoasioita käsittelevä aikakauslehti » Suomen heimo ».}
\kat{Toukomies}{01.03.1925}{P5}{Nämä ovat kuitenkin vasta esitöitä, joihin läheisesti liittyy tutkijan esitys karjalaisille läheisistä Vepsän ohjoisista etujoukoista Kielettären 4 ja 5 vihkossa 1872 ja 1873, vaikka näiden » lyydiläisten » kieli luetaankin vepsänkieleen kuuluvaksi. Mutta samana vuonna hän Suomikirjan II jakson 8 osassa julkaisee varsin laajan — 75-sivuisen — » kertomuksen Suojärven pitäjäästä ja matkustuksistani siellä v. 1867 », joka, lyhyttä esipuhetta mainitsematta, on kokonaan kielellinen esitys. Lehden päätoimittajaksi ryhtyi viime syksystä tohtori Väinö Salminen, jonka johdolla lehti on ottanut voimakkaan suomalaiskansallisen suunnan ja on tässä hengessä miehekkäissä kirjoituksissaan pohtinut monia dots}
\kat{Toukomies}{01.03.1925}{P6}{95 p. ja jälki vaatimus varoja 674,261 mk. — Suojärven rautatieasema -suurituloisimpia asemia Karjalassa. Lähetettyä karjaa, hevosia ja ajoneuvoja 774 kpl.}
\kat{Toukomies}{01.03.1925}{P6}{Valituksi tuli kolmannessa äänestyksessä maalaisliiton ehdokas, Viipurin läänin maaherra Lauri Kristian Relander. Suojärven aseman toiminnasta viime vuonna on » Laatokka » saanut seuraavia tietoja: Yhteensä kannettu liikennevaroja 4,956,445 mk. Sitä varten oli viime tammikuulla toimitettava niiden valitsijamiesten vaali, joiden tuli valita valtakunnan presidentti uudeksi valtakaudeksi.}
\kat{Toukomies}{01.04.1925}{P16}{1. H.; Otteita evankeliuuneista karjalankielellä; Kibunoi, suomalaisia runoja aunukselaisille, sivut 9-16. — Kuvia: Kalevalan heimoa kansallisilla kesäjuhlillaan; Kaksi karjalaisia esittävää kuvaa Olaus Magnin teoksesta v. 1567; Myöhemmän ajan karjalaisia kauppamatkoilla; Kauppias Ivan Mitropanof f: Wanhalan emäntä ja tytär ja Wanhalan talo Wuokkiniemellä. — Karjalan Sivistysseuran apurahat Kontro \& Kuosmasen lahjoitusrahaston viimevuotisten korkojen Suomen puolelle jaettavista varoista on jaettu seuraavain kansakoulujen köyhille oppilaille: Ilomantsin Marjovaaran kk. 250 mk., Ilomantsin Möhkön kk. 400 mk., Korpiselän Saaroisten 350 mk., Salmin Ignoilan kk. 300 mk., Salmin Työmpäisten kk. dots}
\kat{Toukomies}{01.04.1925}{P16}{1. H.; Otteita evankeliuuneista karjalankielellä; Kibunoi, suomalaisia runoja aunukselaisille, sivut 9-16. — Kuvia: Kalevalan heimoa kansallisilla kesäjuhlillaan; Kaksi karjalaisia esittävää kuvaa Olaus Magnin teoksesta v. 1567; Myöhemmän ajan karjalaisia kauppamatkoilla; Kauppias Ivan Mitropanof f: Wanhalan emäntä ja tytär ja Wanhalan talo Wuokkiniemellä. — Karjalan Sivistysseuran apurahat Kontro \& Kuosmasen lahjoitusrahaston viimevuotisten korkojen Suomen puolelle jaettavista varoista on jaettu seuraavain kansakoulujen köyhille oppilaille: Ilomantsin Marjovaaran kk. 250 mk., Ilomantsin Möhkön kk. 400 mk., Korpiselän Saaroisten 350 mk., Salmin Ignoilan kk. 300 mk., Salmin Työmpäisten kk. dots}
\kat{Toukomies}{01.05.1925}{P7}{Siellä edustivat seuraamme sen puheenjohtaja johtaja A. Mitro ja seuran jäsen johtaja Leo Pankkonen, joista edellinen piti kokoukseen avajaistilaisuudessa tervehdyspuheen seuramme ja rajantakaiskarjalaisten puolesta. Suomenpuolelle jaettavat apurahat on jaettu seuraavasti: Salmin Työmpäisten kansakoululle 350 mk., Salmin Ignoilan kk. 300 mk., Suojärven Jehkilän kk. 400 mk., Suojärven Kaitajärven kk. 500 mk. Kirjallisuuden levittämisen asiaa on seura pitänyt vireillä koko vuoden kartuttamalla varastoaan tulevia kansankirjastoja varten ja sen lisäksi toimivuoden lopulla käsitellyt pari asiaa koskevaa yksityisluontoisempaa kysymystä.}
\kat{Toukomies}{01.05.1925}{P7}{Siellä edustivat seuraamme sen puheenjohtaja johtaja A. Mitro ja seuran jäsen johtaja Leo Pankkonen, joista edellinen piti kokoukseen avajaistilaisuudessa tervehdyspuheen seuramme ja rajantakaiskarjalaisten puolesta. Suomenpuolelle jaettavat apurahat on jaettu seuraavasti: Salmin Työmpäisten kansakoululle 350 mk., Salmin Ignoilan kk. 300 mk., Suojärven Jehkilän kk. 400 mk., Suojärven Kaitajärven kk. 500 mk. Kirjallisuuden levittämisen asiaa on seura pitänyt vireillä koko vuoden kartuttamalla varastoaan tulevia kansankirjastoja varten ja sen lisäksi toimivuoden lopulla käsitellyt pari asiaa koskevaa yksityisluontoisempaa kysymystä.}
\kat{Toukomies}{01.06.1925}{P14}{Pistojärven pravl. kuuluu: Pistojärvi, Koljola, Kontu, Malviainen, Suvanto, Kiimaisjärvi, Makarila, Tuhkala, Katoslampi, Akkala, Korpijärvi, Ohta ja Hirvisalmi. Kontokin pravl.: Kontokki, Kostamus, Ahvenjärvi, Suojärvi, Koivajärvi, Kenttijärvi, Vatasalmi, Vonkajärvi, Lusma, Munankilahti, Tetriniemi, Sappovaara, Akonlahti, Lyttä, Vuokinsalmi ja Niskajärvi. Toukomies alkaa julkaista karjalaisten itsensä kirjoittamia pieniä muistelmia ja kuvauksia heidän kotiseuduistaan, havainnoistaan pakolaisajalla y.m. Tässä julkaistaan aluksi nuoren aunukselaisen Elias Laasosen kuvaus keskiaunukselaisesta kotiseudustaan Munajärven pitäjässä, Petroskoin kihlakunnassa.}
\kat{Toukomies}{01.06.1925}{P22}{— Viipurin hiippakunnan läntisestä vaalipiiristä, johon kuuluvat Vaasan, Turun, Hämeenlinnan ja Helsingin seurakunnat, tuli valituksi herra F. A. Kuklin Helsingistä; saman hiippak. keskisestä vaalipiiristä, johon kuuluvat Kotkan, Haminan, Lappeenrannan, Viipurin ja Sorvalin kappelin seurakunnat, kauppias J. Kotshetoff Viipurista; saman hiippak. itäisestä vaalipiiristä, johon kuuluvat Kyyrölän, Uudenkirkon, Raivolan, Terijoen, Kellomäen ja Petsamon seurakunnat, kirkonisännöitsijä Ivan Kondyrev Raivolasta. Maallikkoedustajiksi siihen ovat valitut: Karjalan hiippakunnan eteläisestä vaalipiiristä, johon kuuluvat Palkealan, Käkisalmen, Tiurulan, Kuopio — Savonlinnan, Sortavalan, Kitelän ja dots}
\kat{Toukomies}{01.09.1925}{P1}{Ilmestyy kerran kuussa kunkin kuun 1 p:nä 12 — 16 siv. numeroina, joista heinäk. nro ilm. kesäkuun ja elok. nro syyskuun nron kera yhdessä sekä samoin tänä vuonna nrot I — 21 — 2 yhdessä. Perustetaan » Karjalan sivistysseura », käydään Venäjän väliaikaisen hallituksen luona, anotaan Suojärven rataa, herätetään ja valmistetaan karjalaisia Venäjän perustavaa kansalliskokousta varten, perustetaan oma lehti » Karjalaisten sanomat » ( 1917 — 1920 ). Tuossa v. 1914 alkaneen maailmansodan uuvuttamassa jättiläisvaltakunnassa alkaa v. 1917 vallankumous, jommoista maailma 2i vielä ollut ennen nähnyt, ja tsaarien rautakourin koossapitämä ulkonaisesti mahtava valtakunta, lahonnut ja sisällisesti dots}
\kat{Toukomies}{01.09.1925}{P7}{Niinpä salmilainen maisteri Jo h. Kujola — synnyltään siis jo karjalankielen, vieläpä aunuksen piiristä — v. 1910 julkaisi ansiokkaan ja tarkasti tieteellisen esityksen » Äänneopillinen tutkimus Salmin murteesta », jonka kielennäytteet ja lukuisat sanaesimerkit ovat merkinnältään niin tarkkaa karjalankielen kirjoitustapaa, että Perusteet ovat samat, mutta » ensiksikin on käytetty hyväksi niitä runsaita sanakirjallisia muistiinpanoja, joita dosentti K. F. Karjalainen, lehtori E. V. Ahtia ja maisteri Juho Kujola ovat keräilleet », edellinen Vienan Karjalan ja Tverin, Ahtia Suojärven ja Kujola Tihvinin karjalaismurteista, jotka juuri ovat laajennuksen ja osaksi käsitystenkin muuttumista dots}
\kat{Toukomies}{01.09.1925}{P13}{Ja se puhdistus viimein täytyy ulottua itse kirkon palveluksessa olevien henkilöiden nimiinkin, varsinkin, missä suvun juuri on suomalainen tai karjalainen. Paul Tulehmo y.m. olivat taas esittäneet, että Suojärven suureen, harvaanasuttuun ja sälöiseen rajaseurakuntaan perustettaisiin kolmannen, kiertävän papin toimi. Varsinaisia kysymyksiä oli seurakunnan valtuustojen täydentämistä koskeva kysymys, josta lakivaliokunta oli jo laatinut mietinnön ja joka mietintö sellaisenaan hyväksyttiin,}
\kat{Toukomies}{01.09.1925}{P21}{Kirkko on vaikeassa tilassa, ellei sille saada sivistyneitä, kansallismielisiä pappeja, ellei jumalanpalvelus tule henkeväksi saarnoihin ja veisuihin nähden, ellei papit kykene pitämään hyviä hartaushetkiä, ellei kirkolliskokouksen päätöstä kirkkomme kansallistuttamisesta tositeossa toteuteta j. n. e. Pitäisi ponnistella kirkkolaulujemme uuteen henkeen pukemisessa edelleen. Kun Suojärven ja Aunuksen murteissa sana hauta nominatiivissä sanotaan haudu ja niin myös siis kylännimi Hautavaara kuuluu Hatiduvuara ( genetiivissä ja muissa sijoissa ei koskaan haudun, hauduksi, vaan hauvan, hauvaksi j. n. e. ), niin on eräs » Kotilieden » kirjoittaja lehden 15 — 16 n:ssa kuvatessaan Suojärven oloja dots}
\kat{Toukomies}{01.09.1925}{P21}{Kirkko on vaikeassa tilassa, ellei sille saada sivistyneitä, kansallismielisiä pappeja, ellei jumalanpalvelus tule henkeväksi saarnoihin ja veisuihin nähden, ellei papit kykene pitämään hyviä hartaushetkiä, ellei kirkolliskokouksen päätöstä kirkkomme kansallistuttamisesta tositeossa toteuteta j. n. e. Pitäisi ponnistella kirkkolaulujemme uuteen henkeen pukemisessa edelleen. Kun Suojärven ja Aunuksen murteissa sana hauta nominatiivissä sanotaan haudu ja niin myös siis kylännimi Hautavaara kuuluu Hatiduvuara ( genetiivissä ja muissa sijoissa ei koskaan haudun, hauduksi, vaan hauvan, hauvaksi j. n. e. ), niin on eräs » Kotilieden » kirjoittaja lehden 15 — 16 n:ssa kuvatessaan Suojärven oloja dots}
\kat{Toukomies}{01.10.1925}{P6}{vat, niin olisi niissä paljonkin lähdettä eurooppalaisenkin silmän sietämälle mielikuvituksellisemmin sommitellulle kirkkorakennukselle. Onhan suomenpuolisen Karjalan maaseuduilla jo suomalaismallisempiakin, eurooppalaistyylisempiä kirkkoja — kuten esim. Suojärven Eräs venäläinen Suslov on julkaissut mielenkiintoisen venäläisen julkaisun » Venäläinen rakennustaide kansantarinain muistotiedon perusteella » ( Pietarissa 1911 ).}
\kat{Toukomies}{01.1926}{P9}{Joka tapauksessa tämä nykyisen Suomen Karjalan alue — ennenhän se on ollut suurempi, työntyen etenkin luoteeseen ja mahdollisesti suoraan pohjoiseenkin niin, että vielä oululaisten juonteita on johdettu karjalaisista — on etenkin kielellisiin ja myös sivistyksellisiin seikkoihin katsoen siinä määrin pieniä vivahduseroavaisuuksia täynnä, että se Onpa jo eroa Sortavalan ja Salmin rantaseutujenkin kansan kesken, vaikka ne elävät niin lähekkäin ja viimemainittu taas puolestaan eroaa Suistamon, Suojärven ja muista sisämaan rajakarjalaisista, joista kaikista taas verrattain lähellä elävät Kiteen, Tohmajärven ja Ilomantsinkin asukkaat ovat jo täysiä » ruotsheja » heihin nähden, Mutta Runolaulaja dots}
\kat{Toukomies}{01.1926}{P21}{[ / tm rata Karjalan rajaa kohti. Läskelään asti on jo valmiina ratakappale Suojärven radalta. Komiteaan, jonka ensi kokoukseen edellämainittujen}
\kat{Toukomies}{01.1926}{P21}{Lähiaikojen toimenpide luonnollisesti tulee olemaan, että uutta rataa Pitkästärannasta jatketaan Salmiin asti. Senlisäksi samassa täysistunnossa myönnettiin varat Suojärven radan jatkamiseksi aina rajalle asti. Suomeen jääneistä pakolaisista katsotaan suurimman osan olevan sellaisia, joiden ei sallita ollenkaan palata Karjalaan.}
\kat{Toukomies}{03.1926}{P6}{Se on käännöskokoelma, jonka on tehnyt kirjailija livo Härkönen, puhtaasti karjalaiselta pohjalta noussut, toistaiseksi ainoa persoonallisuus, joka pyrkii tällä alalla korkeampiin kirjallisiin päämääriin. Vaan kaarrellen toisaalta Suojärven saloja myöten ja tyyntyen Sortavalan takanakin vielä tavalliseksi mäkimaaksi tämä pitkän, yli Karjalan ulottuvan kaarteen pohjoispää alkaa vähitellen taas nousta, ensin yläväksi sisämaaksi. Kun kääntäjän runollinen vaisto on puhdistanut ja täydellistänyt ilmaisua, näemme, etteivät nämä erot ole pelkkiä sanastollisia ja äänteellisiä eroavaisuuksia, vaan että niihin sisältyy erilaisia tunnelaatuja ja uusia värisävyjä.}
\kat{Toukomies}{03.1926}{P6}{Ensiksikin tahtoisin omalta kannaltani vastata erääseen kysymykseen: Sisältyykö karjalankielisiin, jo ilmestyneihin kirjallisuusnäytteihin joitakin erikoisia arvoja, joiden säilyminen olisi toivottava joltakin yleisemmältä, esim. puhtaasti kirjalliselta, esteettiseltä kannalta? Korpiselän kirkonkylä, Lehmivaara, Tolvajärven vaarat, Honkavaara, 190 — 208. — Suojärvellä eräs vaara 219, Suojärven ja Ilomantsin välillä, juuri Kuopion läänin rajalla, pari paikkaa 272 — 304 m. ) Täällä jo jotkut järvetkin lainehtivat 175 metrin korkeudessa, kuten vedenjakajalla olevat Suuret vaarat muodostuvat pitkiksi jonoiksi, jotka muistuttavat pieniä maanselkiä — semmoisia Ilomantsissa, Kiihtelysvaarassa, dots}
\subsubsection*{Mäntyselkä \normalfont\small(13 katkelmaa, 13 osumaa aineistossa)}
\kat{Toukomies}{06.1926}{P12}{Sinä et voinut palata rakastamillesi seuduille, vaikka niin hartaasti sitä toivoit, uusi järjestelmä ei herättänyt sinussa luottamusta. Suojärven radan päätekohtaa ja pohjoisin Kontiovaarasta, Äänis järven pohjoisimmasta perukasta, Mäntyselän ja Juustjärven kautta Poraj arvelle. Saamiemme varmojen tietojen mukaan on Venäjä päättänyt rakentaa kolme kapearaiteista rautatietä Muurmanin radalta länteen päin, Suomen rajan lähelle.}
\kat{Toukomies}{01.09.1927}{P13}{rautatie Mäntyselkä Njuhotskaja}
\kat{Suomen Heimo}{20.12.1931}{P26}{Sitäpaitsi uskonnollisestikin Karjala oli Venäjän vallankumoukseen saakka, ortodoksisen kirkon välityksellä, kiinteässä riippuvaisuussuhteessa mainittuun esivaltaan. Poventsan kihlakuntaan kuuluivat seuraavat' 10 kuntaa: Repola, Porajärvi, Rukajärvi, Lintajärvi, Mäntyselkä, Paatene, Aunuksen kihlakuntaan kuuluivat seuraavat 9 kuntaa: Tulemajärvi, Vieljärvi, Vitele, Kotkajärvi, Keskotjärvi, Riipuskala, Mätusova, Vaasheni ja Njekkula.}
\kat{Toukomies}{01.02.1931}{P18}{Koska Valtion Pakolaisavustuskeskus on kuluvaksi vuodeksi kieltäytynyt myöntämästä rahamäärää postimenojen peittämiseksi ja koska kansliaa hoitava papisto on lain mukaan oikeutettu perimään lunastukset todistuksista henkilöiltä, jotka postitse tilaavat papintodistukset tai muut asiakirjat, huomautetaan niille pakolaisille, jotka tulevat hakemaan postitse tarpeellisia todistuksia, että todistus ilman merkintää aviosuhteista maksua Smk. Majalahden kylästä ja palvelijatar Anni Mari Pekantytär Kiiskinen lisalmen ev. lut. seurakunnasta; 6 ) Heikki Kononinpoika Var hikka Aunuksen Mäntyselästä ja palvelijatar Anna Tarsina Inkerin pakolaisseurakunnasta; 7 ) Leski Timo Mikonpoika Levoskin Repolan dots}
\kat{Toukomies}{01.12.1934}{P18}{Kun alueen turvana oli pieni Pohjois-Aunuksen suojeluskunta, jonka muodostivat rohkeimmat miehet Repolan ja Porajärven pitäjistä päällikkönään Iyähes vuotta myöhemmin, kesäkuussa 1919 seurasivat esimerkkiä Porajärvi ja kohta sen jälkeen LÄndajärvi, s. o. Suomen rajalla oleva osa Mäntyselän pitäjää. Tämä miestemme sankariteko se oli, joka antoi voiman ja rohkeuden koko isänmaalliselle osalle kansastamme ryhtyä vapauttamaan itseään vieraan ikeestä.}
\kat{Karjalan Vapaus}{05.1935}{P17}{Kauppias Ilja Zertijeff, Kauppias F. Ojamo on syntynyt Mäntyselän pit. 1914 — 1919 hän hoiti postinkuljettajan virkaa.}
\kat{Karjalan Vapaus}{05.1935}{P19}{Liikemies Pekka Jukola on synt. 29 / 1 1898 Mäntyselän pit. 1921 — 1922.}
\kat{Karjalan Vapaus}{05.1935}{P19}{Kun 1932 hänen toimesta perustettiin T:ni M. Pedas niminen liike, tuli hän sen johtajaksi. Toiminimi M. Pedaksen johtaja J. Dorojejeff on syntynyt Mäntyselän Karhumäellä 24 / 10 -94 mv. Joensuussa kangastavaroilla tukku- ja vähittäiskauppaa harjoittava kauppias S. Saharoff on syntynyt Rukajärvellä helmik.}
\kat{Suomen Heimo}{31.10.1936}{P17}{Tarkoituksenmukaista sensijaan ei olisi suunnitelman supistaminen senkään vuoksi, että karjalan kieli, jonka kohtalo näyttää tällä hetkellä sangen uhatulta, on vanhojen ominaisuuksiensa vuoksi suomalais-ugrilaiselle tutkimukselle erittäin tärkeä, jonka vuoksi mahdollisimman sisältörikkaan ja laajan sanakirjan aikaansaaminen olisi aivan välttämätön. II = Ilomantsi Pa = Paatene Säj = Säämäjärvi Im = Impilahti Phj = Pyhäjärvi Tj = Tulemajärvi J = Jyskyjärvi Pnj = Paariajärvi Tk =. Tunkua Ks = Korpiselkä Prj = Porajärvi U = Uhtua Kst = Kiestinki Ptj = Pistojärvi Usm = Usmana K t = Kontokki R = Repola V = Vitele Ktj = Kotkat järvi R. = Riipuskala Vj = Viel järvi Ktp = Kontupohja Rj = Rukajärvi dots}
\kat{Karjalaisuuden Ystävä}{12.1938}{P31}{Samasta. syystä ei eri sanojen kohdalla osoiteta niiden taivutusta, vaan taivutusselityikset annetaan loppuun liitettävässä tyyppitaulukossa. Pohjoisin aunukselaismunre on Säämäjärven imurre; sen pohjoispuolella oleva Mäntyselän murre on jo n. s. varsinaista karjalan 'Hakusanaksi on otettu Suojärven murteen äänneasu senvuoksi, että se on äänneominaisuuksiltaan ikäänkuin välimuoto Aunuksen ja pohjoisempien murteiden välillä.}
\kat{Itä-Karjala}{23.03.1939}{P3}{Näistä kirjoista on muodostettu Väinölän kirjasto, jo • ta oppilaat käyttävät. VI „ Koivukari, Paul, Mäntyselkä, Koikari, ( Helsinki )... Päätös kuitenkin tehtiin ja ryhdyttiin suorittamaan tarvittavia järjestelyjä.}
\kat{Itä-Karjala}{20.06.1939}{P4}{Tämä rahasto tulee palvelemaan opintomatkoja ja retkeilyjä, kuten nimikin jo osoittaa. VI „ Koivukari, Paul, Mäntyselkä, Koikari, ( Helsinki ).... Näillä edellytyksillä katsottiin voitavan ottaa 15 pakolaislasta huollettavaksi mainittuun koulukotiin.}
\kat{Karjalaisuuden Ystävä}{10.1939}{P4}{Merkkien selitys: ) — (. ) — (. ) 'Saiomen ja Venäjän välinen raja; *. —. —. —. kunnan ( = pitäjän ) raja; o o o o o o ikarjalais-lyydiläis-en kielialueen ja venäjän kielen 'välinen raja; o — o — o — karjalan kielen ja suomen murteiden välinen summittainen raja; i • • \# karjalais- ja lyytiläismurteiden summittainen raja; xxx x x x karjalan kielen aunukselaismurteiden ja varsinaisten karjalaismurteiden raja. Il — Ilomant6i R — Repola lm — Impilahti Rj — Riipuskala J — Jyskyjärvi Rj — Ruka järvi Ks — Korpiselkä S — Salmi Kst — Kiestinlki Skj — Suikujärvi Kt — Kontokki Sm — Suistamo Ktj — Kotkat järvi Suj — Suo järvi Ktp — Kontupohja Säj — Säämäjärvi Mj — Mund järvi Tj — Tulemajärvi Ms — dots}
\subsubsection*{Säämäjärvi \normalfont\small(25 katkelmaa, 42 osumaa aineistossa)}
\kat{Toukomies}{05.1926}{P13}{Tarvitaan pitkä-aikaista, ennakkoluulotonta kansallista kasvatusta ja herätystyötä ennenkuin noiden seutujen asukkaat perinpohjin, juurta jaksain itsensä, » löytävät », tuntevat suomalaisiksi. Säämäjärven ja Pyhäniemen kylistä kotoisin olevat pakolaiset selittävät, että jo ennenkin he ovat käyneet Suoj arvella, mutta eivät koskaan Suomessa saakka — tarkoittaen Korpiselkää. Korpiselällä, Ilomantsissa, Pielis järvellä ei kansa enää pidä itseänsä venäläisenä, mutta käsitteet suomalaisesta kansallisuudesta ovat näil akin tienoilla vielä varsin sekavat ja heikot.}
\kat{Suomen Heimo}{25.05.1929}{P7}{Kevättalvi v. 1919 oli ollut erinomaisen leuto ja Ensimmäiset kelirikon oireet alkoivat ilmestyä jo neuvottelujen aikana maaliskuun puolivälissä. Orusjärven ryhmä hiihtäisi Tulema- ja Videljärven kautta Rihmakylään ( Prjäshaan ), yhtyen siellä Säämäjärven kautta hiihtävään Suojärven ryhmään. 5 miljoonaa kiväärin ja kk:n panosta, 600 shrapnellia, 200 granaattia, 11,000 käsigranaattia y.m. Yhteysvälineetkin olivat hyvin täydellisiä.}
\kat{Suomen Heimo}{25.05.1929}{P9}{Pienempiä joukkoosastoja eteni vielä Aleksanteri Svirin luostarille ja Syvärin välittömään läheisyyteen. Yhtämittaa taistellen sitten edettiin Säämäjärven kautta Prjäshaan, jossa -jo II ryhmä sijaitsi. 2 " p! nä luovuttava Aunuksesta, mutta kaupunki vallattiin taas, ja. toisen kerran siitä luovuttiin toukok.}
\kat{Suomen Heimo}{20.06.1929}{P8}{Huhtikuun 6 päivänä saapuivat ensimmäiset 60 miestä, ja tämän jälkeen päivittäin 100 — 200 miestä, niin että 11. 4. Sortavalassa oli vasta n. 600 miestä. Etenemistä Suo järveltä Säämäjärven yli Rihmakylään oli pidettävä ainoastaan apuliikkeenä, jonka tarkoituksena oli suojata pohjoisen ryhmän vasenta sivustaa yllätyksiltä. Vihollinen, joka olisi uskaltanut jättäytyä molempien etenevien ryhmien väliin, olisi ehdottomasti joutunut saarroksiin, joten vastarintaa sen puolelta ei ollut odotettavissa operatioteiden välisellä alueella.}
\kat{Suomen Heimo}{20.12.1931}{P26}{Edellä esitettyjen kihlakuntien hallinnollisten rajojen mukaan Kauko-Karjala, Kuollaa lukuunottamatta, käsitti 136,300 km 2:n pinta-alan, josta Vienan-Karjalan osalle tuli 58,000 km 2 ja Aunuksen Karjalan osalle 78,300 km 2. Petroskoin kihlakunnassa oli 12 kuntaa: Ostretshini, Soutjärvi ( asuu melkein yksinomaan vepsäläisiä ), Ladva, Pyhäjärvi, Suoju, Säämäjärvi, Munjärvi, Kontupohja, Tepenitsa, Jalguba, Suurlahti ja Talvoja. Siitä lähtien on Kauko-Karjala, joka maantieteellisesti, luonnonhistoriallisesti ja kansallisesti on osa Suomea, kuulunut ja kuuluu edelleen valtiollisesti Venäjän yhteyteen ja sen taloudelliseen, sivistykselliseen ja poliittiseen vaikutuspiiriin.}
\kat{Karjalan Vapaus}{09.1934}{P4}{Siellä yhtyvät Länsi unuksesta tulevat useammat tiet toisiinsa ja jatkuvat sitten edelleen itää kohti yhtenä yllämainittuna kulkuväylänä halki suuren erämaan. Näiden alueiden välimaasto, suunnilleen linjalta Aunus — Prjäsha — Säämäjärvi itään, on suunnatonta, miltei ihmisjalan ja -käden koskematonta erämaata. Koko rajantakainen Itä-Karjala on suurin piirtein katsoen jylhää metsä- ja suoerämaata, jota vain vesistöt sekä hyvin harvalukuiset, huonot tiet ja polut halkovat.}
\kat{Karjalan Vapaus}{09.1934}{P4}{Tämän lisäksi kulkee tuon 150 km leveän ja 50 — 100 km syvän koskemattoman erämaavyöhykkeen lävitse vain pari tuskin ratsumiehen kuljettavaa kapulasiltapolkua. Syvärijoen ja Säämäjärven väliseltä n. 150 km:n levyiseltä alueelta johtaa itään Äänisjärven rannalle vain yksi käyttökelpoinen ajotie, Prjäsha — Poloviina — Suolusmäki — Petroskoi. Heikkisen johtamassa komppaniassa joukkueen johtajana, oli tilaisuudessa — jouduttuaan heti taistelun jälkeen kesäkuun alkupuolelle saakka suorittamaan etuvartiopalvelusta samaisessa Matrosan kylässä}
\kat{Karjalan Vapaus}{09.1934}{P5}{Huhtikuun 21 ja 28 päivien välisenä aikana oli hän tällä pataljoonalla edennyt, useita kiivaita taisteluita käyden, Orus järven kohdalta rajalta Prjäshan edustalle Riihiselkään saakka, lähes 100 km:n pituisen huonoteisen taipaleen. Samanaikaisesti oli Pohjoisen ryhmän toinen pataljoona ( retkikunnan IV patl. ) edennyt luutnantti Maskulan johdossa Suojärveltä Veskelyksen ja Säämäjärven kautta Kindasovan kylään, minne se oli saapunut huhtikuun 28 p:n illalla. Olihan majuri Talvelalla syytä suurella todennäköisyydellä otaksua, ettei häntä vastassa Prjäshassa vielä 29 päivänä olisi muuta kuin aikaisemmissa taisteluissa lyötyjä ja moraalisesti jo järkytettyjä vihollisvoimia, jotka tuskin enää dots}
\kat{Karjalan Vapaus}{09.1934}{P5}{Vaikka majuri Talvela sai Riihiselkään ilmoituksen luutnantti Maskulalta Kindasovasta vielä samana iltana, että viimeksimainittu voisi pataljoonallaan ottaa osaa hyökkäykseen Prjäshaa vastaan vasta 30 päivänä ja että Niiniselän kylä oli todettu vihollisen lujasti miehittämäksi, niin ei hän kuitenkaan katsonut voivansa lykätä hyökkäystään Prjäshaan edes yhdellä vuorokaudella. Aunukseen keväällä 1919 Vieljärven ja Säämäjärven kautta edenneen » Pohjoisen ryhmän » ( 2 patl. ) ensimmäisenä operatiivisena tavoitteena oli mahdollisimman nopean hyökkäysliikkeen avulla saada haltuunsa Prjäshan tärkeä tienristeys ennenkuin vihollinen ennättäisi keskittää taempaa suurempia voimiaan Prjäshan dots}
\kat{Toukomies}{01.06.1934}{P4}{Huomattava on myös: verkkoriäshä = verkkorääsy Rihmakylä ei Prääsä missään tapauksessa ole; prääsällä ei ole mitään tekemistä rihman kera. Säämj arven — Säämäj arven niinen sanovat itse sääm Järveläiset muodossa Seämärvi — Siämärvi, ollen e ainakin paikoitellen e:n ja i:n väliltä, mieluummin e:tä lähempänä. Kuinka helpoittavaa on pitkän taipaleen jälkeen istahtaa tien varteen, jossa kirkas kiven pinta lyhyin kirjoituksin puhuu esi-isäin työstä ja samalla viittaa tulevaisuuteen.}
\kat{Karjalan Vapaus}{05.1935}{P9}{Kun syksyllä ruvettiin keräämään pakolla puna-armeijaan, oli hän kansalaiskokouksia järjestämässä. Kauppias Vasili Molosofkin on syntynyt Aunuksessa Säämäjärven pitäjän Veshkelyksen kylässä maanvilj. Kauppias V. Molosofkin säännöllisellä elämällään ja rehellisyydellään on esimerkiksi kelpaava monelle.}
\kat{Suomen Heimo}{15.03.1935}{P14}{Kashgarin hallituksen kukistuessa hänen onnistui paeta. Kaatunut Säämäjärven talonpoika oli vapausliikkeen ensimmäisiä. Armeija varustettiin nyt mitä täydellisimniin.}
\kat{Suomen Heimo}{15.03.1935}{P14}{Kun nuo onnettomat uhrit olivat kuolleet, hautasivat bolshevikit uhrinsa suohon. Säämäjärven kunnan asukkaat ovat Keski - Aunuksen valistuneinta väkeä ja henkisesti hyvin valveutuneita. Dungaanien ilmoitetaan olevan myöskin hyvin aseistettuja ja saavan tukea englantilaisilta.}
\kat{Suomen Heimo}{15.03.1935}{P14}{Kersantti Helge Halminen oli ollut mukana Viron retkellä, oli pelkäämätön, uljas nuorukainen ja paloi sankari-intoa. Senpä vuoksi komppanian päällikkö lähetti kersantti Helge Halmisen ja erään komppaniansa säämäjärveläisen sotamiehen Säämäjärven kunnan Harjun kylään. Näitten kahden keskuksen väliset neuvottelutkaan eivät ole johtaneet tuloksiin, minkä vuoksi uusi sisällissota on taas ovella.}
\kat{Suomen Heimo}{15.03.1935}{P14}{Kansallisen hallituksen kukistumisesta hän kertoi, että Neuvostoliitto ryhtyi kohta kansallisen hallituksen muodostamisen jälkeen toimenpiteisiin sen kukistamiseksi. Esimerkkinä mainittakoon tästä väkirikkaat Säämäjärven ja Vieljärven kuntain kuntakokoukset ja kuntain valitsemat lähetystöt Suomen hallituksen puheille vuoden vaihteessa 1918. Miten bolshevikit onnistuivat suunnitelmissaan, sen osoittavat Venäjän Turkestanin kansallisen liikkeen lehden " Jash Turkestanin " tilanteesta viime aikoina saamat tiedot.}
\kat{Suomen Heimo}{01.04.1935}{P11}{Veljiänsä, meitä, siunaa vcljesrakkaudella vain! KUKKIA JA KUOLEMAA Karjalaisten tykistöä Uhmoilassa Säämäjärven taistelussa. Armoas Jumala viel' et vaihtaa voi siis lastesi tuoiniohon.}
\kat{Suomen Heimo}{15.04.1935}{P17}{Alueeseen kuuluu myös koko Laatokan rannikko yli 10 km leveydeltä. Näistä Säämäjärven, Matroosan ja Kontupohjan huomattavasti muita suuremmat. ( SH7 / 35. ) 167. Neuvosto- Venäjän sisäasiainkomissariaatti ( ent.}
\kat{Suomen Heimo}{15.04.1935}{P17}{1:10 39:60 2. metsäteollisuuspiirien varusteluosastojen ruokaloissa ja kattilakunnissa:. a. vihannes- ja. ryynipäivällinen o: 65 23:40 b. liha- tai kalapäivällinen 1:05 37: 80 3. kouluaamiaiset 0:30 10:80 Leipää tarjoillaan eri maksusta. ( SH 7 / 35. ) 95. Maaliskuussa 1935 oli seuraavat traktoriasemat eli baa'sit: Olympian, Säämäjärven, Internatsionalen, Matroosan, Vilkan, Koverin ja seuraavat kuorma-autoasemat: Kontupohjan, Motorinskajan, Lososinan ja Vilkan- Suojun. Kolhoosien luku oli 900. ( SH 7 / 35. ) 83. Kollektivisoidun maatalouden luonnetta valaisee se, että Neuvosto-Karjalassa keväällä 1935 oli 648 " erikoisvalmennuksen omaavaa kölhoösituotannon komentajaa ", kaikesta päättäen dots}
\kat{Suomen Heimo}{20.07.1935}{P4}{Kauppoja on perustettu tiheään, mutta harvoinpa niissä on muuta kuin viinaa ja makeisia myytävänä. — Samaan aikaan kuin Petroskoissa valmistauduttiin juhlaan, täyttyi Säämäjärven kirkonkylä toisenlaisesta juhlarahvaasta. Mutta sittenkin karjalakerhoissa on ollut jotakin, joka on tehnyt ne toisenlaisiksi kuin tavalliset nuorisoseurat.}
\kat{Suomen Heimo}{20.07.1935}{P4}{Tuhannet viettivät sitä jo karkoitettuina, tuhannet nääntyen 15-vuotisjuhlan kunniaksi määrätyissä ylimääräisissä " vapaaehtoisissa " rynnäköissä pakkotyömailla. Säämäjärven kirkonkylästä oli tehty keskitysleiri, johon rajaseudun asukkaita koottiin, sitten suurin joukoin rautatielle ja edelleen tuntemattomiin päämääriin kuljetettavaksi. Kuvaavaa sille, mitä huono omatunto voi teettää, on, että hänen puheensa tätä kohtaa oli vaikeata radiossa kuulla ilmeisesti tahallisen häiriön takia.}
\kat{Suomen Heimo}{15.10.1935}{P7}{Ja ennenkuin Karjalan sotien kunnollinen historia voidaan kirjoittaa, on tällaisia niin sanoakseni monografioita oltava ilmestyneinä tai ainakin kirjoitettuina huomattavasti enemmän, kuin niitä vielä toistaiseksi on. — Lopuksi tekijä vielä kuvaa joukkonsa osallistumista Aunuksen eteläisen rintaman tapahtumiin aina Vitelen katastrofiin asti sekä viimeksi komppaniansa — joksi komennuskunta oli muutettu — vaiheita Säämäjärven rintamalla. — Kokonaisuutena Marttinan kirja antaa erittäin kauniin kuvan siitä uhrimielestä ja auttajahengestä, jonka valtaamina niin monet nuorukaisemme vuosina 1918 — 21 läksivät taisteluun heimomme perivihollista vastaan.}
\kat{Suomen Heimo}{15.10.1935}{P12}{Rajavyöhykkeeseen sisältyvät varsinkin Aunuksessa juuri tiheimmin asutut karjalaisseudut. Säämäjärven kyläneuvoston alueelta oli 25 talon asukkaat määrätty karkoitettaviksi. Minulle kerrottiin, että vangitut talonpojat oli määrätty kar koitettaviksi.}
\kat{Suomen Heimo}{15.10.1935}{P12}{Petroskoissa veimme vangit asemalle ja puhuttiin, että heidät lähetetään Siperiaan. Säämäjärven kirkonkylässä oli kirkko ja kaksi taloa täynnä pidätettyjä, kaikkiaan noin 400 henkilöä. Huhtikuun alkupäivinä 1935 olin 8 kuorma-autoa käsittävän kolonnan mukana ruokatavaroiden kuljetuksessa Palalahteen.}
\kat{Suomen Heimo}{15.10.1935}{P12}{Vähitellen tuli rajan yli salateitse muutamia pakolaisia, jotka tiesivät kertoa kokonaisten kylien tyhjentämisestä kauempanakin rajasta sekä suurista keskitysleireistä Säämä järven kirkolla, Porajärvellä, Mundjärven Lahdessa, Palalahdella, Nuosjärvellä ja Aunuksen kaupungissa. Samoihin aikoihin olivat NKVD:n joukot keränneet väestöä Korsan, Ruptsoilan, Jessoilan, Uhmoilan, Kiskoilan, Ängellahden ym. kylistä kuljettaen niistä vanhukset, naiset ja lapset Säämäjärven kirkonkylään ja työkykyiset miehet Uhmoilan kylään, missä heitä pidettiin tarkan valvonnan alaisina. Sotilaiden toimesta oli karkoitetut kylittäin koottu 20 — 30 henkilöä käsittäviksi joukkueiksi, kuuluen joukkoon lisäksi noin 10 dots}
\kat{Suomen Heimo}{15.10.1935}{P13}{Palalahden keskitysleiriltä viety parina päivänä huhtikuun alussa 1935 160 henkeä pois.. Säämäjärvi ja ympäristö: tammi — helmikuussa 193 1 karkoitettu 85 talollista, toukok. Kannanlahden piiri: tietoja puuttuu, mutta jotkut asutukset ovat tyhjinä.}
\subsubsection*{Jukola \normalfont\small(40 katkelmaa, 56 osumaa aineistossa)}
\kat{Suomen Heimo}{15.08.1929}{P24}{Kun maalaisliitto nyt ilmeisesti, pian joutuu hallitusta muodostamaan, " toivovat- kaikki suomalaiskansalliset piirit tässä maassa, että se muodostaa hallituksen, joka on sanoissa ja teoissa suomalainen, kuten on suomalainen se laaja Väestö, Joka on puolueen edustajille äänensä antanut. Nyt edistyspuolueen eduskuntaryhmä rynnistää seitsemän miehen ( tai oikeammin kuuden miehen ja Mandi Hannulan ) voimalla kuin Jukolan veljekset, Itsenäisyyden Liiton kannalta tämän tuloksen luulisi olevan erittäin valitettavan, kun liitto teki niin ankarasti agitatiotyötä juuri edistyspuolueen hyväksi, puolueen, jonka ainoana oikeana vapaamielisenä, kansalaiskiihkottomana puolueena olisi liiton mielestä dots}
\kat{Toukomies}{01.05.1929}{P11}{Hänelle viritteli koko kotoinen luonto, salojen puut sanelivat: Sinulle yksin ylenevät latvojemme kerkät ja puutuu runko j emme jälsi punaiseksi hongaksi. Ellet taas ole niiden tarpeessa, vaan kaipaat mielenkevennystä ilottomiin iltoihisi, pane tarinoimaan Jukolan pojat ja muut saman tapaiset aikasi kuluksi ja muiden mielen ratoksi. Laatokka » ja » Peltomies » riippuivat nauloissaan hänen pirttinsä seinällä, ja mikä niissä oli Miitreille ja hänen ympäristölleen tarpeellista, sen hänen valpas silmänsä muistiin kammitsoi.}
\kat{Suomen Heimo}{15.06.1930}{P5}{On laiminlyöty sukulaiset ja eletty vieraan seurassa. Tämä on Jukolan veljesten ja Toukolaisten sovinnon juhla. Näin voi meidän ystävyydestämme rakentua todella lujaa ja kestävää.}
\kat{Suomen Heimo}{15.06.1930}{P12}{Siinä on eräänlaisen kirjallisen toiminnan palkka: vanha, kunnianarvoisaan ikään ehtinyt suomalainen puolueemme alkaa lukea riveihinsä vain köyryselkäukkoja. Ystäväni ja koulutoverini Martti Jukola, tuo tunnettu urheilija ja urheilukirjailija, tunnettu myös runouden kultakaupungissa nimellä Ilmari Marski, on väitellyt toukokuun 26 pnä. Mainittakoon myöskin, että suuresti kunnioittamani rouvasihminen Serp, Uuden Suomen Sunnuntailiitteen tarmokas toimittaja, on kääntänyt merkillisen sotaromaanin, Sherriffin " Matkan pään ".}
\kat{Toukomies}{01.05.1930}{P17}{Suomen Kuvalehden tilaajille on erikoisetu tämän kirjan hankkimisessa. Teokseen liittyy alkulauseena maisteri Martti Jukolan kirjoittama J. H. Erkon elämäkeita. Otava julkaisi pari vuotta sitten Aleksis Kiven Valitut teokset yhtenä kookkaana niteenä.}
\kat{Toukomies}{01.06.1930}{P22}{Näitä on seuraavissa lajeissa: ■ — Martti Jukola: Juhana Heikki Erkko. Runsaasti kuvitettu teos maksaa 50 mk, sid.}
\kat{Suomen Heimo}{20.12.1931}{P38}{N:o 11 — 12 184 Jukola, Martti: Me uskomme urheiluun. Toppila, Heikki: Päästä meitä pahasta.}
\kat{Suomen Heimo}{30.09.1932}{P12}{'' j-\textasciicircum{}\textasciicircum{}\textasciicircum{}\textasciicircum{}\textasciicircum{}\textasciicircum{} S\textasciicircum{}\textasciicircum{}\textasciicircum{}. Tutustuessanne Linguaphone-kurssiin huomaatte Tekin BnnuJ « >\textasciicircum{}\textasciicircum{}f V iP\textasciicircum{}\textasciicircum{} ett\textasciicircum{} Linguaphone-opiskelun tuloksiin täytyy olla tyyty ff J\textasciicircum{}:\textasciicircum{} — - " " \textasciicircum{} Suomen Joutsenen miehis- HJfl j| / ' / \textasciicircum{}\textasciicircum{}\textasciicircum{}\textasciicircum{} = \textasciicircum{}rl\textasciicircum{}\textasciicircum{}\textasciicircum{}\textasciicircum{}\textasciicircum{} / -li kopurjehdukselle, saa tällekin matkalle mukaan eng- II lanninkielen Linguaphone-kursseja, jotka ylvään koulu I\textasciicircum{}\textasciicircum{}l lU\textasciicircum{}tfÉ H f / J\textasciicircum{}J f\textasciicircum{}\textasciicircum{}-Z\textasciicircum{}5 e3> 7\textasciicircum{}W\#\textasciicircum{}l kivamme edellisellä retkellä jo ovat perehdyttäneet ffPPfflLil\textasciicircum{}Éi\textasciicircum{}H v2\textasciicircum{}iijl ' * i B\textasciicircum{}W\textasciicircum{} ff\textasciicircum{}\textasciicircum{}mfJl meripojat englannintaidon salaisuuksiin. Tohtori Martti Jukola kertoo, etta pojat suuresti I\textasciicircum{}HPrniUlfl PJI \textasciicircum{}t Bffl AfflßAJlP Alm AfflMf li innostuivatkin helppoon, viehättävään ja tehokkaaseen <3\textasciicircum{} X WtofåWwßW x y\textasciicircum{}\textasciicircum{}mV®w n dots}
\kat{Suomen Heimo}{06.11.1932}{P18}{Kososen Piha on mukana samantapaisessa. Martti Jukola: Me uskomme urheiluun. Mutta sittenkin kirja vaikuttaa epä valmiilta.}
\kat{Suomen Heimo}{06.11.1932}{P18}{— Syvemmin ja vakuuttavammin on kuvattu toimittaja Lami. Martti Jukola on urheilukirjaili jäimme ensi rivin miehiä. Hän näki asioiden läpi kuin olisivat ne olleet lasia.}
\kat{Suomen Heimo}{06.11.1932}{P18}{Näin hänen esseensä saavuttavat yleispätevyyttä ja muodostuvat tavallaan jonkinlaisiksi urheilututkielmiksi. Pyrähdyksiä kentällä ja maastossa " on Jukolan uusin, Suomen olympialaispojille omistettu teos. Ja on lisättävä, että ne samalla ovat suurelta osaltaan taiteellisestikin merkittäviä saavutuksia.}
\kat{Suomen Heimo}{06.11.1932}{P19}{Sitten kuvataan oikeaan urheiluun liittyviä kieltämättömiä henkisiä arvoja. Tässä ei ole tilaisuus lähemmin esitellä Jukolan mielenkiintoista kirjaa. O. Y. HERMAN LINDELL P. ESPLANAADINKATU 37 • HELSINKI}
\kat{Suomen Heimo}{06.11.1932}{P19}{Päähenkilössä on paljon yhteistä Soinin tuttujen naistyyppien kanssa: " Unen " tarmokasta keski-ikäistä opiskelijanaista, " Jumalten suosikkien " Ameriikassa matkustanutta itsenäistä virkanaista. Martti Jukolan merkitys urheilukirjani jana onkin mielestämme siinä, että hän on pyrkinyt aina paljastamaan niitä sisäisempiä kulttuuriarvoja, joiden tukena oikeassa mielessä suoritettu urheilutyö on. Se on vaimon selvitys eronneelle miehelleen, itsetunnustus ja kertomus suuresta rakkaudesta, joka 20- vuotisen avioliiton jälkeen oli heitetty pois, jota ei enää huolittu.}
\kat{Suomen Heimo}{17.12.1932}{P24}{Levyt on valmistettu myöskin parhaista aineista, niin että ne ovat siksi kestäviä, että niitä voi " veivata " mielin määrin. Me emme ehdi Jukolan veljesten vanhalla silmäpuoli Valkolla junalle 23.30, vaan pitää ottaa auto His Master Six, kahdeksan sylinteriä, malli 1934. Kaikista onnellisinta kielenopinnoissa on, jos kohta alusta omaksuu hyvän ääntämisen, mikä voi käydä laatuun tarkoin seuraten Linguaphone-levyistä kuultavaa puhetta.}
\kat{Suomen Heimo}{19.12.1934}{P24}{Koskenniemen mielestä ei voi kyllin ihailla sen rakenteellista mestariutta, selvänäköistä psykologiaa, loistavia situatioita ja näkemyksen lujaa varmuutta. Lieneekin totta, että tämä kyläkraatarin onnettoman kihlauksen kuvaus säilyy kuolemattomana Eskon ja Jukolan veljesten hahmojen rinnalla. Hänen syvin runomaailmansa ei ole hyvän ja pahan dualistinen maailma, se on ikuisen hyvyyden maailma lähellä luontoa, elämän suurta alkukotia.}
\kat{Toukomies}{01.05.1934}{P12}{Lopuksi keskusteltiin juoksevista asioista, julkaisutoiminnasta ja painatustöiden painattamisesta. Juho Mustonen ja pöytäkirjaa piti sihteeri kaupanhoitaja Paavo Jukola. Molemmat osat maksavat 20 mk. nid., 28 mk. sid.}
\kat{Toukomies}{01.12.1934}{P40}{( 5 kuvaa. ) 97 Karjalan Sivistysseuran ja sen edeltäjän Vienan Karjalaisten Liiton johtokunnan jäsenet ja toimihenkilöt. ( J. Maksola ); Pitkänrannan Karjalakerhon vuosikokous 41 Paavo \{ Jukola ): Neuvosto-Karjalan kuulumisia 42. Kirjallisuutta. Näytenro 4: Sivu Aari Surakka: Laulu rajalta, runo 15 Vallas H omanen: Matkahavaintoja kerhojen piiristä....}
\kat{Toukomies}{01.12.1934}{P40}{15 Kerhojen työmaalta: Helsingin Karjalakerhon ( V. Homanen ); Martinniemen Karjalakerhon piiristä ( J. Maksimainen ); Lieksan Karjalakerho. Paavo \{ Jukola ): Suojärven kuulumisia 25 Heimotyömaalta: Karjalakerhojen Keskusliitto ( V. Homanen ); Kuollut karjalainen ja Karjalan ystävä. S. -seuran adressi 58 Runolaulajat tekevät kunniaa satavuotiaalle ( 3 kuvaa ) 58 Ilmari Kianto: Viluvirsi Wienan Karjalalle, runo....}
\kat{Toukomies}{01.12.1934}{P41}{( Kansik. ) 137 Kosti Uuksu ( K. Toivonen ): Lintujen lähtiessä, runo 141 livo Härkönen: Kuulkaa Karjalan itku! puhe 142 Kemin karjalaisjuhlat 143 W. Keynäs: Kalevalan pohjalta, puhe 144 ( livo Härkönen ): Runolaulajajuhla Ilomantsissa. ( Kuva. ) 2 Paavo ( Jukola ): Neuvosto-Karjalan oloja 3 V. Homanen: Kerhotoiminnan tarkoituksenmukaistuttaminen 4 Pekka Metsäsissi: Veitsiluodon kerhon tervehdys heimoveljille ja kerhotovereille, runo 4 V ( allas ) H ( omanen ): Mietteitä elämän vakavuudesta, osia puheesta. ( 2 kuv. ) 154 livo Härkönen: Raja-Karjalan ja Aunuksen vanhasta kansanrunoudesta 155 Vaassila Karjalaisten osuus neljännessä suoma ( W. Keynäs ): laisugrilaisessa kulttuurikongressissa dots}
\kat{Karjalan Vapaus}{01.1935}{P59}{F. KUSMIN ka-ipaan\textasciicircum{}auppa Kangas-, Lyhyttavara- ja Valmiiden K A T P A A Vaatteiden Kauppa Puhelin 27 Liha-, herkku-, ruoka- ja KUOPIO - Kauppakatu 28 „, „.. porslunitav arotta Puhelin 842 - Yksityispuhelin 555 Sähköosoite: KUSMIN Sivumyymälät: Kaipaassa ja Annantehtaalla MATTI VÄNSKÄ S\textasciicircum{}\textasciicircum{}\textasciicircum{} ( nuorempi ) AIKKONEN OY. sekatavarakauppa S O R T AVA L A KOVERO - Puhelin Koveron keskus Sähköos. 223 Oraa keskus yhdis, osastoihin MOLOSOFKIN SUOJÅRVI, VÄLIKYLÄ Puhelin 28 Kangas-, lanka-, käsityö- ja muotiliike Villakutomo, r J -. -. -, „ SEKATAVARAKAUPPA Leninki- ja alusvaateompelimo Myy: Ruoka-, kangas- y.m. tavaroita edullisesti Edullinen ostopaikka kuluttaiille ia „ -- \_. „ >. „. \_. „. r J J Ostaa: dots}
\kat{Karjalan Vapaus}{05.1935}{P19}{Juho Dorojejeffin poikana. Pekka Jukola Samana syksynä muutti Suomeen.}
\kat{Karjalan Vapaus}{05.1935}{P19}{Otti osaa karjal. vapausl. Liikemies Pekka Jukola on synt. 29 / 1 1898 Mäntyselän pit.}
\kat{Karjalan Vapaus}{05.1935}{P19}{1921 — 1922. Kauppias Pekka Jukola Jo 19-vuotiaana hän joutui maailmansotaan.}
\kat{Karjalan Vapaus}{05.1935}{P19}{Sen vuoksi kauppias S. Saharoff 1917 tuli Suomeen, missä hän edelleenkin harjoitti samanlaista kauppaa Joensuussa m.m. tunnetulla nimellä Dobrinin \& Saharoff aina vuoteen 1935 tammikuun 1 p:ään, jolloin se lopetettiin. Myöhemmin v. 1930 aloitti itsenäisen liikkeen perustamalla Suojärven Lapinjärvelle sahatavarakaupan sekä harjoittaen puutavaraliikettä, mikä liike v. 1931 muutettiin Veljekset Jukola O / y-nimiseksi osakeyhtiöksi, jonka toimitusjohtajana on edelleenkin. Mutta sen sijalle hän samalle paikkakunnalle perusti taas omalla nimellään kangastavarain tukku- ja vähittäiskaupan, joka kauppias Saharoffin kauppa-alaan harjaantuneella älyllä, ja erinomaisella liikemisneroudella ja dots}
\kat{Karjalan Vapaus}{05.1935}{P19}{Mutta sen sijalle hän samalle paikkakunnalle perusti taas omalla nimellään kangastavarain tukku- ja vähittäiskaupan, joka kauppias Saharoffin kauppa-alaan harjaantuneella älyllä, ja erinomaisella liikemisneroudella ja järjestelykyvyllä alusta alkaen on saanut yhä laajenevan asiakaspiirin sekä jälleenmyyjistä että kuluttajista. Paavo Jukola Puh.}
\kat{Karjalan Vapaus}{05.1935}{P20}{Myy edullisesti SEKATAVARAKAUPPA, SUOJÄRVEN UUSI KE MI KALI KAUPPA Lapinjärvi / Puh. SUOJÄRVI / PUHEL. 17 WELJEKSET JUKOLA OT. • Suu -> hiuB; i a hajuvesiä Suojärvi / Puh. Osakeyhtiö W. A. Ipatof f TUKKU- JA VÄHITTÄISKAUPPA Nurmeksen joensuu Puhelin 73 Kauppa-Osakeyhtiö s i v v m y y m a i « u ENOSSA, ILOMANTSISSA, NURMES POLVIJÄRVELLÄ, TUUPOVAARASSA ja VIINIJÄRVELLÄ v. MOLOsoFKiN Kalpaan Kaappaosakeyhl SUOMBVI. VÄLIKYLÄ KilpAA, pDH „ Puhelin 28 + ♦ Liha-, Herkku-, Ruoka-, Porsliini-, Kangas- ja Lyhyttavaroita SEKATAVARAKAUPPA = \textasciicircum{}\textasciicircum{} \_ = \_ = = \textasciicircum{}\textasciicircum{} = \textasciicircum{}\textasciicircum{} = = MYY: SIVUMYYMÄLÄT: Ruoka-, kangas- y.m. tavaroita edullisesti Kaipaassa ja Annantehtaalla P Piihnnon Kaipaan Talouskauppa 1 • IILIItUILZ dots}
\kat{Viena-Aunus}{01.02.1935}{P22}{( Kansikuva. ) 1 I ( ivo ) H ( ärkönen ): Kymmenvuotismuisto 3 Jussi Sihvo: Liittolaisten Muurmannin retki 1918 — 1919. ( 2 kuv. ) 31 Toivo Nehvonen: Elämänruusu ja orjanruusu 32 Karjalaisjuhlat Kemissä 32 Salojen kävijä ( Paavo Jukola ): Kotiseutukuvauksia. ( Jatkoa. ) 38 Valter Päivinen: Eräs tullikahakka Aunuksen rajalla 20 vuotta sitten 39 I ( ivo ) H ( ärkönen ): Idän kirjoja.}
\kat{Karjalan Vapaus}{03.1936}{P43}{90 PUUTAVARALIIKE konttori Suojärvi as. VELJEKSET JUKOLA OY. Suojärvi * Puh. Sih loppih miehen elaiga.}
\kat{Viena-Aunus}{20.01.1937}{P5}{Hyvin usein nähtiin » batjuska » vieraana internaatissa. Aikamiehetkin silloin Jukolan veljesten tavoin istuivat Johannes Messenius — Suomen kronikan kirjoittaja.}
\kat{Karjalaisuuden Ystävä}{07.1938}{P3}{iSe imielen arreh olgah armahaime, igäine tunnus ] käis meän rahvahan, jeäv kuni yksgi moal viel mies vai naine, kel rinnas herjgi eläv Karjalan! TOIMITUSKUNTA: E. V. Ahtia, Ville Helve, Paavo Jukola ja Nikolai Ruotsi Otto Tossavainen.}
\kat{Karjalaisuuden Ystävä}{07.1938}{P4}{KARJALAISUUDEN YSTÄVÄ Sihteeri, kauppias Paavo Jukola. Osa Karjalaisuuden Ystäväin Liiton johtohenkilöitä}
\kat{Karjalaisuuden Ystävä}{07.1938}{P4}{Aktiivisen toiminnan tielle päästiinkin jo kuluvan v:n alussa, jolloin 13. 2. pidettiin neuvottelukokous kreikkalaiskatolisen kirkkounnan juhlasalissa Sortavalassa. Perustavaa kokousta välittömästi seuranneessa liiton varsinaisessa kokouksessa valittiin liiton puheenjohtajaksi apteekkari Ville Helve ja hallituksen jäseniksi kauppias Paavo Jukola, mv. 'Kautta vuosien jopa vuosikymmenienkin varsinkin rajan läheisyydessä asuvien 'karjalaisten keskuudessa kytenyt ajatus oman kansallisen. heimojärjestön aikaansaamiseksi on vihdoinkin saatu toteutetuksi.}
\kat{Karjalaisuuden Ystävä}{07.1938}{P4}{Sääntökomitea sai työnsä valmiiksi ja kutsui liiton 'perustavan kokouksen koolle Sortavalan Seurahuoneeseen 20. 3. Kun todettiin, että parikymmentä henkilöä oli noudattanut kutsua ja suunnilleen saman verran saatu ilmoituksia muilta asian kannattajilta, päätti kokous, jota johti maanviljelijä Pekika Kyöttinen, yksimielisesti ja innostuksen vallitessa perustaa » Karjalaisuuden Ystäväin Liitto. » nimisen iheimojärjestön ylempänä mainituin ohjelmin. Jukola ja xahastonhoitajaiksi liikkeenedustaja -edust.}
\kat{Karjalaisuuden Ystävä}{07.1938}{P5}{Liittoimime yirittääkin kohdistaa erikoisen harrastuksensa mainittuun osaan ohjelmastaan. Sihteeri: kauppias Paavo Jukola, os.: Rajasuvanto Rahastonhoitaja: teht. Se voi kompastellaikin opetellessaan astumaan, imutta se uskoo varttuvansa ajan ikera.}
\kat{Karjalaisuuden Ystävä}{07.1938}{P5}{Tällöin tietenkin edellytämme, että Ikieli. kehittyy oleellisesti karjalaiseksi eikä sellaiseksi sekamelsikaiksi, joksi se irajan takana iryssäläisellä ymppäyksellä aiotaan tietyssä tarkoituksessa muuntaa. Varsinaiset jäsenet: Liiton puheenjohtaja apteekkari V. Helve, varapuheenjohtaja opettaja J. Mustonen, kauppias ( Paavo Jukola, maanviljelijä Pekka Kyöttinen, luutnantti Paul Marttina, kauppias Vasili Molosovkin, johtaja I. Okulov. Uskomme on, ettei kalevalainen kieli vielä jouda arkistoihin homehtumaan, vaan että sillä on, kuten sen kannattajissakin, voimaa elää ja kehittyä ja että se kerran omalla luontaisella sitkeydellään on murtava ne kahleet, joihin sitä rajan taikana yritetään dots}
\kat{Karjalaisuuden Ystävä}{07.1938}{P24}{Emme osunekaan ihainhaan pääteissämme, että jos se työ, mikä nyt on tässä suhteessa telhty, olisi voitu suorittaa 1:5 vuotta ennen iheimosotia, olisi niiden tuloksetkin ehkä olleet toiset. Liiton puheenjohtajaksi valittiin tunnettu Karjalan ystävä apteekkari Ville Helve ja johtokunnan jäseniksi maanviljelijä P. Kyöttinen, opettaja J. Mustonen, kauppias P. Jukola, johtaja I. Okulov, kauppias V. Molosovkin ja luutnantti P. Marttina. Sikäli kuin heimosodat aikoinaan epäonnistuivat perimmäisten päämääriensä suhteen, oltiin tietyllä taholla valmiita väittämään, että ne olivat olleetkin vain edesvastuuttomia seikkailuretkiä, joiden ainoaksi tulokseksi muka jäisi se kärsimys, jota ne olivat dots}
\kat{Karjalaisuuden Ystävä}{07.1938}{P27}{S U O S I TELLAAN Kahvila ja Ruokala KOTI Veljekset Järvenpää Om. VELJEKSET JUKOLA OY. H. Saulamo puutavaraliike Sekatavarakauppa Suojärvi — Puh. Konttori 2, Myymälä 20 K E. S K U S * KAH V I L A Om.}
\kat{Karjalaisuuden Ystävä}{09.1938}{P3}{Kevät kägiey kezän tuomah, hengettömän eläväizeks luomaih; kovan kahlien kuoreh TOIMITUSKUNTA: E. V. Ahtia, Ville Helve, Paavo Jukola ja N. Ruotsi Kukat kaunehet kukoistetäh, ruohot muas yläih urgevutah, lemahus ladvois lembie.}
\kat{Karjalaisuuden Ystävä}{09.1938}{P27}{Naistenjarvi — Puhelin 7. SIVULIIKE MYY: Saha- ja höylätavaraa Kaitajärvi — Puhelin 4. Torajoen Saha — Puhelin 3. Edullisimmin kaikkia alaan kuuluvia tavaroita SUOJÄRVEN OSUUSLIIKE r. l. Suojärvi / Puh. WÄRZ OY. VELJEKSET JUKOLA OY. Kangas-, lyhy ». ja siirtomaatavarain PUUTAVARALIIKE tukku- ja vähittäiskauppa Suojärvi — Puhelin 90. SORTAVALA Puhelimet: konttoriin 525 RUOKATAVARAKAUPPA myymälään 173 Välikylässä - Puh. 4. Sortavala Suuri valikoima Salmin Kirjakauppa Salmi Halvat hinnat Suojärven Kirjakauppa Suojärvi A. PALONIEMI — i Puutavaraliike N. V. M A RUSKAIN EN SALMI \& KUM P P. OY Puhelimet 104 ja 105. SEKATAVARAKAUPPA OSTAA: Pystymetsiä ja Kr. „,.. \_ hankintapuita.}
\kat{Karjalaisuuden Ystävä}{12.1938}{P3}{95. c,., buuri valikoima SIVULIIKE TT,, naivat hinnat Rajasuvanto, Lapinjärvi — Puh. VELJEKSET JUKOLA OY. H. SAULAMO puutavaraliike Sekatavarakauppa Suojärvi — Puhelin 90. Suojärvi, Jehkilä — Puh. 38 Lieksa Puhelimet: Ruokatavarakauppa 143 SEKATAVARAKAUPPA Leipomo ja konttori 163 Runsas ja hyvä valikoima Hyviä vesirinkeliä ja korppuja a l aan kuuluvia tavaroita, tukuttain ja vähittäin.}
\subsection{Lisähaut}
\subsubsection*{Dorojejeff}
\kat{Karjalan Vapaus}{05.1935}{P19}{13 ja Oikok. 19 Juho Dorojejeffin poikana. Pekka Jukola Samana syksynä muutti Suomeen.}
\kat{Karjalan Vapaus}{05.1935}{P19}{Jo 1890, siis 22 vuotiaana, perusti hän oman tukkuliikkeen, joka myi kangastavaroita. Kun 1932 hänen toimesta perustettiin T:ni M. Pedas niminen liike, tuli hän sen johtajaksi. Toiminimi M. Pedaksen johtaja J. Dorojejeff on syntynyt Mäntyselän Karhumäellä 24 / 10 -94 mv. Joensuussa kangastavaroilla tukku- ja vähittäiskauppaa harjoittava kauppias S. Saharoff on syntynyt Rukajärvellä helmik. Tämän sahan palon jälkeen v. 1928 muutti Viipuriin ja sieltä v. 1929 sahanhoitajaksi Lapinjärven sahalle Suojärvelle.}
\subsubsection*{Karjalaisuuden Ystäväin Liitto ja E. V. Ahtia}
\kat{Toukomies}{12.1925}{P9}{Talon rakennuksista ovat muutamat jo hyvin vanhoja. sia ei näin pitkän matkan päähän voi kunnolla eroittaa. Vanhan karjalaisuuden jäljillä Suomen Pohjois\textasciicircum{}Karjalassa. Vanhin kaikista on tuolla suuren pihamaan äärimmäisessä Varsinainen Kajoon karjalaiskylä on täältä luoteeseen jylhien}
\kat{Toukomies}{12.1925}{P9}{Hänen nimensä mukaan on tämä talonpaikka saanut Timovaaran nimen ja tuota vaaran alla olevaa järveä sanotaan Rauanjärveksi. Raskas ja jylhän vankkumaton on tämän takamaan luonto, samanlaisen tuli olla sen raivaajien elinpaikankin. Tässä katkelmassa kuvattu Timovaaran vanha sukutalo oli viimeisiä ja lujimpia karjalaisuuden tyyssijoja Pielisen vaiheilla. Kaukana, kreivillisen Pielisen pokostan läntisellä äärellä on Kajoon laaja takamaa ja tämän takamaan napana on Timovaaran vanha sukutalo. Hyvin se nimensä ansaitseekin, jos ei lienekään niin tarumaisen rautarikas kuin tämän tienoon ruotsinuskoisten tulokkaiden kesken}
\kat{Toukomies}{01.1926}{P18}{Tähän » Toukomiehen » toimituksen asettamaan kysymykseen tahtoisi allekirjoittanut vastata muutamin sanoin. Onko karjalankielen kirjallinen viljelys suotava ja tarpeellinen — sitä asiaa voipi tarkastaa sekä itse karjalaisten että suomalaisten kannalta. Toiselta puolen vaatii myös karjalaisten henkisen erikoisuuden — karjalaisuuden — suojeleminen kielellistä omintakeisuutta ja siten omaa sivistyskieltä ja kirjallisuutta. Karjalaisista tuskin voinee olla yhdentekevää, tuleeko heidän kielellinen erikoisuutensa säilymään, tulevatko he jäämään kielellisestikin eikä vain nimellisesti karjalaisiksi. Nykyään ei mikään kansa, joka pyrkii edistymään ja nousemaan sivistyksessä, voi, ajan pitkään dots}
\kat{Toukomies}{01.1926}{P18}{Ja mitä olisi suomenkielen ( ja siis suomalaisuudenkin ) tulevaisuus ilman suomalaista kirjallisuutta, joka yksin säilyttää ja yhä uusille sukupolville välittää perittyä kansallista sivistystä ja koko henkistä elämää! Sillä jokaisella kansalla oma kieli ei ole ainoastaan sen henkisen omituisuuden uskollisin ja tarkin ilmaus, vaan myös sen välttämätön suojelija ( missä eivät rotu tai erikoislaatuinen maantieteellinen asema eroita sitä naapurikansoistaan ). Samoin jos karjalaisuuden on mieli säilyä henkisenä, vaikuttavana tekijänä ja omalaatuisuutena, niin se voi tapahtua vain siten, että karjalankieli saapi arvonsa ja tulee välineeksi Kalevalan kansan omintakeisuuden ylläpitämisessä ja dots}
\kat{Toukomies}{01.1926}{P19}{Sulautua sensijaan läntisiin heimolaisiinsa hedelmöittää suomalaisuutta omalla erikoisuudellaan, se on karjalaisten oikea ja ainoa tehtävä. Ilmoitetaan myös, että kehityskykyisemmän nuorison jatko-opetuksesta on huolehdittu maanviljelyskouluissa, seminaareissa y.m. ammattioppilaitoksissa. Eikö päinvastoin todennäköinen tulos olisi senkin tehon heikkeneminen, mikä tähän asti on voinut syntyä karjalaisuuden ollessa vielä alkuperäisessä koskemattomuudessaan? Asianomaiselta taholta on ilmoitettu että Suomessa nykyjään Venäjältä saapuneita pakolaisia on yhteensä n, 21,000. Puolet näistä pakolaisista on varsinaisia venäläisiä. Karjalankielen sana- ja muotorikkaus, äänteellinen sulavuus, dots}
\kat{Toukomies}{01.1926}{P19}{Viime vuosina on näiden pakolaisten luku, kuten tunnettua, huomattavasti vähentynyt, heitä on verrattain paljon palannut takaisin kotiseuduilleen. jonka talon rakentamisen on arvioitu tulevan maksamaan noin io miljoonaa markkaa, on Joensuun kaupungin valtuusto päättänyt ottaa omana asiana ajaakseen. poistaisi tulevaisuudessa kaikki edellytykset vuorovaikutukseen kahden heimokansan välillä, mikä muuten epäilemättä yhä paremmin jatkuisi ja edistyisi karjalaisuuden päästessä omille jaloilleen. Enin osa näistä pakolaisista asuu Viipurin läänissä, jossa niitä lasketaan olevan 7 — 8,000. Uudenmaan läänissä enimmäkseen Helsingissä on venäläisiä pakolaisia n. 2,000. Aivan toinen merkitys dots}
\kat{Toukomies}{01.1926}{P19}{Päinvastoin pitävät enimmät suomalaiset äänet sitä karjalaisille turhana, jopa vahingollisena, se kun on heille muka ylivoimaista, heidän heimohenkeään heikentävää ja sivistystään hidastuttavaa. Inkeriläiset, jotka melkein kaikki, vanhuksia ja orpolapsia lukuunottamatta, huolehtivat itse toimeentulostaan, asuvat pääasiallisesti Kannaksella, Raudun, Metsäpirtin, Kivennavan y.m. pitäjissä. Valitettavasti ei ole suomalaisten puolella osattu edellämainitusta kiittävästä arvolauseesta vetää oikeaa johtopäätöstä — karjalaisuuden jatkuvan elämisen oikeutusta ja sen kielen ja kirjallisuuden tarpeellisuutta sen kannattajana. Karjalaiset pakolaiset, joita on n. 7,500 henkeä, asuvat etupäässä pitkin dots}
\kat{Toukomies}{03.1926}{P7}{Vertailu puhtaaseen suomeen kääntyy venäläisiä lausetapoja ja -sanoja täynnänsä vilisevälle, osaksi suorastaan teennäisesti konstruoidulle tai konstruoitavalle seka- tai murrekielelle epäedulliseksi. Se on enimmäkseen n. s. graniittigneissiä — Karjalan alkuvuorta, sanoo J. E. Rosberg, ellei hän sen nuorempia muodostumia kunnioita aivan kalevalaisvuorilajien nimellä, Kalevalan runojen mukaan. Tämä päätös ja sen seuraukset eivät suinkaan tarvitse sisältää karjalaisuuden kuolemista ja mahdollista » savolaistumista », vaan se voi olla päinvastoin pohjana kansan vähittäiselle entistä korkeammalle sisäiselle itsenäistymiselle. Ja mitä kieliainekseen itseensä tulee, havaitsemme sen äänteissä, dots}
\kat{Toukomies}{03.1926}{P17}{— Toukomiehen maaliskuun numero sisältää: Italian kirkas kevättalvi: näköala Ficsolen luostaripihalta, kansikuva; esittely ja uutisia; Talviajatuksia, kirj. livo Härkönen; Talvilämpöä: Väinö Hämäläinen: Lieden ääressä, taidekuva; Karjalainen ääni, kirj. Se tarvitsee koulutaloja, kirkkoja, kirjastoja, lukukeskuksia; sen rautatiet ja maantiet ovat vasta alullaan; sen oma sanomalehdistö on luomatta; sen teollisuusyritykset ja -laitokset ovat vähäisiä ja aivan olemattomiakin; sen maanviljelys ja luonnonrikkaudet odottavat käyttäjäänsä, kehittäjäänsä. — Kirjailija Simo Erosen mieltäkiinnittävä kaunokirjallinen kuvaus » Kuhasalon luostarin synty ja häviö, vanhan karjalaisuuden kohtaloa », tohtori dots}
\kat{Toukomies}{04.1926}{P9}{» Totikomiehelle » kertonut Ilja munkin työ kantoikin hyviä hedelmiä. Vanhan karjalaisuuden kohtaloa. Tätä vähäistä saarta nimitettiin entisaikaan Kuhasaloksi. Kuhasalon luostarin synty ja häviö.}
\kat{Toukomies}{04.1926}{P15}{— Tähtisarjassa on ilmestynyt A. C. Askewin kirjoittama romaani » Naisen kunniasana », englantilainen tarina nuoresta tytöstä. Nykyinen savolaisperäinen väestö kai unohtaa mielellään ne tarinat, joissa kerrotaan heidän esivanhempiensa karjalaisveljille tekemistä teoista. ja 20-vuotisjuhla; Karjalaisaiheista kirjallisuutta; Kuhasalon luostarin synty ja häviö, vanhan karjalaisuuden kohtaloa, hist. kertomus, kirj. — Kirsikka-sarjassa on uutena numerona ilmestynyt jännittävä seikkailu- ja rakkausromaani Amerikan Suuresta Lännestä Harold Bindlosin kirjoittama » Karjakuninkaan tytär ». Simo Eronen; kuvia Kuolan niemimaalta, aukeama, 10 kuvaa; Kibunoi, II sarjan sivut 9 — 16; Karjalainen esitelmä dots}
\kat{Toukomies}{12.1926}{P28}{Suomen Karjala, yleisesitys, kirj. livo Härkönen. Rajantakaisen Karjalan porteilta ja Petroskoilta, 4 kuvaa. Kuhasalon luostarin synty ja häviö, vanhan karjalaisuuden koh- Menneelle työlle ja työntekijöille, osa livo Härkösen puheesta Talvilämpöä: Väinö Hämäläinen: Lieden ääressä, taidekuva.}
\kat{Toukomies}{01.02.1927}{P12}{Niin kirjat sanovat, viisaat ja oppineet kirjat, mutta minä en usko niitä. Minä kaavin sen puhtaaksi ja sitten sovitan jalkani jyrkkäreunaisesti kovertuneihin astuimiin. Mutta sitä suuremmalla hartaudella syvennyn hävinneen karjalaisuuden muistoihin. Siellä he saivatkin elostaa kenenkään ahdistamatta metsänriistaa pyytäen ja vedenviljaa kalastaen. Tämä Jumalan suuri ihmetyö sai pyhän miehen, kokonaan unohtamaan äskeisen mielenmasennuksensa.}
\kat{Suomen Heimo}{18.11.1930}{P4}{Kansanliike on jälleen saatava kansan liikkeeksi. Sillä sen verran Lapuan liike on varmasti jo tehty puolueasiaksi. Mutta karjalaisuuden ei tarvitse nähdä suomalaisuudessa kilpailijaa, saatikka vastakohtaa. Nämä epäluulot ovat todella häirinneet yhteisymmärrystä karjalaisten ja suomalaisten välillä. Samaten ovat jotkut karjalaiset vanhemman polven miehet suhtautuneet suomalaisiin epäluuloisesti.}
\kat{Suomen Heimo}{18.11.1930}{P5}{Joka väittää, että suomalaiset ovat käyttäneet asemaansa väärin, sen on myös todistettava väitteensä. Tämä ansaitseminen ei tapahdu — se on syytä muistaa — pelkällä heimolaisuuteen ja avuntarpeeseen vetoamisella. Kummallisinta onkin, että monet tämän karjalaisuuden innokkaimmista ajajista ovat lähtöisin Suomen puolelta. Tässä on samalla kysymyksessä periaate eikä suinkaan valtapyrkimys, kuten karjalaisille tahdotaan uskotella. On myönnettävää, että karjalaisilla rajan kummaltakin puolen on enemmän yhteistä kuin rajantakaisilla ja muilla suomalaisilla.}
\kat{Suomen Heimo}{18.11.1930}{P5}{Sillä Suomen ja Karjalan suhde on veljen ja veljen suhde, Suomella ei ole oikeutta eikä syytä ottaa Karjalaan nähden kiitollisuudenvelkaa karhuavan äitipuolen asennetta. Kerhotyössä ja valistustoiminnassa ovat suomalaiset tukeneet ja auttaneet pakolaisveljiensä itsetoiminnallisia pyrkimyksiä, enkä tiedä, että heidän keskensä olisi syntynyt asiallista erimielisyyttä taikka suurempia väärinkäsityksiä. Kun karjalaiset esim. luottamushenkilöitä valitessaan pitävät nimenomaan kiinni asianomaisten karjalaisuudesta, mutta lukevat siihen myöskin Suomen puolella esiintyvän maakunnallisen karjalaisuuden, paisutetaan tarpeettomasti veljien välistä eroavaisuutta. Kun joillakin johtavilla karjalaisilla dots}
\kat{Toukomies}{01.09.1930}{P19}{Vasta viime toimikauden aikanahan saavutettiin kerhojen välille yhteisymmärrys, mikä oli toteutettavissa vain siten, että Liiton toimintaa ohjattiin ankarasti sen sääntöjen rajoissa, eikä satunnaisissa ja karjalaisten pyrkimyksille vieraissa tarkoitusperissä. Me uskomme siitä huolimatta, vaikka vastustajamme nimittävät karjalaisia » laumaksi, jonka ajattelu ei kanna oman reppanan ulkopuolelle » ( hra Takalan sanat ), että karjalaisissa on jo senverran oman arvon tuntoa ja itsekuria, että he kykenevät kylmästi torjumaan moiset eksytykset. Karjala-kerhojen Keskusliiton johtokunta ei muuta suuntaa eikä aja mitään muuta tarkoitusperää kuin karjalaisuuden voimistumista ja yhtenäistymistä, dots}
\kat{Toukomies}{01.09.1930}{P23}{Onkohan laittomuudesta puhujain kädet vain niin puhtaat? Ketkä nuo viisi olivat, ne kyllä hyvin tunnetaan. Mistä te olette saanut karjalaisuuden elinkautisen valtakirjan? Mitä tulee vaalin laittomuuteen, on siitä vastaus säännöissä. Mutta täytyy kysyä, mihin tuollainen valheellinen puhe perustuu.}
\kat{Toukomies}{01.04.1931}{P26}{Hänestähän, poikasesta, miehet puhuvat ja pappi sanoo lähtiessään: » Poika on hyvä oppilas koulussani ja minä toivon hänestä jotakin muuta kuin tavallista karjalaista miestä. Siksi puoltanevat paikkaansa tässä karjalaisten sivistysrientojen merkkipäiväksi ilmestyvässä juhlavihkossa nämäkin tämäntapaiset hajanaiset muistelmat, joista kuvastunee piirre tai toinen asiamme eri vaiheista. Pienet aivot työskentelevät edelleen ja ajatukset saavat omituisen käänteen: » Joudun venäläiseen kouluun, venäläiseen ympäristöön, venäläiseksi opettajaksi, venäläisten uskotuksi ja luottamusmieheksi, tämän rahvaan venäläistyttäjäksi ja — oman kansallisuuteni, karjalaisuuden tappajaksi. » — yhdellä dots}
\kat{Toukomies}{01.09.1931}{P12}{Tämän runon » sammon taonta » -nimisestä osasta lausuu kuuluisa tutkija Ohrt eräässä teoksessaan, ettei tätä runoa ole koskaan tavattu Vienan läänin ulkopuolelta. Orjat lietsoi löyhytteli, Väkipuolet väännätteli Kolme päiveä kesäistä Ja kolme kesäistä yötä: Kivet kasvoi kantapäihin, Vahat varvasten sijoihin. Käsittäen karjalaisten kuntoisuudessa piilevän koko Karjalan tulevaisuuden taian, en saata olla jatkamatta edelleen tätä ajattelua, karjalaisuuden itsetutkistelua. Sillä onhan ilmeistä, että sen kansallisten voimien ja mahdollisuuksien aliarviointi varmaan johtaisi peloittavaan heikontumiseen, samoin kuin yliarviointi veisi helposti pahoihin pettymyksiin. Syventyminen omaan itseensä, dots}
\kat{Toukomies}{01.09.1931}{P12}{Näin ollen kehoituksemme odottaa ei tuntunekaan enää niinkään tyhjältä ja tarkoituksettomalta, varsinkin, kun se sisältää samalla eräänlaisen lämpimän, toivorikkaan tulevaisuuden Pettämättömimpänä peruskirjana kalevalaisen kansan luonteenominaisuuksien tutkimuksessa on pidettävä kansanhengen omaa välittömintä tuotetta, Kalevalaa, jonka jokainen sivu puhuu siitä selvää, ainutlaatuista kieltään. Takaperoista on, mutta totta: Maammekirja lienee ollut monelle meikäläiselle melkein ainoa lähdekirja karjalaisuuden tuntemiseen kymmenisen vuotta takaperin; karjalaisten kirjalliset tietolähteet omasta Muistanemme elävästi, että n. s. sortoajat olivat omansa kokoamaan ja lujittamaan suomalaisten dots}
\kat{Toukomies}{01.02.1932}{P3}{Vienan vetten äärille Kantalahdessa, Sumassa ja Ouegajoen suilla ja pieniksi saarekkeiksi Kuolankin niemelle ulotti heimo asutuksensa pohjoisessa. Ja niin taas kerran oli koko itäinen karjalainen maa vanhan vallanpitäjän hallussa ja raja kulki siinä, mihin, se oli Stolbovan rauhassa saneltu. Aina Vienanjoen takaisille maille sen eräretkeilijät ulottivat matkansa idässä, ja toisaalta Savon yli edeten ovat antaneet aihetta etsiä karjalaisuuden jälkiä aina Pohjanlahden rannoilta. Tämä heimous on siitä paljon supistunut. uMianmeren länsirannan ja Äänisen länsirannan muodostaman viivan toisella puolen on jo kolmisen vuosisataa ollut kaikumatta karjalainen Valdain harjuilta lähtien levisi se dots}
\kat{Toukomies}{01.03.1933}{P13}{Ja niinpä hän päättää teoksensa seuraavalla lauseella: Suomenpuoleinen Karjala esittäytyy ensin ja sitten seuraa rajantakainen Karjala. Okkosen kirjoitus teoksessa osoittaa karjalaisuuden vaikutusta Suomen kuvaamataiteeseen. Kuivankin elämäkerta-aineiston hän on osannut elävöittää miellyttäväksi ja tempaavaksi. Karjala ja karjalaisuus on yhteissuomalaiseen sivistykseen tuonut huomattaviakin tunnusmerkkejä.}
\kat{Toukomies}{01.04.1933}{P10}{Karjalan Sivistysseura on järjestänyt uuden vaikutusalueilleen lähetettäviä kirjastojaan ja tehdyn anomuksen johdosta myöntänyt sen äskenperustetulle lisveden Karjalakerholle, joka on liittynyt seuran alaosastoksi. Suomenkielinen jumalanpalvelus, suomenkielinen uskonnollinen kirjallisuus, kansallinen mieli papiston keskuudessa, uskonnon yhä suurempi lähentyminen suomenkieliseen kansaan — ne kaikki olivat lähellä hänen harrastuksiaan ja tarkoitusperiään. Sitä suurempi syy on pitää yllä niitä julkisen sanan elimiä, jotka ovat jaksaneet pysyä pystyssä ja pitää vireillä karjalaisuuden tulta ja sen pohjalta noussutta karjalaista mieliä, omaa kansallistuntoa, yhteistä isänmaallista henkeä. dots}
\kat{Karjalan Vapaus}{10.1934}{P14}{Sillä varmaa on, ettei meitä kukaan auta, ellemm me itse ole uhrivalmiita ja yritä parastamme asiamme hyväksi ja sen ratkaisemiseksi. Sillä tällaisen uuden, onnellisempaa tulevaisuutta kuvastavan aatteen synnyttävät aikain kuluessa mielestäni vain kulloinkin vallitsevat epämieluisat olot. Aunus on seisonut raskaammassa venäläispuristuksessa kuin muu osa Karjalaa, on se sittenkin alkuperäisen karjalaisuuden kuvastin niin kielellisesti kuin mielellisestikin. Kansallisuusliikkeen kehitysmahdollisuudet karjalaisilla olivat etupäässä varsin löysän ortodoksisen uskonnon suojassa, uskonnon, joka sinne tuotiin yli 700 vuotta sitten. Tällä toimenpiteellään he tietämättään ja tahtomattaan kasvattivat dots}
\kat{Toukomies}{05.02.1925}{P7}{Myös kerrotaan, kuinka v. 1773 saapui Pietariin eräs papinehdokas Poventsan Karjalasta, tullakseen vihityksi papinvirkaan. Näin kertoo eräs kirje, josta tuonnempana enemmän, mutta nähtävästi tämäkin toimenpide jäi tähän. Maisteri E. Ahtia kirjoittaa mielenkiintoisesti Karjalan Kirjassa erityisestä karjalan- ja aunuksenkielisestä kirjallisuudesta. Nyt saattoi tulla Pelguille ajatus kukistaa samalla kertaa sekä ryöppyilevä Novgorodin valta että veneillä vennosti soutava Ruotsin ristivältä. Ja niin on Karjala ja Inkeri saatu ikipäiviksi vapaaksi tai ainakin tästä pahalta näyttävästä ristituliryntäyksestä pelastetuksi.}
\kat{Toukomies}{05.02.1925}{P8}{Matteuksen evankeliumi Tverin Karjalan murteella. V. iByo ilmestyi kolmas karjalankielinen kirja. Ylläolevaan tapaan maisteri Ahtia. — Kirjasen hinnaksi on merkitty rupla. Eräs Tihonov julkaisi karjalais-venäläisen rukouskirjan aunuksen murteella.}
\kat{Toukomies}{01.09.1925}{P6}{29-sivuisena ei se paljoa annakaan, mutta alkuopastajan tarkoituksenhan se voisi täyttää, jos se olisi asiantuntemuksella toimitettu; pääasiassa se kuitenkin voisi olla pienoinen, alkeellinen venäläiskarjalainen sanakirja. ja kohta seuraavina vuosina aletaan samaa rajantakaisen Karjalan osaa varten julkaista kaikkia evaneliumejakin, joista ensimmäinen, Matteuksen evankeliumi, julkaistaan v. 1896 90-sivuisena ja viimeinen, Johanneksen evankeliumi, ilmestyy v. 1900 72-sivuisena. Näistä tutkiskelijoista suomalaisella taholla on ennenkaikkea mainittava harras karjalan kielen ystävä ja myöhemmin jo itsekin karjalankielisen kirjallisuuden kartuttamiseen ryhtyvä maisteri E. V. Ahtia, jonka dots}
\kat{Toukomies}{01.09.1925}{P7}{Vielä ennen kuolemaansa ennenaikaisesti poismennyt tutkija kerkisi laatia ensin julkaisemastaan teoksesta » Karjalan äänneoppi » uuden laajennetun laitoksen, joka ilmestyi v. 1918 ja josta tekijä itse sanoo, että se on kolme kertaa edeltäjäänsä laajempi ja siten siitä melkoisesti eroava. Niinpä salmilainen maisteri Jo h. Kujola — synnyltään siis jo karjalankielen, vieläpä aunuksen piiristä — v. 1910 julkaisi ansiokkaan ja tarkasti tieteellisen esityksen » Äänneopillinen tutkimus Salmin murteesta », jonka kielennäytteet ja lukuisat sanaesimerkit ovat merkinnältään niin tarkkaa karjalankielen kirjoitustapaa, että Perusteet ovat samat, mutta » ensiksikin on käytetty hyväksi niitä runsaita dots}
\kat{Toukomies}{01.09.1925}{P7}{Vielä ennen kuolemaansa ennenaikaisesti poismennyt tutkija kerkisi laatia ensin julkaisemastaan teoksesta » Karjalan äänneoppi » uuden laajennetun laitoksen, joka ilmestyi v. 1918 ja josta tekijä itse sanoo, että se on kolme kertaa edeltäjäänsä laajempi ja siten siitä melkoisesti eroava. Niinpä salmilainen maisteri Jo h. Kujola — synnyltään siis jo karjalankielen, vieläpä aunuksen piiristä — v. 1910 julkaisi ansiokkaan ja tarkasti tieteellisen esityksen » Äänneopillinen tutkimus Salmin murteesta », jonka kielennäytteet ja lukuisat sanaesimerkit ovat merkinnältään niin tarkkaa karjalankielen kirjoitustapaa, että Perusteet ovat samat, mutta » ensiksikin on käytetty hyväksi niitä runsaita dots}
\kat{Toukomies}{01.1926}{P25}{Venäjänikiel. Aunuksenkiel. E. Ahtia. 3. Rajantak. sen loppuvuosina.}
\kat{Toukomies}{01.06.1928}{P17}{Tähän on hänelle annettu asianmukainen vastaus. Näin siis arvovaltaiset suomenkielen ja kirjallisuuden vaalijat. Hra Ahtia oli sitä mieltä, että karja- Kokeiluja, hartaita, kauas tähtääviä yrityksiä olisi jatkettava. Kielitieteen, kotiseutututkimuksen ja kauneusarvojen kannalta}
\kat{Toukomies}{01.06.1928}{P17}{Lisättäköön tähän vielä, mitä professori Koskenniemi aikoinaan » Ajassa » lausui silloin ilmestyneiden livo Härkösen karjalankielisten käännösten johdosta, jotka häner mielestään » ikäänkuin laulavat itsensä ' ) Tämäkin kirjoitus oli alkumuodossaan tarkoitettu lisäselvi- tykseksi saman inkeriläisen kynäntuotteen johdosta, mutta ei pääs- syt päivänvaloon saman lehden palstoina, jossa mainittu kysy- mys oli esitetty. Senj aikeen ovat tätä kysymystä käsitelleet Toukomiehessä lehtori E. V. Ahtia ( Toukomies n:o I — 21 — 2 v. 1926 ), professori Onni Okkonen sekä hra P. Patokoski, selostaen kukin kysymystä omalta näkökannaltaan. Ei lähestulkoonkaan vielä ole pelkoa karjalankielisen kirjallisuuden dots}
\kat{Toukomies}{01.06.1928}{P18}{lukumäärä voisi kohota samoin kuin esim. Provencen ja Catalonian kielimurteiden näytteet Ranskassa ja Espanjassa sekä Sisilian Italiassa, kuten professori Okkonen mainitussa lausunnossaan sanoo olevan. Jos suuren kansalliseepoksemme ja muiden kansanrunousteostemme lisäksi saadaan luoduksi uusia, vaikkapa murrekielisiä kansanrunouskokoelmia, on niiden ilmestyminen vain rikastuttava, uusia arvoja luova ja uusia lähteitä avaava. Tässä, mielessä olisi ennen kaikkea saatava talteen karjalanmurteinen kansanrunous, joka on monipuolinen ja väririkas ja josta maisteri E. V. Ahtia on julkaissut m. m. mielenkiintoisen kokoelman » Rahvahan Kandeleh. » Kielen varsinainen perusta on sen kirjallisuus, dots}
\kat{Toukomies}{01.09.1928}{P5}{Jokainen sana ja sen tärkeimmät taivutusmuodot — ja mahdollisesti myös kerääjälle tunnettujen eri murteiden mukaiset toisinnot — samalle lipulle, se on tämän työn ensimmäinen sääntö ja se helpottaakin ja säännöllistyttää sitä suuressa määrin. Mutta varsinainen, suurempi merkityksensä olisi sellaisella sanakirjalla, josta yksinkertaisesti joka murteen sana tai paremmin sanoen sanatoisinto voitaisiin löytää — yhteisen hakusanan kohdalla mahdollisestikin, mutta kuitenkin niin, että jokainen samankin sanan ja käsitteenilmauksen eri muoto tavattaisiin. joista mahdollisesti vielä toiste. ) Karjalan- » kielistä » kirjallisuutta on jo meillä joku verta, kuten tässä lehdessä olemme aikaisemmin jo dots}
\kat{Toukomies}{01.12.1932}{P8}{Hänen vahvuutensa on lähinnä Suojärven murie, jonka lisäksi hän aikaisempina aikoina kävi läheisiä Suojärven-takaisia heimomaita tutkimassa. Tämän alkuun panemansa erikois- ja harvinaislaatuisen työn innostamana ryhtyi hän v:n 1907 — 1908 vaiheilla kääntämään ja saattamaan suomalaisia kirjoituksia aunus-suoj arven murteelle. Tämän jälkeen maisteri Ahtia 1909 julkaisi pienen vihkosen » Opastus Livviköil », jossa oli suomen ja venäjän kirjaimet ja kymmenkunta sivua suomalaisten sanojen selityksiä aunuksenkielellä. Viimeksimainittu oli luonnollisesti päämääränä, sillä käytännöllisesti ja taloudellisesti katsoen ei ollut mahdollista, että lopullinen, kokonainen aunuksenmurteinen kirjallisuus dots}
\kat{Toukomies}{01.12.1932}{P8}{Pian siirtyi h.?. a suurempaan työhön kääntäen ja saattaen karjalaksi julki koko Luukkaan evankeliumin: » Hospodan meijän lisus Kristan Pyhä Jevanheli Lukaz Karjalan ( livvin ) kielel » ( Sortavala, 1912 ). Kohta julkaisi hän joitakuita raamatun kohtia kahdessa samansisältöisessä pikku vihkosessa: » Bibliiz, Jumalan sanas, lugietah » ja » Hospodan meijän lisusan Kristan Pyheä jevanhelie », joista viimeksimam. painettiin Englannissa. Maisteri Ahtian käännöstyö on liikuttava osoitus sekä Karjalan rakkaudesta että yrityksestä kääntää venäläistyttäj äin samoihin alkama kansan vieraannuttamiseksi tarkoitettu karjalankielinen julkaisutyö vaikutukseltaan päinvastaiseksi ja kohdistaa * sen aseet dots}
\kat{Toukomies}{01.12.1932}{P8}{Kohta julkaisi hän joitakuita raamatun kohtia kahdessa samansisältöisessä pikku vihkosessa: » Bibliiz, Jumalan sanas, lugietah » ja » Hospodan meijän lisusan Kristan Pyheä jevanhelie », joista viimeksimam. painettiin Englannissa. Maisteri Ahtian käännöstyö on liikuttava osoitus sekä Karjalan rakkaudesta että yrityksestä kääntää venäläistyttäj äin samoihin alkama kansan vieraannuttamiseksi tarkoitettu karjalankielinen julkaisutyö vaikutukseltaan päinvastaiseksi ja kohdistaa * sen aseet itseensä venäläistyttäj iin. Tämän niinsanoaksemme suomalaiskansallisen kar j alai - sen murrekirj allisuuden pani alkuun eräs uuras karjalankielen harrastaj a j a sanaston kerääjä, maisteri E. Ahtia, joka dots}
\kat{Toukomies}{01.12.1932}{P8}{Tämän niinsanoaksemme suomalaiskansallisen kar j alai - sen murrekirj allisuuden pani alkuun eräs uuras karjalankielen harrastaj a j a sanaston kerääjä, maisteri E. Ahtia, joka sukulaisuussiteillä jo varhain kiintyi Karjalaan ja sen rajanläheisiin seutuihin. Hänen työnsä siinä mielessä ■ — josta syystä hänkin käytti venäläisten kirjoitustapaa, venäläisin kirjaimin ■ — on muistettava näiden murros vuosien, suurten unelmain ja monenmoisten isänmaallisten yritysten ajoista, j olioin Karj ala ja Suomi lämpimissä mielissä olivat toisiaan lähem- Välillä julkaistuaan joitakuita pieniä runoja karjalaksi Ahtia v. 1920 sitten julkaisi aunukseksi eri kirjasina Markuksen ja Johanneksen evankeliumit ja dots}
\kat{Toukomies}{01.12.1932}{P8}{Välillä julkaistuaan joitakuita pieniä runoja karjalaksi Ahtia v. 1920 sitten julkaisi aunukseksi eri kirjasina Markuksen ja Johanneksen evankeliumit ja vihdoin vuonna 1922 laajahkon kokoelman maallisia lauluja nimellä » Rahvahan kandeleh » sekä 1924 uskonnollisten runojen käännöskokoelman » Vieron virzie ». ( Näitä ennen, vuodesta 1916, oli jo alkanut allekirjoittaneen kokeilu karjalankielisten runokäännösten alulla, jolloin aluksi julkaisin Eino Leinon toimittamassa » Sunnuntaissa » Jänneksen » Väinölän lapset » ja » Karjalan » aunus-karj alaksi, sen jälkeen jatkoin Karjalan Sivistysseuran » Karjalaisten Sanomissa » sekä runokäännösten että alkuperäisten suorasanaisten kirjoitelmain dots}
\kat{Toukomies}{01.12.1932}{P9}{Venäjän vallan aikuisiin esityksiin painoi leimansa myöskin se, ettei kirjoittaja saanut vapaasti puhua esim. sellaisista asioista kuin venäläistyttämistyöstä Karjalassa ja tämän kansan vapaudenkaipuusta. Siinä esitetään karjalaisten asuinpaikat ja lukumäärä, ulkomuoto, rakennukset, puvusto, elinkeinot, tavat, luonne, kieli, muinaisusko j. n. e. Kun kirjoittajat ovat erikoistuntijoita aloillaan, ovat esitykset luotettavia. Erikoisesti pantakoon merkille Aarne Äyräpään selväpiirteiset esitykset kivi- ja pronssikaudesta ja U. Karttusen, E. Ahtian — J. J. Mikkolan y. m. valaisevat kirjoitukset Suomen Karjalan, Vienan ja Aunuksen historiasta. On ymmärrettävää, ettei v. 1910 ollut sellaisia dots}
\kat{Toukomies}{01.02.1933}{P11}{Läpi historian kulkee eräitä piirteitä, joiden vaikutus kokonaisuuden kehitykseen on epäilemättä lyönyt lähtemättömän leimansa ja joita piirteitä vieläkin on varsin voimakkaasti havaittavissa ja vaikuttamassa asiain kulkuun. Toisena erikoispiirteenä kivikauden lopulta mainitaan Karjalan tuotteissa havaittava koristeellisuus, joka ilmeni hirvenpääkirveinä, vaikka kivikausi muualla jo alkoi ikäänkuin rappeutua ja alkoi siirtyminen pronssi- ja rautakauteen. Karjalan historian myöhemmät vaiheet esitetään kahdessa kirjoituksessa: U. Karttusen kirjoituksessa Suomen Karjalan ja J. J. Mikkolan tarkastamassa Ahtian y. m. aikaisempien kirjoitusten mukaan laaditussa kirjoituksessa rajantakaisen dots}
\kat{Karjalan Vapaus}{10.1934}{P24}{J. Maksimainen. Karjalaisuusliikkeen vaiheita, kirj. E. V. Ahtia. A. O. Pajari. Karjalan värit ja vaakuna, kirj.}
\kat{Karjalan Vapaus}{01.1935}{P19}{I Päivällinen 8: — lipuilla. E. V. Ahtia. vaikkeivät olleet toisilleen suvun sukuakaan. » Ei nimi miestä pahenna », sanotaan.}
\kat{Karjalan Vapaus}{01.1935}{P57}{'kanslisti P. Harju 37 Pogostan komun Leninets, kirj. 42 Karjalan ja Länsipohjan kosketuksia, kirj. asemapääll. Viesti Vepsanmaasta, kirj. maisteri E. V. Ahtia... Kuollut oli Karjalan kansan viimeinenkin toivo.... V. O. Jehrimäini 28 Neittä neuvotaan, esittänyt Nasto Huotarinen...}

\vspace{1.2em}{\color{leima}\hrule height 1pt}\vspace{0.5em}
{\footnotesize\color{harmaa}
\textbf{Menetelmästä.} Haut Kielipankin Korp-palvelun kautta Kansalliskirjaston
digitoituun lehtiaineistoon (KLK\_FI), vuosikorpukset 1925--1939. Aineisto on
vapaasti käytettävissä vuoteen 1939 asti. Kaikki hakulausekkeet liitteessä A.

\textbf{Luonteesta.} Tämä on työaineisto mahdollista kirjaa varten, ei valmis tutkimus.
Luvut ovat koneellisesti tuotettuja; tulkinnat ovat tekijän eivätkä ole vertaisarvioituja.

\textbf{Kirjoittajasta.} Juri Jukola on teollisuuden operatiivisen johdon tehtävissä yli
25 vuotta toiminut johtaja, diplomi-insinööri ja kauppatieteiden maisteri sekä hyväksytty
hallituksen jäsen (HHJ). Hän ei ole historioitsija. Romaanit ilmestyvät kirjailijanimellä
J. Lintumies.}

\end{document}
